\documentclass[11pt]{article}
\usepackage[margin=2.6cm]{geometry}
\usepackage{booktabs}
\usepackage{graphicx}
\usepackage{siunitx}
\sisetup{group-digits=integer}
\usepackage{amsmath}
\usepackage{hyperref}
\usepackage{seqsplit}
\graphicspath{{figures/}}
\setlength{\emergencystretch}{3em}
\newcommand{\runid}[1]{\texttt{\small\seqsplit{#1}}}
\newcommand{\dsid}[1]{\texttt{\small #1}}
\newcommand{\best}[1]{\textbf{#1}}
\newcommand{\medape}{medAPE}
\newcommand{\rsq}{$R^2$}
\newcommand{\Rp}{\ensuremath{R'}}
\newcommand{\Cp}{\ensuremath{C'}}

\title{Spotify Engagement from Metadata: Experiment Synthesis}
\author{}
\date{\today}

\begin{document}
\maketitle

\begin{abstract}
This document is the living synthesis of the taste-fit prediction campaigns;
the markdown reports in \texttt{experiments/} are the primary record. The
active target is a binary \emph{rotation} signal --- did a track earn a repeat
listen, $\text{play\_count}\ge 2$ --- predicted from a track's \emph{intrinsic
metadata} alone (genre, artist, artist/track popularity, follower count,
release year, release type) over the user's listening corpus. The headline
metric is the area under the ROC curve (AUC); the engagement-magnitude
regression target that the lab opened on ($\text{play\_count}\times
\text{completion\_ratio}$, $\log1p$-transformed) is retained as a documented
scientific field but is no longer delivered. The behavioral co-listening
embedding and all playback/curation fields are deliberately excluded from the
feature set. Five Jun-2026 campaigns are folded in below; they converge on a
single finding --- categorical artist/genre/album identity read by
gradient-boosted trees beats every dense modality and every dense
distance/margin family --- and they crowned the project's first registered,
promotable champion, \texttt{xgboost-classifier-rotation}, at AUC
\num{0.7429} on the original DS-META corpus. Three Jun-12 campaigns then
re-ran the family question on a larger 2026 corpus (a full-64-d,
\num{23529}-row dataset): a four-family classifier bake-off in which
\texttt{xgboost-classifier} again wins (AUC \num{0.7204}), \texttt{catboost}
follows in-band, and a dense \texttt{pyramid-mlp} is last; a six-geometry
architecture scan that closes the dense-net family for AUC by \emph{refuting}
the budget hypothesis; and a topology bake-off of a new Rust/\texttt{burn}
\texttt{burn-deep-classifier} that, by feeding the identity one-hots
\emph{raw}, repairs the dense net's calibration completely (log-loss
$2.0\to 0.62$) and lifts its AUC $\sim$2 bands to \num{0.7038} --- yet still
lands $\sim$1--2 bands short of the tree tier, with learned entity embeddings
starved by the low identity cardinality. A seventh campaign then closes that
thread: a paired re-run of the topology bake-off on a \emph{higher}
identity-cardinality matrix (artist cap $60\to300$, genre $40\to120$) flips the
ranking --- \texttt{embeddings} now \emph{beats} \texttt{mlp} (a $+0.0130$ swing
that reverses the prior result, confirming the starvation hypothesis) --- and
lifts the net into the in-dataset tree tier (AUC \num{0.7192}, within $\sim$1
band of \texttt{catboost}), the first ANN to clear the rotation save bar, saved
as \texttt{burn-deep-embeddings-highvocab}. An eighth campaign then probes the inverse lever and refutes it: the classic Facebook GBDT$\to$LR trick generalized to GBDT$\to$MLP --- training a fresh head on the frozen champion tree's one-hot leaf indices --- caps \num{1.8}--\num{2.4} noise bands \emph{below} the bare xgboost champion (best head AUC \num{0.6978}) with badly broken calibration, a seventh convergent confirmation that learned representations layered on the trees do not help. A ninth campaign then probes how far the standing identity-resolution lever can go: a vocab-cap scan on the bare GBDT, holding the frozen champion config and moving \emph{only} the artist/genre one-hot cardinality, finds the lever \emph{largely spent} --- the first jump ($60/40\to300/120$) carries the gain, and past 300/120 each step pays only sub-band returns; xgboost reaches a new in-lineage single-split best AUC \num{0.7349} at 1000/300 (only $+0.005$, $\sim$0.4 band, over the anchor --- no save fired), \texttt{catboost} tracks it, and the full-vocabulary point could not be built (host memory ceiling). Section~\ref{sec:anns} consolidates
why artificial neural nets structurally underperform on this one-hot-heavy
feature matrix, how the representation-aware net resolves the diagnosis, and how
enough identity resolution finally lets the embedding topology pay off. A parallel line of work --- \emph{next-track} sequence recommendation, scored on a leak-free session split by Recall@10 --- is also consolidated (Sections~\ref{sec:nexttrack}--\ref{sec:nexttrack-phase2}): a GRU$\times$Markov rank-fusion blend beats the first-order co-listening bar by $\sim$$+61\%$ R@10, and adding a content-kNN third leg on a rich PCA-192 item space (R+M+C) is the current best-on-record row at R@10 \num{0.1859}; a later literature-fit campaign then supersedes it with a supervised \emph{learned content projection} (\Rp+M+\Cp) at R@10 \num{0.2117}, seed-robust and two-split hardened --- the first next-track ranking model registered, promoted and servable (Sections~\ref{sec:nexttrack-projection}--\ref{sec:nexttrack-music}). Its task, metrics and champion are unrelated to the rotation boards and must never be read against them. Two closing next-track architecture campaigns are then folded in, both nulls: a jointly-trained \emph{dual-tower} GRU over paired causal views of the session prefix fails to beat a \SI{0.098}{\percent} parameter-matched single tower (the fourth independent learned-combiner null, and the one that answers the ``train the fusion end-to-end'' rebuttal), and a learned per-step \emph{pre-encoder} is not merely useless but actively \emph{harmful} in a single tower ($-\num{0.0140}$ R@10, CI below zero) --- a loss traced by a counterfactual interpretability read to half-wave rectification of a zero-centred latent deleting exactly its negative half (\texttt{frac\_exact\_zero} \num{0.5327} $=$ \texttt{frac\_negative\_coords} \num{0.5327}). A measured power analysis then reframes all five architecture nulls: at $n=\num{1431}$ the paired-bootstrap CI half-width is \num{0.0124} R@10 while a whole architecture family spans only \num{0.028}, so large harms were detected reliably while moderate wins were never detectable --- and \emph{holisticness@10} exists in three mutually non-comparable code generations, which is why the next-track board stays denominated in R@10 (Sections~\ref{sec:nexttrack-dualtower}--\ref{sec:nexttrack-power}).
Six further next-track campaigns are folded in: a native \emph{MMR} $\lambda$
frontier that is the first confirmed holisticness lever on record and resolves its
mechanism question unanimously (sonic diversification de-concentrates artists in 8
of 8 arms); a purpose-built \emph{rollout} metric returning a sixth architecture
null in the deployment regime, where per-seed variance dominates model identity; a
research memo whose survivors are all \emph{evidence-conditioning} rather than
capacity changes, reported at a lower register because nothing in it was run;
\textbf{\texttt{markov\_gate}}, a nine-line score-time change making the Markov leg
abstain on the \SI{55.6}{\percent} of queries with no bigram evidence, worth
$+\num{0.0203}$ / $+\num{0.0196}$ Recall@10 on two disjoint splits (R@10
\num{0.2320}, the best on record) with the warm evidence buckets byte-for-byte
identical; the refutation-with-mechanism of the training-time anti-eager
regularizer, whose margins were calibrated from the wrong distribution and are
provably inert; the first de-eagering lever that \emph{works}, a hard per-artist
cap that cannot be outscored and trades literal for graded relevance
(\texttt{artist\_conc} \num{0.692}$\to$\num{0.113} with \texttt{music@10}
\emph{rising}, hardened on two splits); and a stage-stack depth campaign closing
the serial architecture branch on \emph{both} evaluation surfaces, whose
best-fitting arm holds the lowest validation loss in its batch and retrieves
\num{55} fewer hits.
\end{abstract}

\section{Introduction}
This document is the living synthesis of the experiment campaigns; the
markdown reports in \texttt{experiments/} are the primary record. Each section
transcribes (never recomputes) the quantitative results of a campaign,
preserves its run and dataset identifiers verbatim for traceability, and adds
the rationale and interpretation that the terse reports omit.

\paragraph{Task and target.} The lab predicts \emph{taste fit} from a track's
intrinsic metadata over the user's listening corpus. It opened on an
engagement-magnitude regression target ---
$\text{play\_count}\times\text{completion\_ratio}$, trained against a
$\log1p$-transform and mapped back with $\text{expm1}$ --- but on 2026-06-08
the active target pivoted to a binary \emph{rotation} signal:
$\text{play\_count}\ge 2$, i.e.\ did the track earn a repeat listen
(Section~\ref{sec:rotation-pivot}). Rotation is the actionable input to the
Playlist Pathfinder's per-track desirability term, it is far more learnable
than engagement magnitude (trees reach $\text{AUC}\approx 0.73$ on rotation
versus \rsq{}$\approx 0.3$ on engagement), and with a base rate near
\SI{44.5}{\percent} it is naturally spread for path search. Engagement is
retained as a documented scientific field, but rotation is the delivered
target and every leaderboard below is on it.

\paragraph{Metrics.} Because the target is binary, the headline metric is the
area under the ROC curve (AUC, higher-better), the quantity the
classification leaderboard (Table~\ref{tab:lb-auc}) sorts on and the
best-models group auto-promotes by; it is reported alongside log-loss,
accuracy and the Brier score. Two distinct metric regimes nonetheless coexist
in the rotation campaigns. The campaigns opened with a \emph{regression} read
--- mean absolute error (MAE), RMSE and \rsq{} on the $[0,1]$ P(rotation)
scale --- and only later adopted the honest \emph{classification} read once
the framework gained a binary mode. MAE is retained on the historical
regression board (Table~\ref{tab:lb-mae}) for same-family band significance,
but it is documented there as \emph{degenerate} on a binary target: an
$\epsilon$-tube SVR can post a deceptively low MAE while being worse than
predicting the mean (Section~\ref{sec:svm}), an illusion the classification
metrics dissolve. The absolute-percentage metrics (MAPE/\medape{}) are
undefined on rotation, since half the labels are exactly zero. The historical
engagement board, where MAE \emph{is} the natural headline, is preserved in
\texttt{PROJECT-FACTS.md}; its numbers live on the $\log1p$ engagement scale
and are not comparable to anything below.

\paragraph{Features.} The feature space is a track's \emph{intrinsic metadata}
alone: genre, artist, artist/track popularity, follower count, release year,
credited-artist count and release type. The dense embedding modalities
(co-listening word2vec, content-text, acoustic, and the song-autoencoder
latent of Section~\ref{sec:ae-blend}) are \emph{behavioral or descriptive}
signals, not categorical identity, and the payload-based predictors over the
metadata one-hots are the canonical models here. Critically, the recipe is
\emph{leakage-free}: the playback fields that mechanically constitute or proxy
the target --- \texttt{play\_count}, \texttt{completion\_ratio},
\texttt{skip\_rate}, \texttt{shuffle\_rate}, \texttt{distinct\_sessions}, and
the demoted \texttt{engagement} field itself --- and the behaviorally
entangled curation fields are all excluded by default.

\paragraph{Three feature matrices.} The first three campaigns ran on the
original \emph{DS-META} corpus (\num{12256} rows, metadata-only). The Jun-12
campaigns ran on a larger 2026 corpus, the \emph{full-64-d} dataset
(\dsid{ds-20260609-204419-p64-s42}, \num{23529} rows, 177 columns: 64 song-AE
PCA latent $+$ 7 numeric $+$ 106 identity one-hots at vocab caps artist 60 /
genre 40 --- the \emph{60/40-vocab} matrix). A final Jun-12 campaign
(Section~\ref{sec:anns-highvocab}) re-ran on a \emph{high-vocab} variant
(\dsid{ds-20260612-143014-p64-s42}, \emph{same} \num{23529} rows / split seed 42
/ PCA64, differing \emph{only} in identity cardinality: artist cap $60\to300$,
genre $40\to120$; 497 columns including 426 one-hots --- artist 301, genre 121,
album-type 4). \textbf{These are three non-comparable matrices}: DS-META differs
from the full-64-d corpus in both size (\num{12256} versus \num{23529} rows) and
feature set (metadata-only versus $+$64 AE latent), and the 60/40 and high-vocab
matrices differ in identity cardinality (177 versus 497 columns). AUC numbers
must therefore never be read across them --- not against the \num{0.7429}
DS-META champion, and not between the 60/40 and high-vocab boards. Each matrix
carries its own leaderboard (Tables~\ref{tab:lb-auc},~\ref{tab:lb-auc-p64}
and~\ref{tab:lb-auc-highvocab}); nothing regressed when the 60/40 xgboost lands
at \num{0.7204}, and every family scores \emph{higher} on the wider high-vocab
matrix because the extra identity granularity is genuine signal.

\section{Background: model families and metrics}
\label{sec:background}
This section gives the minimum theory for a reader who knows ML but not this
repo. It is keyed to \texttt{PROJECT-FACTS.md} and grows only when a new
predictor family or metric enters the reports.

\paragraph{The families in play.}
\begin{description}
  \item[Gradient-boosted trees (\texttt{xgboost}, \texttt{xgboost-classifier},
  \texttt{catboost-classifier}, \texttt{ngboost-classifier}).]
  An additive ensemble of shallow trees, each fit to the residual gradient of
  the previous stage. As a \emph{regressor} its output is a leaf average,
  bounded by the range of training targets; as a \emph{classifier} it emits a
  calibrated probability through a logistic link, the form that wins on
  rotation. Trees split on individual columns, so they read a track's sparse
  artist/genre/album identity one-hots \emph{losslessly} and capture non-linear
  interactions without manual feature crosses --- the reason they dominate this
  corpus. \texttt{catboost} is an ordered-boosting GBDT variant with symmetric
  (oblivious) trees and native categorical handling; on this corpus it is the
  legitimate second tree family, co-leader \emph{in-band} with xgboost.
  \texttt{ngboost} is a natural-gradient boosting variant that fits a full
  predictive \emph{distribution}; its probabilistic machinery, paid for with
  shallow base learners, costs it raw ranking power here. The cost of the GBDT
  family generally is a non-deterministic fit (subsampling and column sampling
  carry an internal RNG seed), so two runs at different seeds differ slightly.
  The registered champion \texttt{xgboost-classifier-rotation} is this family.
  \item[Logistic regression.] Linear log-odds passed through a sigmoid: a
  deterministic, $\ell_2$-regularized linear classifier. It is the natural
  \emph{linear floor} for a binary target --- the analogue of ridge in the
  regression world --- and serves as a cheap, reproducible control on the
  harness. Being a single linear boundary in a dense feature space, it blurs
  the sparse categorical identity that trees exploit, and it lands well below
  them on rotation.
  \item[Kernel SVM / SVC (\texttt{svc} with rbf / linear / poly / sigmoid
  kernels).] A maximum-margin classifier that operates on kernel-induced
  distances between points. As a \emph{regressor} (SVR) on a binary target the
  $\epsilon$-insensitive tube degenerates --- the source of the binary-MAE trap
  of Section~\ref{sec:svm}; as a proper probability \emph{classifier} (SVC) the
  divergence vanishes but every kernel still loses to the trees, because a
  dense-distance / large-margin model cannot read sparse identity one-hots the
  way an axis-aligned tree split does. The sigmoid kernel is the weakest and is
  the documented refutation echo.
  \item[Pyramid-MLP (\texttt{pyramid-mlp-classifier}).] A feed-forward
  artificial neural network with a tapering (``pyramid'') stack of hidden
  layers, ReLU activations and an Adam optimizer, fit on the StandardScaler-ed
  feature matrix. It is the project's only \emph{dense} (every-input-to-every-
  hidden-unit) learner on rotation. As Section~\ref{sec:anns} develops at
  length, a dense net is exactly the regime that one-hot-heavy taste data
  fights: it must learn dense weights over a mostly-zero input, and it lands
  last in the bake-off with a blown calibration that an architecture scan
  proves is structural, not a budget artifact.
  \item[Representation-aware net (\texttt{burn-deep-classifier}).] A
  full-Rust/\texttt{burn} feed-forward classifier built to repair the
  pyramid-MLP's pathologies. It splits the feature matrix into a
  \emph{standardized} continuous block (64 PCA $+$ 7 numeric) and a set of
  \emph{raw} (un-standardized) one-hot identity groups (artist, genre,
  album-type), so the scaler never turns rare categories into outlier spikes
  (Section~\ref{sec:anns-scaler}). A \texttt{topology} enum selects among a
  plain dense \texttt{mlp} over the concatenated inputs, learned per-group
  \texttt{embeddings}, and a \texttt{wide\_deep} split (a wide linear path over
  the one-hots beside a deep path over the continuous block). As
  Section~\ref{sec:anns-burn} reports, the raw-one-hot representation fixes the
  calibration completely and lifts AUC $\sim$2 bands over the pyramid, but the
  net still sits $\sim$1--2 bands below the tree tier, and the embedding
  variants gain nothing because the identity cardinalities here are too low to
  feed them.
  \item[AE-transplant net (\texttt{burn-ae-classifier}).] A two-phase
  Rust/\texttt{burn} network. Phase 1 pretrains an autoencoder to reconstruct
  the \emph{full} assembled feature vector --- the standardized continuous
  block beside the \emph{raw} identity one-hots, $D{=}177$ --- by mean-squared
  error; Phase 2 transplants the trained encoder into a deep softmax classifier
  (encoder $\to$ latent $\to$ head $\to$ logits) and fine-tunes it end-to-end
  with cross-entropy. A \texttt{freeze\_encoder} flag chooses a fixed feature
  extractor versus fine-tuning. It tests whether an \emph{unsupervised}
  reconstruction pretrain hands the classifier a latent that raw one-hots do
  not already give a from-scratch net; as Section~\ref{sec:anns-ae} reports it
  does not --- the reconstructive latent is task-agnostic, and the family ties
  the plain \texttt{burn-deep} \texttt{mlp} without reaching the tree tier.
  \item[Autoencoder-XGBoost blend (\texttt{autoencoder-xgboost}).] A weighted
  ensemble $\hat{y}=w\cdot\hat{y}_{\text{ae}}+(1-w)\cdot\hat{y}_{\text{xgb}}$
  that combines a regressor over a precomputed 64-dimensional
  song-autoencoder latent with a metadata-XGBoost model, the weight $w$ tuned on
  a train holdout. Its purpose was to test whether the dense AE latent carries
  rotation signal the metadata one-hots miss; it does not, and the tuner
  correctly drives $w\to 0$, collapsing the ensemble to metadata-only XGBoost
  (Section~\ref{sec:ae-blend}).
\end{description}

\paragraph{Why a binary target makes MAE degenerate.} For a continuous target,
mean absolute error is a faithful summary of how far predictions sit from the
truth. On a $\{0,1\}$ target it is not: a model can minimize MAE by hugging the
labels regardless of how it \emph{ranks} or \emph{calibrates} them. An SVR with
an $\epsilon{=}0.1$ tube snaps its outputs to a two-level step near
$\{0.1,0.9\}$, sitting only $\sim 0.1$ off every label --- a low MAE --- while
its ranking is worthless (negative \rsq{}). The metrics that matter for a
binary taste-fit signal are therefore the classification ones:
\begin{description}
  \item[AUC (area under the ROC curve).] The probability that a random positive
  is ranked above a random negative; \num{0.5} is chance, \num{1.0} is perfect.
  It is threshold-free and scale-free, depends only on the \emph{ordering} of
  the predicted probabilities, and is exactly the property the Pathfinder needs
  from its desirability term. It is the headline metric and the leaderboard
  sort key.
  \item[Log-loss.] The negative log-likelihood of the labels under the
  predicted probabilities; it punishes \emph{confident} mistakes severely and so
  rewards calibration, not just ranking.
  \item[Brier score.] The mean squared error between predicted probability and
  the $\{0,1\}$ label --- a calibration-and-refinement measure on the natural
  $[0,1]$ scale, lower-better.
\end{description}
A model that wins on AUC \emph{and} log-loss \emph{and} Brier, as the champion
does, is winning on ranking and calibration at once; the binary-MAE artifact
that flatters a degenerate SVR shows up on none of them. The same triad later
exposes the opposite pathology in a dense net (Section~\ref{sec:anns}): a model
can post a respectable AUC --- genuine ranking --- while its log-loss and Brier
are catastrophic, i.e.\ it ranks acceptably but is \emph{confidently wrong}.

\paragraph{A note on the historical log-target.} The lab's opening regression
target, engagement, is heavy-tailed and non-negative, so it was trained against
a $\log1p$-transform and mapped back with $\text{expm1}$. That transform pair
carries a hazard --- a model unbounded in log-space can produce a large
positive log prediction that $\text{expm1}$ amplifies into an enormous value,
the ``$\text{expm1}$ blow-up'' --- which is why bounded-by-construction families
(trees, rbf kernels) are preferred over unbounded ones (linear models, MLPs) on
that target. Engagement remains a documented field and its leaderboard lives in
\texttt{PROJECT-FACTS.md}, but it is no longer the delivered target, so the
blow-up hazard is now a historical footnote rather than a live design
constraint.

\section{The rotation pivot: from engagement magnitude to a binary taste-fit signal}
\label{sec:rotation-pivot}
On 2026-06-08 the active target pivoted away from engagement magnitude to a
binary \emph{rotation} signal, defined as $\text{play\_count} \ge 2$ --- did
the track earn a repeat listen? Rotation is the actionable input to the
Playlist Pathfinder's \texttt{W\_FIT} term, where a per-track $[0,1]$
desirability biases the A* path; a calibrated rotation probability is a
cleaner \texttt{fit\_norm} than a $\log1p$-then-$\text{expm1}$ engagement
estimate. Rotation also turns out to be far more learnable than engagement
magnitude (gradient-boosted trees reach $\text{AUC}\approx 0.73$ on rotation
versus \rsq{}$\approx 0.3$ on engagement) and is naturally spread for path
search, with a base rate near \SI{44.5}{\percent}. Engagement is retained as a
documented scientific field but is no longer the delivered target.

The pivot is decisive for how every number below must be read. Rotation lives
on the $[0,1]$ scale, so its MAE/\rsq{} are \emph{not} comparable to the
engagement leaderboard preserved in \texttt{PROJECT-FACTS.md}; and because half
the rotation labels are exactly zero, the absolute-percentage metrics
(MAPE/\medape{}) are undefined for those rows. Two distinct metric regimes
therefore coexist in the rotation campaigns: an early \emph{regression} read
(MAE, RMSE, \rsq{}, units P(rotation)) and the later, honest
\emph{classification} read (AUC, log-loss, accuracy, Brier). As the campaigns
below show, the choice of metric is not cosmetic --- under MAE a kernel SVR can
\emph{look} competitive while being worse than predicting the mean, an illusion
that the classification metrics dissolve.

\section{Rotation leaderboards}
\label{sec:leaderboard}
Four leaderboards. The first two are on the original DS-META corpus, one per
metric regime; the third and fourth are on the larger 2026 full-64-d corpus.
\textbf{None of the three AUC boards is comparable to the others} --- they sit
on three non-comparable feature matrices (DS-META metadata-only; the full-64-d
\emph{60/40-vocab} matrix; and the full-64-d \emph{high-vocab} matrix, which
shares the 60/40 corpus, rows and split but raises the identity cardinality
caps). Each board carries its own in-dataset references and must never be read
across (see the Introduction's ``Three feature matrices'' note). The DS-META
AUC board carries the project's first registered, promotable rotation champion;
the 60/40 board carries the new-corpus registered xgboost definition; the
high-vocab board carries the first ANN to clear the rotation save bar
(Section~\ref{sec:anns-highvocab}). Finally, \textbf{every AUC on all four
boards is measured on a \emph{random} train/test split}; a chronological
(train-on-earlier, test-on-later) split tells a sharply different story --- the
numbers below are upper bounds on same-era ranking, and honest
forward-prediction skill on this corpus is $\sim$0.11 AUC lower
(Section~\ref{sec:drift}).

\begin{table}[ht]
  \centering
  \caption{Rotation, \emph{classification} regime, \textbf{DS-META corpus} ---
  ranked by AUC (higher-better). The champion is the only registered,
  auto-promoted definition. Datasets: DS-META
  \dsid{ds-20260609-141622-p1-s42} (115 cols, metadata-only, pca floor); the
  178-col row adds the 64 song-AE latent columns
  (\dsid{ds-20260608-204402-p64-s42}). AUC noise band $2\sigma=0.0126$.}
  \label{tab:lb-auc}
  \begin{tabular}{llcccc}
    \toprule
    Model & Dataset & AUC & log-loss & accuracy & Registered? \\
    \midrule
    \best{xgboost-classifier} (champion) & DS-META & \best{0.7429} & \best{0.5883} & 0.6846 & \best{yes (rank 1)} \\
    xgboost-classifier (178, +AE latent) & p64-s42 & 0.7350 & 0.6045 & 0.6818 & no \\
    svc (poly $d{=}3$, best non-tree)    & DS-META & 0.6855 & 0.6452 & 0.6450 & no \\
    logistic ($C{=}1$, linear floor)     & DS-META & 0.6762 & 0.6329 & 0.6491 & no \\
    svc (rbf)                            & DS-META & 0.6722 & 0.6469 & 0.6426 & no \\
    svc (linear)                         & DS-META & 0.6363 & 0.6535 & 0.6373 & no \\
    svc (sigmoid, worst)                 & DS-META & 0.6255 & 0.6673 & 0.6267 & no \\
    \bottomrule
  \end{tabular}
\end{table}

\begin{table}[ht]
  \centering
  \caption{Rotation, \emph{classification} regime, \textbf{new full-64-d
  corpus} --- ranked by AUC (higher-better). Dataset
  \dsid{ds-20260609-204419-p64-s42} (\num{23529} rows, 177 cols: 64 song-AE
  PCA $+$ 7 numeric $+$ 106 identity one-hots). \textbf{Not comparable to
  Table~\ref{tab:lb-auc}}: different corpus size and feature set; the only
  valid in-corpus reference is the \num{0.7204} xgboost point, which the
  bake-off reproduced byte-identically. The winner's default config is saved
  as the new-corpus definition. The single-seed sweep best (depth 8, lr 0.03,
  \num{0.7261}) is $+0.0057$ over the default --- inside the band --- so it is
  flagged for a 3-seed confirm, not promoted. The pyramid-mlp log-loss is its
  documented calibration ceiling (Section~\ref{sec:anns}); the
  \texttt{burn-deep-classifier} rows are the Rust/\texttt{burn}
  representation-aware net that repairs it (log-loss $\to\sim$0.62) yet still
  sits below the tree tier, so neither is registered
  (Section~\ref{sec:anns-burn}). AUC band $2\sigma\approx 0.0126$ (old-corpus
  estimate, approximate here).}
  \label{tab:lb-auc-p64}
  \begin{tabular}{llcccl}
    \toprule
    Model & max\_depth/lr & AUC & log-loss & accuracy & Note \\
    \midrule
    \best{xgboost-classifier} (winner, saved)      & 6 / 0.05  & \best{0.7204} & 0.6075 & 0.6585 & registered (new corpus) \\
    xgboost-classifier (sweep best, single-seed)   & 8 / 0.03  & 0.7261 & \best{0.6044} & \best{0.6651} & flagged for 3-seed confirm \\
    catboost-classifier (2nd tree, in-band)        & 6 / 0.05  & 0.7111 & 0.6132 & 0.6560 & $-0.0093$, inside band \\
    burn-deep-classifier (mlp, best tune)          & lr 0.0005 / [128,64] & 0.7038 & 0.6221 & 0.6521 & best ANN; calib.\ fixed \\
    burn-deep-classifier (mlp, Phase-1 winner)     & lr 0.001 / [128,64]  & 0.7001 & 0.6211 & 0.6470 & dense net; calib.\ fixed \\
    burn-ae-classifier (AE-transplant, best)       & enc[128,64]/clf[128,64] & 0.6972 & 0.6216 & 0.6496 & $-0.0066$ vs burn mlp \\
    ngboost-classifier                             & ---       & 0.6917 & 0.6293 & 0.6392 & $\sim$2.3 bands back \\
    pyramid-mlp-classifier (arch-scan best)        & [512,256] & 0.6792 & 2.0584 & 0.6326 & dense net; calib.\ ceiling \\
    pyramid-mlp-classifier (bake-off, 300 iter)    & [256,128,64] & 0.6781 & 2.4085 & 0.6326 & dense net; calib.\ ceiling \\
    \bottomrule
  \end{tabular}
\end{table}

\begin{table}[ht]
  \centering
  \caption{Rotation, \emph{classification} regime, \textbf{high-vocab full-64-d
  matrix} --- ranked by AUC (higher-better). Dataset
  \dsid{ds-20260612-143014-p64-s42} (\num{23529} rows, 497 cols: 64 song-AE
  PCA $+$ 9 numeric $+$ 426 identity one-hots --- artist 301, genre 121,
  album-type 4). \textbf{Not comparable to Tables~\ref{tab:lb-auc}
  or~\ref{tab:lb-auc-p64}}: it shares the 60/40 corpus, rows and split but
  raises the identity cardinality (artist cap $60\to300$, genre $40\to120$), so
  every family re-ran here for a fresh in-dataset ranking. The trees moved
  \emph{up} on the wider matrix (xgboost $0.7204\to\num{0.7299}$, catboost
  $0.7111\to\num{0.7252}$). The saved \texttt{burn-deep-classifier}
  (\texttt{embeddings}) is the first ANN to reach the in-dataset tree tier and
  clear the rotation save bar (Section~\ref{sec:anns-highvocab}); xgboost is the
  in-dataset leader but no tree save rule fired this campaign. AUC band
  $2\sigma\approx 0.0126$ (old-corpus estimate, approximate here).}
  \label{tab:lb-auc-highvocab}
  \begin{tabular}{llcccl}
    \toprule
    Model & config & AUC & log-loss & accuracy & Note \\
    \midrule
    \best{xgboost-classifier} (in-dataset leader)     & d6 / lr0.05      & \best{0.7299} & \best{0.6022} & \best{0.6719} & not registered \\
    catboost-classifier (2nd tree)                    & defaults         & 0.7252 & 0.6068 & 0.6649 & $-0.0047$, in band \\
    burn-deep-classifier (embeddings, Phase-2 best)   & ed8 / lr0.0005   & 0.7228 & 0.6105 & 0.6645 & tree tier; sub-band gain \\
    burn-deep-classifier (embeddings, \emph{saved})   & ed16 / lr0.001   & 0.7192 & 0.6087 & 0.6694 & registered (high-vocab) \\
    burn-deep-classifier (mlp)                        & lr0.001          & 0.7129 & 0.6166 & 0.6617 & $-0.0063$ vs embeddings \\
    burn-deep-classifier (wide\_deep)                 & lr0.001          & 0.7061 & 0.6256 & 0.6526 & laggard, both regimes \\
    \bottomrule
  \end{tabular}
\end{table}

\begin{table}[ht]
  \centering
  \caption{Rotation, \emph{regression} regime, \textbf{DS-META corpus} ---
  ranked by MAE on the $[0,1]$ P(rotation) scale (lower-better). The
  metadata-only XGBoost leader is the \texttt{mae\_xgb} branch reproduced
  inside every blend run; it was never saved as a standalone definition. The
  SVM-poly row is listed for completeness and \emph{refuted}: its lower MAE is
  the $\epsilon$-tube artifact of Section~\ref{sec:svm}, its \rsq{} is
  negative. MAE noise band $2\sigma=0.0086$.}
  \label{tab:lb-mae}
  \begin{tabular}{lcccl}
    \toprule
    Model & MAE & RMSE & \rsq{} & Source \\
    \midrule
    \best{metadata-only XGBoost} (AE latent dropped) & \best{0.3879} & \best{0.4475} & \best{0.191} & ae-blend / svm \\
    xgboost on latent+metadata combined              & 0.3965 & --- & 0.166 & rotation pivot \\
    xgboost on AE-latent only (64 cols)              & 0.4098 & --- & 0.132 & rotation pivot \\
    svm (poly, $\epsilon$-tube artifact --- REFUTED) & 0.3814 & 0.5303 & $-0.137$ & svm \\
    \bottomrule
  \end{tabular}
\end{table}

\section{Campaigns}

\subsection{Autoencoder-latent + XGBoost blend: a dead modality, isolated}
\label{sec:ae-blend}
\emph{Rationale.} This was the first campaign on the pivoted rotation target.
The hypothesis was that a precomputed 64-dimensional song-autoencoder latent
(from the Qdrant collection \texttt{spotify\_tracks\_song\_ae}) carries a
track's \emph{identity} in a way the metadata one-hots do not, and that
blending an AE-latent regressor with a metadata-XGBoost model ---
$\hat{y}=w\cdot\hat{y}_{\text{ae}}+(1-w)\cdot\hat{y}_{\text{xgb}}$, with the
weight $w$ tuned on a train holdout --- could beat metadata-XGBoost alone. All
six arms ran on \dsid{ds-20260608-204402-p64-s42} (64 AE latent + 114 metadata
$=$ 178 cols, 12{,}256 rows, 9{,}805/2{,}451 at seed 42), varying only the AE
branch (kNN at $k\in\{10,25,50,100,200\}$, plus a ridge-on-latent arm).

\emph{Results.} All six runs (\runid{run-20260609-131545-e34e4} and siblings)
succeeded, and all returned byte-identical blended metrics --- MAE 0.38789,
\rsq{} 0.191, RMSE 0.4475 --- because the weight tuner drove $w\to 0$ in every
arm. The only axis that moved was the AE branch's standalone error
$mae_{\text{ae}}$, which grew monotonically with the neighborhood size, from
0.4259 at $k=10$ to 0.4702 at $k=200$ (the ridge arm landed at 0.4623, between
$k=100$ and $k=200$). Figure~\ref{fig:ae-knn} shows the curve against the three
horizontal references that matter: even the best AE read ($k=10$, 0.4259) never
reaches the metadata-only tree (0.3879), nor even the AE-latent-only
\emph{xgboost} on record (0.4098).

\begin{figure}[ht]
  \centering
  \includegraphics[width=0.8\textwidth]{ae_latent_knn}
  \caption{The song-AE latent carries no rotation signal the metadata
  one-hots do not already carry better. The kNN standalone error
  $mae_{\text{ae}}$ rises monotonically as the neighborhood widens (the
  latent's signal is strongest in tight local neighborhoods and washes toward
  the base rate), and even its best point sits above both the AE-latent-only
  tree and the metadata-only tree. The weight tuner correctly drives the blend
  weight to zero, so the ensemble collapses to metadata-only XGBoost.}
  \label{fig:ae-knn}
\end{figure}

\emph{Interpretation.} The collapse to $w=0$ is the correct answer, not a bug
to fix: given an AE branch strictly worse than the XGB branch on the holdout,
the optimal convex weight \emph{is} zero. The architecture's one real value was
that it \emph{isolated} the latent into a branch the tuner could discard,
recovering a cleaner result than feeding the latent into a unified XGBoost ---
which actually \emph{dilutes} the tree (MAE 0.3965 with the 64 latent columns
present versus 0.3879 with them dropped). The song-AE latent thus joins
co-listening, content-text and acoustic features as a dense modality that loses
to categorical identity one-hots. The one caveat at the time was honesty about
significance: the $0.3879$ versus $0.3965$ edge is only $0.0086$ MAE on a
$[0,1]$ target, and the rotation noise band was then unmeasured --- so the
campaign reported the number but asserted no significant win and promoted
nothing.
\emph{Source: \texttt{2026-06-09-ae-xgboost-blend.md}.}

\subsection{Kernel SVM on rotation: a refutation, and the binary-MAE trap}
\label{sec:svm}
\emph{Rationale.} With the latent ruled out, this campaign tested whether any
SVM kernel (rbf / linear / poly / sigmoid) produces a rotation predictor
competitive with the metadata-only XGBoost leader, and --- only if one passed
--- whether it added orthogonal signal in a grid-weighted blend. The metadata
matrix was reduced to the pca floor in a new dataset
(\dsid{ds-20260609-141622-p1-s42}, 115 cols), with two split-seed siblings
(seeds 17 and 101) built to measure the noise band, and the champion reproduced
as a control (\runid{run-20260609-141717-9a8ba-autoencoder-xgboost}, blend
weight 0). A secondary outcome was the campaign's central deliverable: the
\textbf{first measured rotation noise band}.

\emph{Results.} Across the three split seeds of the metadata-only XGBoost
control (s42 0.387891, s17 0.393072, s101 0.396417) the sample standard
deviation is $\sigma_{\text{split}}=0.00430$, giving a $2\sigma$ significance
band of \textbf{0.0086 MAE} on the $[0,1]$ scale. This retroactively settles
the previous campaign's open question: the $0.3879$ versus $0.3965$ edge is
$\Delta 0.0086 \approx 1{\cdot}\text{band}$ (likely better, not confident),
while the $0.3879$ versus $0.4098$ gap (AE-latent-only) is $\sim
2.5{\cdot}\text{band}$ and therefore real. As for the SVM arms, rbf, linear and
poly all posted MAE inside the control band (0.381--0.390), \emph{lower} than
the champion in the poly case (0.3814) --- yet every one of them is
worse-than-mean on \rsq{} ($-0.13$ to $-0.20$), with RMSE 0.53--0.54 against
the tree's 0.4475. Sigmoid SVR diverged outright: MAE 0.624, RMSE 1.531, \rsq{}
$-8.48$, with predictions reaching 24.7 on a $[0,1]$ target --- the classic
unbounded sigmoid-SVR blow-up. Figure~\ref{fig:svm} contrasts the two metric
views directly.

\begin{figure}[ht]
  \centering
  \includegraphics[width=0.95\textwidth]{svm_mae_vs_r2}
  \caption{On a binary target, low SVR MAE is an artifact, not skill. With
  $\epsilon=0.1$ on a $\{0,1\}$ target the SVR snaps its outputs to the
  $\epsilon$-tube edges (a hard two-level step, $p_{50}\approx 0.10$,
  $p_{90}\approx 0.90$), so each row sits only $\sim 0.1$ off its label and the
  \emph{MAE} (left) reads competitive. The honest \emph{$R^2$} (right) exposes
  it: every kernel is worse than predicting the mean. Sigmoid is off-scale at
  $R^2=-8.48$.}
  \label{fig:svm}
\end{figure}

\emph{Interpretation.} The low SVR MAE is a degeneracy of the metric on a
binary target. With an $\epsilon$-tube of width $0.1$ on $\{0,1\}$ labels, the
SVR collapses to a two-level step function pinned near the right classes ---
each row is about $0.1$ off its label, so the MAE is small, while calibration
and ranking are destroyed (negative \rsq{}). The metric that matters for the
Pathfinder \texttt{fit\_norm} (ranking and calibration, i.e.\ \rsq{}/RMSE/AUC)
says the SVM loses decisively. Worse, the two stage-2 blends
(\runid{run-20260609-143114-528ef-blend} and its lgbm twin) fitted weights of
exactly $[0.0, 1.0]$ --- the grid, optimizing the same degenerate validation
MAE, \emph{mis-selected pure SVM} and discarded the tree, the opposite of the
self-disabling seen in the blend campaign. The lesson is sharp: a grid-by-MAE
blend objective is itself unsafe on a binary target. Nothing was registered or
promoted; kernel SVM became the fifth dense-regime family to lose to one-hot
identity on this corpus.
\emph{Source: \texttt{2026-06-09-svm-kernel-rotation.md}.}

\subsection{The honest AUC bake-off and the first AUC champion}
\label{sec:auc}
\emph{Rationale.} The SVM campaign exposed that MAE is the wrong lens for a
binary target. The framework had meanwhile gained a classification mode
(\texttt{domain.toml} \texttt{task = binary}, primary metric AUC,
higher-better; the best-models group ranks by AUC and auto-promotes the top
run). This campaign re-ran the comparison on that honest footing: xgboost-classifier
against logistic and a four-kernel SVC scan, all on DS-META, plus a
178-column arm with the AE latent re-added to settle the latent-dilution
question on the AUC axis, and three control seeds to measure the AUC noise
band. The hypothesis --- a sixth test of the convergent finding --- was that
xgboost-classifier wins and every dense family loses.

\emph{Results.} All nine runs succeeded. The champion config
(\runid{run-20260609-172920-627d6-xgboost-classifier}, control hyperparameters,
seed 42) reached \textbf{AUC 0.7429} (log-loss 0.5883, accuracy 0.6846, Brier
0.2017) and was registered as \texttt{xgboost-classifier-rotation} and
auto-promoted to best-models rank 1 --- the project's first saved, AUC-scored,
promotable rotation champion. The first measured \emph{AUC} noise band comes
from the three control seeds (s42 0.742934, s17 0.730659, s101 0.734469, mean
0.736020): $\sigma_{\text{AUC}}=0.00628$, $2\sigma=\textbf{0.0126}$
(Figure~\ref{fig:auc-band}). This band is a different metric from the 0.0086
MAE band and does not replace it. Every dense family lost by many bands ---
svc-poly 0.6855, logistic 0.6762, svc-rbf 0.6722, svc-linear 0.6363,
svc-sigmoid 0.6255 --- a spread of 4.6 to 9.3 noise bands below the champion
(Figure~\ref{fig:auc-bakeoff}). The 178-column arm (AE latent re-added) scored
0.7350, $-0.0079$ below the metadata-only champion --- the same dilutive
direction as the regression result, but inside the AUC band, so
neutral-leaning-dilutive at a single seed.

\begin{figure}[ht]
  \centering
  \includegraphics[width=0.8\textwidth]{auc_bakeoff}
  \caption{The honest classification metric confirms the convergent finding a
  sixth time. xgboost-classifier reads the sparse artist/genre/album identity
  one-hots losslessly; the dense distance/margin families (logistic and all
  four SVC kernels) blur that identity in a dense feature space and land 4.6 to
  9.3 noise bands lower. The shaded strip is the $2\sigma=0.0126$ band below
  the champion; no challenger comes within a band of it.}
  \label{fig:auc-bakeoff}
\end{figure}

\begin{figure}[ht]
  \centering
  \includegraphics[width=0.7\textwidth]{auc_noise_band}
  \caption{The first measured rotation AUC noise band. Three data-split seeds
  of the registered champion config give $\sigma_{\text{AUC}}=0.0063$ and a
  $2\sigma$ significance band of $0.0126$. In relative terms this is actually
  tighter than the MAE band (0.85\% versus 1.10\% coefficient of variation),
  as expected for a rank metric on 2{,}451 test rows. The band only matters for
  the sub-band latent-dilution call; champion and challengers are separated by
  many bands.}
  \label{fig:auc-band}
\end{figure}

\emph{Interpretation.} Reframing rotation as classification did two things at
once. It dissolved the SVR illusion of Section~\ref{sec:svm} --- under AUC,
log-loss \emph{and} Brier, every SVC kernel is now visibly worse than the
trees, all in the same direction, so the $\epsilon$-tube artifact that made
SVR's MAE look competitive simply vanishes. And it confirmed that the metric
was only ever a problem for the \emph{fairness of the comparison}; the dense
family was always the loser. Notably svc-sigmoid does not collapse to chance
under classification (it holds AUC 0.6255 rather than blowing up): a proper
probability objective removes the unbounded divergence that wrecked the
SVR-sigmoid run. The convergent finding now stands on an honest metric: the
taste-fit signal lives in artist/genre/album \emph{identity}, captured by
sparse one-hots and exploited by trees, not in dense distance or margin models.
The remaining lever is identity \emph{resolution} --- raising the vocabulary
caps so the \texttt{artist=\_\_other\_\_} tail gets real columns --- not new
modalities or families. The next planned campaign is a champion hyperparameter
scan now that the family and metric are settled.
\emph{Source: \texttt{2026-06-09-rotation-classification.md}.}

\subsection{Four-family bake-off on the new full-64-d corpus, and an xgboost depth/lr sweep}
\label{sec:bakeoff-p64}
\emph{Rationale.} The first three campaigns ran on the original DS-META corpus.
This campaign re-asked the family question on the larger 2026 corpus --- ``all
of the 64-d embedding'', the full 177-column matrix
(\dsid{ds-20260609-204419-p64-s42}, \num{23529} rows, 64 song-AE PCA $+$ 7
numeric $+$ 106 identity one-hots) --- by putting four classifier families head
to head at registry defaults, seed 42: \texttt{xgboost-classifier},
\texttt{catboost-classifier}, \texttt{ngboost-classifier} and a dense
\texttt{pyramid-mlp-classifier}. The latter three are \emph{new families on
this project}. The winner then earned a \texttt{max\_depth}$\times$
\texttt{learning\_rate} sweep. The only valid in-corpus reference is a
pre-existing run on this dataset (\runid{run-20260609-204726-439d8-xgboost-classifier},
AUC \num{0.7204}); nothing here may be read against the DS-META \num{0.7429}
champion (different corpus and feature set).

\emph{Results.} All four bake-off arms succeeded (Table~\ref{tab:bakeoff-p64};
Figure~\ref{fig:bakeoff-p64}). \texttt{xgboost-classifier}
(\runid{run-20260612-115909-840c0-xgboost-classifier}) won at AUC \num{0.7204},
reproducing the in-corpus reference byte-identically (\num{0.7204383403}) ---
the determinism gate passed. \texttt{catboost-classifier}
(\runid{run-20260612-115909-c7ef6-catboost-classifier}) followed at \num{0.7111},
$-0.0093$ behind --- \emph{inside} the $\sim$0.0126 band, so co-leader and the
legitimate second tree family, not a loss. \texttt{ngboost-classifier}
(\runid{run-20260612-115909-d725b-ngboost-classifier}) landed at \num{0.6917}
($\sim$2.3 bands back), and the dense \texttt{pyramid-mlp-classifier}
(\runid{run-20260612-115909-ea698-pyramid-mlp-classifier}) was last at
\num{0.6781} --- with a log-loss of \num{2.4085} and Brier \num{0.3251} far
above the trees' $\sim$0.60--0.63 / $\sim$0.21, the diagnostic dissected in
Section~\ref{sec:anns}.

\begin{table}[ht]
  \centering
  \caption{Phase 1 --- the four-family bake-off on the new full-64-d corpus
  (\dsid{ds-20260609-204419-p64-s42}, seed 42, registry defaults). All four
  succeeded. AUC higher-better; log-loss and Brier lower-better. Winner bolded;
  best cell per metric bolded.}
  \label{tab:bakeoff-p64}
  \begin{tabular}{llcccc}
    \toprule
    Arm / run id & Predictor & AUC & log-loss & accuracy & Brier \\
    \midrule
    \best{C0} (\runid{run-20260612-115909-840c0-xgboost-classifier}) & \best{xgboost-classifier}   & \best{0.7204} & \best{0.6075} & \best{0.6585} & \best{0.2109} \\
    A1 (\runid{run-20260612-115909-c7ef6-catboost-classifier})       & catboost-classifier         & 0.7111 & 0.6132 & 0.6560 & 0.2132 \\
    A2 (\runid{run-20260612-115909-d725b-ngboost-classifier})        & ngboost-classifier          & 0.6917 & 0.6293 & 0.6392 & 0.2193 \\
    A3 (\runid{run-20260612-115909-ea698-pyramid-mlp-classifier})    & pyramid-mlp-classifier      & 0.6781 & 2.4085 & 0.6326 & 0.3251 \\
    \bottomrule
  \end{tabular}
\end{table}

\begin{figure}[ht]
  \centering
  \includegraphics[width=0.8\textwidth]{bakeoff_p64}
  \caption{The family ordering survives the corpus change. On the larger 2026
  full-64-d matrix the two gradient-boosted-tree families lead
  (\texttt{xgboost} winning, \texttt{catboost} co-leader inside the shaded
  $2\sigma$ band below the winner); \texttt{ngboost} sits $\sim$2.3 bands back;
  the dense \texttt{pyramid-mlp} is last, the same convergent finding seen on
  DS-META in Figure~\ref{fig:auc-bakeoff}.}
  \label{fig:bakeoff-p64}
\end{figure}

The winner earned a Phase-2 sweep over \texttt{max\_depth}$\in\{4,6,8\}\times$
\texttt{learning\_rate}$\in\{0.03,0.05,0.1\}$ (Table-form data in the source
report; Figure~\ref{fig:xgb-heatmap}). Depth is the dominant axis: every
depth-4 cell is at the bottom (0.7074--0.7119), and the depth-8 / lr-0.03 corner
tops the grid at \num{0.7261} (\runid{run-20260612-120207-2dd3d-xgboost-classifier},
log-loss \num{0.6044}). High learning rate (0.1) consistently hurts both AUC and
calibration. But the best-vs-default gap (+0.0057) is \emph{inside} the band, so
the default depth-6/lr-0.05 config was saved as the new-corpus definition and the
depth-8/lr-0.03 cell was flagged for a 3-seed confirmation rather than promoted.

\begin{figure}[ht]
  \centering
  \includegraphics[width=0.62\textwidth]{xgb_depth_lr_heatmap}
  \caption{The xgboost \texttt{max\_depth}$\times$\texttt{learning\_rate} sweep
  on the new corpus. Depth is the lever (the AUC climbs monotonically up the
  depth axis at the slow learning rates), and a fast learning rate hurts. The
  best cell (depth 8, lr 0.03, boxed) reaches \num{0.7261} but only $+0.0057$
  over the default --- inside the $2\sigma$ band --- so the default config
  remains saved and the best cell awaits a 3-seed confirm.}
  \label{fig:xgb-heatmap}
\end{figure}

\emph{Interpretation.} The convergent finding survives the corpus change intact:
on a different, larger corpus with the AE latent present, the two tree families
again lead and the dense net is again last. \texttt{catboost} earns its place as
a legitimate second tree family (co-leader in-band) --- a tree-diversity /
second-opinion member, not an AUC gain over xgboost. \texttt{ngboost} pays for
its distributional machinery: shallow base learners (\texttt{base\_max\_depth} 3)
under-fit the 177-column one-hot-heavy matrix. The pyramid-MLP's last-place AUC
\emph{plus} its blown log-loss are the dense-net signature; the bake-off read the
log-loss as a 300-iteration under-training artifact and prescribed more budget ---
a hypothesis the next campaign tests and refutes.
\emph{Source: \texttt{2026-06-12-classifier-bakeoff-p64.md}.}

\subsection{Pyramid-MLP architecture scan: refuting the budget hypothesis}
\label{sec:mlp-scan}
\emph{Rationale.} The bake-off flagged the pyramid-MLP's last place (AUC
\num{0.6781}, log-loss \num{2.4085}) as a 300-iteration-cap under-training
artifact, prescribing ``raise \texttt{max\_iter} so the net converges.'' This
follow-up tests that mechanism directly: six pyramid geometries, all with
\texttt{max\_iter} raised to \textbf{800}, alpha \num{1e-4}, learning rate
\num{1e-3}, seed 42, on the same dataset --- so the comparison reflects
\emph{architecture}, not budget. Does any geometry, given the budget to
converge, lift the dense family to a competitive AUC?

\emph{Results.} Refuted on every count (Table~\ref{tab:mlp-scan};
Figures~\ref{fig:mlp-arch} and~\ref{fig:mlp-calib}). \textbf{No arm hit the 800
cap}: every geometry self-stopped at \textbf{73--100 iterations} via sklearn's
\texttt{n\_iter\_no\_change} loss-plateau tolerance --- far below even the old
300 cap, so the budget was never the binding constraint. The decisive tell: the
re-baseline arm M0 ([256,128,64], the bake-off geometry now at \texttt{max\_iter}
800) reproduced the under-trained run's AUC to all digits ---
\num{0.6780926} versus \num{0.6780926} --- because it ran the identical 73
iterations. Raising the cap from 300 to 800 changed nothing. The best converged
geometry was M1 [512,256] at AUC \num{0.6792} (\runid{run-20260612-123155-0ace5}),
a statistical tie with the under-trained default ($+0.0011 \approx$ zero) and
$\sim$3.3 bands below xgboost. The whole family lands in a 0.0080-wide AUC strip
(0.6712--0.6792), entirely inside one noise band --- the geometry ranking is not
resolvable at a single seed --- with the log-loss (2.06--2.78) and Brier
(0.32--0.34) blown at the same level as the under-trained default across every
arm. No definition was saved (below the 0.683 save bar).

\begin{table}[ht]
  \centering
  \caption{The six-geometry pyramid-MLP architecture scan
  (\dsid{ds-20260609-204419-p64-s42}, \texttt{max\_iter} 800, alpha \num{1e-4},
  lr \num{1e-3}, seed 42), sorted by AUC. All six succeeded; \emph{none} hit
  the 800-iteration cap --- each self-stopped at the listed iteration count.
  Best cell per metric bolded.}
  \label{tab:mlp-scan}
  \begin{tabular}{llcccc}
    \toprule
    Arm & Hidden layers & AUC & log-loss & Brier & Iters (of 800) \\
    \midrule
    \best{M1} wide-shallow & [512, 256]            & \best{0.6792} & \best{2.0584} & \best{0.3216} & 75 \\
    M0 re-baseline         & [256, 128, 64]        & 0.6781 & 2.4085 & 0.3251 & 73 \\
    M4 gentle taper        & [384, 307, 246]       & 0.6777 & 2.7764 & 0.3292 & 87 \\
    M2 deep pyramid        & [256, 128, 64, 32]    & 0.6737 & 2.7631 & 0.3303 & 100 \\
    M3 steep taper         & [512, 169, 56]        & 0.6736 & 2.7031 & 0.3365 & 89 \\
    M5 wide-deep           & [512, 307, 184, 111]  & 0.6712 & 2.6457 & 0.3348 & 84 \\
    \bottomrule
  \end{tabular}
\end{table}

\begin{figure}[ht]
  \centering
  \includegraphics[width=0.95\textwidth]{mlp_arch_scan}
  \caption{Both knobs of the dense family are exhausted. \emph{Left:} the six
  geometries land inside a single $2\sigma$ band around the under-trained
  default --- architecture is not the binding constraint. \emph{Right:} every
  geometry self-stops at 73--100 iterations, well below even the old 300 cap
  (let alone the raised 800), so the training \emph{budget} was never the
  binding constraint either. The optimizer reaches a loss plateau it cannot
  escape at this learning rate and halts. The binding constraint is the input
  \emph{representation}.}
  \label{fig:mlp-arch}
\end{figure}

\emph{Interpretation.} The refutation is clean and twofold: neither the training
budget (every arm self-stops far below the cap; M0 reproduces the under-trained
AUC to all digits) nor the architecture (all six geometries inside one band)
moves the dense family. The high log-loss is therefore not a truncated training
trajectory but a genuine \emph{calibration ceiling} at these hyperparameters ---
the optimizer's chosen plateau. With the dense family's own knobs exhausted, the
binding constraint must lie elsewhere: in the input \emph{representation}, the
subject of Section~\ref{sec:anns}. The directional within-family trend ---
wider-and-shallower mildly favored (M1 top), depth/capacity mildly harmful (M5
wide-deep last) --- is consistent with a dense net struggling to route signal
through 106 sparse one-hots: extra layers add capacity to memorize and blur,
not to separate. The family is closed for AUC on this feature set; do not
re-scan geometry or budget.
\emph{Source: \texttt{2026-06-12-pyramid-mlp-arch-scan.md}.}

\section{Why artificial neural nets underperform on one-hot-heavy taste data}
\label{sec:anns}
The dense \texttt{pyramid-mlp} is the only artificial neural net in the
rotation campaigns, and across two Jun-12 campaigns it lands last with a
distinctive failure pattern. This section consolidates, rigorously, \emph{why}
--- grounding every claim in the campaign evidence. The short version is the
project's documented \emph{convergent finding}: on this feature matrix the
binding constraint is the input \emph{representation}, and a dense net is the
one model class structurally mismatched to it. We build the argument in five
parts and close with the motivated next step.

\subsection{The feature matrix is mostly sparse identity}
\label{sec:anns-matrix}
The full-64-d matrix has 177 columns, and their \emph{composition} is the whole
story. \textbf{106 of the 177 columns are sparse identity one-hots}: artist
($\times 61$), \texttt{genre\_primary} ($\times 41$) and \texttt{album\_type}
($\times 4$). Only 71 columns are continuous --- the 64 song-AE PCA latents and
7 numeric metadata fields (popularities, follower count, release year,
credited-artist count). For any given track exactly one artist column, one genre
column and one album-type column is hot; the rest of the 106 are zero. The
matrix is, by construction, dominated by a mostly-zero block of mutually
exclusive indicators.

A gradient-boosted tree reads this block \emph{losslessly and cheaply}: a single
axis-aligned split on column \texttt{artist=X} partitions the data on exactly
that identity, with no interaction among the one-hots forced and no cost paid for
the columns that are zero. Capturing ``tracks by artist $X$ in genre $Y$ rotate''
costs the tree two nested splits. A dense net cannot do this. It must learn a
\emph{dense} weight from every one of the 106 mostly-zero inputs into every
hidden unit, then back out a non-linear combination --- spending capacity and
gradient signal modeling the structure of a block that the tree handles with one
comparison. This is precisely the regime in which trees structurally dominate
dense nets, and it is the project's \emph{convergent finding}, now observed
across four feature modalities and five model families: the taste-fit signal is
carried by artist/genre identity, captured by one-hots and exploited by trees,
not by descriptive attributes or distance/margin/dense models. The pyramid-MLP
is simply the dense family's turn to lose to it.

\subsection{Standardization turns rare categories into outlier spikes}
\label{sec:anns-scaler}
The dense net is fed a StandardScaler-ed matrix, and standardization interacts
pathologically with sparse one-hots. A one-hot column for a category present in
a fraction $p$ of rows has mean $p$ and standard deviation $\sqrt{p(1-p)}$.
Standardizing it maps the ``on'' rows to
\[
  z_{\text{on}} = \frac{1-p}{\sqrt{p(1-p)}} = +\sqrt{\frac{1-p}{p}},
  \qquad
  z_{\text{off}} = \frac{0-p}{\sqrt{p(1-p)}} = -\sqrt{\frac{p}{1-p}}.
\]
For a \emph{rare} category --- the long tail of artists and genres, where $p$ is
small --- the ``on'' value $+\sqrt{(1-p)/p}$ becomes a large positive spike (for
$p=0.01$ it is $\approx +9.9$), while the ``off'' value sits at a tiny negative
$-\sqrt{p/(1-p)}\approx -0.1$. Every rare identity therefore arrives at the net
as an occasional, enormous outlier on its own axis. ReLU units with Adam at
learning rate \num{1e-3} swing hard on these spikes --- a single rare-category
activation drives a large pre-activation that the optimizer chases --- and the
softmax/sigmoid head saturates, emitting near-0 or near-1 probabilities on thin
evidence. Trees are immune: a split threshold on a one-hot is invariant to any
monotone rescaling, so the scaler's spike does nothing to them. This is a direct
mechanistic contributor to the blown calibration documented next.

\subsection{The diagnostic signature: genuine ranking, broken calibration}
\label{sec:anns-signature}
The dense net's failure has a precise, reproducible signature in the metrics ---
and it is \emph{not} the failure of a weak model. The pyramid-MLP posts AUC
$\approx 0.68$ (0.6781 in the bake-off, 0.6792 best in the arch scan): genuine
ranking signal, well above the 0.5 chance line, only $\sim$3 bands below the
trees. Yet \emph{at the same time} its log-loss is $\approx 2.0$--$2.8$ and its
Brier $\approx 0.32$--$0.34$, against the trees' log-loss $\approx 0.60$--$0.63$
and Brier $\approx 0.21$ (Figure~\ref{fig:mlp-calib}). A constant base-rate
predictor (always outputting $p\approx 0.445$) scores a log-loss of only
$\approx 0.69$. \textbf{The net's log-loss of 2.0--2.8 is thus far worse than the
random baseline, while xgboost's 0.61 is well below it.}

\begin{figure}[ht]
  \centering
  \includegraphics[width=0.95\textwidth]{mlp_calibration_signature}
  \caption{The dense-net diagnostic: the net is confidently wrong, not merely
  weak. \emph{Left:} on AUC the pyramid-MLP genuinely ranks ($\approx 0.68$,
  well above chance, $\sim$3 bands below the trees) --- there is real signal in
  its ordering. \emph{Right:} but its log-loss ($\approx 2.0$--$2.8$) is far
  \emph{worse} than the random base-rate baseline ($\approx 0.69$, dashed),
  while the trees sit well below it ($\approx 0.61$). A model that ranks
  acceptably yet scores worse-than-random log-loss is emitting confident,
  miscalibrated probabilities.}
  \label{fig:mlp-calib}
\end{figure}

This combination --- acceptable ranking, worse-than-random log-loss --- means the
net is \emph{confidently wrong}: it pushes many predictions to the saturated
extremes (near 0 or near 1) on the wrong side of the label, and log-loss punishes
each such confident error severely. It is miscalibrated, not weak: the ordering
is there, but the probabilities attached to that ordering are unusable as a
\texttt{fit\_norm}. The mechanism is exactly the scaler-spike saturation of
Section~\ref{sec:anns-scaler}. Contrast the trees, whose leaf-frequency outputs
are calibrated by construction: a leaf reporting probability $q$ is the empirical
rotation frequency of the training rows that fall in it, so the tree's
probabilities are honest by design --- the reason it wins on log-loss and Brier as
well as AUC.

\subsection{Optimization versus representation: the binding constraint}
\label{sec:anns-binding}
A natural objection is that the net is merely under-optimized --- too few
iterations, the wrong geometry. The architecture scan (Section~\ref{sec:mlp-scan})
forecloses both. On the \emph{budget} axis, every one of the six geometries
self-stops at 73--100 iterations, well below even the old 300 cap (let alone the
raised 800): the optimizer reaches a loss plateau it cannot escape, so more budget
buys nothing --- the re-baseline arm M0 reproduced the under-trained AUC to all
digits because it ran the identical 73 iterations. On the \emph{architecture}
axis, all six geometries land inside a single noise band, and the best converged
geometry reproduces the under-trained AUC to four digits (0.6792 versus 0.6781).
Both knobs internal to the dense family are therefore exhausted. By elimination,
the binding constraint is neither budget nor architecture but the input
\emph{representation} --- the raw one-hots of Section~\ref{sec:anns-matrix}, which
the net cannot consume the way a tree does. This is the diagnostic that points at
representation-aware topologies (entity embeddings, wide-and-deep) as the only
untested lever, since those are designed precisely to replace raw one-hots with a
learned dense representation.

\subsection{PCA has already pre-empted the net's nonlinear job}
\label{sec:anns-pca}
There is a second reason the dense net has little to do. The 64 song-AE PCA
latents are already a \emph{linear, decorrelated} compression of a
1536-dimensional embedding: PCA has extracted the dominant linear structure and
handed the net orthogonal components with most of the easy linear variance
already separated. Whatever nonlinear interaction a hidden layer might discover
in that block has largely been pre-empted; little nonlinear structure remains for
the net's depth to exploit. The evidence is direct: a plain logistic regression
reaches AUC 0.676 on this corpus and the converged MLP reaches 0.679 --- a dead
heat. A multi-layer ReLU network buys essentially nothing over a single linear
boundary, which is exactly what one expects when the continuous block is already
linearly pre-conditioned and the categorical block is the part that carries
signal the net cannot read.

\subsection{The open question and the motivated next step}
\label{sec:anns-next}
The diagnosis is now sharp: the dense net's ceiling is set by its inability to
consume raw one-hot identity, not by budget, architecture, or the continuous
block. That isolates a single, testable open question --- \emph{can a
representation-aware topology close the gap to the trees?} The two ingredients
that target the diagnosis directly are (i) \textbf{entity embeddings} for the
identity groups (artist, genre, album-type), which replace each sparse one-hot
block with a learned low-dimensional dense vector --- the standard fix for
high-cardinality categoricals in nets, removing both the dense-weights-over-zeros
problem of Section~\ref{sec:anns-matrix} and the scaler-spike problem of
Section~\ref{sec:anns-scaler}; and (ii) a \textbf{wide-and-deep} split --- a wide
linear path over the one-hots (so identity is read directly, as the tree reads
it) alongside a deep path over the continuous block. A Rust/\texttt{burn}
implementation of such a representation-aware network --- the
\texttt{burn-deep-classifier} --- was built precisely to answer this question.
Section~\ref{sec:anns-burn} reports its results, and they resolve the diagnosis
cleanly: the calibration pathology was indeed a representation artifact and is
fixed, but the AUC gap to the trees only \emph{narrows}, it does not close.

\emph{Sources: \texttt{2026-06-12-classifier-bakeoff-p64.md},
\texttt{2026-06-12-pyramid-mlp-arch-scan.md}.}

\subsection{Resolved: the burn-deep-classifier fixes calibration but not the AUC gap}
\label{sec:anns-burn}
\emph{Rationale.} The preceding subsections left a single sharp open question:
the dense net's ceiling is set by its inability to consume raw one-hot identity,
so \emph{can a representation-aware topology close the gap to the trees?} The
\texttt{burn-deep-classifier} --- a full-Rust/\texttt{burn} network --- was built
to answer it. Its design targets the two diagnosed pathologies directly: it keeps
the identity one-hots \emph{raw} (standardizing only the 64 PCA $+$ 7 numeric
continuous block), so the scaler-spike mechanism of Section~\ref{sec:anns-scaler}
never fires, and it exposes a \texttt{topology} enum over a plain dense
\texttt{mlp}, learned entity \texttt{embeddings}, and a \texttt{wide\_deep} split.
Phase 1 compared the three topologies (epochs 250, patience 20, early-stop on
validation log-loss, hidden $[128,64]$, embed\_dim 8, lr \num{1e-3}, batch 256);
Phase 2 tuned the winner. All nine runs ran on the same full-64-d dataset
(\dsid{ds-20260609-204419-p64-s42}) and succeeded cleanly.

\emph{Results.} The headline is the calibration recovery
(Figure~\ref{fig:burn-topology}). Every topology lands at log-loss
\num{0.621}--\num{0.628} and Brier \num{0.217}--\num{0.220} --- squarely in the
trees' $\sim$0.60--0.63 / $\sim$0.21 regime --- against the sklearn pyramid's
log-loss \num{2.0}--\num{2.4} and Brier \num{0.33}. \textbf{The calibration
pathology is fixed completely.} On the topology axis (Table~\ref{tab:burn-topo}),
the plain \texttt{mlp} wins at AUC \num{0.7001}, with \texttt{embeddings}
(\num{0.6934}) and \texttt{wide\_deep} (\num{0.6918}) both inside the
$2\sigma=0.0126$ band of it and of each other --- \texttt{embeddings} edges
\texttt{wide\_deep} by a sub-band $+0.0016$. The Phase-2 \texttt{mlp} tune
(Table~\ref{tab:burn-tune}) tops out at AUC \num{0.7038} (lr \num{5e-4},
$[128,64]$): the whole 6-cell grid is a \num{0.0081}-wide strip
(\num{0.6957}--\num{0.7038}), every $[128,64]$ cell beats every $[256,128]$ cell,
and the best config is only $+0.0037$ (inside the band) over the Phase-1 winner,
so the Phase-1 \texttt{mlp} config stands. At \num{0.7038} the net sits
$-0.0073$ below \texttt{catboost} ($\sim$0.6 band) and $-0.0166$ below
\texttt{xgboost} ($\sim$1.3 bands); below the $\sim$0.71 save bar, so no
definition was persisted.

\begin{table}[ht]
  \centering
  \caption{Phase 1 --- the \texttt{burn-deep-classifier} topology bake-off
  (\dsid{ds-20260609-204419-p64-s42}, held params: epochs 250, patience 20,
  hidden $[128,64]$, embed\_dim 8, dropout 0.1, lr \num{1e-3}, batch 256, relu;
  the only axis is \texttt{topology}). All three succeeded. AUC higher-better;
  log-loss and Brier lower-better. Winner bolded; best cell per metric bolded.}
  \label{tab:burn-topo}
  \begin{tabular}{llcccc}
    \toprule
    Arm / run id & Topology & AUC & log-loss & accuracy & Brier \\
    \midrule
    \best{T-mlp} (\runid{run-20260612-140116-41fc9-burn-deep-classifier}) & \best{mlp}        & \best{0.7001} & \best{0.6211} & \best{0.6470} & \best{0.2169} \\
    T-emb (\runid{run-20260612-140116-0eda4-burn-deep-classifier})        & embeddings        & 0.6934 & 0.6251 & 0.6436 & 0.2187 \\
    T-wd (\runid{run-20260612-140116-e85be-burn-deep-classifier})         & wide\_deep        & 0.6918 & 0.6281 & 0.6407 & 0.2198 \\
    \bottomrule
  \end{tabular}
\end{table}

\begin{table}[ht]
  \centering
  \caption{Phase 2 --- tuning the winning \texttt{mlp} topology, a
  \texttt{lr}$\times$\texttt{hidden} grid (other params held). All six
  succeeded, sorted by AUC; winner bolded, best cell per metric bolded. The
  whole grid is inside one noise band, every $[128,64]$ cell beats every
  $[256,128]$ cell, and the best gain over the Phase-1 winner ($+0.0037$) is
  sub-band --- so the Phase-1 config is the config of record.}
  \label{tab:burn-tune}
  \begin{tabular}{lllcccc}
    \toprule
    Run id & lr & hidden & AUC & log-loss & accuracy & Brier \\
    \midrule
    \best{\runid{run-20260612-140211-6575a-burn-deep-classifier}} & \best{0.0005} & \best{[128,64]}  & \best{0.7038} & 0.6221 & 0.6521 & 0.2172 \\
    \runid{run-20260612-140211-21ad5-burn-deep-classifier}        & 0.002         & [128,64]         & 0.6997 & 0.6227 & \best{0.6532} & 0.2174 \\
    \runid{run-20260612-140211-18c9e-burn-deep-classifier}        & 0.001         & [128,64]         & 0.6996 & \best{0.6213} & 0.6447 & \best{0.2170} \\
    \runid{run-20260612-140211-fecb0-burn-deep-classifier}        & 0.0005        & [256,128]        & 0.6966 & 0.6232 & 0.6424 & 0.2179 \\
    \runid{run-20260612-140211-4c359-burn-deep-classifier}        & 0.002         & [256,128]        & 0.6963 & 0.6245 & 0.6460 & 0.2181 \\
    \runid{run-20260612-140211-0f4e1-burn-deep-classifier}        & 0.001         & [256,128]        & 0.6957 & 0.6252 & 0.6424 & 0.2187 \\
    \bottomrule
  \end{tabular}
\end{table}

\begin{figure}[ht]
  \centering
  \includegraphics[width=0.95\textwidth]{burn_topology}
  \caption{The representation-aware net resolves the open question of
  Section~\ref{sec:anns-next}. \emph{Left:} the three topologies are inside one
  $2\sigma$ band of each other; the plain \texttt{mlp} wins and the entity
  \texttt{embeddings} barely edge \texttt{wide\_deep}, because at these identity
  cardinalities (artist 61, genre 41, album-type 4) learned embeddings are
  starved --- they pay off at hundreds-to-thousands of categories, not dozens (a
  hypothesis Section~\ref{sec:anns-highvocab} confirms by raising the caps).
  \emph{Right:} the headline. Keeping the one-hots \emph{raw} drops the
  log-loss from the sklearn pyramid's $\approx 2.0$ (confidently miscalibrated,
  worse than the random baseline) to $\approx 0.62$, into the tree regime --- a
  clean confirmation that the calibration pathology was a representation
  artifact, not an intrinsic property of nets on this data.}
  \label{fig:burn-topology}
\end{figure}

\emph{Interpretation.} The result resolves Section~\ref{sec:anns-next}'s open
question in two parts, and the two answers point in opposite directions. First,
the \emph{calibration} verdict is overturned: the pyramid-MLP's worse-than-random
log-loss of Section~\ref{sec:anns-signature} was a \emph{representation} artifact
--- raw one-hots crushed into the standardizer --- not an intrinsic property of
dense nets on taste data, and a net that keeps the one-hots raw is calibrated in
the trees' regime. The earlier ``broken calibration is a fixed ceiling'' reading
holds only for the sklearn-pyramid representation. AUC also recovered: the burn
net's \num{0.7038} is $+0.025$ over the best converged pyramid (\num{0.6792}),
$\sim$2 bands of genuine ranking-power gain, so the prior conclusion that ``the
MLP family is closed for AUC at $\approx 0.68$'' is revised --- a raw-one-hot
Rust net reaches $\sim$0.70. Second, and against the optimistic reading, the
representation fix does \emph{not} reach the tree tier: at \num{0.7038} the net is
still $\sim$1 band below \texttt{catboost} and $\sim$1.3 below \texttt{xgboost}.
The gap narrowed dramatically (from the pyramid's $\sim$3.3 bands to $\sim$1) but
did not close, so the convergent finding survives intact --- trees still read the
sparse identity one-hots a hair more effectively than even a raw-one-hot net. The
topology comparison sharpens this: the plain \texttt{mlp} beating both the
embedding and wide-deep variants is the \emph{same} ``the identity is already in
the one-hots'' lesson the tree campaigns taught, now reproduced inside the ANN
family --- learned per-group embeddings have little to add over 61 artist and 41
genre one-hots, and the standing lever for the net (as for the trees) is identity
\emph{resolution}: raise the vocabulary caps and the embedding topology should
finally start to pay off --- a prediction that the high-vocab re-run of
Section~\ref{sec:anns-highvocab} then confirms, flipping \texttt{embeddings}
above \texttt{mlp} and lifting the net into the tree tier. No definition was
saved here (below the $\sim$0.71 bar); the
net is a strong, now-honest dense baseline and a natural ensemble member for its
decorrelated inductive bias, not a champion.
\emph{Source: \texttt{2026-06-12-burn-deep-topology.md}.}

\subsection{Confirmed: raising identity cardinality lifts the embedding net into the tree tier}
\label{sec:anns-highvocab}
\emph{Rationale.} Section~\ref{sec:anns-burn} closed on a sharp, testable
prediction: the entity-\texttt{embeddings} topology lost to plain \texttt{mlp}
\emph{only} because the identity cardinalities (artist 61, genre 41) were too
low to feed a learned embedding, so raising the vocabulary caps should flip the
ranking. This campaign is the paired re-run that tests it directly. A new matrix
(\dsid{ds-20260612-143014-p64-s42}) was built \emph{identical} to the 60/40
full-64-d dataset of Section~\ref{sec:anns-burn} --- same \num{23529} rows, same
split seed 42, byte-identical PCA64 --- differing in \emph{one} axis only:
identity cardinality, with the artist cap raised $60\to300$
(\SI{50.4}{\percent} coverage) and \texttt{genre\_primary} $40\to120$
(\SI{78.8}{\percent}). It is a \emph{third} non-comparable matrix (497 columns:
64 PCA $+$ 9 numeric $+$ 426 one-hots --- artist 301, genre 121, album-type 4 ---
versus the 60/40 matrix's 177), so all metrics here are fresh in-dataset
references and must never be read against the 60/40 or DS-META boards. The
headline question: with cardinality $5\times$ (artist) and $3\times$ (genre)
larger, does \texttt{embeddings} now beat \texttt{mlp}? Phase 0 re-ran the two
trees for an in-dataset tree tier; Phase 1 re-ran the topology bake-off (the
embeddings arm now at embed\_dim 16, since higher cardinality justifies a wider
table); Phase 2 tuned the winning topology over embed\_dim $\times$ lr.

\emph{Results.} The ranking \textbf{flipped} (Table~\ref{tab:highvocab-topo};
Figure~\ref{fig:burn-highvocab}). \texttt{embeddings}
(\runid{run-20260612-143411-e8db6-burn-deep-classifier}) wins at AUC
\num{0.7192}, with \texttt{mlp} at \num{0.7129} ($+0.0063$) and \texttt{wide\_deep}
last at \num{0.7061} ($+0.0131$). On the 60/40 matrix \texttt{mlp} had beaten
\texttt{embeddings} by $+0.0067$; here \texttt{embeddings} beats \texttt{mlp} by
$+0.0063$ --- a \textbf{swing of $+0.0130$} in the embedding's favour, crossing
the band, the qualitative reversal the starvation hypothesis predicted. Phase 0
shows both trees also moved \emph{up} on the wider matrix: xgboost
$0.7204\to\num{0.7299}$ ($+0.0095$, \runid{run-20260612-143411-02fda-xgboost-classifier},
log-loss \num{0.6022}) and catboost $0.7111\to\num{0.7252}$ ($+0.0141$,
\runid{run-20260612-143411-8f15e-catboost-classifier}), so high vocab is genuine
signal for \emph{every} family --- it does not hand the embedding a unique
advantage, it \emph{unlocks} a topology that was previously starved. The Phase-2
embed\_dim$\times$lr tune (best embed\_dim 8 / lr \num{5e-4}, AUC \num{0.7228},
\runid{run-20260612-143508-3310f-burn-deep-classifier}) shows the win came from
\emph{enabling} embeddings plus a slower learning rate, not from a wider table:
AUC is monotone-decreasing in embed\_dim at the slow lr (ed8 \num{0.7228} $>$
ed16 \num{0.7210} $>$ ed32 \num{0.7200}, a sub-band $\sim$\num{0.0028} spread),
so embed\_dim is already saturated at 8. The best tuned cell beats the Phase-1
winner by only $+0.0036$ (sub-band), so the decision rule saved the \emph{Phase-1}
config (embeddings, embed\_dim 16, lr \num{1e-3}) as
\texttt{burn-deep-embeddings-highvocab}. A determinism twin reproduced the
Phase-1 run to four digits ($0.71923$ versus $0.71923$). At \num{0.7192} the
saved net sits $-0.0060$ below catboost ($<$ 1 band) and $-0.0107$ below xgboost,
inside the band of the second tree family --- the first ANN to reach the
in-dataset tree tier and clear the rotation save bar.

\begin{table}[ht]
  \centering
  \caption{Phase 1 --- the \texttt{burn-deep-classifier} topology bake-off on
  the \textbf{high-vocab} matrix (\dsid{ds-20260612-143014-p64-s42}, held
  params: epochs 250, patience 20, hidden $[128,64]$, embed\_dim 16, dropout
  0.1, lr \num{1e-3}, batch 256, relu; the only axis is \texttt{topology}). All
  three succeeded. The ranking \emph{reverses} the 60/40 result of
  Table~\ref{tab:burn-topo}: \texttt{embeddings} now wins. AUC higher-better;
  log-loss and Brier lower-better. Winner bolded; best cell per metric bolded.}
  \label{tab:highvocab-topo}
  \begin{tabular}{llcccc}
    \toprule
    Topology / run id & Topology & AUC & log-loss & accuracy & Brier \\
    \midrule
    \best{T-emb} (\runid{run-20260612-143411-e8db6-burn-deep-classifier}) & \best{embeddings} & \best{0.7192} & \best{0.6087} & \best{0.6694} & \best{0.2110} \\
    T-mlp (\runid{run-20260612-143411-01d7b-burn-deep-classifier})        & mlp               & 0.7129 & 0.6166 & 0.6617 & 0.2143 \\
    T-wd (\runid{run-20260612-143411-7d117-burn-deep-classifier})         & wide\_deep        & 0.7061 & 0.6256 & 0.6526 & 0.2183 \\
    \bottomrule
  \end{tabular}
\end{table}

\begin{table}[ht]
  \centering
  \caption{Phase 2 --- tuning the winning \texttt{embeddings} topology, an
  embed\_dim$\times$\texttt{lr} grid (hidden $[128,64]$, other params held),
  sorted by AUC; winner bolded, best cell per metric bolded. AUC is
  monotone-decreasing in embed\_dim at lr \num{5e-4} (the table is already
  saturated at width 8), and lr \num{5e-4} beats lr \num{1e-3} at every width.
  The best gain over the Phase-1 winner ($+0.0036$) is sub-band, so the Phase-1
  config is the config of record.}
  \label{tab:highvocab-tune}
  \begin{tabular}{llcccc}
    \toprule
    embed\_dim & lr & AUC & log-loss & accuracy & Brier \\
    \midrule
    \best{8}  & \best{0.0005} & \best{0.7228} & 0.6105 & 0.6645 & 0.2114 \\
    16        & 0.0005        & 0.7210 & \best{0.6088} & \best{0.6711} & \best{0.2108} \\
    32        & 0.0005        & 0.7200 & 0.6100 & 0.6657 & 0.2114 \\
    8         & 0.001         & 0.7193 & 0.6092 & 0.6643 & 0.2112 \\
    16        & 0.001         & 0.7192 & 0.6087 & 0.6666 & 0.2110 \\
    32        & 0.001         & 0.7183 & 0.6097 & 0.6668 & 0.2114 \\
    \bottomrule
  \end{tabular}
\end{table}

\begin{figure}[ht]
  \centering
  \includegraphics[width=0.95\textwidth]{burn_highvocab}
  \caption{Identity resolution closes the ANN-pitfalls thread. \emph{Left:} the
  paired \texttt{embeddings}-versus-\texttt{mlp} comparison across the two
  cardinalities. On the 60/40 matrix \texttt{mlp} wins by $+0.0067$; on the
  high-vocab matrix \texttt{embeddings} wins by $+0.0063$ --- a $+0.0130$ swing
  that reverses the ranking and confirms the starvation diagnosis: with artist
  cardinality $5\times$ and genre $3\times$ larger the embedding table finally
  has enough distinct identities to learn a useful low-rank geometry instead of
  collapsing near-constant. \emph{Right:} the five-family in-dataset AUC ranking
  on the high-vocab matrix. The saved \texttt{embeddings} net (\num{0.7192})
  reaches the bottom of the tree tier (shaded: from a band below \texttt{catboost}
  up to \texttt{xgboost}) --- the first ANN to do so --- while \texttt{xgboost}
  (\num{0.7299}) remains the in-dataset leader.}
  \label{fig:burn-highvocab}
\end{figure}

\emph{Interpretation.} This resolves Section~\ref{sec:anns-burn}'s open question
\emph{confirmed}: the embedding topology's earlier loss was a starvation
artifact of low cardinality, not an intrinsic limit, and raising the caps flips
the ranking exactly as predicted. Two honest caveats keep the claim calibrated.
First, the single-split emb-vs-mlp margin ($+0.0063$) is itself sub-band, so the
\emph{magnitude} of the reversal needs a 3-split confirm; what is decisive is
the \emph{direction} --- the whole ranking reordered, with the embedding arm
moving from the bottom of the 60/40 ranking to the top of this one. Second, high
vocab helps \emph{everyone}: the trees gained as much as (catboost) or more than
the embedding arm, so the correct reading is ``higher resolution is better
signal for every family, and it specifically unlocks the embedding topology'',
not ``embeddings now win outright'' --- \texttt{xgboost} at \num{0.7299} is still
the in-dataset leader, $\sim$\num{0.0107} ($\sim$1 band) above the saved net.
But the qualitative arc that ran through Section~\ref{sec:anns} now closes
cleanly: a dense net that consumes raw one-hots is calibrated in the tree regime
(Section~\ref{sec:anns-burn}), and given \emph{enough} identity resolution its
embedding topology reaches the tree tier --- the standing ``identity resolution
is the lever'' lever, demonstrated. Logloss ($\sim$0.61) and Brier ($\sim$0.21)
stay in the trees' calibrated regime throughout, so there is no calibration
penalty for the win. One data-hygiene note: the high-vocab manifest's redundancy
report flags 21 near-perfect genre$\leftrightarrow$artist one-hot collinearities
(e.g.\ \texttt{genre\_primary=australian psych} $\equiv$ \texttt{artist=Tame
Impala}, correlation \num{1.000}; 19 more $\ge\num{0.957}$), an artifact of
admitting rare genres carried by a single newly-admitted artist; the trees and
the net's regularizers absorb it with no observed run pathology, but it is worth
knowing before reading individual one-hot importances on this matrix. The saved
\texttt{burn-deep-embeddings-highvocab} is flagged for best-model-selector as
the first tree-tier ANN and a natural family-diversity candidate.
\emph{Source: \texttt{2026-06-12-high-vocab-topology.md}.}

\subsection{A second representation-aware net: the AE-transplant confirms the ceiling}
\label{sec:anns-ae}
\emph{Rationale.} Section~\ref{sec:anns-burn} established that a net which keeps
the identity one-hots \emph{raw} is calibrated in the tree regime yet still caps
$\sim$1 band below the trees. A natural next question is whether a \emph{learned
unsupervised representation} of the same inputs --- rather than the raw inputs
themselves --- lifts the net further. The \texttt{burn-ae-classifier} tests it:
Phase 1 pretrains an autoencoder to reconstruct the full assembled feature
vector (standardized continuous block beside the raw one-hots, $D{=}177$) by
MSE, Phase 2 transplants the encoder into a softmax classifier and fine-tunes.
This campaign is the family's first characterization --- a 13-config scan on the
60/40 full-64-d matrix (\dsid{ds-20260609-204419-p64-s42}, seed 42) over latent
size, head depth/width, activation, and the \texttt{freeze\_encoder} flag, read
against the in-corpus xgboost champion (\num{0.7204}) and the comparable
\texttt{burn-deep} \texttt{mlp} (\num{0.7038}).

\emph{Results.} \textbf{Capacity does not matter for this family on rotation}
(Table~\ref{tab:burn-ae}; Figure~\ref{fig:burn-ae}). The 12 fine-tuned configs
span a single \num{0.0054}-wide strip (\num{0.69175}--\num{0.69719}), entirely
inside the $2\sigma=0.0126$ band --- a flat tie across a $4\times$ latent range
(32 to 128), across head depth and width, and across activation. The nominal
best (cfg3, enc$[128,64]$/clf$[128,64]$, AUC \num{0.69719}) beats the
\num{0.69225} baseline by only $+0.00494$ (sub-band), and even the bare linear
probe (cfg2, clf$[]$, \num{0.69175}) is only $-0.00544$ below it: a single linear
layer on the fine-tuned latent already captures essentially all the family's
skill. The one out-of-band effect is a \emph{trap}: freezing the MSE-pretrained
encoder is safe under a shallow head (cfg11, \num{0.69530}, in-band) but
\emph{collapses} under a deep head (cfg10, \num{0.68210}, $-0.0151$ $\approx$
1.2 bands, the worst row). relu edges gelu and silu by a sub-band margin.
Calibration is healthy throughout (log-loss $\sim$0.62, Brier $\sim$0.22 --- the
trees' regime), so this is a legitimate calibrated dense member, just not an AUC
contender: the best config sits $\sim$1.8 bands below \texttt{xgboost} and is
\emph{statistically equal} to the plain \texttt{burn-deep} \texttt{mlp}
($-0.0066$, sub-band). No definition was saved.

\begin{table}[ht]
  \centering
  \caption{The \texttt{burn-ae-classifier} capacity + head-architecture scan
  (\dsid{ds-20260609-204419-p64-s42}, seed 42; held: ae\_epochs 80, ae\_lr
  \num{1e-3}, epochs 200, lr \num{1e-3}, dropout 0.1, patience 25). All 13 runs
  succeeded. AUC and accuracy higher-better; log-loss lower-better. Winner
  (cfg3) and best cell per metric bolded; the frozen-deep-head collapse (cfg10)
  is the only out-of-band row. Sorted by AUC.}
  \label{tab:burn-ae}
  {\small\begin{tabular}{@{}p{3.3cm}p{5.0cm}ccc@{}}
    \toprule
    Run id & enc / clf / freeze / act & AUC & log-loss & acc \\
    \midrule
    \best{\runid{run-20260615-152359-21f11-burn-ae-classifier}} & \best{[128,64] / [128,64] / ft / relu} & \best{0.69719} & \best{0.6216} & \best{0.6496} \\
    \runid{run-20260615-152359-05bce-burn-ae-classifier} & [256,128] / [128] / ft / relu   & 0.69712 & 0.6254 & 0.6485 \\
    \runid{run-20260615-152359-e5a75-burn-ae-classifier} & [128,32] / [64] / ft / relu     & 0.69657 & 0.6246 & 0.6451 \\
    \runid{run-20260615-152359-94f8f-burn-ae-classifier} & [128,64] / [64] / ft / relu (baseline re-run) & 0.69645 & 0.6264 & 0.6496 \\
    \runid{run-20260615-152359-7a459-burn-ae-classifier} & [128,64] / [256,128] / ft / relu & 0.69561 & 0.6235 & 0.6434 \\
    \runid{run-20260615-152359-cb93b-burn-ae-classifier} & [384,192] / [192,96] / ft / relu & 0.69541 & 0.6224 & 0.6426 \\
    \runid{run-20260615-160112-aa1db-burn-ae-classifier} & [256,128] / [128] / \emph{frozen} / relu & 0.69530 & 0.6242 & 0.6409 \\
    \runid{run-20260615-160112-4a764-burn-ae-classifier} & [128,64] / [128,64] / ft / gelu & 0.69340 & 0.6256 & 0.6411 \\
    \runid{run-20260615-152359-3b006-burn-ae-classifier} & [256,128,64] / [64] / ft / relu & 0.69302 & 0.6279 & 0.6481 \\
    \runid{run-20260615-160113-7567b-burn-ae-classifier} & [128,64] / [128,64] / ft / silu & 0.69242 & 0.6230 & 0.6436 \\
    \runid{run-20260615-152359-86fd5-burn-ae-classifier} & [256,128] / [256,128] / ft / relu & 0.69240 & 0.6275 & 0.6392 \\
    \runid{run-20260615-152359-dbb6a-burn-ae-classifier} & [128,64] / [] (linear probe) / ft / relu & 0.69175 & 0.6282 & 0.6405 \\
    \runid{run-20260615-160112-17da7-burn-ae-classifier} & [128,64] / [128,64] / \emph{frozen} / relu (collapse) & 0.68210 & 0.6378 & 0.6373 \\
    \bottomrule
  \end{tabular}}
\end{table}

\begin{figure}[ht]
  \centering
  \includegraphics[width=0.8\textwidth]{burn_ae_capacity}
  \caption{The AE-transplant net confirms the dense-net ceiling on two axes.
  \emph{Left:} the 12 fine-tuned configs collapse onto a single
  \num{0.0054}-wide strip inside the noise band, far below the comparable
  \texttt{burn-deep} \texttt{mlp} (dashed, \num{0.7038}) and the xgboost
  champion (green, \num{0.7204}) --- latent size, head depth/width and
  activation are all spent levers; capacity is saturated at the smallest config
  tested. \emph{Right:} the one real effect. Freezing the MSE-pretrained encoder
  is safe under a (near-)linear head (cfg11, in-band) but \emph{collapses} the
  model under a deep head (cfg10, $-0.0151 \approx 1.2$ bands): a deep head
  stacked on a frozen, task-agnostic representation just underfits the label,
  because the gradients that would have re-shaped the encoder toward the target
  are blocked. End-to-end fine-tuning (the registry default) avoids the trap.}
  \label{fig:burn-ae}
\end{figure}

\emph{Interpretation.} The AE pretrain buys a fast, stable init but \emph{no
AUC} over a from-scratch raw-one-hot net: because the reconstruction objective
is purely descriptive (MSE on the assembled vector), the latent it hands the
classifier is task-agnostic, and no amount of encoder/latent/head capacity turns
it into tree-tier ranking power. This is the same ``a learned representation
downstream of the identity one-hots adds nothing the trees do not already
have'' lesson --- the AE-latent dilution of Section~\ref{sec:ae-blend} and the
entity-embedding starvation of Section~\ref{sec:anns-burn} --- now reproduced in
a third dense family. The freeze-collapse is a citable pitfall: an MSE-optimal
feature extractor is not label-optimal, so a deep head over a frozen encoder
stacks unhelpful capacity and underfits. If the family is ever revisited the
lever is the pretraining \emph{objective} (supervised or contrastive), not its
size. The family is closed for AUC on rotation as configured.
\emph{Source: \texttt{2026-06-15-burn-ae-classifier-capacity-scan.md}.}

\subsection{Refuted: a GBDT-leaf-embedding head cannot recover the tree that generated its features}
\label{sec:leaf-head}
\emph{Rationale.} The ANN thread above closed on a clean lesson --- a learned
representation fed downstream of the identity one-hots does not buy ranking
power the trees do not already have (the AE-latent dilution of
Section~\ref{sec:ae-blend}; entity embeddings starved unless cardinality is
high, Section~\ref{sec:anns-burn}). This campaign stress-tests that lesson at
its hardest point with the classic Facebook GBDT$\to$LR trick, generalized to
GBDT$\to$MLP. The frozen champion tree --- the registered
\texttt{xgboost-classifier-rotation-p64} config, \emph{hard-frozen}, not a
tunable --- emits, for each row, the index of the leaf it falls into in every
one of its 500 boosted trees; those indices are one-hot-encoded into the tree's
own learned axis-aligned partition ($\sim$10k leaf-membership columns), and a
fresh head (an $\ell_2$-logistic linear head, or a shallow MLP) is trained on
top. The input to the head is thus the tree's \emph{own} partition; the
hypothesis --- a seventh test of the convergent finding --- was that even this
would \emph{not} beat the bare tree, the expected refutation being the citable
finding. All arms ran on the 60/40 full-64-d matrix
(\dsid{ds-20260609-204419-p64-s42}, seed 42); the only valid reference is the
in-dataset \num{0.7204} champion, never the \num{0.7429} DS-META or
\num{0.7299} high-vocab numbers (different matrices).

\emph{Results.} Refuted exactly as predicted (Table~\ref{tab:leaf-head};
Figure~\ref{fig:leaf-head}). The control C0
(\runid{run-20260615-151551-f2911-xgboost-classifier}) reproduced the champion
to all 16 digits (\num{0.7204383403100314}), so the determinism gate passed and
the treatments are interpretable. Every head then capped \emph{below} the bare
tree by more than one noise band: GBDT$\to$linear at AUC \num{0.6978}
($-0.0227$, $\sim$1.8 bands), GBDT$\to$MLP$[256,128]$ at \num{0.6921}
($\sim$2.2 bands), GBDT$\to$MLP$[128]$ at \num{0.6897} ($\sim$2.4 bands) --- the
best of them is itself \num{0.0100} below the refutation threshold
(champion $-$ band $=$ \num{0.7078}). The in-band-tie branch (which would have
fired a 3-seed spot-check) and the win branch ($>\num{0.7330}$) never triggered;
no definition was saved. The linear head --- the literal Facebook recipe --- is
the best of the three, and adding MLP capacity over the leaf one-hots does not
recover the tree, consistent with the recipe's original framing as a cheap
\emph{linear} trick. Worse than the AUC deficit, \emph{every head is badly
mis-calibrated}: log-loss \num{0.8162}--\num{1.3119} and Brier
\num{0.2603}--\num{0.2951} against the tree's \num{0.6075} / \num{0.2109}. The
$\ell_2$-logistic head at the default inverse-$C{=}1$ is under-regularized over
$\sim$10k sparse leaf indicators --- it drives a handful of leaf weights to
large magnitudes and emits over-confident probabilities, the same
sklearn-pyramid over-confidence signature (log-loss far above the trees' while
AUC is only modestly behind) seen in Section~\ref{sec:anns-signature}, and a
calibration red flag for best-model-selector should any of these runs ever be
considered for the group.

\begin{table}[ht]
  \centering
  \caption{The GBDT-leaf-head ceiling probe
  (\dsid{ds-20260609-204419-p64-s42}, seed 42; the leaf-emitting GBDT is the
  frozen champion config in every treatment, only the head varies). All four
  runs succeeded. AUC higher-better; log-loss and Brier lower-better. Control
  (winner) bolded; best cell per metric bolded; treatment rows sorted by AUC.
  Refutation threshold (champion $-2\sigma$) $=\num{0.7078}$: every head is
  below it.}
  \label{tab:leaf-head}
  \begin{tabular}{llcccc}
    \toprule
    Arm / run id & Head & AUC & log-loss & accuracy & Brier \\
    \midrule
    \best{C0} (\runid{run-20260615-151551-f2911-xgboost-classifier}) & \best{bare tree (champion)} & \best{0.7204} & \best{0.6075} & \best{0.6585} & \best{0.2109} \\
    T1 (\runid{run-20260615-151551-b4db5-gbdt-leaf-head-classifier}) & linear ($\ell_2$ logistic, l2 1.0) & 0.6978 & 1.3119 & 0.6441 & 0.2944 \\
    T3 (\runid{run-20260615-151551-44b1d-gbdt-leaf-head-classifier}) & mlp $[256,128]$, head\_lr \num{1e-3} & 0.6921 & 1.1813 & 0.6396 & 0.2951 \\
    T2 (\runid{run-20260615-151551-9d655-gbdt-leaf-head-classifier}) & mlp $[128]$, head\_lr \num{1e-3} & 0.6897 & 0.8162 & 0.6356 & 0.2603 \\
    \bottomrule
  \end{tabular}
\end{table}

\begin{figure}[ht]
  \centering
  \includegraphics[width=0.8\textwidth]{leaf_head_ceiling}
  \caption{The leaf-embedding stack is refuted on both axes. \emph{Left:} every
  GBDT$\to$head arm caps below the bare champion tree on AUC --- the best
  (GBDT$\to$linear, \num{0.6978}) is itself a full band below the refutation
  threshold (the dashed line at champion $-2\sigma=\num{0.7078}$; the shaded
  strip is the band). \emph{Right:} and every head is mis-calibrated, with
  log-loss \num{0.82}--\num{1.31} far above the tree regime (\num{0.61}, dotted)
  and past the random base-rate baseline (\num{0.69}, dashed) for the two
  worst. Learning a head over the tree's own one-hot leaf indices is strictly
  lossier than reading the tree's leaf-value sums directly.}
  \label{fig:leaf-head}
\end{figure}

\emph{Interpretation.} The refutation is decisive and mechanistic. The tree's
additive leaf-value scoring already weights every leaf \emph{jointly} across 500
boosting rounds with shrinkage; re-encoding those leaves as one-hot membership
indicators and fitting a head \emph{post-hoc} discards the boosting structure
and re-estimates it worse. The leaf one-hots are a strictly lossier view of a
tree that already scores better directly --- the dense head cannot even recover
the tree that generated its own features. This is the same family as the
AE-latent dilution (Section~\ref{sec:ae-blend}) and the dense-head ceiling
(Section~\ref{sec:anns}), pushed to its limit: here the head's \emph{input} is
the tree's optimal partition, and it still caps $\sim$2 bands below. Re-tuning
the head regularization could tame the over-confidence but cannot close an
$\sim$1.8-band AUC deficit, so the leaf-head family is closed for rotation. This
is the seventh convergent confirmation that the durable lever is identity
\emph{resolution} (raising vocab caps, which already lifted xgboost to
\num{0.7299} on the high-vocab matrix), not new heads or representations layered
on the tree.
\emph{Source: \texttt{2026-06-15-gbdt-leaf-head-ceiling-probe.md}.}

\subsection{The identity-resolution vocab-cap scan: a largely-spent lever}
\label{sec:vocab-cap}
\emph{Rationale.} Every campaign above closes on the same standing lever ---
identity \emph{resolution}: raising the artist/genre vocabulary caps so the
\texttt{\_\_other\_\_} tail gets real one-hot columns, which already lifted
xgboost from \num{0.7204} (60/40) to \num{0.7299} (300/120,
Section~\ref{sec:anns-highvocab}) and flipped the embedding net into the tree
tier. But only \emph{two} points existed on that cap$\to$AUC curve, and the
high end was the in-lineage record. This campaign maps where the lever stops
paying. It holds the registered champion \texttt{xgboost-classifier-rotation-p64}
config \emph{frozen} (\texttt{catboost} defaults as a secondary arm), holds
\emph{everything else} fixed --- same \num{23529}-row corpus, same split seed 42
(\num{18823} train / \num{4706} test), same PCA64, same field block --- and
moves the \emph{single} axis \texttt{vocab\_top\_n}, adding three higher-cap
points: 500/200 (777 cols), 1000/300 (1377 cols) and the full vocabulary,
5594/962 (6633 cols). The hypothesis to discriminate: \emph{saturation} (gain
per step $<$ one band) versus \emph{reversal} (AUC drops as singleton-artist
columns inject memorizable variance).

\emph{Results.} All eight launched runs succeeded; the controls reproduced both
anchors to four digits (60/40 xgboost \num{0.72044} versus reference
\num{0.7204}, $\Delta+0.00004$; 300/120 xgboost \num{0.72988} versus
\num{0.7299}, $\Delta-0.00002$), so the whole curve sits under one consistent
measurement (Table~\ref{tab:vocab-cap}; Figure~\ref{fig:vocab-cap}). xgboost
rose to a new in-lineage single-split best of \textbf{\num{0.73490}} at 1000/300
(\runid{run-20260615-163803-e75be-xgboost-classifier}), but the step structure
tells the real story: the jump $60\to300$ was $+0.00944$ ($+0.75$ band), then
$300\to500$ was \emph{flat} ($-0.00022$, within noise) and $500\to1000$ was
$+0.00524$ ($+0.42$ band). The best point is only $+0.00502$ ($\sim$0.40 band)
over the 0.7299 anchor --- a \emph{sub-band} improvement, well short of the
$+0.0126$ a confident beat would require. \texttt{catboost} tracks xgboost
qualitatively at every cap (\num{0.71105}$\to$\num{0.72516}$\to$\num{0.72372}%
$\to$\num{0.72633}), staying $\sim$0.005--0.009 below it throughout and rising
just as gently after 300/120 --- its native categorical handling does
\emph{not} scale better at high cardinality. The top-of-curve all-vocab point
(6633 cols) \textbf{could not be built}: two POST attempts each returned a
build id but died mid-materialization against the host memory ceiling
($\sim$2.3 GB free of 15 GB under concurrent load), leaving no dataset and no
artifacts. That is an infrastructure ceiling, \emph{not} a modeling result, and
the curve's top end is recorded as unmeasured.

\begin{table}[ht]
  \centering
  \caption{The vocab-cap scan on the bare GBDT
  (frozen champion xgboost config; \texttt{catboost} defaults; the only axis is
  \texttt{vocab\_top\_n}, all else byte-identical to the 300/120 anchor: same
  \num{23529} rows, split seed 42, PCA64, field block). Sorted by artist/genre
  cap. AUC higher-better; log-loss and Brier lower-better. The new in-lineage
  single-split best run is bolded; best cell per metric bolded. The all-vocab
  row is the failed build (host memory ceiling, two attempts), recorded as an
  infrastructure failure, not a result. AUC band $2\sigma\approx 0.0126$.}
  \label{tab:vocab-cap}
  \begin{tabular}{lllccccl}
    \toprule
    Run id & cap & cols & model & AUC & log-loss & Brier & Note \\
    \midrule
    \runid{run-20260615-163751-de9a4-xgboost-classifier} & 60/40    & 177  & xgb & 0.72044 & 0.60754 & 0.21087 & control \\
    \runid{run-20260615-163803-3391f-catboost-classifier} & 60/40   & 177  & cat & 0.71105 & 0.61324 & 0.21316 & control \\
    \runid{run-20260615-163803-536d5-xgboost-classifier} & 300/120  & 497  & xgb & 0.72988 & 0.60223 & 0.20803 & anchor \\
    \runid{run-20260615-163803-6f170-catboost-classifier} & 300/120 & 497  & cat & 0.72516 & 0.60681 & 0.20994 & anchor \\
    \runid{run-20260615-163803-1a0ba-xgboost-classifier} & 500/200  & 777  & xgb & 0.72966 & 0.60392 & 0.20856 & flat ($-0.00022$) \\
    \runid{run-20260615-163803-1fae3-catboost-classifier} & 500/200 & 777  & cat & 0.72372 & 0.60894 & 0.21068 & dips below own 300/120 \\
    \best{\runid{run-20260615-163803-e75be-xgboost-classifier}} & \best{1000/300} & \best{1377} & \best{xgb} & \best{0.73490} & \best{0.60070} & \best{0.20700} & \best{new in-lineage best (sub-band)} \\
    \runid{run-20260615-163803-9ed65-catboost-classifier} & 1000/300 & 1377 & cat & 0.72633 & 0.60931 & 0.21068 & catboost best on its curve \\
    (build failed) & all 5594/962 & 6633 & --- & --- & --- & --- & NOT built (memory ceiling) \\
    \bottomrule
  \end{tabular}
\end{table}

\begin{figure}[ht]
  \centering
  \includegraphics[width=0.8\textwidth]{vocab_cap_scan}
  \caption{The identity-resolution lever is largely spent. Holding the frozen
  champion config and moving only the artist/genre one-hot cap, the cap$\to$AUC
  curve climbs sharply on the first jump ($60/40\to300/120$, $+0.75$ band) and
  then \emph{plateaus}: $300\to500$ is flat ($-0.00022$, within noise) and
  $500\to1000$ recovers only $+0.42$ band, so the new in-lineage single-split
  best (\num{0.73490} at 1000/300) clears the 300/120 anchor by just $+0.00502$
  ($\sim$0.40 band, inside the shaded $+2\sigma$ envelope and so a \emph{sub-band}
  improvement). \texttt{catboost} (dashed) tracks xgboost at every cap, $\sim$%
  0.005--0.009 below --- native categorical handling buys no high-cardinality
  edge. The full-vocabulary point (6633 cols, red $\times$) is \emph{unmeasured}:
  its build hit a host memory ceiling on both attempts, an infrastructure
  failure rather than a modeling result.}
  \label{fig:vocab-cap}
\end{figure}

\emph{Interpretation.} The verdict is a clean ``not saturated, not reversed,
but diminishing.'' No interior point fell a full band below the anchor (the
lowest, 500/200, is $-0.00022$), so the reversal hypothesis does not fire ---
the singleton-artist columns at high cardinality do \emph{not} hurt the tree,
consistent with the project's no-blow-up prior that axis-aligned splits tolerate
collinear, mostly-zero one-hots by simply not splitting on the useless ones. But
the gains past the first jump are all sub-band and wobble up and down inside the
$2\sigma=0.0126$ envelope: the curve past 300/120 is better read as a flat
plateau with single-split jitter than as a monotone climb, which is exactly why
the decision rule guards a save behind a $\ge$1-band beat plus a 3-seed confirm.
The best point's $+0.00502$ is in the ``likely better, needs confirmation''
window, so \textbf{the 3-seed confirm did not fire and no definition was saved};
the cross-matrix DS-META champion (\num{0.7429}) holds rank 1 unchanged, and the
1000/300 point auto-promoted to best-models rank 3 without displacing any
registered model. The honest read for future work: raw vocab-cap \emph{height}
is a largely-spent lever on this corpus --- the big payoff was the first jump,
already on record. The catboost agreement seals it: two independent tree
families say the marginal columns past 300/120 carry little independent signal.
The residual rotation signal is the behavioral/contextual part no
song-metadata cardinality can capture, so identity work should pivot from
bigger caps to resolution \emph{quality} --- artist-followers / popularity
interactions, genre hierarchy --- and the standing ``identity resolution is the
lever'' guidance should now read ``weak and plateauing past 300/120, with the
top end still unmeasured pending a wide build on a quiescent host.''
\emph{Source: \texttt{2026-06-15-vocab-cap-scan-gbdt.md}.}

\subsection{Reading the nets: per-model sparse autoencoders}
\label{sec:anns-sae}
\emph{Rationale.} Every result so far has been read off the leaderboard --- AUC,
log-loss, Brier. But two nets can post identical AUC while carving the input
into very different internal representations, a distinction the leaderboard
cannot see. This campaign opens that box. It first ports a \emph{per-model
sparse autoencoder} (\texttt{POST /api/interp/model-sae}) into the
\texttt{burn-deep-classifier} family (the SAE head previously existed only for a
regressor), so that promotion now auto-queues an interpretability read: a
$2\times$-overcomplete SAE (l1 \num{0.0015}, 40 epochs) is fit on each hidden
layer's activations, its atoms are probed for interpretable \emph{concepts}, and
the linear \emph{decodability} of rotation is measured layer by layer. The first
use compares two geometries at fixed AUC on the canonical full-64-d matrix
(\dsid{ds-20260609-204419-p64-s42}, seed 42): a tapering \texttt{pyramid}
$[256,128,64]$ against a uniform \texttt{rect} $[128,128,128]$, differing
\emph{only} in hidden shape.

\emph{Results.} The two nets are \textbf{statistically tied} on AUC (pyramid
\num{0.69857}, rect \num{0.69937}; $\Delta\num{0.0008}\ll$ band) --- both in the
on-record \texttt{burn-deep} \texttt{mlp} tier and $\sim$1.7 bands below the
xgboost champion. They also capture the \emph{same amount} of decodable signal:
linear decodability of rotation barely lifts above the raw input ($\sim$0.111
$\to\sim$0.12) for both. The difference is purely in \emph{layout}
(Figure~\ref{fig:sae-pyramid-rect}). The pyramid's interpretable concepts peak
in its middle 128-d layer (72) and then compress $\sim$40\% out at the 64-d
bottleneck (42); the uniform net spreads them more evenly and \emph{retains}
more to the end ($17\to66\to54$), because it has the width to. Neither net is
starved --- both fully utilize every SAE atom (0 dead / 0 rare), give
essentially every identity segment ($\sim$103--105/105) a dedicated separating
atom, and reconstruct $>96\%$ of activation variance --- so the concept-count
gap is genuine allocation, not capacity. A no-early-stop control confirmed a
registered pitfall: run to the full 100 epochs with patience 0, both 3-layer
nets overfit (AUC $-\num{0.015}$, log-loss blows to \num{0.96}--\num{1.22}); the
$\sim$0.699 / 0.625 calibrated regime is recovered only \emph{because} the
production runs early-stop. \textbf{Deeper \texttt{burn-deep} nets need
patience.}

\begin{figure}[ht]
  \centering
  \includegraphics[width=0.72\textwidth]{sae_pyramid_rect}
  \caption{At equal AUC, two geometries lay out the same signal differently ---
  a representational difference invisible to the leaderboard. The tapering
  \texttt{pyramid} $[256,128,64]$ (blue) peaks its SAE-interpretable concepts in
  the middle 128-d layer (72) and squeezes $\sim$40\% of them out at the 64-d
  neck (42); the uniform \texttt{rect} $[128,128,128]$ (grey) spreads its
  concepts and retains more to the output layer ($17\to66\to54$). Both nets
  fully utilize every atom and decode rotation to the same $\sim$0.12 ceiling,
  so the pyramid \emph{consolidates} its interpretable structure into a
  compressed bottleneck while the uniform net keeps a broader concept set ---
  same AUC, different allocation.}
  \label{fig:sae-pyramid-rect}
\end{figure}

\emph{Interpretation.} At equal ranking power the pyramid consolidates its
interpretable structure into a narrowing bottleneck while the uniform net
distributes it, and both merely re-express the raw input's $\sim$0.11
decodability rather than manufacturing new taste signal. The SAE thus confirms,
mechanistically, what the leaderboard only implies: the dense nets are not
finding structure the trees miss; they are re-packaging the same identity signal
in a representation that AUC cannot distinguish. This ``same signal, different
allocation'' pattern recurs across families at higher vocab
(Section~\ref{sec:anns-capacity}). \emph{Source:
\texttt{2026-07-12-pyramid-vs-rect-sae.md}.}

\subsection{Capacity is a spent lever at high vocab, and what the SAE sees}
\label{sec:anns-capacity}
\emph{Rationale.} The high-vocab matrix (Section~\ref{sec:anns-highvocab}) was
where the embedding net first reached the tree tier, but it was never
depth/width-tuned. This campaign asks the direct question: does \emph{capacity}
--- width or depth --- lift the ANN family across the remaining $\sim$0.2--0.6
band gap to the trees on that matrix (\dsid{ds-20260612-143014-p64-s42}, 497
cols)? Two one-axis geometry ladders, lr held at each topology's on-record best:
an \texttt{embeddings} ladder (ed8, lr \num{5e-4}) and an \texttt{mlp} ladder
(lr \num{1e-3}), each stepping $[128,64]$ (control) $\to[256,128]$ (wider)
$\to[256,128,64]$ (deeper). Each promoted arm also gets the per-model SAE of
Section~\ref{sec:anns-sae}, for a cross-family representation read at high vocab.

\emph{Results.} \textbf{CEILING --- capacity is saturated for both topologies}
(Table~\ref{tab:highvocab-capacity}; Figure~\ref{fig:highvocab-capacity}).
Neither widening nor deepening beats its own $[128,64]$ control by more than one
band; the best config stays the on-record \texttt{embeddings}
ed8/lr\num{5e-4}/$[128,64]$ (reproduced here at \num{0.72274}, $\Delta$
$-\num{0.0001}$). Deepening \texttt{embeddings} to three layers actively
\emph{hurts} ($-\num{0.0152}\approx$ 1.2 bands, the only out-of-band move);
deepening \texttt{mlp} helps a sub-band touch ($+\num{0.0057}$). Crucially, the
best \texttt{mlp} (B6, \num{0.71800}) closes to within $+\num{0.0047}$ of the
best \texttt{embeddings} (B1, \num{0.72274}): at their best geometries the two
topologies tie within a band, so \textbf{capacity, not the embeddings mechanism,
was the lever} --- the earlier ``embeddings wins at high vocab'' reading is an
in-band tie, and the ANN family remains $\sim$0.2--0.6 band below the tree tier
(catboost \num{0.7252}, xgboost \num{0.7299}). The saved
\texttt{burn-deep-embeddings-highvocab} definition and the champion are both
unchanged; the best arm was promoted for the record as
\texttt{emb-128-64-ed8-highvocab}.

\begin{table}[ht]
  \centering
  \caption{The \texttt{burn-deep-classifier} capacity scan on the high-vocab
  matrix (\dsid{ds-20260612-143014-p64-s42}, 497 cols; \texttt{activation=relu},
  dropout 0.1, epochs 250, patience 20; only the hidden geometry varies within a
  topology). All six succeeded. AUC and accuracy higher-better; log-loss and
  Brier lower-better. Winner (B1) and best cell per metric bolded; sorted by
  AUC. The deep \texttt{embeddings} arm (B3) is the only out-of-band move.}
  \label{tab:highvocab-capacity}
  {\small\begin{tabular}{@{}p{3.3cm}p{2.9cm}cccc@{}}
    \toprule
    Run id & topology / hidden & AUC & log-loss & acc & Brier \\
    \midrule
    \best{\runid{run-20260712-171008-065bd-burn-deep-classifier}} & \best{embeddings / [128,64]} & \best{0.72274} & \best{0.61043} & 0.66426 & \best{0.21140} \\
    \runid{run-20260712-171008-7c302-burn-deep-classifier} & mlp / [256,128,64]  & 0.71800 & 0.61266 & 0.66065 & 0.21276 \\
    \runid{run-20260712-171008-cf746-burn-deep-classifier} & embeddings / [256,128] & 0.71798 & 0.61323 & 0.66511 & 0.21246 \\
    \runid{run-20260712-171008-fc3e0-burn-deep-classifier} & mlp / [256,128]     & 0.71576 & 0.61220 & \best{0.66638} & 0.21249 \\
    \runid{run-20260712-171008-c093f-burn-deep-classifier} & mlp / [128,64]      & 0.71229 & 0.61694 & 0.65852 & 0.21447 \\
    \runid{run-20260712-171008-786b4-burn-deep-classifier} & embeddings / [256,128,64] & 0.70750 & 0.61671 & 0.66001 & 0.21444 \\
    \bottomrule
  \end{tabular}}
\end{table}

\begin{figure}[ht]
  \centering
  \includegraphics[width=0.8\textwidth]{highvocab_capacity}
  \caption{Capacity is a spent lever for the dense net at high vocab. Holding lr
  at each topology's best and stepping the hidden geometry
  $[128,64]\to[256,128]\to[256,128,64]$, neither the \texttt{embeddings} (blue)
  nor the \texttt{mlp} (grey) ladder clears its own control by a band; a third
  \texttt{embeddings} layer \emph{hurts} by \num{0.0152} ($\sim$1.2 bands, the
  only out-of-band move), because the extra depth dilutes an already
  weaker-conditioned learned entity-embedding input. The best \texttt{mlp} closes
  to within $+\num{0.0047}$ of the best \texttt{embeddings}, so the two families
  tie within a band --- capacity, not the embeddings mechanism, was the lever ---
  and the whole ANN tier stays $\sim$0.2--0.6 band below the in-dataset tree
  tier (shaded, catboost \num{0.7252} up to xgboost \num{0.7299}).}
  \label{fig:highvocab-capacity}
\end{figure}

\emph{Interpretation.} The per-model SAE explains \emph{why} both moves fail.
The AUC-tied cross-family pair --- \texttt{embeddings} $[128,64]$ (\num{0.72274})
and \texttt{mlp} $[256,128,64]$ (\num{0.71800}) --- plateau at the \emph{same}
decodability ceiling (peak layer test $R^2_{\text{target}}$ \num{0.149} versus
\num{0.141}) and both give all 245 identity segments a dedicated separating
atom; the difference is layout only. The winner reaches the ceiling with few,
monotonically-built concepts ($33\to75$, $\Sigma108$), while the deeper/wider
\texttt{mlp} forms nearly $2\times$ the concepts ($55\to113\to32$, $\Sigma200$)
that peak in its middle layer and \emph{collapse} through its neck --- the same
consolidation-versus-proliferation pattern of Section~\ref{sec:anns-sae}, now
reproduced across families and invisible to AUC. And the $-\num{0.0152}$
deep-embeddings penalty is dilution, not starvation: that arm has 0 dead atoms
and full utilization but a weaker learned entity-embedding \emph{input}
($R^2_{\text{target}}$ \num{0.123} versus the winner's \num{0.133}), and its
extra depth never recovers it (decodability dips mid-trunk). This is the
mechanistic face of the capacity ceiling: a third layer over a learned embedding
table dilutes ranking rather than sharpening it. Both this scan and the GBDT
vocab-cap scan (Section~\ref{sec:vocab-cap}) point the same way --- identity
\emph{resolution quality}, not the raw capacity of the model or the vocabulary,
is the remaining lever. \emph{Source:
\texttt{2026-07-12-high-vocab-capacity-scan.md}.}

\section{Taste drift: the leaderboard is an upper bound on forward prediction}
\label{sec:drift}
\emph{Rationale.} Every AUC in this document is measured on a \emph{random}
train/test split, which lets a track discovered in a given period land on both
sides of the split. That answers ``can the model rank tracks from the same era
as its training data?'' --- but the deployed question is forward-looking:
\emph{will I replay a song I discover next?} This campaign asks it honestly by
adding a first-class chronological-split capability to the builder
(\texttt{split\_order\_field} on the build request): rows are ordered by a
payload field (here \texttt{first\_played}), the earliest go to train and the
latest to test, and PCA, vocabulary and imputation are all fit on the
early-train rows only, so the split is leak-free. Four datasets were built from
the active \texttt{spotify\_tracks\_song\_ae} collection --- full and intrinsic,
each random and chronological --- and the champion plus a model sweep were
evaluated on each dataset's own native split. If AUC collapses under the
chronological split, taste drifts and the leaderboard overstates real forward
skill.

\emph{Results.} \textbf{Taste is not stationary} (Table~\ref{tab:drift};
Figure~\ref{fig:taste-drift}). The full-random build reproduces the on-record
champion to five digits (\num{0.73490}), validating the fresh builds; the
chronological build of the \emph{same} features drops it to \num{0.61894} --- a
$+\num{0.116}$ collapse ($\approx$9 noise bands). The drop is \num{0.109}--%
\num{0.116} for \emph{every} model family (\texttt{catboost} \num{0.72889}$\to$%
\num{0.61294}, \texttt{logistic} \num{0.72141}$\to$\num{0.61234}), so it is the
data/taste, not a model quirk. Identity one-hots barely help forward:
\emph{dropping} them (the intrinsic build) costs only $\sim$0.02 chronological
(\num{0.61894}$\to$\num{0.59858}) versus $\sim$0.025 random --- they mostly
memorize same-era artists. A recency-weighted \texttt{xgboost} (up-weighting
recent discoveries) recovers only $+\num{0.003}$ (sub-band), so drift is not a
simple recency effect. And the walk-forward read --- train on
\texttt{first\_played}~$<Y$, test on year $Y$ --- lands at \num{0.59}--\num{0.67}
for every single year from 2019 to 2026, \emph{never} near the \num{0.735}
random baseline, and more training history does not close the gap. An
off-framework re-split of the pre-existing random-built features agrees to
$\sim$0.007 AUC (\num{0.626}), and controlling for right-censoring (restricting
the chronological test to tracks with $\ge$180 days of replay opportunity)
barely moves it (\num{0.62579}$\to$\num{0.62319}) --- so the penalty is genuine
non-stationarity, not the base-rate shift of mechanically-less-rotated recent
tracks.

\begin{table}[ht]
  \centering
  \caption{The random-versus-chronological AUC gap, definitive framework builds
  (all \texttt{spotify\_tracks\_song\_ae}, \num{23529} rows, seed 42; each model
  on its dataset's native split). ``full'' carries the identity one-hots,
  ``intrinsic'' drops them (64 PCA $+$ 7 numeric only). The champion row is
  bolded; the drift column is the random$-$chronological drop on the full
  matrix. Full-random reproduces the on-record champion (\num{0.73490}),
  validating the builds. Datasets: full-random \dsid{ds-20260712-183615},
  full-chrono \dsid{ds-20260712-183629}, intrinsic-random
  \dsid{ds-20260712-183630}, intrinsic-chrono \dsid{ds-20260712-183631}. AUC
  band $2\sigma\approx 0.0126$.}
  \label{tab:drift}
  \begin{tabular}{lccccc}
    \toprule
    Model & full $\cdot$ random & full $\cdot$ chrono & drift & intr.\ $\cdot$ random & intr.\ $\cdot$ chrono \\
    \midrule
    \best{xgboost} (champion) & \best{0.73490} & \best{0.61894} & \best{$+0.116$} & 0.71020 & 0.59858 \\
    catboost                  & 0.72889 & 0.61294 & $+0.116$ & 0.70172 & 0.59556 \\
    logistic                  & 0.72141 & 0.61234 & $+0.109$ & 0.64836 & 0.59016 \\
    \bottomrule
  \end{tabular}
\end{table}

\begin{figure}[ht]
  \centering
  \includegraphics[width=0.9\textwidth]{taste_drift}
  \caption{The leaderboard AUC is an upper bound on forward prediction.
  \emph{Left:} moving from a random to a chronological (train-earlier /
  test-later) split drops AUC by $\sim$0.11 for \emph{every} model family --- the
  penalty is the data, not any one model. \emph{Right:} the walk-forward read
  (train on \texttt{first\_played}~$<Y$, test on year $Y$) lands at
  \num{0.59}--\num{0.67} every year and never approaches the \num{0.735} random
  baseline (green), tracking the $\sim$\num{0.619} chronological level (red
  dashed) instead; more accumulated history does not close the gap. Forward
  ``will I replay a song I discover next'' skill on this corpus is
  $\sim$0.62--0.63, $\sim$0.11 below the random-split leaderboard.}
  \label{fig:taste-drift}
\end{figure}

\emph{Interpretation.} The mechanism is identity churn, and it ties the whole
document together. Rotation is driven by artist/genre \emph{identity} one-hots
(the convergent finding of every campaign above); in a chronological split the
test period is full of newly discovered artists that were unseen in training and
fall into \texttt{\_\_other\_\_}, so the model must fall back on the weaker
intrinsic acoustic/metadata signal --- which tops out near \num{0.60}. That is
exactly why raising vocab caps (Section~\ref{sec:vocab-cap}) and adding net
capacity (Section~\ref{sec:anns-capacity}) are spent levers for forward skill:
they only sharpen same-era \emph{memorization}, which does not transfer. It is
also the same-signal story the per-model SAEs tell --- linear decodability
$\sim$0.14 with no clean monosemantic taste atom (Sections~\ref{sec:anns-sae},
\ref{sec:anns-capacity}): remove identity memorization and little generalizable
taste signal remains. The practical reading: every on-record AUC is an upper
bound on same-era ranking, not forward skill; a deployed model needs periodic
retraining because taste drifts; and the only lever that would move
\emph{forward} AUC is richer generalizable intrinsic signal (better acoustics,
cross-artist structure), since identity resolution and model capacity are both
spent and same-era-only. This does not overturn any leaderboard number --- the
random-split AUCs are valid measurements of same-era ranking --- but it reframes
what they mean. \emph{Source:
\texttt{2026-07-12-taste-drift-temporal-split.md}.}

\section{Next-track sequence recommendation: the task, the Markov bar, and the fusion champion}
\label{sec:nexttrack}
Every campaign above targets the binary \emph{rotation} signal. The remaining
sections fold in a \emph{separate} line of work --- \emph{next-track} sequence
recommendation --- whose task, metrics and champion are unrelated to the
rotation leaderboards and must never be read against them. This section
consolidates the four campaigns that established the task, its baseline bar and
its fusion champion; the two sections that follow
(Sections~\ref{sec:nexttrack-compression} and~\ref{sec:nexttrack-phase2})
develop the item-representation and model-family axes and crown the current
best-on-record row.

\paragraph{Task and split.} The task is next-item \emph{retrieval}: given a
listening \emph{session} (an ordered sequence of tracks), rank the catalog and
recover the next distinct track the user actually played. It is scored on a
single leak-free chronological \emph{next-distinct} split of the sequence
artifact \dsid{seq-20260713-014226} (\num{7154} sessions / \num{19402} items;
\num{5723} train / \num{1431} test at cut 2024-08-24). The \textbf{cold-item
rate is 0.55} --- more than half of test-session targets never appear in any
training session --- so absolute recall is modest \emph{by construction} and
every result is read against the incumbent baseline, not an absolute bar.

\paragraph{Metrics.} The headline metric is \textbf{Recall@10} (R@10, the
fraction of test sessions whose held-out next track appears in the top ten
retrieved), reported beside mean reciprocal rank (MRR) and, since 2026-07-14,
the first-class \emph{graded-relevance} diagnostics: \emph{artist@10} /
\emph{genre@10} (a top-ten candidate shares the truth's artist / primary genre)
and \emph{artist\_mrr}. All confidence intervals here are bootstrap 95\% CIs
over \num{2000} resamples of the \num{1431} test sessions, and the decisive
statistic is the \emph{paired} per-session $\Delta$ against a named reference
(the Markov bar, a same-space GRU anchor, or the champion). The graded metrics
are the honest scorecard: the champion lands the right \emph{artist} in the
top ten $\sim$32\% of the time and the right \emph{genre} $\sim$40\%, far more
useful than the $\sim$17\% exact-recall number the cold-item rate caps. Since
2026-07-20 the board also carries \textbf{music@10}, a continuous
``sounds-alike'' graded metric (Section~\ref{sec:nexttrack-music}): the mean best
cosine, in a balanced intrinsic musical-distance space, between the true next
track and the top-ten candidates --- a smooth superset of Recall@10, display-only
(the primary metric stays Recall@10). Since 2026-07-22 the board also carries a
display-only \emph{session-holisticness} suite (artist\_adj@10 / mood\_coh@10 /
ild@10; Section~\ref{sec:nexttrack-holisticness}), which does not reorder the board
but fingerprints the champion as the most artist-eager row.

\paragraph{The Markov bar.} The reference every arm is measured against is the
first-order \emph{Markov} co-occurrence model (\texttt{B2}: a train-only
\texttt{track\_bigram} with artist/genre back-off, the transition term the
sibling app already ships) at R@10 \num{0.1069}, MRR \num{0.0694}. Popularity
(\texttt{B0}) and recency (\texttt{B1}) floors sit near zero (R@10
\num{0.0007}). A board \emph{win} requires beating the reference on
\emph{both} R@10 and MRR beyond the paired-$\Delta$ / marginal-CI tests, then a
reproduction confirm.

Table~\ref{tab:nexttrack-lb} is the standing next-track best-on-record board.
The current champion \emph{row} --- and, since 2026-07-18, a \emph{registered,
promoted and servable} model --- is a supervised \emph{learned content
projection}, \emph{\Rp+M+\Cp}, on the PCA-192 item space
(Sections~\ref{sec:nexttrack-projection}--\ref{sec:nexttrack-projection-scan}):
the GRU and content-kNN legs retrieve in an InfoNCE metric map of the item space
fit toward next-track adjacency. It supersedes the three-leg \emph{R+M+C} z-blend
on the same space (Section~\ref{sec:nexttrack-modelsweep}), which had dethroned
the two-leg \emph{R+M} z-blend on the AE-64 space
(Section~\ref{sec:nexttrack-blend}). Since 2026-07-18 the sequence family carries
a ranking serving contract, so the champion is servable via
\texttt{POST /api/models/blend-gru-markov-content-proj/predict} (best-models rank
1); the earlier blends below it, and the leg standalones, remain off-server
references. Since 2026-07-25 the board also carries three \emph{refuted}
architecture datapoints --- the widest single-GRU capacity control, the best
dual-tower arm and the best single-tower pre-encoder arm
(Sections~\ref{sec:nexttrack-dualtower}--\ref{sec:nexttrack-preencoder}). They
are listed because they \emph{bound} their axes, not because they are
candidates: none beats the single-GRU baseline by the paired statistic and none
was registered. Since 2026-07-27c the board's top row is the \emph{evidence-gated}
champion (Section~\ref{sec:nexttrack-markovgate}): making the Markov leg abstain on
the \SI{55.6}{\percent} of queries where it has no bigram evidence lifts Recall@10
to \num{0.2320} (\num{332} of \num{1431}), paired CI above zero on two disjoint
splits, with MRR and both graded metrics up. It is the best recall figure in this
line of work, and it is a promotion \emph{recommendation} only --- nothing was
registered or promoted --- so the servable champion below it is unchanged.
On the server-effective primary metric, and on why this board
stays denominated in Recall@10 across three holisticness generations, see
Section~\ref{sec:nexttrack-power-generations}.

\begin{table}[ht]
  \centering
  \caption{Next-track best-on-record board (RANKING regime), on the leak-free
  next-distinct split (\num{1431} test sessions; bootstrap 95\% CI over
  \num{2000} resamples). R@10, MRR, artist@10, genre@10, music@10 higher-better;
  the Markov bar is the incumbent reference. music@10 values are the offline
  backfill (Section~\ref{sec:nexttrack-music}); rows without a mapped run are
  left blank. Champion row and prior champions bolded; best cell per metric
  bolded. The three refuted architecture rows added 2026-07-25/26 are shown to
  \emph{bound} their axes and are excluded from the best-cell bolding; they carry
  no marginal CI, only a paired $\Delta$ against the single-GRU baseline
  (Tables~\ref{tab:nexttrack-dualtower} and~\ref{tab:nexttrack-preencoder}).
  The top row, added 2026-07-27c, is the evidence-gated champion of
  Section~\ref{sec:nexttrack-markovgate}: a promotion \emph{recommendation} rather
  than a registered model, so its R@10 is bolded as the best on record but it is
  otherwise excluded from the best-cell bolding, and it carries a paired $\Delta$
  (Table~\ref{tab:nexttrack-markovgate-paired}) rather than a marginal CI.
  \textbf{Not comparable to any rotation board above.} Item spaces: AE-64 (\dsid{spotify\_tracks\_song\_ae}); PCA-$k$
  (\dsid{spotify\_tracks\_song\_pca$k$}), all fit from the leak-safe
  \dsid{spotify\_tracks\_content} source.}
  \label{tab:nexttrack-lb}
  {\footnotesize\setlength{\tabcolsep}{4pt}\begin{tabular}{lcccccl}
    \toprule
    Model & R@10 [95\% CI] & MRR & art@10 & gen@10 & mus@10 & Note \\
    \midrule
    \best{\Rp+M+\Cp\ $+$ gate} (best recall) & \best{0.2320} & 0.1335 & 0.367 & 0.463 & 0.493 & not promoted \\
    \best{\Rp+M+\Cp, PCA-192} (champion, servable) & \best{0.2117 [0.191, 0.233]} & \best{0.1220} & 0.333 & 0.401 & 0.454 & new champion \\
    R+M+C, PCA-192 (prior fusion champ.)  & 0.1859 [0.166, 0.207] & 0.1025 & 0.331 & 0.394 & 0.449 & content additive \\
    R+M+C, PCA-128                        & 0.1824 [0.162, 0.203] & 0.0993 & 0.332 & 0.396 & --- & content additive \\
    \best{R+M, AE-64} (prior champion)    & 0.1726 [0.154, 0.192] & 0.0963 & 0.322 & 0.397 & 0.441 & prod.\ anchor; 3-seed \\
    R+M+C, AE-64 (graded-preferred)       & 0.1740 [0.154, 0.194] & 0.1014 & \best{0.330} & 0.400 & --- & content only tied \\
    R+M, PCA-192                          & 0.1712 [0.151, 0.192] & 0.0952 & 0.329 & 0.398 & 0.443 & $\approx$ prior champion \\
    R+M, PCA-128                          & 0.1670 [0.148, 0.187] & 0.0920 & 0.329 & 0.398 & 0.444 & GRU gain not carried \\
    R+A+M, AE-64                          & 0.1684 & 0.0943 & --- & --- & --- & ANN leg redundant \\
    GRU single h512, PCA-192          & 0.1300 & --- & --- & --- & 0.458 & width tie \#7; not sig.\ \\
    tower bank lat$\times$4, PCA-192   & 0.1300 & --- & --- & --- & 0.459 & bank $=$ width; closed \\
    LSTM single, PCA-128                  & 0.1272 [0.110, 0.145] & 0.0495 & 0.335 & \best{0.458} & 0.451 & best single; $\ll$ blend \\
    GRU single h454, PCA-192          & 0.1251 & 0.0539 & 0.338 & 0.454 & 0.457 & width closed \\
    GRU single, PCA-192                   & 0.1230 [0.106, 0.140] & 0.0534 & \best{0.338} & 0.450 & 0.456 & dim point-optimum \\
    dual-tower GRU f0, PCA-192        & 0.1188 & 0.0563 & 0.337 & 0.456 & 0.455 & dual tower: refuted \\
    GRU single, PCA-128 (base)            & 0.1146 [0.099, 0.132] & 0.0489 & 0.333 & 0.452 & \best{0.457} & confirmed base space \\
    pre-MLP$\to$GRU, PCA-192          & 0.1090 & 0.0446 & 0.324 & 0.469 & 0.450 & pre-encoder harmful \\
    B2 first-order Markov (bar)           & 0.1069 [0.092, 0.124] & 0.0694 & 0.298 & 0.389 & 0.407 & app incumbent \\
    ANN pooled-MLP, PCA-128               & 0.0901 [0.076, 0.105] & 0.0394 & 0.299 & 0.435 & 0.426 & loses to the GRU \\
    XGB stacker (M+G / M+G+A)             & 0.080--0.087 & 0.042--0.049 & 0.31 & 0.38 & 0.41 & combiner underperforms \\
    content-kNN (song-AE)                 & 0.0636 [0.050, 0.077] & 0.0321 & 0.272 & 0.405 & --- & the additive C leg \\
    item2vec (train-only SGNS)            & 0.0517 & 0.0214 & 0.166 & 0.312 & --- & drags the blend \\
    \bottomrule
  \end{tabular}}
\end{table}

\subsection{The sequence sweep: a GRU-infonce ties the Markov bar}
\label{sec:nexttrack-sweep}
\emph{Rationale.} The instance's opening campaign asked the founding question:
does a session-context sequence model --- a GRU/LSTM over the ordered prefix,
predicting the next track's 64-dim song-AE latent and retrieving by cosine ---
beat the first-order Markov transition model the sibling app already ships,
\emph{forward}, on the leak-free session split? Four arms crossed the recurrent
cell (GRU, LSTM), capacity (\texttt{hidden} 128, 256) and, decisively, the
training loss (cosine regression versus \texttt{InfoNCE} contrastive), all on
\dsid{seq-20260713-014226} at 60 epochs.

\emph{Results.} No arm beat the bar. The decision rule required beating
\texttt{B2} on \emph{both} R@10 and MRR beyond non-overlapping CIs; the best arm
(M3, GRU \texttt{hidden}128 \texttt{InfoNCE}) reached a statistical \emph{tie}
with the Markov bar on R@10 (\num{0.0894}, paired-$\Delta$ CI includes zero) but
\emph{lost} on MRR (\num{0.0388} versus \num{0.0694}) and Recall@20. The
dominant lever was the loss objective: \texttt{InfoNCE} was $\approx 2.3\times$
cosine on R@10 (\num{0.0894} versus \num{0.0391}), while capacity (\texttt{h256})
and the LSTM cell barely moved the cosine arms (all $\sim$0.038--0.041).

\emph{Interpretation.} A latent-plus-single-cosine-hop is a harder route to the
\emph{exact} next track than a co-listening bigram: the bigram directly
memorizes ``what follows what,'' while the GRU must land near the right point in
a 64-d content space and then retrieve, losing on precise ranking (MRR). The
finding is consistent with the parent lab's taste-drift result --- \num{0.55}
cold items, identity/adjacency-driven behavior --- and it set the agenda:
\texttt{InfoNCE} is the base loss, and the natural next arm is a \emph{blend} of
the GRU's candidate scores with the Markov bigram.
\emph{Source: \texttt{2026-07-13-nexttrack-sequence-sweep.md}.}

\subsection{Topology scan: depth, width and direction are a spent lever}
\label{sec:nexttrack-topology}
\emph{Rationale.} With the loss fixed at \texttt{InfoNCE}, this scan isolated
network \emph{topology}: does depth (stacked layers), width (\texttt{hidden}
256), an inter-layer residual stack or bidirectionality lift the sequence model
over the Markov bar? The \texttt{seq-nexttrack} predictor gained
\texttt{num\_layers}, \texttt{residual} and \texttt{bidirectional}; because a
bidirectional RNN under per-step teacher forcing would leak future items, the
bidirectional arm used a leak-free last-item objective (encode the full prefix,
predict only the held-out next item).

\emph{Results.} Refuted. No arm beat the bar on both metrics. The best arm
(width \texttt{h256}, R@10 \num{0.0957}) is a \emph{tie} with the bar on R@10
and loses on MRR (\num{0.0411} versus \num{0.0694}). Depth is flat 1$\to$2
layers (\num{0.0894}$\to$\num{0.0922}, in-band) and \emph{hurts} at 3 layers
(\num{0.0825}, out-of-band below the bar); the residual stack over depth-2/h256
added nothing. The bidirectional arm collapsed (\num{0.0482}), but this is an
\emph{objective} artifact, not a directionality verdict: the leak-free last-item
objective supplies $\sim$1 training target per session versus teacher forcing's
$\sim$13 next-step pairs, so it trained on $\sim$13$\times$ less signal.

\emph{Interpretation.} Architecture is not the lever for next-track on this
corpus. The transferable latent-retrieval signal caps at R@10 $\sim$0.09--0.10
across depth, width, residual and direction, below the co-listening bigram ---
consistent with the whole arc that the loss objective moved the needle and
capacity/topology do not. The one remaining unexplored lever was the candidate
framing / fusion, which the next campaign pulled decisively.
\emph{Source: \texttt{2026-07-13-nexttrack-topology-scan.md}.}

\subsection{The GRU$\times$Markov fusion champion, and why more ensembling fails}
\label{sec:nexttrack-blend}
\emph{Rationale.} The sweep and the topology scan agreed that the best single
GRU only \emph{ties} the Markov bar on R@10 while losing on MRR. This campaign
asked whether \emph{combining} models beats the bar, and how far ensembling
goes. Three legs were defined, each emitting full-vocab candidate scores:
\textbf{M} = the train-only Markov bigram (the bar), \textbf{R} = the
GRU-infonce-h256 recurrence, and \textbf{A} = an ANN bag-of-prefix MLP (a
feed-forward twin of R). They were combined by a weighted z-normalized blend
($\alpha\cdot z(\text{leg})+\dots$) and, as a learned alternative, an XGB
stacker (a GBDT meta-learner over per-candidate features).

\emph{Results.} The \textbf{R+M z-blend} ($\alpha{=}0.5$) \emph{beats the bar
decisively}: R@10 \num{0.1726} $[0.154, 0.192]$, MRR \num{0.0963}, a paired
$\Delta$ over Markov of $+\num{0.0657}$ $[+0.049, +0.083]$ (CI entirely above
zero) --- $\sim$$+61\%$ R@10, $+39\%$ MRR --- on a robust plateau across
$\alpha \in [0.25, 0.75]$ and reproduced across three seeds (\num{0.1705} /
\num{0.1726} / \num{0.1726}). It was registered as the \texttt{seq-blend}
predictor and the \texttt{blend-gru-markov} definition, the first next-track
champion. Further ensembling then \emph{failed} to improve on it: R+A+M
(\num{0.1684}) is slightly worse, R+A ($\approx$ GRU alone, \num{0.0894}) is
much worse, and both XGB stacker variants (M+R \num{0.1020}; M+R+A \num{0.0887})
barely clear Markov alone.

\begin{table}[ht]
  \centering
  \caption{The fusion campaign: the R+M z-blend beats the Markov bar
  decisively, and no richer ensemble improves on it. R@10 and MRR
  higher-better; champion bolded. \emph{Not comparable to any rotation board.}}
  \label{tab:nexttrack-blend}
  {\small\setlength{\tabcolsep}{5pt}\begin{tabular}{lccl}
    \toprule
    Model & R@10 [95\% CI] & MRR & Note \\
    \midrule
    \best{R+M z-blend $\alpha{=}0.5$} (champion) & \best{0.1726 [0.154, 0.192]} & \best{0.0963} & $\Delta$ vs Markov $+0.066$ [$+0.05,+0.08$] \\
    R+A+M z-blend                                & 0.1684 & 0.0943 & ANN leg redundant with GRU \\
    B2 Markov bar (M)                            & 0.1069 [0.092, 0.123] & 0.0694 & the reference \\
    XGB stacker M+R                              & 0.1020 & 0.0561 & learned combiner $\ll$ z-blend \\
    GRU alone (R)                                & 0.0957 & 0.0411 & ties the bar on R@10 only \\
    R+A z-blend                                  & 0.0894 & 0.0377 & $\approx$ GRU alone \\
    XGB stacker M+R+A                            & 0.0887 & 0.0572 & barely clears Markov \\
    \bottomrule
  \end{tabular}}
\end{table}

\emph{Interpretation.} The win is about signal \emph{diversity}, not model
count. The GRU (latent similarity) and Markov (co-occurrence counts) rank
\emph{different} correct next-tracks, so fusing them roughly doubles the hit
rate; piling on more of the \emph{same} kind of signal --- another neural
latent predictor (A), redundant with R --- adds nothing. The learned XGB
stacker \emph{underperforms} the simple z-blend, counter to the usual
``stacking beats blending'' intuition: its pointwise binary-with-negatives
training does not optimize the ranking objective as well as directly fusing
normalized scores. The remaining headroom is bounded by the \num{0.55}
cold-item rate and requires \emph{genuinely new signal}, not more model fusion
--- the hypothesis the graded-relevance campaign tests next.
\emph{Source: \texttt{2026-07-14-nexttrack-blend-and-ensemble.md}.}

\subsection{Graded-relevance metrics and new base learners: content is the only additive signal}
\label{sec:nexttrack-graded}
\emph{Rationale.} Two threads in one campaign. (A) Make the evaluation credit
landing the right \emph{band/vibe} even when the exact track is wrong --- the
observation that ``the first choice is usually not the exact track, but the
artist is usually right'' --- by promoting \emph{artist@k}, \emph{genre@k} and
\emph{artist\_mrr} to first-class metrics wired end-to-end
(\texttt{seq\_common.eval\_from\_scores} $\to$ Rust \texttt{Metrics} /
\texttt{MetricsSpec} $\to$ \texttt{domain.toml} $\to$ UI). (B) Try genuinely new
leak-free base learners --- \textbf{C} = content-kNN (cosine of each candidate's
song-AE latent to the prefix, no training), \textbf{M2} = a second-order Markov
trigram with back-off, and \textbf{I} = a train-only item2vec SGNS --- scored on
both the exact and the new graded metrics.

\emph{Results.} The exact-metric path was byte-identical (purely additive
keys); the champion reproduced R@10 \num{0.1726} / MRR \num{0.0963} exactly. No
new base learner beat the bar standalone (M2 \num{0.1076} $\approx$ M
\num{0.1069}; the trigram falls through to the bigram --- sessions are too short
to fire; item2vec \num{0.0517} is the weakest leg). Blended into the champion,
\textbf{content is the only additive signal}: \textbf{R+M+C} (adding the content
leg) is the top blend on MRR (\num{0.0963}$\to$\num{0.1014}), artist@10
(\num{0.3215}$\to$\num{0.3298}) and artist\_mrr, and ties the champion on exact
R@10 (paired $\Delta$ $+\num{0.0014}$, CI straddles zero). Item2vec \emph{drags}
the blend (R+M+I \num{0.1349}); the trigram is neutral.

\emph{Interpretation.} Content-kNN ranks \emph{different} near-correct tracks
(latent vibe similarity) than the GRU (session recurrence) and Markov
(co-occurrence), exactly the diversity the fusion campaign predicted new signal
would need. On the AE-64 space the content leg only \emph{tied} the champion on
exact recall while lifting the graded metrics --- so it was carried forward as
the graded-preferred third leg, not a new exact champion. That ``tie'' becomes a
\emph{win} once the item space is richer (Section~\ref{sec:nexttrack-phase2}):
on PCA-192 the same R+M+C construction is the strongest on record. Co-listening
item2vec and the trigram Markov added nothing on this data.
\emph{Source: \texttt{2026-07-14-nexttrack-graded-relevance-and-new-legs.md}.}

\section{Next-track sequence recommendation: the item-representation compression bake-off}
\label{sec:nexttrack-compression}
Every campaign above targets the binary \emph{rotation} signal. This section
folds in the first report from a \emph{separate} line of work --- \emph{next-track}
sequence recommendation --- whose task, metrics and champion are unrelated to the
rotation leaderboards and must never be read against them. The task is
retrieval: given a listening \emph{session} (an ordered sequence of tracks),
rank the catalog and recover the next distinct track the user actually played.
It is scored on a leak-free chronological \emph{next-distinct} split (\num{7154}
sessions / \num{19402} items; \num{5723} train / \num{1431} test at cut
2024-08-24; cold-item rate 0.55), and its headline metric is \textbf{Recall@10}
(R@10, the fraction of test sessions whose held-out next track appears in the
top ten retrieved), reported beside mean reciprocal rank (MRR) and the
category-hit diagnostics artist@10 / genre@10 / artist\_mrr. All confidence
intervals here are bootstrap 95\% CIs over \num{2000} resamples of the
\num{1431} test sessions. The standing champion of this track --- the reference
this campaign is read against --- is a rank-fusion \emph{R+M z-blend}
(\texttt{blend-gru-markov}, $\alpha{=}0.5$) of a sequence GRU with a first-order
Markov co-occurrence model on the 64-dim multimodal song-autoencoder space
(\texttt{spotify\_tracks\_song\_ae}, ``AE-64''), at R@10 \num{0.173}
$[0.154, 0.192]$, MRR \num{0.096}, three-seed confirmed.

\emph{Rationale.} The champion's item space --- a 64-dim AE latent --- was one
arbitrary early choice, and the sequence GRU predicts \emph{into} and retrieves
\emph{over} whatever geometry that space imposes. This campaign asks which
\emph{compression} is the best space for the GRU: linear PCA versus the
nonlinear autoencoder, each at latent dim 32 / 64 / 128. Five alternative spaces
were baked from the \emph{same} leak-safe 384-dim content source
(\texttt{spotify\_tracks\_content}), each fit unsupervised over all \num{23529}
items using only intrinsic content/acoustic/catalog fields (text-384 $+$
z-scored numerics $+$ \texttt{af\_*} acoustics $+$ genre/album-type multi-hot),
never a behavioral or target field. Each space was then scored under a
\emph{frozen} GRU readout (\texttt{gru-infonce-h256}: \texttt{arch=gru},
\texttt{loss=infonce}, \texttt{hidden=256}, one layer, dropout 0.1, 40 epochs,
lr \num{1e-3}, batch 128, $k{=}10$), so the item latent space is the only axis
that moves. The design bet was that the AE's equal-weight-per-block
reconstruction loss would \emph{preserve} the small acoustic/catalog blocks that
PCA's text-dominated global variance would starve; the null was that dim, not
the compressor, sets the ceiling. Every artifact passed an identical parity gate
(same \num{7154} sessions / \num{19402} items / split / cut), so the comparison
is apples-to-apples, and the paired-$\Delta$ is the per-session R@10 difference
against the in-batch AE-64 control.

\emph{Results.} The parity gate passed: the AE-64 control reproduced its
on-record R@10 \num{0.0957} $\in [0.081, 0.112]$, so the ranking is valid. All
six spaces trained cleanly (Table~\ref{tab:nexttrack-compression};
Figure~\ref{fig:nexttrack-compression}). \textbf{PCA-128}
(\dsid{spotify\_tracks\_song\_pca128}, artifact \texttt{seq-20260715-030509};
run \runid{run-20260715-042100-a1d9d}) is the single-GRU winner at R@10
\num{0.1146}, a paired-$\Delta$ of $+\num{0.0189}$ $[+0.0042, +0.0335]$ over the
AE-64 control --- a CI \emph{entirely above zero}, the only arm that clears the
control significantly --- and it also tops MRR (\num{0.0489}), artist@10
(\num{0.3326}) and artist\_mrr (\num{0.2644}). PCA-64 is the runner-up
(\num{0.1041}, paired-$\Delta$ straddling zero: a tie). The two 32-dim spaces
(AE-32 \num{0.0545}, PCA-32 \num{0.0412}) sit far below the 64-dim control, both
with paired-$\Delta$ CIs entirely below zero. The training-free content-kNN
probe (a pure cosine retrieval in each space, last column of the report) is flat
at $\sim$0.06 across all six spaces and does \emph{not} track the GRU ranking:
PCA-128's probe R@10 (\num{0.0636}) exactly ties the AE-64 control
(\num{0.0636}), and AE-32 tops the probe (\num{0.0720}) despite ranking
near-worst under the GRU.

\begin{table}[ht]
  \centering
  \caption{The item-representation compression bake-off: six item spaces (linear
  PCA vs the nonlinear AE $\times$ latent dim 32 / 64 / 128), each scored under
  the \emph{same} frozen \texttt{gru-infonce-h256} readout on the next-distinct
  split (\num{1431} test sessions; bootstrap 95\% CI over \num{2000} resamples;
  paired-$\Delta$ vs the in-batch AE-64 control). All six succeeded. R@10, MRR,
  artist@10 and genre@10 higher-better. Winner (PCA-128) bolded; best cell per
  metric bolded. \emph{Not comparable to any rotation board.}}
  \label{tab:nexttrack-compression}
  {\small\begin{tabular}{llccccc}
    \toprule
    Space & Run id & R@10 [95\% CI] & paired-$\Delta$ vs AE-64 & MRR & artist@10 & genre@10 \\
    \midrule
    \best{PCA-128} (winner) & \best{\runid{run-20260715-042100-a1d9d}} & \best{0.1146 [0.099, 0.132]} & \best{$+0.0189$ [$+0.004$, $+0.034$]} & \best{0.0489} & \best{0.3326} & 0.4521 \\
    PCA-64 (runner-up)      & \runid{run-20260715-042100-8851d} & 0.1041 [0.088, 0.121] & $+0.0084$ [$-0.005$, $+0.022$] (tie) & 0.0416 & 0.3117 & 0.4598 \\
    AE-64 (control)         & \runid{run-20260715-042059-dd261} & 0.0957 [0.081, 0.111] & --- (control) & 0.0411 & 0.2942 & 0.4500 \\
    AE-128                  & \runid{run-20260715-042059-e96ae} & 0.0929 [0.078, 0.108] & $-0.0028$ [$-0.015$, $+0.009$] (tie) & 0.0398 & 0.3026 & \best{0.4640} \\
    AE-32                   & \runid{run-20260715-042059-d3f9e} & 0.0545 [0.043, 0.066] & $-0.0412$ [$-0.055$, $-0.027$] & 0.0274 & 0.2376 & 0.3955 \\
    PCA-32                  & \runid{run-20260715-042059-b1242} & 0.0412 [0.031, 0.052] & $-0.0545$ [$-0.069$, $-0.041$] & 0.0224 & 0.2187 & 0.3816 \\
    \bottomrule
  \end{tabular}}
\end{table}

\begin{figure}[ht]
  \centering
  \includegraphics[width=0.8\textwidth]{nexttrack_compression}
  \caption{Latent dimension, not the compressor, is the lever. Both series climb
  steeply from dim 32 to 64; PCA (blue) then keeps climbing to 128 and
  \emph{overtakes} the AE, while the AE (grey) plateaus --- its 128-dim point
  (\num{0.0929}) actually dips below its own 64-dim control (\num{0.0957}). At
  dim 32 the AE edges PCA, a low-capacity crossover that reverses once dim is
  generous. The AE-64 control (red ring) is the champion's item space; the
  PCA-128 winner (green star) is the only arm whose paired-$\Delta$ CI clears
  zero. Error bars are bootstrap 95\% CIs; the neighbouring CIs overlap, so the
  win rests on the \emph{paired} test, not the marginal intervals.}
  \label{fig:nexttrack-compression}
\end{figure}

The single-GRU win, however, does \textbf{not} survive into the production
blend. Re-fusing the winning PCA-128 GRU leg with the Markov leg
(\texttt{seq-blend}, $\alpha{=}0.5$, seed 1337;
run \runid{run-20260715-044920-36dce}) scores R@10 \num{0.1670}
$[0.148, 0.187]$, MRR \num{0.0923} --- \emph{below} the AE-64 champion blend
(\num{0.173}, MRR \num{0.096}) on both metrics, with overlapping CIs. By the
decision rule this is not a board event: nothing was promoted, the AE-64
champion blend stands, and PCA-128 is carried forward only as the sanctioned
Phase-2 base space.

\emph{Interpretation.} Four things follow. First, \textbf{latent dimensionality
is the dominant lever and the compressor is secondary}: for PCA the trend is
monotone (\num{0.041} $\to$ \num{0.104} $\to$ \num{0.115} across dim 32/64/128),
and 32 dims starve the space regardless of method. Second, the design's
\textbf{acoustic-preservation hypothesis is refuted}: at generous dim the linear
PCA beats the nonlinear AE (PCA-128 $>$ AE-128, PCA-64 $\ge$ AE-64), and although
AE-32 preserves the numeric/acoustic blocks near-perfectly it ranks near the
bottom --- so acoustic preservation is irrelevant to next-track retrieval. What
helps is a \emph{variance-ordered, decorrelated} target geometry (PCA-128
captures cumulative EVR \num{0.982}) that the GRU finds easier to regress into
than the AE's entangled equal-weight latents; PCA's global-variance ordering is
a feature here, not the liability the hypothesis assumed. Third, the AE's edge
is a purely \emph{low-capacity} phenomenon (AE-32 $>$ PCA-32): the nonlinear
bottleneck packs more usable structure into 32 dims, but the advantage
evaporates and reverses once dim is generous. Fourth, the flat content-kNN probe
that fails to track the GRU ranking is a \emph{flagged divergence}: the
intrinsic retrieval richness of these spaces is similar, so PCA-128's gain comes
from the \emph{learnability} of its target geometry for the trained GRU, not
from a demonstrably richer space in the absolute --- which is exactly why the
decision rule flagged rather than auto-promoted it. The load-bearing caution for
downstream work is the blend result: a clearly stronger GRU leg (\num{0.1146}
vs \num{0.096}) bought \emph{nothing} at the blend, because the Markov leg is
representation-independent and the extra tracks the better GRU recovers evidently
overlap what Markov already supplies. \textbf{Representation gains measured on
the single GRU do not automatically translate to the production blend}, so the
compression axis is genuine progress for the single model but essentially spent
for the blend. \emph{Source:
\texttt{2026-07-15-nexttrack-compression-representation-bakeoff.md}.}


\section{Next-track Phase 1b/2: the dimension ceiling, the model-family sweep, and the R+M+C champion}
\label{sec:nexttrack-phase2}
The compression bake-off (Section~\ref{sec:nexttrack-compression}) found latent
\emph{dimension} to be the dominant lever and left PCA still climbing at 128.
Two follow-on campaigns located the dimension ceiling, then swept the model
family on the chosen spaces --- and, in the process, crowned a new
best-on-record row.

\subsection{The compression dimension ceiling: PCA peaks at 192}
\label{sec:nexttrack-dim}
\emph{Rationale.} With dimension established as the lever, this leg extended the
PCA sweep upward (PCA-192, PCA-256) to locate the real optimum before Phase 2
committed a base space, and added AE-192 to test whether PCA $\ge$ AE holds at
higher dimension. The protocol was identical: all reps fit from the leak-safe
\dsid{spotify\_tracks\_content} source, scored under the \emph{same} frozen
\texttt{gru-infonce-h256} readout, each new sequence artifact passing the exact
parity gate; the reference is PCA-128 (\runid{run-20260715-042100-a1d9d}), the
Phase-1 winner, and the paired per-session $\Delta$ is measured against it.

\emph{Results.} The PCA R@10-versus-dimension curve \textbf{peaks at
$\sim$192} (R@10 \num{0.1230}) and \emph{turns over} by 256 (\num{0.1167}); it
was not still climbing. The full trend is
$32\!\to\!\num{0.041}$, $64\!\to\!\num{0.104}$, $128\!\to\!\num{0.115}$,
$\mathbf{192\!\to\!\num{0.123}}$, $256\!\to\!\num{0.117}$
(Figure~\ref{fig:nexttrack-dim}). PCA-192 is the point-optimum, but its paired
$\Delta$ versus PCA-128 is $+\num{0.0084}$ $[-0.0007, +0.0189]$ --- it misses
the ``CI entirely above zero'' bar by \num{0.0007}, a whisker --- so by the
mechanical rule \textbf{PCA-128 stands as the Phase-2 base}, with PCA-192
flagged as the confirm-and-promote candidate. PCA $\ge$ AE holds strongly and
the gap is widest at 192 (PCA-192 \num{0.123} versus AE-192 \num{0.098}); the AE
plateaus flat at $\sim$0.09--0.10 from dim 64 upward.

\begin{figure}[ht]
  \centering
  \includegraphics[width=0.8\textwidth]{nexttrack_dim_curve}
  \caption{The dimension ceiling is found. Extending the compression sweep of
  Figure~\ref{fig:nexttrack-compression} to PCA-192/256 and AE-192, the PCA
  curve (blue) peaks at $\sim$192 (R@10 \num{0.123}) and turns over by 256, while
  the AE (grey) plateaus flat at $\sim$0.09--0.10 and never catches it --- the
  PCA$\,\ge\,$AE gap is widest at 192. PCA-192 is the point-optimum (green star)
  but its paired $\Delta$ over PCA-128 grazes zero, so PCA-128 (the marked base)
  stays the Phase-2 single-model base. Error bars are bootstrap 95\% CIs over
  the \num{1431} test sessions.}
  \label{fig:nexttrack-dim}
\end{figure}

\emph{Interpretation.} Extra variance past $\sim$192 (PCA-256 captures
cumulative EVR \num{0.997}) is low-signal / noise dimensions the GRU cannot
exploit and that slightly dilute the target geometry. The AE's flat plateau
reconfirms the bake-off's refuted hypothesis at higher dimension: its
equal-weight reconstruction objective keeps spending capacity on the tiny
numeric/acoustic blocks (reconstruction $R^2 \approx 0.999$ at 192) rather than
the high-variance text/catalog structure the GRU needs. As before, the
training-free content-kNN probe stays flat ($\sim$0.064--0.070) and does not
reproduce the GRU's peak, so the dimension optimum is about target
\emph{learnability}, not intrinsic retrieval richness.
\emph{Source: \texttt{2026-07-15b-nexttrack-compression-dim-extension.md}.}

\subsection{The model-family sweep and the new champion, R+M+C on PCA-192}
\label{sec:nexttrack-modelsweep}
\emph{Rationale.} With the compression axis settled (PCA-128 base, PCA-192
candidate), this campaign swept the model \emph{family} on the chosen spaces:
does a different recurrent cell (LSTM), a non-recurrent pooled MLP (ANN), or a
\emph{learned} XGB stacker beat the fixed GRU$\times$Markov z-blend --- and does
adding the content-kNN third leg (R+M+C) on the \emph{rich} PCA space finally
beat the champion? A frozen held config ran across all arms; only the model
family and the representation varied, and the server was serialized
(\texttt{-{}-max-runs 1}) to avoid the thread-oversubscription trap.

\emph{Results.} A \textbf{new best-on-record row}: the \textbf{R+M+C z-blend on
PCA-192} (GRU $\times$ train-only Markov $\times$ content-kNN, equal thirds)
reaches R@10 \num{0.1859} $[0.166, 0.207]$, MRR \num{0.1025}, a paired $\Delta$
over the prior champion of $+\num{0.0133}$ $[+0.0007, +0.0259]$ --- entirely
above zero, with no MRR regression (Table~\ref{tab:nexttrack-phase2};
Figure~\ref{fig:nexttrack-phase2}). The content leg is genuinely additive on the
rich space, beating same-space R+M by $+\num{0.0154}$ (PCA-128) / $+\num{0.0147}$
(PCA-192), both CIs above zero --- \emph{stronger} than the AE-64 R+M+C that
only tied (Section~\ref{sec:nexttrack-graded}), because the richer, higher-EVR
space makes the content leg pay off. The learned \textbf{XGB stacker badly
underperforms} the fixed z-blend (\num{0.080}--\num{0.087}, losing even to its
own single GRU base leg), reconfirming the fusion campaign; the \textbf{LSTM is
the best single model} (PCA-128 \num{0.1272}, paired $\Delta$ over the GRU
$+\num{0.0126}$, CI above zero) but far below the blends; and the Phase-2 single
base \textbf{stays PCA-128} (GRU-192's edge still grazes zero).

\begin{table}[ht]
  \centering
  \caption{Phase-2 model-family sweep on the rich PCA spaces, ranked by R@10
  (higher-better). The fixed z-blend is the combiner of record; the learned XGB
  stacker underperforms it badly. Champion row and prior champion bolded; best
  cell per metric bolded. Paired $\Delta$ is per-session against the named
  reference. \emph{Not comparable to any rotation board.}}
  \label{tab:nexttrack-phase2}
  {\footnotesize\setlength{\tabcolsep}{4pt}\begin{tabular}{lp{2.4cm}ccp{2.9cm}}
    \toprule
    Model & Run id / source & R@10 [95\% CI] & MRR & paired $\Delta$ \\
    \midrule
    \best{R+M+C, PCA-192} (champ.\ row) & \runid{analyze\_rmc} (A4) & \best{0.1859 [0.166, 0.207]} & \best{0.1025} & vs champ $+0.0133$ [$+0.001,+0.026$] \\
    R+M+C, PCA-128                        & \runid{analyze\_rmc} (A2) & 0.1824 [0.162, 0.203] & 0.0993 & vs champ $+0.0098$ (tie) \\
    \best{R+M, AE-64} (prior champion)   & \runid{run-20260714-170928-645ae} & 0.1726 [0.154, 0.192] & 0.0963 & --- (reference) \\
    R+M, PCA-192                          & \runid{run-20260715-142440-96f53} & 0.1712 [0.151, 0.192] & 0.0952 & vs champ $-0.0014$ (tie) \\
    R+M, PCA-128                          & \runid{run-20260715-044920-36dce} & 0.1670 [0.148, 0.187] & 0.0920 & vs champ $-0.0056$ (tie) \\
    LSTM single, PCA-128                  & \runid{run-20260715-142440-546d0} & 0.1272 [0.110, 0.145] & 0.0495 & vs GRU-128 $+0.0126$ ($>0$) \\
    GRU single, PCA-192                   & \runid{run-20260715-142440-89d6a} & 0.1230 [0.106, 0.140] & 0.0534 & vs GRU-128 $+0.0084$ (tie) \\
    ANN pooled-MLP, PCA-128               & \runid{run-20260715-142440-7e528} & 0.0901 [0.076, 0.105] & 0.0394 & vs GRU-128 $-0.0245$ \\
    XGB stacker M+G, PCA-128              & \runid{run-20260715-142440-371b1} & 0.0867 [0.073, 0.101] & 0.0493 & vs GRU-128 $-0.0280$ \\
    XGB stacker M+G+A, PCA-192            & \runid{run-20260715-142440-f111d} & 0.0853 [0.071, 0.099] & 0.0490 & vs GRU-192 $-0.0377$ \\
    \bottomrule
  \end{tabular}}
\end{table}

\begin{figure}[ht]
  \centering
  \includegraphics[width=0.85\textwidth]{nexttrack_phase2_families}
  \caption{The Phase-2 model-family sweep. The fixed z-blend is the combiner of
  record: R+M+C on PCA-192 (green) is the new champion row, the content leg is
  genuinely additive on the rich space, and the learned XGB stacker (red) badly
  underperforms --- it loses even to its own single GRU base leg. LSTM edges GRU
  as the best single model (grey) but sits far below the blends. The first-order
  Markov bar (\num{0.107}) is the incumbent reference; the champion row is at
  \num{0.186}.}
  \label{fig:nexttrack-phase2}
\end{figure}

\emph{The crown, and its two honest caveats.} The board-event candidate was
confirmed by reproducing the GRU-192 leg five times (the original plus four
reruns) and recomputing R+M+C each time; the paired $\Delta$ over the prior
champion was $+\num{0.0133}$ $[+0.0007, +0.0259]$ \emph{every} rerun, with the
MRR lead intact, so R+M+C on PCA-192 was recorded as the new champion row. Two
caveats travel with the crown and must not be omitted. First, \textbf{the
reproduction spread was exactly zero} because \texttt{seq-nexttrack} training is
deterministic on this CPU --- all five GRU-192 checkpoints were byte-identical
--- so the confirm verified \emph{determinism, not seed-robustness}; the only
uncertainty quantified is the single-split bootstrap. Second, \textbf{the paired
$\Delta$ lower bound is $+\num{0.0007}$}: the crown clears zero by a whisker on a
single split. The crown is real by the pre-agreed criterion, but the margin over
the prior champion is thin, and hardening it needs a second leak-free test split
(or seed-varied legs once the blend is generalized to three legs).

\emph{Interpretation.} The compression-plus-model-family thread is otherwise
closed: PCA-128 is the single-model base, the LSTM is the best single model, the
learned stacker is dead as a combiner, and the fixed z-blend is the combiner of
record. The load-bearing result is that \emph{content is the one leg that adds
genuinely new ranking signal, and it pays off more on the higher-EVR PCA
space} --- the exception to the standing rule that single-model / representation
gains do not carry to the blend (the R+M pair on the rich PCA spaces only ties
the AE-64 champion). The productization gap remains: the champion is train-only
and off-server, so the prior champion R+M on AE-64 stays the deployable anchor
until the three-leg blend is registered with a \texttt{predict} subcommand.
\emph{Source: \texttt{2026-07-15c-nexttrack-phase2-model-sweep.md}.}

\section{Next-track: the learned content projection, registration, and two-split hardening}
\label{sec:nexttrack-projection}
The R+M+C champion (Section~\ref{sec:nexttrack-modelsweep}) carried two thin
caveats --- a determinism-only reproduction confirm and a paired-$\Delta$ lower
bound of only $+\num{0.0007}$ --- and its headroom was capped by the \num{0.55}
cold-item rate, which only \emph{new} signal can attack. Three late-July campaigns
then crowned, registered, and hardened a genuinely new champion. This section folds
in the first two (the win and its productization); the item-vector and music@10
threads follow in Sections~\ref{sec:nexttrack-representation} and~\ref{sec:nexttrack-music}.

\subsection{The literature-fit campaign: a supervised content projection wins}
\label{sec:nexttrack-projection-win}
\emph{Rationale.} Guided by a survey of the established session-recommendation
literature (the Jannach--Ludewig session-kNN line, sequential rules / factorized
Markov, GRU4Rec training tricks, rank fusion and learning-to-rank, content /
cold-start-aware session rec), this campaign pre-registered three tiers of classic
methods and asked which actually improve the fit to \emph{this} problem --- a
single-user, session-based, next-distinct \emph{retrieval} task with 55\% cold
targets. All arms were measured off-server on the PCA-192 champion artifact
\dsid{seq-20260715-131139}, reusing the sanctioned driver that reproduces the
champion (R+M+C \num{0.1859}) exactly; the decision statistic is the paired
per-session $\Delta$ against that champion.

\emph{Results --- the one win.} A supervised \textbf{learned content projection}
is a new best-on-record row. Fitting a linear metric map $W$ of the PCA-192
content space toward next-track adjacency (InfoNCE over train consecutive pairs,
self-loops dropped) and then retrieving \emph{both} the content leg (\Cp) and the
GRU (\Rp) in that projected space gives \textbf{\Rp+M+\Cp z-blend R@10 \num{0.2117}
$[0.191, 0.233]$, MRR \num{0.1220}} --- a paired $\Delta$ over the champion of
$+\num{0.0259}$ $[+0.012, +0.041]$, entirely above zero
(Table~\ref{tab:nexttrack-projection}). Unlike the R+M+C crown it is
\textbf{seed-robust} (three seeds; paired $\Delta$ $+\num{0.026}$ / $+\num{0.025}$ /
$+\num{0.029}$, every CI $>0$ --- a genuine seed-varied confirm, since the
projection has random init and shuffled batches) and \textbf{cold-reaching}: the
projected content leg \Cp\ scores cold targets \num{0.152} $\approx$ warm
\num{0.151}, versus Markov's warm \num{0.226} / cold \num{0.010}, so it reaches
never-seen items as well as seen ones --- the one channel that reaches the 55\% cold
majority. The symmetric single map beats the directional two-tower, and rank-192
$=$ rank-64 (the gain is the supervision, not the capacity). As a side effect the
projected single GRU \Rp\ ($\sim$\num{0.176}) nearly matches the \emph{old}
champion blend and crushes the prior best single (LSTM \num{0.127}).

\begin{table}[ht]
  \centering
  \caption{Tier A1 of the literature-fit campaign: the learned content projection
  on PCA-192, versus the R+M+C champion (\num{1431} test sessions; bootstrap 95\%
  CI; paired $\Delta$ per session vs the champion). Winner (\Rp+M+\Cp) bolded; best
  cell per metric bolded. \emph{Not comparable to any rotation board.}}
  \label{tab:nexttrack-projection}
  {\footnotesize\setlength{\tabcolsep}{4pt}\begin{tabular}{lccccl}
    \toprule
    Model & R@10 [95\% CI] & MRR & art@10 & gen@10 & paired $\Delta$ vs champion \\
    \midrule
    R+M+C champion (thirds)           & 0.1859 [0.166, 0.207] & 0.1025 & 0.331 & 0.394 & --- (reference) \\
    \best{\Rp+M+\Cp (both projected)} & \best{0.2117 [0.191, 0.233]} & \best{0.1220} & 0.333 & 0.401 & \best{$+0.0259$ [$+0.012,+0.041$] $>0$} \\
    R+M+C+\Cp (4-leg)                 & 0.2061 [0.185, 0.227] & 0.1148 & \best{0.345} & \best{0.410} & $+0.0203$ [$+0.013,+0.029$] $>0$ \\
    R+M+\Cp (swap C$\to$\Cp)          & 0.1971 [0.177, 0.218] & 0.1107 & 0.334 & 0.400 & $+0.0112$ [$-0.000,+0.023$] cand. \\
    \Cp\ proj-content (sym, r192)     & 0.1516 & 0.0708 & 0.305 & 0.401 & vs unsup.\ C 0.0685: $+0.083$ \\
    \Cp\ dir.\ two-tower (r192)       & 0.1349 & 0.0552 & 0.291 & 0.393 & symmetric wins \\
    \bottomrule
  \end{tabular}}
\end{table}

\emph{The lanes that closed.} Every other literature-fit direction failed to beat
the fixed z-blend, exactly as the survey predicted for a single-user, cold-heavy
corpus (Table~\ref{tab:nexttrack-projection-lanes}). Principled Markov smoothing
(absolute-discounting $+$ a Kneser--Ney continuation floor) is a real but tiny
standalone estimator gain (\num{0.1069}$\to$\num{0.1104}) that washes out in the
blend ($-\num{0.0021}$, CI straddles 0). A ranking-aware combiner has no honest
headroom: even the \emph{oracle} weights, fit directly on the test labels, reach
only \num{0.1901} ($+\num{0.0042}$, straddles 0), and RRF rank fusion \emph{hurts}
recall ($-\num{0.029}$, CI$<$0). A gated V-SKNN session-kNN is warm-only against
55\% cold targets and \emph{drags} the blend ($-\num{0.038}$, CI$<$0) --- the
single-user setting strips the cross-user collaborative neighborhood that drives
its multi-user benchmark wins. And GRU4Rec importance negatives give a within-noise
standalone tick ($+\num{0.005}$) that also washes out ($-\num{0.0014}$, tie).
\textbf{The only additive gain comes from new signal, not smarter fusion.}

\begin{table}[ht]
  \centering
  \caption{The literature-fit lanes that closed: every borrowed method fails to
  beat the fixed z-blend champion (\num{1431} test sessions; blend paired $\Delta$
  per session vs the champion).}
  \label{tab:nexttrack-projection-lanes}
  {\footnotesize\setlength{\tabcolsep}{5pt}\begin{tabular}{llcl}
    \toprule
    Lane & Standalone R@10 & Blend $\Delta$ vs champion & Verdict \\
    \midrule
    Markov KN-smoothing (Tier A2)         & 0.1069$\to$0.1104 & $-0.0021$ (straddles 0) & null in blend \\
    Ranking combiner, \emph{oracle} (B1)  & 0.1901 (blend) & $+0.0042$ (straddles 0) & no honest headroom \\
    RRF rank fusion (B1)                   & 0.1572 (blend) & $-0.029$ (CI$<$0) & hurts recall \\
    V-SKNN session-kNN (B2)                & 0.0671 & $-0.038$ (CI$<$0) & drags (warm-only) \\
    GRU importance negatives (C)           & 0.1230$\to$0.1279 & $-0.0014$ (tie) & washes out \\
    \bottomrule
  \end{tabular}}
\end{table}

\emph{Interpretation.} The supervised content projection is the first
genuinely-additive new signal since the content leg itself, and it works for the
reason the survey flagged: content is the only channel with cold reach, and
replacing the unsupervised, variance-ordered PCA geometry with a target-fit one
lifts \emph{both} legs that share it. The flashy borrowed methods are dead ends
here because they either re-shuffle the same three legs (no headroom over the
z-blend, even for the oracle) or import a collaborative assumption the single-user
corpus violates.
\emph{Source: \texttt{2026-07-18-nexttrack-literature-fit-campaign.md}.}

\subsection{Registration, the objective / rank / $\tau$ scan, and second-split hardening}
\label{sec:nexttrack-projection-scan}
\emph{Rationale.} The projection win landed as an off-server driver. This
follow-through implemented it as a first-class capability of the \texttt{seq-blend}
predictor (hyperparameters \texttt{projection}, \texttt{projection\_rank},
\texttt{projection\_objective}, \texttt{projection\_tau}, \dots), then reproduced
the win on-server, registered $+$ promoted it, and ran a pre-approved ablation scan.

\emph{Registration.} The verify run
(\runid{run-20260718-191539-1b89c-seq-blend}, definition default $=$ InfoNCE /
full-rank / $\tau$ \num{0.07}) reproduced the off-server driver to four decimals ---
R@10 \num{0.21174}, MRR \num{0.12199} --- and was promoted via the API as
\texttt{blend-gru-markov-content-proj}, the projection matrix baked into the model
directory for serving. \texttt{GET /api/best-models} confirms it at \textbf{rank
1} (\num{0.21174}), above the prior servable champion
\texttt{blend-gru-markov-content} (rank 2, \num{0.18588}); it is servable via
\texttt{POST /api/models/blend-gru-markov-content-proj/predict}.

\emph{The scan.} The decision rule --- re-promote only a config beating the default
on \emph{both} R@10 and MRR by more than the $\sim$$\pm\num{0.015}$ noise band ---
confirms the default (Table~\ref{tab:nexttrack-projection-scan}). \textbf{InfoNCE
$\gg$ BPR} is the sharpest axis: both BPR cells (\num{0.1579}, \num{0.1523}) fall
$\sim$\num{0.055} below the InfoNCE default and land \emph{below even the old R+M+C
champion} (\num{0.1859}). \textbf{Full-rank $\ge$ rank-64} ($-\num{0.006}$, a tie).
And \textbf{$\tau$ is mildly sensitive and monotone up} (\num{0.05}$\to$\num{0.1929},
\num{0.07}$\to$\num{0.2117}, \num{0.10}$\to$\num{0.2229}) --- but the nominal
$\tau$-\num{0.10} top cell beats the default by only $+\num{0.0112}$ on recall,
\emph{inside} the noise band, so it is a tie and the default \textbf{HELD}. Unlike
the reweighting axis of Section~\ref{sec:nexttrack-projection-win}, the graded
metrics here \emph{track} recall, because the scan moves the \emph{signal quality}
rather than re-fusing fixed legs.

\begin{table}[ht]
  \centering
  \caption{The on-server objective / rank / $\tau$ ablation scan (6 cells, sorted
  by recall@10). All overlay the promoted definition; hit@10 $=$ recall@10 by
  construction. Promoted champion (default, held) bolded; best cell per metric
  bolded (all held by the sub-noise $\tau$-\num{0.10} cell).}
  \label{tab:nexttrack-projection-scan}
  {\footnotesize\setlength{\tabcolsep}{4pt}\begin{tabular}{lccccp{3.0cm}}
    \toprule
    Config & recall@10 & MRR & art@10 & gen@10 & run id \\
    \midrule
    S2 infonce / full / $\tau$ \num{0.10}          & \best{0.2229} & \best{0.1273} & \best{0.3424} & \best{0.4032} & \runid{run-20260718-193111-a781c} \\
    \best{infonce / full / $\tau$ \num{0.07} (promoted, held)} & \best{0.2117} & \best{0.1220} & 0.3326 & 0.4011 & \runid{run-20260718-191539-1b89c} \\
    S3 infonce / rank 64 / $\tau$ \num{0.07}        & 0.2055 & 0.1169 & 0.3326 & 0.3990 & \runid{run-20260718-193650-66c96} \\
    S1 infonce / full / $\tau$ \num{0.05}          & 0.1929 & 0.1111 & 0.3354 & 0.3997 & \runid{run-20260718-192443-57432} \\
    S4 bpr / full / $\tau$ \num{0.07}              & 0.1579 & 0.0939 & 0.3187 & 0.3913 & \runid{run-20260718-194228-4cc06} \\
    S5 bpr / rank 64 / $\tau$ \num{0.07}            & 0.1523 & 0.0909 & 0.3194 & 0.3906 & \runid{run-20260718-194832-d070f} \\
    \bottomrule
  \end{tabular}}
\end{table}

\emph{Second-split hardening.} The crown was measured on a single leak-free split;
a validation-only pass re-ran it (plus the two reference blends) on a
\textbf{second, disjoint, leak-free split} \dsid{seq-20260718-211238} --- an
earlier-holdout 60--80\% chronological band on the same PCA-192 space, item-space
byte-identical, test window disjoint in index and time from the champion split's
80--100\% test, with a \textbf{cold-item rate of \num{0.64}} (vs \num{0.55}). By
construction absolute R@10 lands below the first-split \num{0.212}; the verdict is
the paired $\Delta$ (Table~\ref{tab:nexttrack-projection-hardening};
Figure~\ref{fig:nexttrack-projection-hardening}). The decision rule (paired
$\Delta$ vs R+M+C with CI entirely $>0$ and MRR not regressing) is met: paired
$\Delta$ \textbf{$+\num{0.0426}$ $[+0.0280, +0.0573]$}, MRR $+\num{0.0234}$. The
margin is \emph{wider} here than on the first split ($+\num{0.0259}$) --- on a
colder split the crown's Markov / content-kNN legs reach fewer targets and the
projection's cold-reaching \Cp\ leg carries proportionally more of the win. The
crown is now hardened on two disjoint leak-free splits, resolving both R+M+C
caveats.

\begin{table}[ht]
  \centering
  \caption{Second-split hardening on the disjoint 0.64-cold earlier-holdout split
  \dsid{seq-20260718-211238} (\num{1431} test sessions; validation-only). Point
  metrics (top) and the paired per-session $\Delta$ (bottom, 2000-resample
  bootstrap, seed 1337). Projection champion bolded.}
  \label{tab:nexttrack-projection-hardening}
  {\footnotesize\setlength{\tabcolsep}{4pt}\begin{tabular}{lccccp{3.0cm}}
    \toprule
    Config & recall@10 & MRR & art@10 & gen@10 & run id \\
    \midrule
    \best{\Rp+M+\Cp projection champion} & \best{0.1705} & \best{0.1050} & \best{0.2991} & \best{0.3417} & \runid{run-20260718-212639-50e1d} \\
    R+M+C crown (primary reference)      & 0.1279 & 0.0816 & 0.2956 & 0.3410 & \runid{run-20260718-213156-cf1dd} \\
    R+M 2-leg (context reference)        & 0.1174 & 0.0791 & 0.2844 & 0.3368 & \runid{run-20260718-214114-93bad} \\
    \midrule
    \multicolumn{6}{l}{\emph{Paired $\Delta$ recall@10 (projection $-$ reference):}} \\
    PROJ vs R+M+C (primary) & \multicolumn{5}{l}{$+\num{0.0426}$ $[+0.0280, +0.0573]$ --- entirely $>0$; MRR $+\num{0.0234}$ (no regression)} \\
    PROJ vs R+M (context)   & \multicolumn{5}{l}{$+\num{0.0531}$ $[+0.0391, +0.0678]$ --- entirely $>0$} \\
    \bottomrule
  \end{tabular}}
\end{table}

\begin{figure}[ht]
  \centering
  \includegraphics[width=0.8\textwidth]{nexttrack_projection_hardening}
  \caption{The projection lead \emph{widens} on the colder split. On the canonical
  0.55-cold split the projection champion beats the R+M+C crown by paired $\Delta$
  $+\num{0.0259}$; on a disjoint, colder 0.64-cold earlier-holdout split the margin
  \emph{grows} to $+\num{0.0426}$ (both CIs entirely $>0$). Absolute R@10 is lower
  on the colder split by construction; the verdict is the paired $\Delta$. The
  cold-reaching \Cp\ leg carries proportionally more of the win where more targets
  are cold.}
  \label{fig:nexttrack-projection-hardening}
\end{figure}

\emph{Source: \texttt{2026-07-18b-nexttrack-projection-registration-and-scan.md}.}

\section{Next-track: item-vector representation exploration}
\label{sec:nexttrack-representation}
\emph{Rationale.} The compression axis (Sections~\ref{sec:nexttrack-compression}
and~\ref{sec:nexttrack-dim}) found latent \emph{dimension} to dominate but never
asked whether the \emph{composition} of the 192-d item vector --- the
equal-weighted 4-block PCA (text / numeric / acoustic / categorical) --- is itself a
handicap. An explained-variance-ratio (EVR) preflight said it might be: the raw
input variance budget is 44\% acoustic $+$ 43\% numeric versus only 2.5\% text (text
is L2-normalized, so its per-dimension scale is tiny), so full-matrix PCA-192 spends
\textbf{87\% of its variance budget} reconstructing the numeric $+$ acoustic blocks
to $R^2 = \num{1.000}$ while the signal-carrying text block reaches only $R^2 =
\num{0.717}$; per-column standardize (whitening) flips the captured budget to 73\%
text / 21\% categorical (text $R^2$ \num{0.905})
(Figure~\ref{fig:nexttrack-evr-budget}). Eight alternative item-latent spaces were
built parity-gated to the champion split (only \texttt{item\_latents.f32} differs)
and each evaluated under \emph{two} readouts: \textbf{(A)} the champion projection
\texttt{blend-gru-markov-content-proj} (the space's ceiling under the winning
recipe) and \textbf{(C)} the frozen \texttt{gru-infonce-h256} single (raw
representation quality). Validation-only; nothing promoted.

\begin{figure}[ht]
  \centering
  \includegraphics[width=0.78\textwidth]{nexttrack_evr_budget}
  \caption{Why the item vector is a handicap for a raw model --- and why the
  projection is the fix. The equal-weighted PCA-192 orders dimensions by raw
  variance, so 87\% of its captured budget reconstructs the behaviorally-inert
  numeric $+$ acoustic blocks (popularity / era / loudness) while the
  signal-carrying text block (2.5\% of raw input variance, L2-normalized) is
  under-preserved ($R^2$ \num{0.717}). Per-column whitening flips the captured
  budget to 73\% text / 21\% categorical (text $R^2$ \num{0.905}). The learned
  InfoNCE projection learns this same reweighting from data.}
  \label{fig:nexttrack-evr-budget}
\end{figure}

\begin{table}[ht]
  \centering
  \caption{Item-vector representation exploration: eight alternative 192-d item
  spaces vs the PCA-192 champion, under readout A (champion projection) and readout
  C (frozen GRU). Paired $\Delta$ on exact hit@10 (2000-resample bootstrap, rng
  1337) vs the same-readout baseline; a lift needs the CI entirely $>0$. Baseline
  rows and the best readout-C lift (V5) bolded. \emph{Not comparable to any rotation
  board.}}
  \label{tab:nexttrack-representation}
  {\scriptsize\setlength{\tabcolsep}{3pt}\begin{tabular}{llcccc}
    \toprule
    & & \multicolumn{2}{c}{Readout A (projection)} & \multicolumn{2}{c}{Readout C (frozen GRU)} \\
    \cmidrule(lr){3-4}\cmidrule(lr){5-6}
    Variant (space) & dataset id & R@10 & $\Delta$ hit@10 [CI] & R@10 & $\Delta$ hit@10 [CI] \\
    \midrule
    \best{full --- PCA-192 (baseline)} & \dsid{seq-20260715-131139} & \best{0.2117} & (baseline) & \best{0.1230} & (baseline) \\
    V1 noaco               & \dsid{seq-20260718-222742} & 0.2089 & $-0.0028$ [$-0.013,+0.008$] & 0.1481 & $+0.0252$ [$+0.009,+0.041$] \\
    V2 std (whiten all)    & \dsid{seq-20260718-222754} & 0.2082 & $-0.0035$ [$-0.015,+0.008$] & 0.1530 & $+0.0300$ [$+0.011,+0.049$] \\
    V3 std-noaco           & \dsid{seq-20260718-222807} & 0.2166 & $+0.0049$ [$-0.008,+0.017$] & 0.1572 & $+0.0342$ [$+0.015,+0.055$] \\
    V4 textcat             & \dsid{seq-20260718-222820} & 0.1950 & $-0.0168$ [$-0.029,-0.006$] & 0.1565 & $+0.0335$ [$+0.015,+0.051$] \\
    \best{V5 balanced-pca192} & \dsid{seq-20260718-222834} & 0.2159 & $+0.0042$ [$-0.006,+0.015$] & \best{0.1586} & \best{$+0.0356$ [$+0.019,+0.052$]} \\
    V7 std-txtcatup        & \dsid{seq-20260718-222848} & 0.2055 & $-0.0063$ [$-0.019,+0.006$] & 0.1488 & $+0.0259$ [$+0.006,+0.046$] \\
    V10 text-only          & \dsid{seq-20260718-222901} & 0.1999 & $-0.0119$ [$-0.025,+0.001$] & 0.1377 & $+0.0147$ [$-0.005,+0.034$] \\
    V6 sonic-64 (64-d)     & \dsid{seq-20260718-222905} & 0.1936 & $-0.0182$ [$-0.031,-0.005$] & 0.1216 & $-0.0014$ [$-0.017,+0.013$] \\
    \bottomrule
  \end{tabular}}
\end{table}

\begin{figure}[ht]
  \centering
  \includegraphics[width=0.8\textwidth]{nexttrack_representation}
  \caption{The representation handicap is real for the \emph{raw} GRU but
  \emph{absorbed} by the champion projection. Six of eight whitening / acoustic-drop
  / metric spaces lift the frozen GRU (readout C, green; all six with paired-$\Delta$
  CI entirely $>0$), but none lifts the champion projection (readout A, blue; the two
  nominally-positive spaces straddle 0, and V4/V6 significantly regress). The learned
  InfoNCE map already re-weights away the loud acoustic / numeric directions:
  whitening and the projection are the same fix.}
  \label{fig:nexttrack-representation}
\end{figure}

\emph{Verdict.} \textbf{Readout A: no variant beats the PCA-192 champion with paired
$\Delta$ CI entirely $>0$} (V3 std-noaco $+\num{0.0049}$ and V5 balanced
$+\num{0.0042}$ straddle 0; V4 textcat $-\num{0.0168}$ and V6 sonic-64 $-\num{0.0182}$
regress). \textbf{Readout C: six of eight variants beat the frozen GRU with paired
$\Delta$ CI entirely $>0$} ($+\num{0.025}$ to $+\num{0.036}$). So the raw geometry
\emph{is} a handicap --- suppressing the loud directions lifts the raw GRU by
$\sim$20--29\% relative --- \textbf{but the gain is fully absorbed by the
projection}, which already re-weights those directions from data. Whitening and the
learned projection are the same fix, and the champion has already pulled it. The
graded axis moves more cleanly than exact recall (whitening lifts the frozen GRU's
genre@10 to $\sim$\num{0.47}; a text-only space posts the readout-A artist@10 /
genre@10 maxima), and the continuous music@10 metric
(Section~\ref{sec:nexttrack-music}) corroborates: the projection champion leads on
it (\num{0.4538}) over R+M+C (\num{0.4491}), and it is flat across the eight
representations ($\sim$\num{0.447}--\num{0.460}). One caution: dropping
numeric$+$acoustic \emph{entirely} (V4) hurts the champion even though it helps the
raw GRU, so keep numeric and whiten / drop only acoustic.
\emph{Source: \texttt{2026-07-18c-nexttrack-representation-exploration.md}.}

\section{Next-track: the music@10 musical-distance metric}
\label{sec:nexttrack-music}
\emph{Rationale.} Exact Recall@10 is unforgiving on a 55\%-cold single-user corpus.
This campaign introduces a \emph{continuous} ``sounds-alike'' graded metric,
\textbf{music@10}, and back-fills it across every next-track run. It is the mean
over test cases of the best (max) cosine, in a \emph{balanced} intrinsic
musical-distance space (\dsid{spotify\_tracks\_content\_metric}: a block-weighted,
partial-whitened fusion of the whitened 64-d song-AE latent, the 384-d name/tag text
embedding, five metadata numerics, and a per-genre AE prototype), between the true
next track and the top-ten retrieved candidates. A top-ten containing the exact
truth scores \num{1.0}, so music@10 is a \emph{smooth superset} of Recall@10; it is
display-only (the primary metric stays Recall@10).

\emph{Results.} music@10 was recomputed offline and non-destructively across all 53
next-track runs (from each run's stored \texttt{predictions.json} and its dataset
\texttt{items.json}, on the shared \num{1431}-case split). It \textbf{correlates}
with Recall@10 ($r = \num{0.609}$) but is far from redundant, and it crowns a
\emph{different} winner (Table~\ref{tab:nexttrack-music-families};
Figure~\ref{fig:nexttrack-music-scatter}). The Recall@10 champion is a
\texttt{seq-blend} run (\num{0.223}), but the \emph{musically-closest} predictions
come from \texttt{seq-nexttrack} --- which regresses the next track's content latent
directly and retrieves by cosine --- at best music@10 \num{0.482} and the highest
family mean (\num{0.448}). The \texttt{seq-popularity} floor is correctly worst
(\num{0.185}), confirming the metric is not trivially saturated.

\begin{table}[ht]
  \centering
  \caption{music@10 by predictor family, from the offline backfill over all 53
  next-track runs (one shared \num{1431}-case split). The recall@10 champion
  (\texttt{seq-blend}) is \emph{not} the music@10 champion (\texttt{seq-nexttrack}).
  Best cell per column bolded.}
  \label{tab:nexttrack-music-families}
  {\footnotesize\setlength{\tabcolsep}{6pt}\begin{tabular}{lcccc}
    \toprule
    Predictor family & runs & best recall@10 & best music@10 & mean music@10 \\
    \midrule
    seq-blend      & 23 & \best{0.223} & 0.461 & 0.441 \\
    seq-nexttrack  & 24 & 0.159 & \best{0.482} & \best{0.448} \\
    seq-markov     & 1  & 0.107 & 0.407 & 0.407 \\
    seq-ann        & 1  & 0.090 & 0.426 & 0.426 \\
    seq-stacker    & 3  & 0.087 & 0.418 & 0.411 \\
    seq-popularity & 1  & 0.001 & 0.185 & 0.185 \\
    \bottomrule
  \end{tabular}}
\end{table}

\begin{figure}[ht]
  \centering
  \includegraphics[width=0.82\textwidth]{nexttrack_music_scatter}
  \caption{music@10 versus Recall@10 across all 53 next-track runs, one point per
  run, coloured by predictor family. The two axes correlate ($r = \num{0.609}$) but
  crown different winners --- the Recall@10 champion (a \texttt{seq-blend} run at
  \num{0.223}) is not the musically-closest, which is a \texttt{seq-nexttrack} run
  (music@10 \num{0.482}) that regresses the content latent and retrieves by cosine.
  music@10 is display-only; Recall@10 stays the primary metric.}
  \label{fig:nexttrack-music-scatter}
\end{figure}

\emph{Interpretation.} music@10 is a genuinely distinct, non-saturated view of
next-track quality: it rewards \emph{sonic} proximity, which the exact-hit metric
ignores. The recall-versus-music decoupling has a clean mechanism --- a model that
regresses the content latent (\texttt{seq-nexttrack}) lands \emph{near} the truth in
content space even when it misses the exact next track, so its misses ``sound
right'' more often than the blend's. This is why the metric is display-only and must
not silently become a selection objective; promotion still ranks by Recall@10. It
inherits a mild circularity from the \emph{balanced} dial (the text block embeds the
artist name; the genre block is a genre prototype), so it partly rewards
same-artist / genre proximity.
\emph{Source: \texttt{2026-07-20-musical-distance-metric.md}.}

\section{Next-track: session holisticness metrics and the anti-eager levers (MMR $\lambda$ / eager $\beta$)}
\label{sec:nexttrack-holisticness}
\emph{Rationale.} A recurring qualitative complaint about the top blends is that
their predictions feel ``too album-eager / bigram-like'' --- they parrot the
current artist rather than curating a holistic continuation. This campaign
quantifies that symptom with a newly-added \emph{display-only session-holisticness}
metric suite, tests the user's hypothesis that \emph{the GRU beats the LSTM because
it is more eager}, and measures whether two new Phase-2 levers reduce eagerness at
an acceptable recall cost: an eval-time MMR re-rank (\texttt{mmr\_lambda}) and an
anti-eager training regularizer (\texttt{eager\_beta}). It is a measurement /
exploratory pass: nothing is promoted, the primary metric stays Recall@10, and the
holisticness objective decision is explicitly deferred to the user. Every arm runs
on the \emph{canonical} PCA-192 champion split \dsid{seq-20260715-131139}
(\num{1431} test sessions; cold-item rate \num{0.55}), and the champion reproduced
its on-record Recall@10 \num{0.2117} exactly (\num{0.21174}).

\paragraph{The new display-only holisticness metrics.} Wired end-to-end alongside
music@10 (\texttt{seq\_common.eval\_from\_scores} $\to$ Rust $\to$
\texttt{domain.toml} $\to$ UI); the primary metric is unchanged.
\begin{description}
  \item[artist\_adj@10 ($\downarrow$).] Fraction of the top-10 sharing the
  \emph{seed} track's artist --- the ``parrots the current artist'' eagerness rate.
  \textbf{Lower is better}.
  \item[album\_adj@10.] The same on (artist, album); \textbf{omitted on this
  dataset} because its \texttt{items.json} carries no album field (expected, not an
  error --- covered by a separate album-carrying dataset).
  \item[mood\_coh@10 ($\uparrow$).] Mean cosine of the top-10 to the prefix mood
  centroid in the content-metric space; higher is more mood-coherent.
  \item[ild@10 ($\uparrow$).] Intra-list diversity (mean pairwise sonic distance) of
  the top-10; higher is less duplicative.
  \item[suffix\_recall@10 / cont\_prec@10.] Multi-step continuation coverage,
  produced only by the on-demand harness \texttt{seq\_continuation\_eval.py}
  (Phase E below), not by normal runs.
\end{description}
The sanity gate passed: music@10 / mood\_coh@10 / ild@10 populate on all blend and
GRU runs, so the content-metric vectors loaded and the MMR / mood / ild levers are
live, not no-ops. Both new levers are OFF by default and verified byte-identical
when off.

\emph{Results --- the holisticness fingerprint (Phase A).} Table~\ref{tab:nexttrack-holisticness-fingerprint}
fingerprints the model family. \textbf{Eagerness tracks Recall@10 inside the blend
family} --- the champion (Recall@10 \num{0.2117}) is the \emph{most} artist-eager
model measured (artist\_adj@10 \num{0.610}), with R+M next (\num{0.579}) --- but the
eagerness lives in the \emph{Markov} leg, not the RNN: the standalone first-order
Markov bar is already at \num{0.543}, while the single GRU / LSTM sit far lower at
artist\_adj@10 $\approx$ \num{0.345}. So the blends buy their Recall@10 lead by being
artist-adjacent (bigram-like), which is exactly what the holisticness suite flags.
The single RNNs are conversely the \emph{mood-coherent, low-eager} models (highest
mood\_coh@10, lowest real-model artist\_adj@10). Crucially the metrics are \emph{not
trivially gamed}: the popularity floor maxes ild@10 (\num{0.898}) and minimizes
artist\_adj@10 (\num{0.003}) yet scores $\approx 0$ recall and \emph{negative}
mood\_coh@10 --- you maximize them by being random, so the meaningful low-eager /
diverse models are the real single RNNs.

\begin{table}[ht]
  \centering
  \caption{Phase A holisticness fingerprint across the model family, sorted by
  Recall@10 (\num{1431} test sessions). All runs share the prefix
  \protect\runid{run-20260722-165345-}; the id column lists the distinguishing suffix.
  Recall@10 (raw hits of \num{1431}), MRR, artist@10, genre@10, music@10,
  mood\_coh@10, ild@10 higher-better; \textbf{artist\_adj@10 lower-better}. Champion
  row and best cell per metric bolded. \textsuperscript{\dag}~The popularity floor's
  extreme artist\_adj@10 / ild@10 are \emph{degenerate} (it ignores the seed): it
  lands $\approx 0$ recall and negative mood\_coh@10, proving the metrics are not
  gamed low/high by a good model.}
  \label{tab:nexttrack-holisticness-fingerprint}
  {\scriptsize\setlength{\tabcolsep}{3pt}\begin{tabular}{llcccccccc}
    \toprule
    Config & id & R@10 (raw) & MRR & art@10 & gen@10 & mus@10 & artadj@10\,$\downarrow$ & mood@10\,$\uparrow$ & ild@10\,$\uparrow$ \\
    \midrule
    \best{A1 champion \Rp+M+\Cp} & \texttt{82a8f} & \best{0.2117 (303)} & \best{0.122} & 0.333 & 0.401 & 0.454 & 0.610 & 0.398 & 0.406 \\
    A2 R+M (blend)               & \texttt{28dbb} & 0.1712 (245) & 0.095 & 0.329 & 0.398 & 0.443 & 0.579 & 0.399 & 0.403 \\
    A3 GRU alone ($\beta{=}0$)   & \texttt{e39c9} & 0.1230 (176) & 0.053 & \best{0.338} & 0.450 & \best{0.456} & 0.347 & \best{0.479} & 0.454 \\
    B LSTM (else matched to A3)  & \texttt{bb387} & 0.1188 (170) & 0.049 & 0.328 & \best{0.455} & 0.449 & 0.344 & 0.458 & 0.463 \\
    A4 markov (first-order)      & \texttt{047cc} & 0.1069 (153) & 0.069 & 0.298 & 0.389 & 0.407 & 0.543 & 0.310 & 0.462 \\
    A5 popularity                & \texttt{34614} & 0.0007 (1)   & 0.001 & 0.014 & 0.085 & 0.185 & \best{0.003}\textsuperscript{\dag} & $-0.041$ & \best{0.898}\textsuperscript{\dag} \\
    \bottomrule
  \end{tabular}}
\end{table}

\emph{GRU versus LSTM (Phase B) --- the hypothesis is not supported.} With
hyperparameters matched (h256, InfoNCE, 40 epochs, lr \num{1e-3}), the GRU and LSTM
are statistically \emph{tied}: Recall@10 \num{0.1230} versus \num{0.1188} ($\Delta$
$+\num{0.0042}$, inside the $\pm\num{0.015}$ band), and artist\_adj@10 \num{0.347}
versus \num{0.344} (a negligible $+\num{0.003}$). The GRU is not measurably more
eager than the LSTM on this split, so the ``GRU wins because it parrots the current
artist harder'' hypothesis is \textbf{refuted here}.

\emph{The MMR $\lambda$ lever (Phase C) --- the strong knob.} MMR is a pure
eval-time re-rank, so the four originally-launched per-$\lambda$ training runs were
\emph{cancelled} (not failed experiments); the sweep was computed offline from the
single champion checkpoint (the $\lambda{=}$off cell reproduced the champion's
Recall@10 \num{0.2117} exactly). As $\lambda$ drops from off toward \num{0.3}
(Table~\ref{tab:nexttrack-mmr-sweep}; Figure~\ref{fig:nexttrack-mmr-tradeoff}), Recall@10 falls
--- the \emph{only} quantity traded --- while ild@10 rises and artist\_adj@10 falls
(de-eagering). \textbf{$\lambda{=}\num{0.9}$ is essentially free}: the recall
$\Delta$ is $-\num{0.0014}$, well inside the $\pm\num{0.015}$ band (a tie), yet it
already lifts ild@10 \num{0.406}$\to$\num{0.436} and trims eagerness
\num{0.610}$\to$\num{0.603}. \textbf{$\lambda\approx\num{0.7}$} trades $\sim$2 recall
points ($\Delta$ $-\num{0.0181}$, just outside the band) for a large ild@10 gain
(\num{0.406}$\to$\num{0.520}) and a real eagerness drop
(\num{0.610}$\to$\num{0.571}). A notable side effect: \textbf{music@10 rises
monotonically} as you diversify (\num{0.454}$\to$\num{0.492}) --- sonic closeness and
list diversity are not in tension; only exact recall is traded.

\begin{table}[ht]
  \centering
  \caption{Phase C MMR $\lambda$ sweep on the champion, computed offline from its
  checkpoint (\texttt{mmr\_pool} \num{200}; \num{1431} test sessions). Recall@10 (raw
  hits), MRR, artist@10, music@10, mood\_coh@10, ild@10 higher-better;
  \textbf{artist\_adj@10 lower-better}. Best cell per metric bolded; the champion
  default ($\lambda{=}$off) recall / MRR bolded.}
  \label{tab:nexttrack-mmr-sweep}
  {\footnotesize\setlength{\tabcolsep}{4pt}\begin{tabular}{lcccccccc}
    \toprule
    $\lambda$ & R@10 (raw) & $\Delta$ vs off & MRR & art@10 & mus@10 & artadj\,$\downarrow$ & mood\,$\uparrow$ & ild\,$\uparrow$ \\
    \midrule
    off (1.0) & \best{0.2117 (303)} & --- & \best{0.122} & 0.333 & 0.454 & 0.610 & \best{0.398} & 0.406 \\
    0.9 & 0.2103 (301) & $-0.0014$ (TIE) & 0.121 & 0.341 & 0.462 & 0.603 & 0.391 & 0.436 \\
    0.7 & 0.1936 (277) & $-0.0181$ & 0.116 & 0.347 & 0.476 & 0.571 & 0.364 & 0.520 \\
    0.5 & 0.1642 (235) & $-0.0475$ & 0.108 & \best{0.359} & 0.487 & 0.486 & 0.315 & 0.653 \\
    0.3 & 0.1097 (157) & $-0.1020$ & 0.100 & 0.358 & \best{0.492} & \best{0.293} & 0.238 & \best{0.815} \\
    \bottomrule
  \end{tabular}}
\end{table}

\begin{figure}[ht]
  \centering
  \includegraphics[width=0.8\textwidth]{nexttrack_mmr_tradeoff}
  \caption{The MMR $\lambda$ trade-off on the next-track champion. As $\lambda$ drops
  from off toward \num{0.3}, Recall@10 (red, left axis) falls --- the only quantity
  traded --- while ild@10 rises and artist\_adj@10 falls (right axis; more diverse,
  less artist-eager) and music@10 even \emph{rises}. $\lambda{=}\num{0.9}$ sits inside
  the $\pm\num{0.015}$ recall tie band yet already lifts ild@10 by \num{0.030};
  $\lambda\approx\num{0.7}$ buys a large ild@10 gain and real de-eagering for
  $\sim$2 recall points. MMR is the strong holisticness lever; $\beta$
  (Table~\ref{tab:nexttrack-eager-beta}) is weak.}
  \label{fig:nexttrack-mmr-tradeoff}
\end{figure}

\emph{The anti-eager regularizer $\beta$ (Phase D) --- the weak knob.} A
training-time penalty on near-copy prediction of the current track's latent
(\texttt{eager\_margin} \num{0.0}) barely moves the needle
(Table~\ref{tab:nexttrack-eager-beta}): across $\beta$ \numrange{0}{0.2}, artist\_adj@10 shifts
only \num{0.347}$\to$\num{0.340} ($\Delta$ $-\num{0.0074}$) and ild@10 only
\num{0.454}$\to$\num{0.458}, with Recall@10 flat (all four inside $\pm\num{0.015}$).
The single GRU is already low-eager (\num{0.347}), so there is little copy behaviour
left to penalize --- $\beta$ is not a useful holisticness knob at these settings.

\begin{table}[ht]
  \centering
  \caption{Phase D anti-eager regularizer $\beta$ sweep on the single GRU
  (\texttt{eager\_margin} \num{0.0}; \num{1431} test sessions). Recall@10, artist@10,
  mood\_coh@10, ild@10 higher-better; \textbf{artist\_adj@10 lower-better}. Best cell
  per metric bolded.}
  \label{tab:nexttrack-eager-beta}
  {\footnotesize\setlength{\tabcolsep}{4pt}\begin{tabular}{lccccccc}
    \toprule
    run id & $\beta$ & R@10 (raw) & MRR & art@10 & artadj\,$\downarrow$ & mood\,$\uparrow$ & ild\,$\uparrow$ \\
    \midrule
    \runid{run-20260722-165345-e39c9} & 0.0  & 0.1230 (176) & 0.053 & \best{0.338} & 0.347 & \best{0.479} & 0.454 \\
    \runid{run-20260722-165345-353b5} & 0.05 & \best{0.1237 (177)} & 0.053 & 0.337 & 0.346 & 0.478 & 0.455 \\
    \runid{run-20260723-002044-af11a} & 0.1  & 0.1202 (172) & 0.053 & 0.335 & 0.344 & 0.477 & 0.456 \\
    \runid{run-20260723-002044-17cf8} & 0.2  & 0.1202 (172) & 0.051 & 0.336 & \best{0.340} & 0.475 & \best{0.458} \\
    \bottomrule
  \end{tabular}}
\end{table}

\emph{Multi-step continuation (Phase E).} The on-demand harness
(\texttt{seq\_continuation\_eval.py} \texttt{-m 3 -k 10}, \num{1172} continuation
cases) shows the GRU \emph{beats} the Markov bigram on both multi-step metrics
(Table~\ref{tab:nexttrack-continuation}): the RNN's advantage \emph{widens} past the immediate
next track, because the bigram's single-step co-occurrence does not compound into
good 3-step continuations whereas the GRU's learned sequence state does --- the
mirror image of the single-step board, where the bigram-carrying blend wins raw
Recall@10.

\begin{table}[ht]
  \centering
  \caption{Phase E multi-step continuation (harness \texttt{-m 3 -k 10},
  \num{1172} cases) on \dsid{seq-20260715-131139}: the GRU (A3 checkpoint) beats the
  first-order Markov bigram on both metrics (higher-better). Winner bolded.}
  \label{tab:nexttrack-continuation}
  {\footnotesize\setlength{\tabcolsep}{6pt}\begin{tabular}{lcc}
    \toprule
    Model & suffix\_recall@10 & cont\_prec@10 \\
    \midrule
    markov (first-order bigram) & 0.0869 & 0.0236 \\
    \best{GRU (gru-infonce-h256, A3)} & \best{0.0963} & \best{0.0273} \\
    \bottomrule
  \end{tabular}}
\end{table}

\emph{Interpretation and the deferred decision.} The campaign delivers a clean
mechanistic account of the eagerness symptom: it is sourced from the Markov bigram
leg (which parrots the current artist via raw co-occurrence), not the neural leg, so
any future de-eagering should target the bigram / the fusion, not the RNN. Of the two
new levers, the eval-time MMR re-rank is the \emph{strong} knob and the training-time
$\beta$ regularizer is \emph{weak}. Nothing is promoted --- the primary metric stays
Recall@10 and the champion is unchanged --- but \textbf{$\lambda\approx\num{0.9}$ is
flagged as a near-free holisticness upgrade for a future user promotion decision}
(Recall@10 tie within noise, ild@10 \num{0.406}$\to$\num{0.436}, eagerness
\num{0.610}$\to$\num{0.603}), with $\lambda\approx\num{0.7}$ the ``meaningful
de-eagering for $\sim$2 recall points'' operating point; adopting either requires a
user decision plus a 3-seed / native-run confirm. The $\beta$ regularizer is not
recommended as a lever at these settings.
\emph{Source: \texttt{2026-07-22-session-holisticness-and-antieager-levers.md}.}

\section{Next-track: dual-tower fusion refuted --- the fourth combiner null}
\label{sec:nexttrack-dualtower}
\emph{Rationale.} By 2026-07-25 the record held three independent
\emph{learned-combiner} nulls --- the XGB stacker (\numrange{0.080}{0.087},
\emph{below} its own GRU base leg; Section~\ref{sec:nexttrack-modelsweep}), reciprocal-rank
fusion ($-\num{0.029}$) and even \emph{oracle} test-fit blend weights
($+\num{0.004}$, CI straddling zero; Section~\ref{sec:nexttrack-projection-win}). All three
were \textbf{post-hoc, score-level} combiners, which left one obvious rebuttal
standing: \emph{of course} a bolted-on combiner fails --- train the fusion
end-to-end so the parts co-adapt. This campaign runs exactly that rebuttal, at
the \emph{representation} level, jointly, from scratch.

The vehicle is a new predictor, \texttt{seq-dualgru}: two \emph{independent}
recurrent towers (separate weights) read the prefix in parallel over
configurable \emph{causal views}, their per-step hidden states are concatenated,
and one MLP head fuses them into the predicted next-item latent. Three views are
exposed, all strictly causal at step $t$: \texttt{latent} (the raw item vector),
\texttt{delta} (the step-to-step difference, $x_{-1}=0$), and \texttt{cummean}
(the running session centroid --- the quantity the \texttt{seq-ann} pooled-MLP
leg was built from). Crucially it reuses \texttt{seq\_nexttrack.make\_batches}
and \texttt{step\_loss} verbatim, so objective, batching, early stopping and
evaluation are byte-for-byte the family's, and it carries no bidirectional mode
(a backward pass would leak the teacher-forcing target).

The hypothesis was sharp: two towers should beat one \emph{only if they see
genuinely different views}; two towers on the same \texttt{latent} view should
reduce to a capacity / implicit-ensemble effect a wider single GRU can buy more
cheaply. The design therefore crosses the \emph{complete upper triangle} of the
$\{\texttt{latent},\texttt{delta},\texttt{cummean}\}^2$ view matrix (six cells at
\texttt{fusion\_layers}${}=1$) with a \emph{fusion-depth} axis (two cells at
\texttt{fusion\_layers}${}=0$, a pure linear read of the concatenated states) and
--- the axis the earlier topology scan lacked (Section~\ref{sec:nexttrack-topology}) --- a
\textbf{parameter-matched single-GRU capacity control}. Parameter counts were
re-measured by instantiating the real classes at \texttt{latent\_dim}${}=192$:
the control C1 (h425, \num{871017} parameters) matches the six
\texttt{fusion\_layers}${}=1$ dual arms (\num{871872}) to
\SI{0.098}{\percent}, and is \SI{10}{\percent} \emph{larger} than the
\texttt{fusion\_layers}${}=0$ arms (\num{789696}) --- deliberately generous to
the single tower.

\paragraph{The honest bar.} Two boards were pre-registered \emph{separately}, so
that the scientific verdict and the board event could not be traded off after the
fact. On \textbf{Recall@10} the bar for a standalone single model is the
single-model board \num{0.123}, \emph{not} the three-leg blend champion
\num{0.212} --- comparing one trained model against a fusion of three legs would
be a category error, and dual-tower-as-the-\Rp-leg is a \texttt{seq-blend} code
change explicitly out of scope. On \textbf{holisticness@10} the crown rule
carries an anti-degenerate recall floor (see below). All ten runs are on the
canonical split \dsid{seq-20260715-131139} at seed 1337, executed one at a time
on a single remote worker, and the baseline C0 reproduced its on-record
Recall@10 \num{0.12299} (\num{176} of \num{1431}) \emph{exactly}.

\emph{Results.} Table~\ref{tab:nexttrack-dualtower} is the batch. \textbf{No arm is bolded
as a winner because there is no winner}: the single-GRU baseline holds the
batch's best Recall@10, and not one of the eight dual arms beats it by
paired-$\Delta$ CI above zero --- six lose outright and two (the
\texttt{fusion\_layers}${}=0$ pair) tie. The pre-registered refutation rule
fired.

\begin{table}[ht]
  \centering
  \caption{Dual-tower view-pairing $\times$ fusion-depth scan, ten runs on
  \dsid{seq-20260715-131139} (\num{1431} test sessions), sorted by Recall@10.
  All arms share the run-id prefix \protect\runid{run-20260725-143502-}; C0 is
  \protect\runid{run-20260725-143143-97c1e}. Paired $\Delta$ is a per-session
  paired bootstrap against C0 (\num{2000} resamples, rng 1337). Recall@10 raw =
  successes of \num{1431}; CI bounds are given to three decimals, and a CI
  straddling zero is a \emph{tie}. Arm codes carry the view pair
  (lat\,$=$\,\texttt{latent}, del\,$=$\,\texttt{delta},
  cum\,$=$\,\texttt{cummean}) and the fusion depth (\texttt{-f0} $=$ linear
  read, otherwise a 256-wide MLP head). R@10, MRR, mus@10, ild@10
  higher-better; \textbf{artadj@10 lower-better}; H is the \emph{3-factor}
  holisticness generation (Section~\ref{sec:nexttrack-power}). Best cell per metric
  bolded.
  \textsuperscript{\dag}~D-CC's extreme artadj@10 / H are \emph{not} a real
  crown: they are bought at \num{95} of \num{1431} hits, \num{0.0566} below the
  baseline, failing the $\pm\num{0.015}$ recall floor by $3.8\times$.}
  \label{tab:nexttrack-dualtower}
  {\scriptsize\setlength{\tabcolsep}{1.7pt}\begin{tabular}{llrrlccccc}
    \toprule
    Arm & id & params & R@10 (raw) & $\Delta$ vs C0 [95\% CI] & MRR & mus@10 & artadj\,$\downarrow$ & ild\,$\uparrow$ & H\,$\uparrow$ \\
    \midrule
    C0 h256 \emph{base} & \texttt{97c1e} & \num{394944} & \best{0.12299 (176)} & --- & 0.05341 & 0.4558 & 0.3474 & 0.4536 & 0.14172 \\
    C1 h425 \emph{ctrl} & \texttt{be11b} & \num{871017} & 0.12089 (173) & $-0.0021$ [$-.011$, $+.006$] & 0.05215 & \best{0.4597} & 0.3640 & 0.4506 & 0.13850 \\
    D-LL-f0 \emph{lat/lat} & \texttt{18b4a} & \num{789696} & 0.11880 (170) & $-0.0042$ [$-.013$, $+.004$] & \best{0.05631} & 0.4553 & 0.3630 & 0.4506 & 0.13774 \\
    D-LC-f0 \emph{lat/cum} & \texttt{3e5d4} & \num{789696} & 0.11740 (168) & $-0.0056$ [$-.017$, $+.005$] & 0.05087 & 0.4458 & 0.3557 & 0.4630 & 0.13648 \\
    D-LL \emph{lat/lat} & \texttt{12fe6} & \num{871872} & 0.10273 (147) & $-0.0203$ [$-.035$, $-.007$] & 0.04235 & 0.4446 & 0.3092 & 0.4873 & 0.14635 \\
    D-LC \emph{lat/cum} & \texttt{a95c5} & \num{871872} & 0.09713 (139) & $-0.0259$ [$-.039$, $-.013$] & 0.03901 & 0.4380 & 0.2997 & 0.4913 & 0.15000 \\
    D-LD \emph{lat/del} & \texttt{8d214} & \num{871872} & 0.09154 (131) & $-0.0315$ [$-.045$, $-.018$] & 0.04029 & 0.4385 & 0.2999 & 0.4988 & 0.14673 \\
    D-DC \emph{del/cum} & \texttt{e9cfb} & \num{871872} & 0.08665 (124) & $-0.0363$ [$-.050$, $-.022$] & 0.03379 & 0.4125 & 0.2319 & 0.5413 & 0.16340 \\
    D-CC \emph{cum/cum} & \texttt{88a4c} & \num{871872} & 0.06639 (95) & $-0.0566$ [$-.072$, $-.041$] & 0.02664 & 0.3935 & \best{0.1899}\textsuperscript{\dag} & 0.5235 & \best{0.19715}\textsuperscript{\dag} \\
    D-DD \emph{del/del} & \texttt{4324f} & \num{871872} & 0.05800 (83) & $-0.0650$ [$-.082$, $-.049$] & 0.02819 & 0.4016 & 0.2363 & \best{0.5457} & 0.13335 \\
    \bottomrule
  \end{tabular}}
\end{table}

\emph{Capacity is not the missed lever either.} The param-matched control C1
(h425, \num{871017} parameters --- $2.2\times$ C0) scores \num{0.12089} against
C0's \num{0.12299}, a paired $\Delta$ of $-\num{0.0021}$ straddling zero. This
fills a real gap: single-GRU \emph{width} had only ever been scanned on the
64-dim AE space (Section~\ref{sec:nexttrack-topology}), and the rich PCA-192 space was
untested. It behaves the same way, so \textbf{single-GRU width is now closed on
PCA-192 as well as AE-64}. Because C1 is a genuine \SI{0.098}{\percent}
parameter match to the \texttt{fusion\_layers}${}=1$ arms, the null cannot be
re-attributed to ``capacity would have helped'': it demonstrably would not.

\emph{View attribution --- complementarity is real, but lands on the wrong
pair.} The complete view triangle is what makes this mechanically decidable
rather than a judgment call, and the answer is genuinely surprising
(Table~\ref{tab:nexttrack-dualtower-views}, Figure~\ref{fig:nexttrack-dualtower-matrix}).
The pair the design nominated as its best hope --- \texttt{latent} $+$
\texttt{cummean} --- is \emph{not} complementary: it ties its latent/latent
parent and beats only the centroid-only one, i.e.\ the raw-latent tower carries
the cell and the centroid tower is dead weight. Same for \texttt{latent} $+$
\texttt{delta}. But \texttt{delta} $+$ \texttt{cummean} \emph{is} complementary
--- it beats \emph{both} parents with CI above zero. Read together: the
co-adaptation mechanism the hypothesis proposed genuinely \textbf{works}; it is
simply that \textbf{the raw \texttt{latent} view already subsumes everything the
other two carry}. A GRU over raw latents integrates step-to-step movement and
the running mean internally, so pairing it with an explicit computation of
either is redundant --- while two \emph{impoverished} views do have
non-overlapping information to pool, at a hopeless absolute level
(\num{0.08665}). That is a clean mechanistic explanation for the null, not
merely an absence of effect.

\begin{table}[ht]
  \centering
  \caption{View attribution: each mixed cell against \emph{both} of its
  single-view parents (paired $\Delta$ Recall@10, \num{1431} sessions). Only one
  of three mixed pairs is genuinely complementary --- and it is the one with no
  raw-latent tower.}
  \label{tab:nexttrack-dualtower-views}
  {\scriptsize\setlength{\tabcolsep}{3.4pt}\begin{tabular}{llll}
    \toprule
    Mixed cell & vs parent 1 & vs parent 2 & Verdict \\
    \midrule
    D-LC \emph{lat/cum} & D-LL: $-0.0056$ [$-.018$, $+.007$] & D-CC: $+0.0308$ [$+.018$, $+.045$] & \emph{not} complementary \\
    D-LD \emph{lat/del} & D-LL: $-0.0112$ [$-.023$, $+.001$] & D-DD: $+0.0335$ [$+.020$, $+.047$] & \emph{not} complementary \\
    \best{D-DC \emph{del/cum}} & D-DD: $+0.0287$ [$+.015$, $+.043$] & D-CC: $+0.0203$ [$+.008$, $+.034$] & \best{complementary: both CI$>$0} \\
    \bottomrule
  \end{tabular}}
\end{table}

\begin{figure}[ht]
  \centering
  \includegraphics[width=0.8\textwidth]{nexttrack_dualtower_matrix}
  \caption{The complete view-pairing matrix at \texttt{fusion\_layers}${}=1$,
  one run per cell. Every dual cell sits \emph{below} the single-GRU baseline
  (\num{0.12299}), and the ordering is monotone in how much raw \texttt{latent}
  the pair contains: both-latent is best, both-delta worst. The lower triangle
  was never run (the pairing is symmetric by construction) and is masked in grey
  rather than mirrored, so the colour scale reflects only measured cells. The
  gradient is the finding: a recurrent tower over raw latents already computes
  what \texttt{delta} and \texttt{cummean} make explicit.}
  \label{fig:nexttrack-dualtower-matrix}
\end{figure}

\emph{The MLP fusion head is actively harmful, not merely useless.} Both
\texttt{fusion\_layers}${}=0$ arms beat their \texttt{fusion\_layers}${}=1$ twins
by paired-$\Delta$ CI above zero ($+\num{0.0161}$ on latent/latent,
$+\num{0.0203}$ on latent/cummean) at \SI{10}{\percent} \emph{fewer} parameters
--- so removing the head makes the model measurably better in \emph{both} view
regimes. Mechanistically this fits the whole family arc: the objective is
cosine/InfoNCE retrieval \emph{in} the item-latent space, so the head's job is
to land \emph{in} that geometry, and an extra 256-wide non-linear layer between
the recurrent state and a metric read-out just adds a warp the loss must then
undo. The trade-off runs the other way on holisticness (both f0 arms have
\emph{lower} H, $\Delta$H CI below zero): the f1 arms' extra warp diffuses the
top-ten (higher ild, lower artist\_adj) at the cost of exact hits. The rule
adopted is \textbf{default \texttt{fusion\_layers}${}=0$}, later confirmed a
third time (Table~\ref{tab:nexttrack-fusion-depth}).

\emph{Holisticness and recall are anti-correlated, and the recall floor earned
its place.} Five arms lift holisticness@10 by paired-$\Delta$H CI above zero,
topping out at D-CC's \num{0.19715} ($+\num{0.055}$) --- which on the naked crown
metric would have ranked second on the entire board. But D-CC lands \num{95} of
\num{1431} exact hits against the baseline's \num{176}. The ordering of the batch
by H is almost the \emph{reverse} of its ordering by recall
(Figure~\ref{fig:nexttrack-holisticness-antiwin}), so \textbf{all five
$\Delta$H winners fail the pre-registered $\pm\num{0.015}$ recall floor} and the
crown rule never fired. Without that floor the campaign would have ``won the
crown'' with a model \SI{46}{\percent} worse at its actual job --- precisely the
degenerate \texttt{mood-session} failure mode (H \num{0.2332} at Recall@10
\num{0.0238}) the floor was written to block. On the present evidence,
\textbf{unconstrained holisticness is maximized by predicting less well}, so any
future holisticness objective on this corpus must be a \emph{constrained} one.

\begin{figure}[ht]
  \centering
  \includegraphics[width=0.8\textwidth]{nexttrack_holisticness_antiwin}
  \caption{Why the crown rule needs a hard recall floor. Across the ten arms,
  holisticness@10 runs almost exactly \emph{backwards} along Recall@10: the five
  arms that clear the $\Delta$H CI-above-zero leg (green) are the five furthest
  \emph{below} the baseline on exact recall, and all five fall outside the
  $\pm\num{0.015}$ floor (shaded). The open square is the incumbent
  \texttt{mood-session} model that motivated the floor --- the highest
  holisticness on record at \num{34} of \num{1431} correct next tracks. H values
  are the 3-factor generation (Section~\ref{sec:nexttrack-power}).}
  \label{fig:nexttrack-holisticness-antiwin}
\end{figure}

\emph{Interpretation.} Nothing was registered and nothing promoted; the champion,
the leaderboard and the pins are untouched. What the campaign buys is a
\emph{closure}: the combiner null is now robust across the whole combiner design
space on this corpus --- post-hoc score-level (three ways) \emph{and}
representation-level jointly-trained with a parameter-matched control --- which
materially consolidates ``the fixed equal-thirds z-blend is the combiner of
record.'' Two genuine sub-findings survive the null (the harmful MLP fusion head,
and \texttt{delta}$+$\texttt{cummean} complementarity), and the machinery itself
is sound: 10 of 10 runs succeeded, the new predictor trains, converges and
evaluates at exact parity with the family, and D-LL-f0 even posts the batch's
best MRR (\num{0.05631} against C0's \num{0.05341}) --- the one metric on which a
dual arm leads. Per the noise-band discipline that is ``at least equal, likely
neither better nor worse,'' not a win.

\paragraph{A provenance finding that revises two earlier numbers.} Before the
batch, the live remote worker directory was found \emph{stale}: its
\texttt{seq\_common.py} had no \texttt{holisticness\_at\_k} at all,
\texttt{seq\_nexttrack.py} differed from the hub by md5, and three
\texttt{seq\_*.py} modules were missing entirely. After refreshing and verifying
the whole family md5-identical to the hub, C0 returned \num{0.12299},
bit-identical to the hub's on-record value --- \emph{not} the \num{0.1201957}
that two earlier remote runs (\runid{run-20260723-002044-af11a},
\texttt{-17cf8}) had returned. The ``cross-host reduction-order offset of
$-\num{0.0028}$'' those runs had suggested is therefore \textbf{withdrawn: it was
a code difference, not a host effect}. Their recall numbers belong to a
superseded code state and are not code-comparable to later batches; they remain
valid as the $\beta$-eagerness measurements they were taken for
(Table~\ref{tab:nexttrack-eager-beta}), since those conclusions rest on within-batch
comparisons. The general lesson --- a divergent remote checkout produces a
\emph{wrong number}, not a crash --- is now a standing pre-flight gate.
\emph{Source: \texttt{2026-07-25-nexttrack-dual-tower-fusion-scan.md}.}

\section{Next-track: the per-step pre-encoder refuted, and the rectifier mechanism}
\label{sec:nexttrack-preencoder}
\emph{Rationale.} The dual-tower null closed \emph{deterministic} view
asymmetry: \texttt{delta} and \texttt{cummean} are lossy functions of what a
raw-latent tower already holds, so pairing them buys nothing. This campaign asks
the strictly stronger version --- does a \textbf{learned} per-step
re-embedding, which is \emph{not} a function of information the tower already
has, enlarge the hypothesis space usefully? The axis is a nonlinear
re-embedding in front of the recurrence,
$[\,\text{Linear}(d_{\text{in}},w),\ \text{ReLU},\ \text{Dropout}\,]\times
\texttt{pre\_layers}$, applied per timestep (strictly causal), with the GRU's
\texttt{input\_size} becoming $w$.

The hypothesis has an unusually sharp analytic edge. A GRU's own input transform
is \emph{already linear} ($W_i x$ per gate), and a \emph{linear} pre-encoder at
$w \ge d_{\text{in}}$ is exactly as expressive as a bare GRU, since
$W_i(Vx) = (W_iV)x$ with no rank constraint. \textbf{The ReLU is therefore the
only expressivity change}, which means no linear-pre control run is needed and a
null cannot be explained away as ``the pre-encoder was redundant by
construction.'' The design's decisive move is the \textbf{single-tower ablation}
S1 --- a \texttt{MLP256}$\to$\texttt{GRU256} against a parameter-matched plain
GRU --- because without it any dual-tower result is uninterpretable: pre-encoder,
or second tower?

Ten runs, same canonical split, same single remote worker, zero failures. Two
reproduction gates both passed \emph{bit-identically}: C0 returned Recall@10
\num{0.12299091544374564} (\num{176} of \num{1431}, identical to the last digit,
a seventh reproduction), confirming the \texttt{pre\_hidden} edit did not perturb
the family's default code path; and C3, the bare dual-tower arm, returned
\num{0.10272536687631029}, exactly its on-record value, validating every dual
verdict.

\emph{Results.} Table~\ref{tab:nexttrack-preencoder} is the batch and
Figure~\ref{fig:nexttrack-preencoder-forest} is the verdict at a glance. Not one
of the six pre-MLP arms beats the baseline by paired-$\Delta$ CI above zero:
five lose outright, one ties. \textbf{And the decisive ablation does not merely
tie --- it loses}: S1 at \num{0.10901} against C0's \num{0.12299}, $\Delta$
$-\num{0.0140}$ with the CI entirely below zero, \num{156} hits against
\num{176}. Since a linear pre-encoder here is provably expressivity-neutral,
\textbf{the ReLU re-embedding is destroying prefix information that the GRU's own
gated linear input transform preserved}.

\begin{table}[ht]
  \centering
  \caption{Pre-encoder presence $\times$ placement $\times$ dose, ten runs on
  \dsid{seq-20260715-131139} (\num{1431} test sessions), sorted by Recall@10.
  All arms share the prefix \protect\runid{run-20260726-161801-}; C0 is
  \protect\runid{run-20260726-161429-bdcbc}. Paired $\Delta$ is a per-session
  paired bootstrap against C0 (\num{2000} resamples, rng 1337); CI bounds are
  given to three decimals, and a CI straddling zero is a \emph{tie}. Dual-arm
  codes give the two towers' pre-encoder widths (A0/B256 $=$ tower A bare, tower
  B rectified at width 256) and the fusion depth. H here is the
  \emph{4-factor} (grounded) holisticness generation and is \textbf{not}
  comparable to Table~\ref{tab:nexttrack-dualtower}'s 3-factor column
  (Section~\ref{sec:nexttrack-power}). Best cell per metric bolded; \textbf{artadj@10
  lower-better}. \textsuperscript{\dag}~T4's minimal artadj and T1's maximal ild
  are the familiar de-eagering-by-predicting-worse trade: both sit
  $\sim$\num{0.025} below C0 on recall.}
  \label{tab:nexttrack-preencoder}
  {\scriptsize\setlength{\tabcolsep}{1.7pt}\begin{tabular}{llrrlccccc}
    \toprule
    Arm & id & params & R@10 (raw) & $\Delta$ vs C0 [95\% CI] & MRR & mus@10 & artadj\,$\downarrow$ & ild\,$\uparrow$ & H\,$\uparrow$ \\
    \midrule
    C2 h454 \emph{ctrl} & \texttt{7646d} & \num{969936} & \best{0.12509 (179)} & $+0.0021$ [$-.007$, $+.011$] & 0.05389 & 0.4572 & 0.3668 & 0.4466 & 0.045765 \\
    T1-f0 \emph{A0/B256, f0} & \texttt{a6423} & \num{888256} & 0.12369 (177) & $+0.0007$ [$-.012$, $+.013$] & \best{0.05470} & \best{0.4575} & 0.3627 & 0.4616 & 0.044541 \\
    C0 h256 \emph{base} & \texttt{bdcbc} & \num{394944} & 0.12299 (176) & --- & 0.05341 & 0.4558 & 0.3474 & 0.4536 & 0.047086 \\
    C1 h297 \emph{ctrl} & \texttt{7bab3} & \num{494697} & 0.11321 (162) & $-0.0098$ [$-.017$, $-.003$] & 0.05199 & 0.4513 & 0.3602 & 0.4474 & 0.045285 \\
    \best{S1 \emph{preMLP256}} & \texttt{ad53f} & \num{493504} & 0.10901 (156) & $\mathbf{-0.0140}$ [$-.027$, $-.001$] & 0.04463 & 0.4499 & 0.3240 & 0.4815 & \best{0.047132} \\
    C3 \emph{bare dual, f1} & \texttt{0df6f} & \num{871872} & 0.10273 (147) & $-0.0203$ [$-.035$, $-.007$] & 0.04235 & 0.4446 & 0.3092 & 0.4873 & 0.046612 \\
    T2 \emph{A0/B128, f1} & \texttt{dc7d2} & \num{847424} & 0.09993 (143) & $-0.0231$ [$-.037$, $-.009$] & 0.04490 & 0.4415 & 0.3039 & 0.4939 & 0.046431 \\
    T3 \emph{A0/B384, f1} & \texttt{ca0b8} & \num{1093440} & 0.09993 (143) & $-0.0231$ [$-.037$, $-.009$] & 0.04333 & 0.4441 & 0.3068 & 0.4961 & 0.045895 \\
    T4 \emph{A256/B256, f1} & \texttt{bd5b5} & \num{1068992} & 0.09783 (140) & $-0.0252$ [$-.039$, $-.011$] & 0.03829 & 0.4343 & \best{0.2845}\textsuperscript{\dag} & 0.4986 & 0.043985 \\
    T1 \emph{A0/B256, f1} & \texttt{9450b} & \num{970432} & 0.09713 (139) & $-0.0259$ [$-.040$, $-.013$] & 0.04176 & 0.4369 & 0.2962 & \best{0.5016}\textsuperscript{\dag} & 0.045769 \\
    \bottomrule
  \end{tabular}}
\end{table}

\begin{figure}[ht]
  \centering
  \includegraphics[width=0.8\textwidth]{nexttrack_preencoder_forest}
  \caption{The pre-encoder verdict. Paired per-session bootstrap $\Delta$
  Recall@10 against the reproduced baseline, one row per arm; the shaded strip is
  the $\pm\num{0.015}$ practical tie band. Not one pre-encoder arm reaches a CI
  entirely above zero. The decisive row is S1 --- the single tower, the only arm
  whose result is unambiguously attributable to the pre-encoder --- which loses
  outright. Because a \emph{linear} pre-encoder at this width is provably
  expressivity-neutral, the rectifier is the only candidate cause.}
  \label{fig:nexttrack-preencoder-forest}
\end{figure}

\emph{The dose curve is flat, and placement is irrelevant.} Sweeping
\texttt{pre\_hidden\_b} through $\{0,128,256,384\}$ at
\texttt{fusion\_layers}${}=1$ (Table~\ref{tab:nexttrack-preencoder-dose}), every adjacent
pair straddles zero and so do the endpoints --- a \textbf{flat} axis, which by
the pre-registered clause would have blocked registration even had a win rule
fired. But note the \emph{nominal} ordering is downward-then-flat, never upward:
every dosed arm sits below the bare dual. Taken with S1's significant loss, the
sign of this axis is consistently negative across both towers and all four
widths, so there is no interior peak to chase. Symmetry is likewise a
non-factor: T1 (asymmetric, one bare tower) against T4 (both towers rectified)
straddles zero on recall, so keeping tower A bare is not load-bearing at this
dose.

\begin{table}[ht]
  \centering
  \caption{The dose curve, as measured: paired $\Delta$ Recall@10 against the
  bare dual-tower arm C3 at \texttt{fusion\_layers}${}=1$. Flat by the
  pre-registered test --- but nominally downward, never upward.}
  \label{tab:nexttrack-preencoder-dose}
  {\footnotesize\setlength{\tabcolsep}{6pt}\begin{tabular}{rlll}
    \toprule
    \texttt{pre\_hidden\_b} & Arm & R@10 (raw) & $\Delta$ vs C3 (bare) \\
    \midrule
    0 & C3 & 0.10273 (147) & --- \\
    128 & T2 & 0.09993 (143) & $-0.0028$ [$-0.0147$, $+0.0091$] straddles \\
    256 & T1 & 0.09713 (139) & $-0.0056$ [$-0.0189$, $+0.0077$] straddles \\
    384 & T3 & 0.09993 (143) & $-0.0028$ [$-0.0154$, $+0.0105$] straddles \\
    \bottomrule
  \end{tabular}}
\end{table}

\emph{The one solid positive: the fusion-depth regression, measured a third
time and bigger.} T1-f0 beats T1 by $+\num{0.0266}$ with the CI above zero,
against the $+\num{0.016}$ and $+\num{0.020}$ measured in the two dual-tower view
regimes (Table~\ref{tab:nexttrack-fusion-depth}). The mechanism holds and sharpens: every
extra nonlinear layer between the recurrent state and the metric read-out adds a
warp the InfoNCE loss must undo, and \emph{the more nonlinearity already added on
the input side, the more the output-side warp costs}. Note too that T1-f0 --- the
only pre arm that ties the baseline --- is exactly the arm with \emph{both} a bare
tower A \emph{and} \texttt{fusion\_layers}${}=0$: the shortest un-warped path from
raw latents to the read-out. \texttt{fusion\_layers}${}=0$ is the family default,
full stop.

\begin{table}[ht]
  \centering
  \caption{The \texttt{fusion\_layers} regression, three independent
  measurements. Paired $\Delta$ Recall@10 of the linear-read arm minus its
  MLP-head twin; positive means removing the fusion MLP \emph{helps}. The
  $\Delta$H column is available only for the two same-batch 3-factor pairs.
  \emph{Sources: \texttt{2026-07-25-nexttrack-dual-tower-fusion-scan.md} (rows
  1--2), \texttt{2026-07-26-nexttrack-pre-encoder-scan.md} (row 3).}}
  \label{tab:nexttrack-fusion-depth}
  {\scriptsize\setlength{\tabcolsep}{4pt}\begin{tabular}{llll}
    \toprule
    Pair (f0 $-$ f1) & Regime & $\Delta$ R@10 [95\% CI] & $\Delta$H [95\% CI] \\
    \midrule
    D-LL-f0 $-$ D-LL & redundant views, no pre-encoder & $+0.0161$ [$+0.0021$, $+0.0301$] & $-0.0086$ [$-0.0116$, $-0.0055$] \\
    D-LC-f0 $-$ D-LC & complementary views, no pre-encoder & $+0.0203$ [$+0.0063$, $+0.0335$] & $-0.0135$ [$-0.0171$, $-0.0100$] \\
    \best{T1-f0 $-$ T1} & \best{with a per-step pre-encoder} & \best{$+0.0266$ [$+0.0140$, $+0.0391$]} & --- \\
    \bottomrule
  \end{tabular}}
\end{table}

\emph{The mechanism, measured rather than inferred.} The analytic argument named
the rectifier as the only candidate; a follow-up interpretability read then
\emph{measured} it (Figure~\ref{fig:nexttrack-rectifier}). The decisive control
was a counterfactual tap the sparse-autoencoder engine cannot express --- the
model's \emph{own} affine map with the ReLU deleted. Reading next-item genre
decodability at the same step (\num{11184} probe test rows, 25 classes, paired
standard error $\le\num{0.0065}$): the raw latent scores \num{0.4033}, the
affine map \num{0.4084} (\textbf{neutral, ${\sim}0.8\sigma$}, exactly as
$W_i(Vx)=(W_iV)x$ predicts), and the rectified tap \num{0.3700}
(\textbf{${\sim}5.9\sigma$ --- the ReLU is the whole loss}). The smoking gun is
arithmetic: \texttt{frac\_exact\_zero} on the rectified tap is \num{0.5327},
matching the affine tap's \texttt{frac\_negative\_coords} \num{0.5327} to four
digits. The rectifier deletes \emph{precisely} the negative half, leaving
${\sim}120$ live units for 192 signed directions --- \emph{fewer live coordinates
than the input has dimensions}. PCA-192 latents are zero-centred and
near-symmetric by construction, and at $w=256$ (only $1.33\times$ the input
dimension) there is nowhere near enough over-completeness to re-encode the
discarded half-space, unlike the wide embeddings where ReLU MLPs normally live.
The damage survives the recurrence: within one model, \texttt{tower\_b} (fed
rectified) scores \num{0.3929} against its \emph{identical twin}
\texttt{tower\_a} (fed raw) at \num{0.4183}, ${\sim}3.9\sigma$ --- and
\texttt{tower\_a} reproduces the standalone plain GRU (\num{0.4235}), which is
exactly what ``the bare tower \emph{rescues}'' requires.

\begin{figure}[ht]
  \centering
  \includegraphics[width=0.8\textwidth]{nexttrack_rectifier}
  \caption{The rectifier mechanism, measured. Next-item genre decodability at
  each representation tap, including one counterfactual the interpretability
  engine cannot express: the model's own affine map with the ReLU deleted. On the
  input side the affine map is \emph{free} (${\sim}0.8\sigma$) and the ReLU is
  the \emph{entire} cost (${\sim}5.9\sigma$); after the recurrence the damage
  persists (${\sim}3.9\sigma$ between two identical towers that differ only in
  whether their input was rectified), while the raw-fed tower reproduces the
  standalone plain GRU. \texttt{frac\_exact\_zero} on the rectified tap
  (\num{0.5327}) matches the affine tap's \texttt{frac\_negative\_coords}
  (\num{0.5327}) to four digits: half-wave rectification of a zero-centred
  latent deletes exactly the negative half, leaving ${\sim}120$ live units for
  192 signed directions.}
  \label{fig:nexttrack-rectifier}
\end{figure}

\emph{Interpretation, and one honest caveat about the controls.} This is the
fifth consecutive architecture null, and it closes the \emph{input side} as the
view null closed the input-view direction: together with the four
learned-combiner nulls, adding representation machinery \emph{around} a single
GRU over a good latent space does not help on this corpus, whereas every real
gain on record came from improving the \emph{latent space itself} (the InfoNCE
content projection, PCA-192) or from the fixed z-blend. The axis is now
quantitatively bounded --- best pre arm \num{0.12369}, a tie; dose flat to twice
the input dimension; placement irrelevant; the fusion-depth interaction
resolved. The caveat is that the \emph{parameter-matched width control was a
weaker instrument than assumed}: C1 (h297) came in at \num{0.11321}, itself
losing to C0, and the batch's width points run \num{176} $\to$ \num{162} $\to$
\num{173} $\to$ \num{179} hits at h256 / h297 / h425 / h454 --- non-monotone,
spanning \num{0.0119} in recall with no trend. ``Single-GRU width is flat on
PCA-192'' survives as a \emph{trend} claim (h454 ties h256), but per-width jitter
is ${\sim}\pm\num{0.01}$ on a single split, roughly two-thirds of the noise band.
\textbf{A width-matched control that is itself a single unreplicated run is not a
reliable capacity anchor at this effect size} --- which is why it matters that
S1's verdict rests on the C0 comparison, reproduced to the last digit. A second
caveat is un-disentangled: the pre-MLP ends in \texttt{Dropout}, so pre arms
receive GRU-\emph{input} dropout that the bare arms lack, and S1's loss could be
partly over-regularization rather than the rectifier --- separating them needs a
\texttt{pre\_dropout} parameter that does not yet exist.
\emph{Source: \texttt{2026-07-26-nexttrack-pre-encoder-scan.md}.}

\section{Next-track: reading the record --- statistical power and the holisticness generations}
\label{sec:nexttrack-power}
Two cross-cutting facts emerged from the 2026-07-25/26 batches that change how
every number above should be read. Neither is a result about a model; both are
results about the \emph{measuring instrument}, and they belong here rather than
buried in one campaign.

\subsection{The nulls are power-limited, and that is a different claim}
\label{sec:nexttrack-power-stats}
Every next-track verdict in this document rests on a paired per-session
bootstrap over the
same \num{1431} test sessions. That $n$ has a measurable resolution. Averaged
over the nine paired comparisons of the pre-encoder batch the 95\% CI half-width
is \textbf{\num{0.0124}} on Recall@10 (\num{0.0128} over the dual-tower batch's
nine, so the figure is stable), and \textbf{${\sim}\num{0.0018}$} on
holisticness@10. Now compare those to the \emph{entire spread} of an architecture
family: \num{0.028} on Recall@10 and \num{0.0031} on holisticness@10
(Table~\ref{tab:nexttrack-power}). The ratio is the whole story --- a family's full range
is only ${\sim}2.3$ resolution widths wide on recall and ${\sim}1.9$ on
holisticness, so of eight adjacent arms roughly \textbf{two tiers} are separable
on Recall@10 and \textbf{one} on holisticness@10. Resolving a \num{0.005}
Recall@10 difference at this precision would require
$n \approx 1431\cdot(\num{0.0124}/\num{0.005})^2 \approx \num{8800}$ test
sessions, six times what the split provides
(Figure~\ref{fig:nexttrack-power}).

\begin{table}[ht]
  \centering
  \caption{Measured resolution of the test split. Half-widths are means over the
  nine paired-bootstrap comparisons of the pre-encoder batch (\num{2000}
  resamples, $n=\num{1431}$); the family spread is that batch's full range on
  each metric. ``Tiers'' is spread\,$/$\,half-width, i.e.\ how many genuinely
  distinguishable levels eight adjacent arms can occupy.
  \emph{Sources: \texttt{2026-07-25-nexttrack-dual-tower-fusion-scan.md},
  \texttt{2026-07-26-nexttrack-pre-encoder-scan.md}.}}
  \label{tab:nexttrack-power}
  {\footnotesize\setlength{\tabcolsep}{6pt}\begin{tabular}{lcccc}
    \toprule
    Metric & CI half-width at $n=\num{1431}$ & family spread & tiers & $n$ for $\Delta=\num{0.005}$ \\
    \midrule
    Recall@10 & 0.0124 & 0.0280 & $\sim$2 of 8 & $\sim$\num{8800} \\
    holisticness@10 & 0.0018 & 0.0031 & $\sim$1 of 8 & --- \\
    \bottomrule
  \end{tabular}}
\end{table}

\begin{figure}[ht]
  \centering
  \includegraphics[width=0.8\textwidth]{nexttrack_power}
  \caption{The architecture nulls are power-limited. The paired-bootstrap CI
  half-width falls as $1/\sqrt{n}$; at the split's $n=\num{1431}$ it is
  \num{0.0124} on Recall@10, while a whole architecture family spans only
  \num{0.028}. Resolving a \num{0.005} difference would need
  ${\sim}\num{8800}$ test sessions. The honest reading of the five architecture
  nulls follows directly: \emph{large harms were detected reliably, moderate wins
  were never detectable at all}.}
  \label{fig:nexttrack-power}
\end{figure}

\textbf{The honest reading of the five architecture nulls follows directly, and
it is not the reading a casual glance gives.} What these campaigns established
firmly is that the tested architectures do not deliver \emph{large} gains, and
that several of them do real \emph{harm} --- the harms are the well-powered
results, because they are big (S1 $-\num{0.0140}$; the \texttt{fusion\_layers}
regression $+\num{0.0266}$; six of eight dual arms below the baseline). What the
campaigns \emph{cannot} establish is the absence of a moderate gain: a genuine
$+\num{0.005}$ improvement --- which on a \num{0.123} baseline would be a
respectable \SI{4}{\percent} relative lift --- would have been invisible by
construction. So ``dual-tower fusion is refuted'' and ``the pre-encoder is
refuted'' should be read as \emph{no large effect, and no effect of the
hypothesized magnitude}, not as ``these architectures do not work.'' The
practical consequence is a prioritization rule rather than a caveat: on this
split, further architecture arms are a poor use of runs no matter how they come
out, because the design space's whole dynamic range is barely twice the
measurement error. Either enlarge $n$ or move to axes whose effects are large ---
which is, empirically, exactly what the representation and supervision axes have
delivered (the content projection's $+\num{0.0259}$ is twice the half-width, and
the second-split hardening's $+\num{0.0426}$ is three times it).

\subsection{Three generations of holisticness@10, and why the board is denominated in recall}
\label{sec:nexttrack-power-generations}
The composite \emph{holisticness@10} was revised twice in four days, and stored
values are \textbf{not comparable across the cuts}:
\begin{description}
  \item[Gen.\ 1 --- 3-factor (to 2026-07-24).] The bare product
  \[\mathrm{clamp}(\text{mood\_coh},0)\times\text{ild}\times(1-\text{artist\_adj}),\]
  carrying \emph{no correctness term at all}. Values land in
  \numrange{0.063}{0.233} --- the column in Table~\ref{tab:nexttrack-dualtower}.
  \item[Gen.\ 2 --- 4-factor, grounded (2026-07-25).] Multiplies in a
  \texttt{music@10} grounding factor,
  \[\mathrm{clamp}\bigl((\text{music@10}-\num{0.185})/\num{0.815},\,0,\,1\bigr).\]
  Because \texttt{music@10} falls in lockstep with recall, this cancels the
  de-eagering ``wins'' the 3-factor form rewarded; values land in
  \numrange{0.044}{0.050} --- the column in Table~\ref{tab:nexttrack-preencoder}.
  \item[Gen.\ 3 --- concentration-aware (2026-07-26).] Eagerness becomes
  \texttt{max(artist\_adj,} \texttt{artist\_conc)} and the grounding floor
  rises \num{0.185}$\to$\num{0.34}. Values land near \numrange{0.009}{0.025}.
\end{description}
Every facet is stored per row, so any generation is recomputable after the fact
(the helper \texttt{rescore\_holisticness.py} does this for the runs it can
reach),
and Table~\ref{tab:nexttrack-hgen} re-derives all three for a representative slice. Two
consequences matter for reading this document. First, \textbf{a holisticness
column must always be labelled with its generation}, because the generation
effect dwarfs the model effect: every 3-factor value mechanically outranks every
4-factor value, since the grounding factor multiplies by ${\sim}0.3$. Second,
and more importantly, \textbf{the recall and holisticness boards point in
opposite directions at every generation}. The champion blend --- \num{0.21174}
Recall@10, \num{303} of \num{1431} --- ranks \emph{last} in
Table~\ref{tab:nexttrack-hgen} under all three formulas, because it is the most
artist-eager model on record; the degenerate \texttt{cummean/cummean} arm and
\texttt{mood-session} lead. Generation 3 demoted \texttt{mood-session} from
first to a genuine zero, but it did \emph{not} fix the dual-tower degenerate arm,
so the $\pm\num{0.015}$ recall floor remains mandatory on any holisticness claim.

\begin{table}[ht]
  \centering
  \caption{The same models under all three holisticness generations, sorted by
  Recall@10. The generation effect dwarfs the model effect, and the ordering is
  roughly \emph{inverse} to Recall@10 throughout --- which is why the next-track
  leaderboard stays denominated in Recall@10. Dashes are runs whose facets could
  not be re-derived under that generation (remote-only runs with no materialized
  predictions file). \emph{Sources:
  \texttt{2026-07-25-nexttrack-dual-tower-fusion-scan.md},
  \texttt{2026-07-26-nexttrack-pre-encoder-scan.md}.}}
  \label{tab:nexttrack-hgen}
  {\footnotesize\setlength{\tabcolsep}{2.4pt}\begin{tabular}{llccc}
    \toprule
    Model (run) & R@10 (raw) & gen-1 (3-f) & gen-2 (4-f) & gen-3 (current) \\
    \midrule
    \best{\Rp+M+\Cp blend} (\texttt{1c1a3}) & \best{0.21174 (303)} & 0.0632 & 0.020830 & 0.00860 \\
    single GRU h454 (\texttt{7646d}) & 0.12509 (179) & --- & 0.045765 & 0.02399 \\
    single GRU h256 (\texttt{97c1e}, \texttt{bdcbc}) & 0.12299 (176) & 0.14172 & 0.047086 & \best{0.02454} \\
    single GRU h425 (\texttt{be11b}) & 0.12089 (173) & 0.13850 & 0.046677 & 0.02487 \\
    S1 pre-MLP$\to$GRU (\texttt{ad53f}) & 0.10901 (156) & 0.14499 & 0.047132 & 0.02410 \\
    D-CC \emph{cum/cum} (\texttt{88a4c}) & 0.06639 (95) & 0.19715 & \best{0.050436} & --- \\
    \texttt{mood-session} (degen.) & 0.02380 (34) & \best{0.23319} & 0.043881 & 0.00000 \\
    \bottomrule
  \end{tabular}}
\end{table}

\paragraph{What is primary, and what this document sorts on.} For completeness:
the \texttt{[metrics].primary} key in \texttt{domain.toml}, and the best-models
group's \texttt{primary\_metric}, have both read \texttt{holisticness@10} since
2026-07-23, so the \emph{server-effective} primary metric is no longer
Recall@10. Statements
in the earlier next-track sections that ``the primary metric stays Recall@10''
were true when written and are superseded \emph{for the server}. This document
nevertheless continues to denominate the next-track leaderboard
(Table~\ref{tab:nexttrack-lb}) and every campaign verdict in \textbf{Recall@10},
deliberately: it is the record's comparability currency, it is stable across all
three metric generations, and --- as Table~\ref{tab:nexttrack-hgen} shows --- a
holisticness ordering would place the servable champion last. Holisticness is
reported throughout as a labelled, generation-tagged diagnostic column.
\emph{Sources: \texttt{2026-07-25-nexttrack-dual-tower-fusion-scan.md},
\texttt{2026-07-26-nexttrack-pre-encoder-scan.md}.}

\section{Next-track: the MMR $\lambda$ frontier --- the first confirmed crown lever}
\label{sec:nexttrack-mmr}
\emph{Rationale.} Section~\ref{sec:nexttrack-holisticness} left one standing flag: an
eval-time MMR re-rank at $\lambda\approx\num{0.9}$ was ``the only near-free
holisticness lever on record'', measured in an \emph{offline} sweep over a stored
checkpoint and explicitly deferred to a user decision. This campaign is the
native confirm, and it was designed to settle a second, pre-registered
\emph{mechanism} question rather than merely to re-measure the knob: does sonic
diversification de-concentrate \textbf{artists} (H1, ``responsive''), or are
same-artist tracks far enough apart in the content-metric space that the
diversity penalty leaves eagerness untouched (H2, ``sticky'')? The answer decides
whether an explicit artist-diversity penalty is the indicated sequel or a
redundant one.

\paragraph{What is being maximized.} All ten arms are scored under the
\emph{live} four-factor crown, taken verbatim from
\texttt{predictors/seq\_common.py}:
\[
  \begin{aligned}
    H = {}& \max(\texttt{mood\_coh},0)\;\times\;\texttt{ild}\\
          &\times\;\bigl(1-\mathrm{clamp}(\max(\texttt{artist\_adj},
             \texttt{artist\_conc}),0,1)\bigr)\\
          &\times\;\mathrm{clamp}\bigl((\texttt{music@10}-\num{0.34})/\num{0.66},
             \,0,\,1\bigr),
  \end{aligned}
\]
each facet the \emph{mean over rows} and the product taken \emph{of the
aggregates} --- generation 3 of Section~\ref{sec:nexttrack-power-generations}. Two anchors
frame the batch: \textbf{A}, the recall champion \Rp+M+\Cp\ (H \num{0.008603}, its
eagerness factor bound by \texttt{artist\_conc} \num{0.691653}), and \textbf{B},
the best \emph{single} model on the live crown, \texttt{gru-infonce-h256} (H
\num{0.024538}). The bar to beat is the all-time healthy-arm crown maximum
$H_{\mathrm{bar}} = \num{0.024872}$, held by an h425 capacity control
(\runid{run-20260725-143502-be11b-seq-nexttrack}). All arms ride the canonical
split \dsid{seq-20260715-131139}; no dataset work was needed.

\paragraph{Why this batch is unusually sharp.} $\lambda$ is an \emph{eval-only}
parameter: the trained weights are bit-identical across the arms of one model, so
$\lambda=\num{1.0}$ is a bit-identity gate and every $\Delta$ is a maximally
paired comparison. Both controls reproduced their anchors exactly (\num{303} and
\num{176} hits of \num{1431}, MRR agreeing to \num{3.1e-11}), and every arm's $H$
was \emph{recomputed} from the stored \texttt{predictions.json} under one code
path, so the per-row rounding bias cancels in every $\Delta$ (recomputed versus
stored agreed to \num{7.4e-8} and \num{2.1e-7}). A third gate --- MMR-active ---
required each $\lambda<1$ arm's \texttt{ild} to differ from its own control by far
more than \num{1e-6}, which is the registered detector for MMR silently no-opping
when the content-metric index is unreachable. No arm was a no-op.

\emph{Results.} Table~\ref{tab:nexttrack-mmr-arms} is the frontier, and
Figure~\ref{fig:nexttrack-mmr-frontier} plots it. Arms are tiered by a
pre-registered rule: \textbf{TIER-1 FREE} means $\Delta H$ above zero with the
recall $\Delta$ straddling zero and $|\Delta\text{recall}| \ll \num{0.015}$;
\textbf{TIER-2 PRICED} means a real $H$ gain at a recall cost the CI can see;
\textbf{TIER-3 REJECTED} means $\Delta$recall below $-\num{0.03}$.

\begin{table}[ht]
  \centering
  \caption{The MMR $\lambda$ frontier on both anchors, canonical split
  \dsid{seq-20260715-131139} ($n=\num{1431}$), sorted by holisticness@10. Every
  run id is \protect\runid{run-20260726-} plus a launch time and hex suffix, in
  arm order: \texttt{184115-bf395} (B3), \texttt{184115-1dc77} (A4),
  \texttt{184116-20d9b} (B4), \texttt{184115-147a6} (B2),
  \texttt{184115-e1539} (B1), \texttt{183333-9a0b2} (B0),
  \texttt{184115-7a2f7} (A3), \texttt{184115-36edb} (A2),
  \texttt{184115-6cb38} (A1), \texttt{183333-1c361} (A0); the A arms are
  \texttt{seq-blend}, the B arms \texttt{seq-nexttrack}. Every $\Delta$ is a
  paired bootstrap (\num{2000} resamples, rng \num{1337}) against the arm's
  \emph{own} $\lambda=\num{1.0}$ control, with all five facets re-aggregated on
  the same resampled index set before the product; CI bounds are given to three
  decimals for width. Tier codes: \emph{t1} TIER-1 FREE, \emph{t2} TIER-2 PRICED,
  \emph{t3} TIER-3 REJECTED, \emph{ctrl} the arm's own $\lambda=\num{1.0}$
  control. The binding facet --- the one entering the eagerness factor through the
  $\max()$ --- is \texttt{conc} on A0/A1/A2, B0 and \texttt{adj} everywhere else.
  H@10, R@10, \texttt{mus} ($=$ music@10) and \texttt{ild}
  higher-better; \texttt{adj}, \texttt{conc} lower-better. Best cell per metric
  bolded; the promotable arm B4 is bolded in full. \texttt{mood\_coh} is omitted
  for width and read in the interpretation: B0$\to$B3 \num{0.4788}$\to$\num{0.4585}
  ($-\SI{4}{\percent}$) against A0$\to$A4 \num{0.3982}$\to$\num{0.2378}
  ($-\SI{40}{\percent}$), the champion's collapse under diversification.
  \emph{Source: \texttt{2026-07-26-nexttrack-mmr-lambda-frontier.md}.}}
  \label{tab:nexttrack-mmr-arms}
  {\scriptsize\setlength{\tabcolsep}{1.9pt}\begin{tabular}{lllllllll}
    \toprule
    Arm & $\lambda$/pool & H@10 & $\Delta H$ [95\% CI] & R@10 (hits) & $\Delta$R@10 [95\% CI] & mus & ild & adj / conc \\
    \midrule
    B3 \emph{t2} & 0.7 / 200 & \best{0.034171} & $+.0096$ [$+.008$, $+.011$] & 0.11461 (164) & $-.0085$ [$-.016$, $-.001$] & 0.4747 & 0.5314 & 0.313 / 0.271 \\
    A4 \emph{t3} & 0.3 / 200 & 0.031626 & $+.0230$ [$+.021$, $+.026$] & 0.10971 (157) & $-.1019$ [$-.120$, $-.084$] & \best{0.4925} & \best{0.8146} & 0.293 / \best{0.190} \\
    \best{B4} \emph{t1} & \best{0.7 / 50} & \best{0.030355} & \best{$+.0058$ [$+.005$, $+.007$]} & \best{0.12020 (172)} & \best{$-.0029$ [$-.008$, $+.003$]} & 0.4670 & 0.5017 & 0.327 / 0.302 \\
    B2 \emph{t1} & 0.8 / 200 & 0.029458 & $+.0049$ [$+.004$, $+.006$] & 0.11950 (171) & $-.0036$ [$-.009$, $+.001$] & 0.4648 & 0.4961 & 0.330 / 0.307 \\
    B1 \emph{t1} & 0.9 / 200 & 0.026761 & $+.0022$ [$+.002$, $+.003$] & 0.12089 (173) & $-.0021$ [$-.006$, $+.002$] & 0.4597 & 0.4716 & 0.341 / 0.337 \\
    B0 \emph{ctrl} & 1.0 / 200 & 0.024539 & --- & 0.12299 (176) & --- & 0.4558 & 0.4536 & 0.347 / 0.356 \\
    A3 \emph{t3} & 0.5 / 200 & 0.023576 & $+.0150$ [$+.013$, $+.017$] & 0.16422 (235) & $-.0473$ [$-.061$, $-.033$] & 0.4871 & 0.6527 & 0.486 / 0.461 \\
    A2 \emph{t2} & 0.7 / 200 & 0.015289 & $+.0067$ [$+.006$, $+.008$] & 0.19357 (277) & $-.0181$ [$-.029$, $-.008$] & 0.4759 & 0.5201 & 0.571 / 0.608 \\
    A1 \emph{t1} & 0.9 / 200 & 0.010435 & $+.0018$ [$+.001$, $+.002$] & 0.21034 (301) & $-.0013$ [$-.008$, $+.005$] & 0.4625 & 0.4362 & 0.603 / 0.671 \\
    A0 \emph{ctrl} & 1.0 / 200 & 0.008603 & --- & \best{0.21174 (303)} & --- & 0.4538 & 0.4064 & 0.610 / 0.692 \\
    \bottomrule
  \end{tabular}}
\end{table}

\begin{figure}[ht]
  \centering
  \includegraphics[width=0.95\textwidth]{nexttrack_mmr_frontier}
  \caption{\textbf{Left:} the $\lambda$ frontier on the recall $\times$ crown
  plane. Each engine traces a path as $\lambda$ falls from \num{1.0}; colour is
  the pre-registered tier (green TIER-1 FREE, blue TIER-2 PRICED, red TIER-3
  REJECTED, grey control). Three single-GRU arms clear the prior all-time crown
  bar \num{0.024872} at a recall whose paired CI straddles zero --- the first
  genuinely free crown movement on record. The champion's path reaches comparable
  $H$ only after paying recall the CI can see, and its two best-$H$ points are
  rejected as degenerate. \textbf{Right:} the mechanism question. H1 required
  $|\Delta\texttt{artist\_conc}|\ge|\Delta\texttt{artist\_adj}|$; H2 required the
  ratio below \num{0.5} (shaded). Measured \numrange{1.59}{3.06} in \emph{8 of 8}
  arms, so H2 is refuted by a factor of \numrange{3}{6} in the wrong direction ---
  and the ratio shrinks monotonically as $\lambda$ falls, i.e.\ the easy
  de-concentration is bought first and the mechanism saturates.}
  \label{fig:nexttrack-mmr-frontier}
\end{figure}

\emph{Interpretation.} \textbf{MMR is a real crown lever, and the first one on
record.} Every one of the eight $\lambda<1$ arms lifts $H$ with the CI entirely
above zero. What makes the gain legitimate rather than an artefact of the
composite is that \emph{three} of the four factors improve at once: \texttt{ild}
rises (\num{0.4536}$\to$\num{0.5314} on the GRU, \num{0.4064}$\to$\num{0.8146} on
the champion), eagerness falls, and --- the load-bearing part ---
\texttt{music@10} \emph{also} rises (\num{0.4558}$\to$\num{0.4747}). A lever that
raised $H$ by degrading relevance would be cancelled by the grounding factor; this
one is not, which is precisely the distinction the floor was written to enforce
after the \texttt{cummean/cummean} and \texttt{seq-mood} pathologies.

\textbf{H1 is confirmed unanimously and H2 is refuted, which \emph{cancels} the
sequel the previous campaign had pre-registered.} In 8 of 8 arms
$\Delta\texttt{artist\_conc} \le \Delta\texttt{artist\_adj}$ (both negative) at a
ratio of \numrange{1.59}{3.06}. Same-artist tracks \emph{are} near-neighbours in
the content-metric space, so a max-similarity penalty de-concentrates artists as a
side effect, with no explicit notion of artist anywhere in it. Sonic MMR
therefore \emph{is} a de-eagering lever, and porting \texttt{seq-mood}'s explicit
\texttt{artist\_penalty} to the blend --- the pre-registered H2 follow-up --- is no
longer the indicated next move. It stays on file as a \emph{complement}, because
\texttt{artist\_adj} falls ${\sim}\num{2.5}\times$ slower than
\texttt{artist\_conc} and becomes the binding facet at low $\lambda$.

\textbf{The binding facet switches, and that locates the end of the lever.} On the
champion the switch from \texttt{artist\_conc} to \texttt{artist\_adj} happens
between $\lambda\num{0.7}$ and $\lambda\num{0.5}$; on the GRU it has already
happened by $\lambda\num{0.9}$. Once \texttt{artist\_adj} binds, the
$(1-\text{eager})$ factor is the stiff one and each further recall point buys
progressively less crown --- the structural reason the GRU's frontier flattens
where it does, and a prediction that pushing below ${\sim}\lambda\num{0.7}$ on the
GRU is poor value. B3 is the confirmation: the highest $H$ measured anywhere
(\num{0.034171}) at the cost of losing the recall tie.

\textbf{The promotable operating point is the pool probe.} \textbf{B4}
($\lambda\num{0.7}$, \texttt{mmr\_pool}~\num{50}) is the single best TIER-1 arm at
$H$~\num{0.030355} --- $+\SI{22.1}{\percent}$ over $H_{\mathrm{bar}}$ and
\num{5.4} measured half-widths clear of it --- and it beats B2 on $H$ by
$+\num{0.000889}$ with the CI above zero at a recall that straddles, so its rank
is \emph{resolved} rather than nominal. B2 ($\lambda\num{0.8}$) was the design's
pre-registered promotion hypothesis and fired almost exactly as projected
($\num{0.0321}$ predicted against \num{0.029458} measured, with a \emph{smaller}
recall cost than projected). Nothing was promoted: six definitions were
registered, each carrying the note ``NOT PROMOTED --- robustness confirm $+$ user
approval pending''. The recall board is untouched --- \textbf{no arm improves
recall}; MMR only ever costs it.

\textbf{The second axis, \texttt{mmr\_pool}, closes as a fine-tuning knob rather
than an axis.} Capping the candidate pool at \num{50} recovers recall
($\Delta$ vs pool~\num{200} $+\num{0.005584}$, CI above zero, as hypothesized ---
a cap bounds how far MMR can reach for diversity) but gives $H$ back
($-\num{0.003838}$, CI below zero). It does not win at fixed $\lambda$; it moves
along the \emph{same} frontier. What earns it a line in the record is the
cross-comparison above: the cap reaches a slightly \emph{better} point than the
nearest pure-$\lambda$ setting. Prefer $\lambda$ for coarse movement, pool for
fine-tuning at a chosen $\lambda$.

\textbf{The champion/GRU contrast resolves in favour of factor health, not
de-eagering headroom.} The design framed this as the campaign-carrying question:
the champion is $\num{1.9}\times$ more eager, so a given absolute de-eagering is
worth ${\sim}\num{2.1}\times$ more relative $H$ there --- yet it starts
$\num{2.85}\times$ behind, and the gap \emph{widens} under diversification
($H$~\num{0.015289} against \num{0.034171} at $\lambda\num{0.7}$). The facets say
why: the champion's \texttt{mood\_coh} \textbf{collapses} as it diversifies
(\num{0.3982}$\to$\num{0.2378} at $\lambda\num{0.3}$, $-\SI{40}{\percent}$) while
the GRU's barely moves ($-\SI{4}{\percent}$). The champion's recall comes from
artist-adjacency --- the Markov bigram leg, per
Section~\ref{sec:nexttrack-holisticness} --- so diversifying away from the seed artist also
walks it out of the prefix's mood neighbourhood. It pays twice. \textbf{MMR's
value scales with the health of the other three factors, not with room to
de-eager}, and no $\lambda$ closes the champion's crown gap.

\paragraph{Two instrument findings, one of them a warning about the record
itself.} First, \textbf{the shared-weights instrument is
$\numrange{3}{33}\times$ sharper than the architecture-family noise floor}:
measured half-widths run \numrange{0.000382}{0.002476} on $\Delta H$ against the
family's conservative \num{0.0018}, with the number of sessions whose hit@10
changed at all as low as \textbf{\num{9}}. Arm A1 had been pre-registered as
``likely UNRESOLVABLE'' and resolved at \num{4.8} measured half-widths.
Eval-time levers must therefore budget their minimum detectable effect from the
\emph{pairing}, not from the architecture figure of
Section~\ref{sec:nexttrack-power-stats} --- the two regimes differ by more than an order of
magnitude, and conflating them under-powers a whole class of cheap experiments.
The corollary is reassuring in the other direction: the 2026-07-22 \emph{offline}
$\lambda$ sweep projected recall \num{0.2103} / \num{0.1936} / \num{0.1642} /
\num{0.1097} and the native runs returned \num{0.21034} / \num{0.19357} /
\num{0.16422} / \num{0.10971} --- four for four to four decimals, which validates
the offline driver as a \emph{screening} instrument.

Second, a new infrastructure pitfall was found the hard way and it retroactively
qualifies part of the record: \textbf{a remote worker running a
\texttt{(deleted)} binary image silently strips metric keys its struct does not
know}. The metric-completeness gate failed on the first attempt with
\texttt{artist\_conc\_at\_k} missing, although the predictors were md5-identical
to the hub and the worker's own \texttt{metrics.json} carried the field: the
worker \emph{daemon} had started a day before the field existed, \texttt{rsync}
had replaced the file on disk but not the running image, and the old image
deserialized the predictor's metrics into a struct lacking the field, dropped it,
and posted the filtered set \textbf{without erroring}. \texttt{holisticness@10}
was nevertheless correct throughout, because the predictor computes $H$ itself
before the server sees it --- which is exactly what made the defect invisible.
Every remote run in a four-hour window on 2026-07-26 lost
\texttt{artist\_conc\_at\_k}, which is why anchor B's quoted concentration was a
\emph{reconstruction} from \texttt{top\_k\_ids} and the dataset artist map. That
reconstruction is now independently validated: the natively emitted values
(\num{0.6916530786551751} and \num{0.35591272614333413}) match the rescoring
tool's reconstructions to the last digit. Remedy: cycle the daemon whenever a
metric field is added; \texttt{rsync} of the binary is not sufficient.
\emph{Source: \texttt{2026-07-26-nexttrack-mmr-lambda-frontier.md}.}

\section{Next-track: autoregressive rollouts: a generation metric, and five sweeps that say no}
\label{sec:nexttrack-rollout}
\emph{Rationale.} This campaign was provoked by a \emph{user preference}, not by a
metric. On a real 135-track playlist the user strongly preferred one engine's
20-step journey over another's, while \texttt{holisticness@10} scored those two
journeys \num{0.003} apart. That gap is the point: every metric on record ---
Recall@10, holisticness@10 and its four facets, \texttt{music@10}, even the
multi-step continuation evaluator --- scores \emph{one} ranked list produced from a
\emph{real} prefix. None of them ever feeds a model its own output. A 20-step
journey does exactly that, so a trajectory failure --- a rollout whose energy
decays until a party seed has become a comedown --- is structurally invisible to
all of them. This is the first campaign in this document to build an instrument
for the \emph{deployment} regime, and it is the direct ancestor of the walk
surface of Section~\ref{sec:nexttrack-artistcap}.

\paragraph{What was built.} Four things, all off the training path.
\begin{description}
  \item[\texttt{hold@N}, a generation metric.] Per legible axis $a$ in
  \{energy, valence, acousticness, tempo\},
  $\texttt{fade}_a = |\mathrm{mean}(\text{2nd half}) -
  \mathrm{mean}(\text{1st half})|/\texttt{band}_a$, and then
  \[
    \texttt{hold@}N = 1 - \mathrm{clamp}\bigl(\mathrm{mean}_a(\texttt{fade}_a)
      + \texttt{stasis},\,0,\,1\bigr).
  \]
  The per-axis bands are calibrated on \num{2135} real
  \emph{train} sessions (energy fade p10 \num{0.0141}, p50 \num{0.0769}, p90
  \num{0.1963}; tempo \num{1.28}, \num{8.05}, \num{22.43}), so axes with
  incompatible units combine and ``a big fade'' means big \emph{relative to how
  this listener actually listens}. Real sessions are not flat --- the median energy
  fade is \num{0.077} --- so a metric rewarding zero fade would rank an
  unnaturally rigid generator above the listener it imitates.
  \item[\texttt{stasis}, the anti-gaming term.] A pure fade metric is
  degenerate: twenty near-identical tracks fade by \num{0} and score \num{1.0},
  which is the
  same pathology that forces \texttt{holisticness@10} to carry an \emph{external}
  recall floor. \texttt{stasis} puts the floor \emph{inside} the metric by
  penalizing steps smaller than the listener's own p10 jump, and it was verified
  on a synthetic flatline (moves-naturally \num{1.000}, at-your-floor \num{1.000},
  half-floor \num{0.500}, barely-moves \num{0.048}).
  \item[An offline rollout evaluator.] \texttt{seq\_continuation\_eval.py
  --rollout} scores sampled \emph{test} sessions by \textbf{importing} the
  generation loop from the serving path rather than reimplementing it, so an
  offline score describes the journeys the lab actually serves; the refactor was
  verified byte-identical. It is on-demand and sampled because a rollout costs
  \texttt{steps} full-vocabulary passes per session against the primary path's one
  (${\sim}20\times$).
  \item[A playlist-lab readout.] A mood report card, a comparative journey-shape
  chart and per-step \texttt{passed\_over} diagnostics, with the holisticness
  facets demoted to a collapsed detail block.
\end{description}

\emph{Results.} Table~\ref{tab:nexttrack-rollout} is five sweeps, \num{120} paired
rollouts and \num{240} journeys, under a paired bootstrap over \emph{seeds} plus a
leave-one-seed-out jackknife. Rollout policy was frozen throughout at the lab
defaults (\num{20} steps, artist penalty \num{2}, anchor \num{1}, temperature
\num{0} --- deterministic, so one seed is one sample and there is no rollout
variance).

\begin{table}[ht]
  \centering
  \caption{Five \texttt{hold@20} sweeps, paired bootstrap over seeds (\num{2000}
  resamples, rng \num{1337}) with a leave-one-seed-out jackknife (``LOO'' is the
  number of refits preserving the verdict). Positive $\Delta$ favours the
  dual-tower arm. The pre-registered rule was CONFIRMED iff the CI lies entirely
  above zero \emph{and} the verdict survives $\ge\SI{90}{\percent}$ of LOO refits;
  TIE iff the CI straddles. S2 and S3 met that rule and were nonetheless
  \emph{retracted} --- see the interpretation.
  \emph{Source: \texttt{2026-07-27-rollout-hold-metric.md}.}}
  \label{tab:nexttrack-rollout}
  {\footnotesize\setlength{\tabcolsep}{4pt}\begin{tabular}{llrlll}
    \toprule
    \# & Sweep & $n$ & $\Delta$\,\texttt{hold@20} [95\% CI] & LOO & Verdict \\
    \midrule
    S1 & pilot, mood-extreme seeds, 10-track prefix & 8 & $-0.034$ [$-0.138$, $+0.079$] & 0/8 & TIE \\
    S2 & pilot, mood-extreme seeds, 30-track prefix & 8 & $+0.314$ [$+0.107$, $+0.529$] & 8/8 & \emph{retracted} \\
    S3 & replication, \textbf{random} seeds & 39 & $+0.122$ [$+0.044$, $+0.201$] & 39/39 & \emph{confounded} \\
    S4 & \textbf{same-split control} & 39 & $+0.033$ [$-0.030$, $+0.093$] & 0/39 & \best{TIE} \\
    S4b & same-split, other dual checkpoint & 39 & $+0.051$ [$-0.012$, $+0.118$] & 0/39 & \best{TIE} \\
    S4c & \textbf{within-architecture} & 39 & $-0.072$ [$-0.141$, $-0.002$] & --- & significant \\
    S5 & \textbf{the user's own \num{11} playlists} & 11 & $+0.004$ [$-0.088$, $+0.104$] & 0/11 & \best{TIE} \\
    --- & the single judged playlist & 1 & $+0.379$ (\num{0.623} / \num{0.245}) & --- & outlier \\
    \bottomrule
  \end{tabular}}
\end{table}

\emph{Interpretation.} \textbf{Rollout fade does not separate this family.} At
matched training split the two architectures tie \emph{twice}; on the user's own
playlists the pair they preferred ties; and the only surviving effect is
\textbf{checkpoint-level} within one architecture (S4c), which does not transfer to
the playlist regime. Mean \texttt{hold@20} over the identical \num{39} random
seeds runs \num{0.6488} / \num{0.6309} for the two dual-tower checkpoints against
\num{0.5983} / \num{0.5266} for the two single-GRU ones --- so S3's apparent
architecture win decomposes into a non-significant architecture component
($+\num{0.05}$) and a significant \emph{checkpoint} component ($+\num{0.072}$).
Two \texttt{seq-nexttrack} models differ from each other more than dual-tower
differs from single-GRU at matched split, and the checkpoint the playlist lab
pre-selected by default happens to be the worst holder of the four. The user's
original comparison was the best holder against the worst.

Combined with the five prior architecture nulls this is a sixth, and it is the
strongest form available, because the standing rebuttal to those nulls was
``Recall@10 cannot see session quality''. It cannot; and when an instrument built
for the deployment regime looks, there is still nothing there.

\textbf{Per-seed variance dominates model identity, and that is the operational
finding.} Individual gaps on \emph{one} model pair span $+\num{0.39}$ to
$-\num{0.27}$, while the model-level mean difference on the user's playlists is
$+\num{0.004}$. Which seed you hand the generator matters far more than which of
these models generates. Two consequences: prefer serving an \emph{ensemble} and
letting the listener choose over hunting a single better generator, and budget
$\ge\num{40}$ seeds for any future rollout comparison --- \numrange{8}{12} cannot
resolve anything at this variance.

\textbf{\texttt{hold@N} is a working instrument but not a model ranker.} It fired
on a real checkpoint difference, it fired on the judged playlist (\num{0.62}
against \num{0.25}) where \texttt{holisticness@10} saw a \num{0.003} gap, and it
correctly returned TIE three times. The recommendation is therefore to keep it as
a per-journey \emph{lab readout} and \textbf{not} to wire it as a leaderboard
column: its ground truth is listener preference, which stands at $n=1$ blind vote.
Promoting an unvalidated composite to primary is how the crown ended up needing an
external recall floor, and replacing one unvalidated composite with another repeats
the mistake.

\paragraph{Four methodological pitfalls, each of which cost a retracted
positive.} They are worth more than the null. \textbf{(1) A cross-split comparison
masquerading as an architecture result.} The lab's two default pre-selects are
trained on \emph{different} splits with identical vocab and item latents, so
comparing them attributes checkpoint and split differences to architecture; pair
from the same batch when the claim is architectural. \textbf{(2) Mood-extreme seed
selection inflates effects ${\sim}\num{2.5}\times$} (S2's $+\num{0.314}$ became
S3's $+\num{0.122}$ on a random sample). \textbf{(3) An anti-gaming term can be
too weak to do its job:} \texttt{stasis} first normalized the shortfall by the
band \emph{width}, and because this listener's jump band is wide a flatlined
rollout scored \num{0.740} --- \emph{above} a real generator's \num{0.477}.
Normalize by the \emph{floor}; a guard beatable by the degeneracy it targets
launders that degeneracy as a good score. \textbf{(4) A metric that returns
nothing must drop the \emph{pair}} and report the smaller $n$, or the pairing
silently breaks. The jackknife earned its keep at $n<40$ throughout: it killed two
of three per-axis positives at $n=8$ and flagged S4 and S5 as unsupported at 0/39
and 0/11.

\paragraph{What the per-axis read adds, including a failed negative control.} At
$n=\num{39}$ only \emph{energy} fade separates the arms ($-\num{0.039}$
[$-\num{0.069}$, $-\num{0.009}$], favouring dual-tower); valence, acousticness and
tempo all straddle. But one pre-registered \emph{negative} control moved: the
energy \emph{level} deviation from the seed favours the single GRU
($+\num{0.018}$ [$+\num{0.005}$, $+\num{0.032}$]). Those are different properties,
so the claim had to be re-scoped from ``holds the seed better'' to ``is more
stationary'' --- and it shows that \texttt{hold@N}'s deliberate exclusion of a level
term is a choice that \emph{changes the ranking}, not a neutral simplification.
Two other limits travel with the batch: \SI{74}{\percent} of the corpus carries
the acoustic fields the metric needs, so seeds were pre-filtered for
$\ge\SI{60}{\percent}$ prefix coverage (a mild selection effect), and nothing here
has been listened to.

Nothing was promoted, registered or wired; the champion, the crown board and the
best-models pins are unchanged. For the record, the live-crown standing relevant
to the dual-tower question is a \emph{tie} rather than a loss --- the best
legitimate dual arm \num{0.02378} against the single GRU's \num{0.02454}, a
\num{0.0008} gap against a ${\sim}\num{0.0018}$ half-width. The family was refuted
on Recall@10, ties on the live primary, and shows nothing on the rollout axis.
The indicated next axis is therefore not topology but the \emph{training regime}:
all five architecture nulls held one-step teacher forcing byte-identical while the
product rolls twenty steps on its own output, and that exposure-bias mismatch is
the one thing those nulls could not have detected.
\emph{Source: \texttt{2026-07-27-rollout-hold-metric.md}.}

\section{Next-track: a research memo, and what its provenance permits}
\label{sec:nexttrack-memo}
\emph{Status, first.} This section consolidates a document that is \textbf{not a
campaign report}. Nothing was run on the server, nothing was promoted, nothing was
registered; the champion, the leaderboard, the crown board and the best-models
pins are all unchanged by it. It was produced by a research workflow (four scouts,
five proposal lenses, five adversarial screeners, then synthesis) plus
\emph{offline} scoring against frozen checkpoints. Where it reports a number, that
number came from an offline screen over stored artifacts, not from a native run.
It is included here because it is the direct parent of the campaign in
Section~\ref{sec:nexttrack-markovgate} and because its negative results --- the graveyard
below --- are cheaper to reuse than to rediscover. It is reported at a lower
register than the measured campaigns, and its figures are deliberately kept out of
every board in this document.

\paragraph{The provenance partition, kept visible.} The memo itself distinguishes
two classes of number, and that distinction survives into this synthesis.
\begin{description}
  \item[Independently re-derived by the orchestrator.] The share of test queries
  whose last prefix item has \emph{zero} train-bigram support:
  \textbf{\num{795} of \num{1431} $=$ \num{0.5556}}, reproduced by an independent
  script. The ungated control, which is self-validating because it must equal the
  champion: \textbf{\num{0.21174}}, exact. The gated arm: \textbf{\num{0.23201}},
  a paired $\Delta$ of $+\num{0.02027}$ $[+\num{0.01048}, +\num{0.03075}]$ with the
  CI above zero. And its bucket decomposition: the zero-support bucket moves
  \num{0.1660}$\to$\num{0.2025} over \num{795} queries while the two warm buckets
  are \textbf{bit-identical} (\num{0.3116} and \num{0.2326}, unchanged).
  \item[\textbf{Not} verified by the orchestrator --- agent-reported only, treated
  as provisional.] The $+\num{0.01957}$ second-split figure, the replay channel's
  $+\num{0.02865}$, the \texttt{latent\_gamma} numbers, and the recomputed
  \num{0.0706} discordance with its ${\sim}\num{0.0145}$ noise floor. None of these
  appears in Table~\ref{tab:nexttrack-lb} or on any crown board, and the first of them is
  precisely what the follow-up campaign was commissioned to establish natively.
\end{description}

\paragraph{One caveat the memo states about its own headline.} The ten gating
variants were screened \textbf{on the test set}, so quoting the best arm's CI is
optimistic by construction and the screen cannot itself be evidence. Three things
mitigate it --- four mechanistically related variants all land at
\numrange{0.019}{0.020}, the bucket decomposition is a mechanical explanation
rather than a fit, and the warm buckets do not move by a single hit --- but the
mitigation is an argument, not a measurement. The native two-split confirm is the
real gate, and it is Section~\ref{sec:nexttrack-markovgate}.

\paragraph{The honest headline: ``a better topology'' is the wrong question.} The
record had by then answered it six times --- cell and loss on the AE-64 space;
depth, width and direction; width on PCA-192; dual-tower fusion
(Section~\ref{sec:nexttrack-dualtower}); the MLP fusion head, an \emph{active} regression
recovered by deleting it; the per-step pre-encoder, a measured loss with the
mechanism identified down to \texttt{frac\_exact\_zero} $=$
\texttt{frac\_negative\_coords} (Section~\ref{sec:nexttrack-preencoder}); and rollout fade
(Section~\ref{sec:nexttrack-rollout}). Five lenses were pointed at the topology question
and \emph{every} architecture-shaped proposal came back weakened, sub-noise, or
pre-killed. Two came back with measured negative results: a split Markov leg lost
at every point of its own pre-registered grid, and a factored artist softmax head
was killed by its own pre-kill run, which found that top-5 artist gating on the
real checkpoint changes recall by \textbf{exactly \num{0.0}} because the trained
GRU's top-10 already lives inside its own top-5 artists.

\textbf{Where the evidence points instead: the model is not the constraint. What
it is asked to rank against, and which legs are allowed to speak on a given
query, are.} Every proposal that survived adversarial screening is a candidate-set
or evidence-conditioning change --- an evidence-gated Markov leg, a causal
cross-session replay channel, per-leg reliability scalars, a prefix-only negative
mask, the never-swept softmax temperature. None is a topology. This is the same
conclusion the record had reached from the representation side, arrived at from
five independent directions, and it comes with a measurement rationale: eval-time
and score-time levers are $\numrange{3}{33}\times$ cheaper to resolve than
training-time ones (Section~\ref{sec:nexttrack-mmr}), so a shortlist of score-time
proposals is the rational response to the instrument this lab has.

\paragraph{The two survivors worth naming.} \textbf{(1) \texttt{markov\_gate}} ---
gate the Markov leg to abstain where it has no bigram evidence, restoring the
\texttt{trust} term that the blend's z-normalization erases. Twelve lines,
zero learned parameters, and confirmed natively in the next section.
\textbf{(2) A causal cross-session replay channel} --- \SI{76.17}{\percent} of the
held-out targets were played in a \emph{strictly earlier} session and are not in
the current prefix, and nothing in the family can see this, because the Markov leg
is a train-frozen bigram and the recency scorer only ever scored prefix items,
which the evaluator masks to $-\infty$ (its recorded \num{0.001} is dead code, not
a measurement). Its hit-Jaccard with the GRU is \num{0.060} against \num{0.201}
for the GRU/Markov pair that produced a working blend, i.e.\ it is genuinely
different signal. Its catch is \emph{governance}, not code: it is an
information-set amendment that moves the rankable ceiling from \num{0.4493} to
\num{0.7617}, which needs its own board column and an explicit decision before any
number from it is comparable with anything in this document. A third proposal,
\texttt{latent\_gamma} (a one-scalar spectral rescale of the item latents), is
worth recording for its \emph{reframe} rather than its delta: on the offline
screen the champion's entire recall contribution is reproducible by a
\num{192}-number diagonal rescale, which recasts ``the single biggest win on
record'' as a conditioning fix that supervision happened to find, and predicts
that further linear metric learning is near-exhausted. All of these figures are in
the provisional class.

\paragraph{The graveyard, which is the reusable part.} Each entry is a lane that
should not be re-opened without new information. Full-catalogue or
support-restricted cross-entropy is pre-empted by an existing arm at
${\sim}\SI{53}{\percent}$ support coverage (\num{0.1223} against \num{0.1230}) ---
the literature's top candidate does not transfer because the item side is frozen.
A factored artist softmax head is refuted by its own pre-kill (exactly \num{0.0}).
Splitting the Markov leg into track and back-off components was measured at every
grid point and lost at all five ($-\num{0.0217}$ to $-\num{0.0384}$, all CIs below
zero), because separate z-normalization structurally destroys the \texttt{trust}
term inside the warm regime. A pool-matched stage-2 re-ranker is unobtainable: the
train pool's positive density is \num{0.5975} against test's \num{0.3368}, so its
own go-gate passes \emph{because of} the confound. A mixture-of-vMF head loses to
its own oracle decomposition. Diffusion, DreamRec and flow-matching are
externally falsified for this setting, and generative retrieval with semantic IDs
is strictly dominated when exact cosine over \SI{15}{\mega\byte} is already
affordable and the corpus is \SIrange{55}{64}{\percent} cold. Popularity
debiasing has nothing to debias (the popularity baseline is \num{0.0007}).
Hard-negative mining is harmful at low temperature here, because the hardest
negatives are same-album and same-artist tracks --- i.e.\ mining false positives.
And the GRU's pre-normalization output norm is not a confidence signal
(correlation with correctness $+\num{0.0073}$, and the loss is exactly
norm-invariant); the top-1-minus-next-10 score margin is, at $+\num{0.294}$.

\paragraph{The one durable theoretical result.} The memo's cross-cutting
observation that \textbf{the learned-combiner anomaly is mechanically explained}
is the part that changes how the record should be read. A prior oracle search over
\emph{global static scalar} leg weights returned $+\num{0.004}$ straddling zero,
and it is now clear why that is compatible with a real $+\num{0.020}$ from a
gate: dropping the Markov leg \emph{globally} gains in the zero-support bucket and
loses in the warm bucket, and the two cancel almost exactly. No constant can
express an intervention whose sign flips with a per-query data statistic, and
z-normalization is scale-invariant so it provably cannot see one. The registered
``four independent learned-combiner nulls'' pitfall should therefore forbid
\emph{learned score-level fusion} --- fitted weights, stackers, rank fusion,
end-to-end representation fusion --- and \textbf{not} per-query evidence gating,
which is a third thing: nothing is learned, nothing co-adapts, and the condition
is a train-side count.
\emph{Source: \texttt{2026-07-27b-nexttrack-model-topology-research-memo.md}
(research memo; see the provenance partition above).}

\section{Next-track: \texttt{markov\_gate}: a leg that abstains where it has no evidence}
\label{sec:nexttrack-markovgate}
\emph{Rationale.} This is the native two-split confirm of the memo's headline
proposal (Section~\ref{sec:nexttrack-memo}), and it produces the largest single recall gain
in the whole next-track arc from a nine-line change with no new parameters.

\paragraph{The mechanism, and why it is not the refuted combiner lane.} The Markov
leg mixes its bigram row with an artist/genre back-off through a trust weight
$\texttt{trust} = n_u/(n_u+8)$, where $n_u$ is the train-level bigram support of
the query's last prefix item $u$. At $n_u = 0$ the trust is zero and the leg is
\textbf{pure back-off}: it carries no track-level evidence at all. But the blend's
candidate normalizer is \emph{scale-invariant} by construction --- it maps each
leg's score vector to unit variance so that three legs on wildly different scales
are commensurable --- so the evidence-free back-off is re-inflated and speaks at a
full one-third weight anyway. On the canonical split that happens on \textbf{\num{795}
of \num{1431} test queries}, \SI{55.6}{\percent}. The intervention is to make the
leg \emph{abstain} at $n_u = 0$ and renormalize the divisor from three legs to two.
Nothing is learned; the condition is a train-side count. That is exactly the third
category the registered learned-combiner pitfall did not cover, and the bucket
decomposition below is a mechanical proof that no global static weight could find
the effect: dropping the leg globally is a $-\num{0.0070}$ tie, because the
zero-support gain and the warm-bucket loss cancel.

\paragraph{Design and the preflight gate.} Six runs on the hub, all
\texttt{seq-blend} from the champion definition, with the controls
\textbf{re-trained in this batch} rather than substituted from the archive so that
every paired comparison is code-matched to one generation. Two splits: the
canonical \dsid{seq-20260715-131139} and the disjoint earlier-holdout
\dsid{seq-20260718-211238} (\SI{64}{\percent} cold, $n_{\mathrm{test}}$
\num{1431}). The hard pre-registered stop was that the ungated control had to
reproduce Recall@10 \num{0.21174}, exactly \num{303} of \num{1431}, proving the new
flag left the champion bit-identical. \textbf{It did, to the last digit}
(\num{0.21174004192872117}).

\emph{Results.} Table~\ref{tab:nexttrack-markovgate} is the batch,
Table~\ref{tab:nexttrack-markovgate-paired} the paired statistics, and
Table~\ref{tab:nexttrack-markovgate-buckets} the decomposition that carries the argument.
Figure~\ref{fig:nexttrack-markov-gate} plots both.

\begin{table}[ht]
  \centering
  \caption{The \texttt{markov\_gate} batch, six runs from the champion definition
  \texttt{blend-gru-markov-content-proj}. \dsid{ds1} $=$
  \dsid{seq-20260715-131139} (canonical, \SI{55}{\percent} cold), \dsid{ds2} $=$
  \dsid{seq-20260718-211238} (disjoint, \SI{64}{\percent} cold);
  \dsid{ds2}'s absolute recall is lower \emph{by construction}, so read it only
  through the paired $\Delta$. Every run id is \protect\runid{run-20260728-} plus
  the value in the \emph{run} column. The graded metrics track the gate too:
  artist@10 \num{0.333}$\to$\num{0.367} and genre@10 \num{0.401}$\to$\num{0.463}
  from C1 to A1, and \num{0.382} / \num{0.497} on A3.
  Rows sorted by Recall@10. H@10 is generation~3
  (Section~\ref{sec:nexttrack-power-generations}). R@10, MRR, H@10, mus@10, \texttt{ild},
  \texttt{mood}, art@10, gen@10 higher-better; \texttt{adj}, \texttt{conc}
  lower-better. Best cell per metric bolded among the canonical-split arms; the
  new best-recall row A1 is bolded in full.
  \emph{Source: \texttt{2026-07-27c-nexttrack-markov-gate-confirm.md}.}}
  \label{tab:nexttrack-markovgate}
  {\scriptsize\setlength{\tabcolsep}{1.9pt}\begin{tabular}{llllllllllll}
    \toprule
    id & run & ds & overrides & R@10 (hits) & MRR & H@10 & mus@10 & ild & adj & conc & mood \\
    \midrule
    \best{A1} & \texttt{005349-507bb} & ds1 & \best{gate} & \best{0.23201 (332)} & \best{0.13352} & 0.018699 & 0.49319 & 0.43377 & 0.57477 & 0.59821 & \best{0.46226} \\
    A3 & \texttt{012030-6d1b7} & ds1 & gate $+$ $\lambda$0.7 & 0.21314 (305) & 0.12852 & \best{0.030980} & \best{0.51649} & \best{0.56770} & \best{0.51279} & \best{0.47325} & 0.41885 \\
    C1 & \texttt{004444-b32ea} & ds1 & --- (champion) & 0.21174 (303) & 0.12199 & 0.008603 & 0.45379 & 0.40637 & 0.60971 & 0.69165 & 0.39822 \\
    C3 & \texttt{011245-bfc67} & ds1 & $\lambda$0.7 only & 0.19357 (277) & 0.11560 & 0.015289 & 0.47592 & 0.52010 & 0.57135 & 0.60823 & 0.36433 \\
    A2 & \texttt{010658-39e20} & ds2 & gate & 0.19008 (272) & 0.11801 & 0.008414 & 0.42953 & 0.41401 & 0.58672 & 0.64847 & 0.42618 \\
    C2 & \texttt{010113-e2599} & ds2 & --- (champion) & 0.17051 (244) & 0.10498 & 0.003392 & 0.39616 & 0.40177 & 0.61013 & 0.71945 & 0.35364 \\
    \bottomrule
  \end{tabular}}
\end{table}

\begin{table}[ht]
  \centering
  \caption{Paired bootstrap over the \num{1431} test sessions (\num{2000}
  resamples, rng \num{1337}, paired by row id). $\Delta H$ is computed
  \emph{product-of-means per resample} --- all five facets re-aggregated on the same
  resampled index set for both arms, then multiplied --- not as a mean of per-row
  products. $n_{disc}$ counts sessions whose hit@10 changed at all, split by sign.
  A verdict of WIN means the recall CI lies entirely above zero, TIE that it
  straddles; bounds are abbreviated to four decimals for width.
  \emph{Source: \texttt{2026-07-27c-nexttrack-markov-gate-confirm.md}.}}
  \label{tab:nexttrack-markovgate-paired}
  {\scriptsize\setlength{\tabcolsep}{4pt}\begin{tabular}{llllr}
    \toprule
    Comparison & $\Delta$R@10 [95\% CI] & Verdict & $\Delta H$ [95\% CI] & $n_{disc}$ ($+$/$-$) \\
    \midrule
    A1 $-$ C1 (gate, ds1) & \best{$+.02027$ [$+.0105$, $+.0308$]} & \best{WIN} & $+.010096$ [$+.0085$, $+.0117$] & 57 (43/14) \\
    A2 $-$ C2 (gate, ds2) & \best{$+.01957$ [$+.0105$, $+.0301$]} & \best{WIN} & $+.005022$ [$+.0039$, $+.0062$] & 56 (42/14) \\
    A3 $-$ C3 (gate, $\lambda$0.7) & \best{$+.01957$ [$+.0091$, $+.0301$]} & \best{WIN} & $+.015691$ [$+.0136$, $+.0179$] & 64 (46/18) \\
    A3 $-$ C1 (gate$+$MMR / champ.) & $+.00140$ [$-.0119$, $+.0147$] & TIE & \best{$+.022377$ [$+.0199$, $+.0249$]} & 96 (49/47) \\
    C3 $-$ C1 (MMR $\lambda$0.7 alone) & $-.01817$ [$-.0294$, $-.0077$] & costs recall & $+.006686$ [$+.0057$, $+.0077$] & 60 (17/43) \\
    \bottomrule
  \end{tabular}}
\end{table}

\begin{table}[ht]
  \centering
  \caption{The bucket decomposition, reproduced \emph{natively} through the real
  transition builder and bucketed by the train-level bigram support $n_u$ of each
  query's last prefix item --- the exact quantity the Markov scorer computes.
  Recall@10 within bucket, hits in parentheses. In both warm buckets the two arms'
  top-$k$ lists are \textbf{byte-for-byte identical}, so the gate is a surgically
  pure $n_u=0$ intervention with literally zero warm-side cost, on three
  independent control/arm pairs.
  \emph{Source: \texttt{2026-07-27c-nexttrack-markov-gate-confirm.md}.}}
  \label{tab:nexttrack-markovgate-buckets}
  {\footnotesize\setlength{\tabcolsep}{5pt}\begin{tabular}{llrllll}
    \toprule
    Pair & Bucket & $n$ & control & gated & $\Delta$ & top-$k$ identical? \\
    \midrule
    \multicolumn{7}{l}{\emph{Canonical split, A1 vs C1}} \\
    & $n_u = 0$ & 795 & 0.1660 (132) & 0.2025 (161) & \best{$+0.0365$} & no \\
    & $n_u$ 1--10 & 292 & 0.3116 (91) & 0.3116 (91) & $+0.0000$ & \best{YES} \\
    & $n_u > 10$ & 344 & 0.2326 (80) & 0.2326 (80) & $+0.0000$ & \best{YES} \\
    \midrule
    \multicolumn{7}{l}{\emph{Second split, A2 vs C2}} \\
    & $n_u = 0$ & 892 & 0.1244 (111) & 0.1558 (139) & \best{$+0.0314$} & no \\
    & $n_u$ 1--10 & 322 & 0.2578 (83) & 0.2578 (83) & $+0.0000$ & \best{YES} \\
    & $n_u > 10$ & 217 & 0.2304 (50) & 0.2304 (50) & $+0.0000$ & \best{YES} \\
    \midrule
    \multicolumn{7}{l}{\emph{Canonical split under MMR $\lambda$0.7, A3 vs C3}} \\
    & $n_u = 0$ & 795 & 0.1472 (117) & 0.1824 (145) & \best{$+0.0352$} & no \\
    & $n_u$ 1--10 & 292 & 0.2945 (86) & 0.2945 (86) & $+0.0000$ & \best{YES} \\
    & $n_u > 10$ & 344 & 0.2151 (74) & 0.2151 (74) & $+0.0000$ & \best{YES} \\
    \bottomrule
  \end{tabular}}
\end{table}

\begin{figure}[ht]
  \centering
  \includegraphics[width=0.95\textwidth]{nexttrack_markov_gate}
  \caption{\textbf{Left:} the bucket decomposition on all three control/arm pairs
  (\texttt{ds1} canonical, \texttt{ds2} the \SI{64}{\percent}-cold second split,
  \texttt{ds1 $\lambda$0.7} the same canonical pair under an MMR re-rank). The
  entire effect sits in the zero-support bucket; in both warm buckets the arms'
  top-$k$ lists are byte-for-byte identical, which is why a global leg weight
  provably cannot express this intervention. \textbf{Right:} the same result on the
  recall $\times$ crown plane. MMR alone (blue) is a \emph{priced} lever on the
  champion --- it buys crown and loses recall with the CI below zero. The gate
  (green) buys \emph{both}. Composed, gate$+$MMR lands at champion recall
  ($\Delta$ $+\num{0.00140}$, straddling) with $\num{3.6}\times$ the champion's
  holisticness, so the gate changes MMR's tier on the champion from PRICED to FREE.}
  \label{fig:nexttrack-markov-gate}
\end{figure}

\emph{Interpretation.} \textbf{The defect is real, it belongs to the normalizer,
and it is worth \num{29} hits.} Scale-invariance is exactly what makes the z-blend
robust to legs on incommensurable scales, and exactly what destroys their
\emph{confidence}. A Markov leg that has never seen the current track emits a pure
artist/genre back-off; the normalizer re-inflates it to the same footprint as a GRU
that has actually modelled the prefix, and the blend averages the two as equals.
Turning that third of the vote off where it is uninformed converts \num{43} misses
into hits and \num{14} hits into misses --- net $+\num{29}$, $+\SI{9.6}{\percent}$
relative, on the best model on record. Recall@10 \num{0.23201} (\num{332} of
\num{1431}) is \textbf{the best recall figure in this document}, with the paired CI
above zero on two disjoint splits and MRR up on both. For scale: the champion's own
coronation margin over the prior crown was $+\num{0.0259}$ canonical and
$+\num{0.0426}$ on the second split; this gate adds $+\num{0.0203}$ /
$+\num{0.0196}$ \emph{on top of that}, at no additional training cost.

\textbf{The gate is also a crown lever, which was not predicted.} $\Delta H$ is
above zero in all three paired comparisons, and the facets say why: abstaining
kills the \emph{bigram-echo} component of the recommendation. Artist concentration
falls \num{0.69165}$\to$\num{0.59821} and adjacency \num{0.60971}$\to$\num{0.57477}
--- the list stops repeating one artist --- while \texttt{mood\_coh} \emph{rises}
\num{0.39822}$\to$\num{0.46226} and \texttt{music@10} rises
\num{0.45379}$\to$\num{0.49319}. All four crown factors move the right way at
once. This is the first lever on record that de-eagers the champion \emph{without}
paying in mood coherence, and it resolves the exact double-payment
Section~\ref{sec:nexttrack-mmr} identified as structural: there, de-eagering was a re-rank
fighting the model's own scores, so walking away from the seed artist also walked
the blend out of the prefix's mood neighbourhood. Here the de-eagering is the
removal of a vote that had no business being cast.

\textbf{Gate and MMR are additive, not substitutive, and their composition is the
headline for the crown board.} The gate-alone $\Delta$ ($+\num{0.02027}$) and the
gate-under-$\lambda\num{0.7}$ $\Delta$ ($+\num{0.01957}$) are
indistinguishable, and the buckets show why: MMR re-ranks the top of a scored list,
the gate changes which list gets scored, and they act on disjoint parts of the
pipeline. Standing alone, MMR $\lambda\num{0.7}$ on the champion is TIER-2 PRICED
($-\num{0.01817}$ recall, CI below zero). Add the gate and the recall cost is paid
back exactly: A3 is $+\num{0.00140}$ $[-\num{0.01188}, +\num{0.01468}]$, a
statistical tie with the champion, at $H$~\num{0.030980} --- $\num{3.6}\times$ the
champion's \num{0.008603}, comfortably above the mandatory \num{0.108} recall
floor. On the crown board that sits at \textbf{rank 2}, behind only the
$\lambda\num{0.7}$ single GRU, which scores its \num{0.034171} at recall
\num{0.11461}: A3 delivers \SI{96}{\percent} of the incumbent leader's crown at
$\num{1.86}\times$ its recall.

\textbf{The honest caveat, and why the out-of-sample structure redeems it.} The
memo screened ten gating variants on the test set, so its banked
$+\num{0.02027}$ was selection-optimistic and could not be evidence. Four things
redeem it here: the effect reproduces on a \emph{second, disjoint,
\SI{64}{\percent}-cold split that was not screened on}, within \num{0.0007} of the
canonical figure; it reproduces under a third condition ($\lambda\num{0.7}$) that
was not screened either; four mechanistically related variants in the memo all
landed at \numrange{0.019}{0.020}, so the screen was measuring a plateau rather
than selecting a lucky corner; and the warm buckets do not move by a single hit,
which is a \emph{structural} prediction of the mechanism rather than a fitted
outcome. One residual risk remains: the $n_u = 0$ \emph{boundary} itself was chosen
with test knowledge. The boundary-free version --- scaling the leg by
\texttt{trust} \emph{before} normalization instead of hard-abstaining --- is the top
follow-up, and it is strictly more general. Nor is any of this a fluke of one
split's cold rate: the second split is \SI{64}{\percent} cold, its zero-support
bucket grows to \SI{62.3}{\percent} of queries while the per-query gain shrinks
slightly, and the net $\Delta$ comes out nearly identical. The lever's value scales
with how much of the query distribution has no bigram evidence, which is the
majority on any split of this corpus.

\paragraph{Governance, and an infrastructure finding that nearly ate the result.}
Nothing was promoted, registered or tagged; by design boundary this campaign
carries a promotion \emph{recommendation} only, and the board row in
Table~\ref{tab:nexttrack-lb} is labelled accordingly. The server's deterministic
top-12-by-holisticness recompute nevertheless auto-promoted four of the batch's
runs, which is server behaviour rather than a campaign action and left the group
needing re-curation. Separately: the remote training worker was live with a
\textbf{stale checkout containing zero occurrences of the new flag}. Had any gated
arm been dispatched there it would have executed the \emph{ungated} path and
silently returned a null --- the registered ``remote checkout can silently diverge''
pitfall, in its worst form, because the failure is a wrong number rather than a
crash. It could not happen here only because remote dispatch is opt-in and every
run in this batch was posted without it (confirmed by an unclaimed marker on all
six). Until the checkout is freshened, any queued \texttt{seq-blend} run silently
ignores the gate.

\paragraph{Record amendments.} Two, both narrowing rather than overturning.
\textbf{(1)} The learned-combiner pitfall should forbid \emph{learned score-level
fusion} and not \emph{per-query evidence gating}; the prior oracle null and this
$+\num{0.020}$ win are both true and are reconciled by the bucket decomposition,
since a global static scalar cannot express a gate whose sign flips with $n_u$.
\textbf{(2)} The normalizer's erasure of per-leg \emph{confidence} is a
\textbf{general} defect, of which the Markov leg's zero-support back-off is one
instance. Any leg with a query-dependent notion of ``I have no evidence here'' is
mis-weighted the same way, and the content and GRU legs have not been audited for
it.
\emph{Source: \texttt{2026-07-27c-nexttrack-markov-gate-confirm.md}.}

\section{Next-track: the eagerness regularizer at a calibrated margin: refuted, with a mechanism}
\label{sec:nexttrack-eagermargin}
\emph{Rationale.} Section~\ref{sec:nexttrack-holisticness} left \texttt{eager\_beta} --- the
only \emph{training-time} de-eagering lever in the repository --- on record as ``a
weak knob, not recommended''. That verdict deserved a re-test, because the probe
behind it used a margin of zero at a dose worth \SIrange{1}{3}{\percent} of the
loss. This campaign asked whether the lever works when the margin is placed in the
measured \emph{right tail} of the transition-cosine distribution at
$\numrange{15}{60}\times$ that dose. The stake was deployability: at the time,
\texttt{mmr\_lambda} entered only through the evaluation path and never reached
serving, so the training-time regularizer was the one anti-eager candidate that
could actually ship.

\paragraph{The design was written so that a null would still be informative.} The
pre-registered primary was the paired $\Delta$ on \texttt{artist\_conc@10} against
an in-batch control, with holisticness@10 explicitly \emph{demoted} to secondary
because at the operating point $\partial H/\partial\texttt{music@10} =
+\num{0.2062}$ against $\partial H/\partial\texttt{eager} = -\num{0.0378}$ --- a
$\num{5.46}\times$ ratio that would have scored a \emph{successful} de-eagering as
approximately zero. A relevance floor was pre-registered in two parts: absolute
\texttt{music@10} $\ge\num{0.4300}$ \emph{and} a paired $\Delta$\texttt{music@10}
lower bound above $-\num{0.030}$. Eight \texttt{seq-dualgru} runs on the canonical
split \dsid{seq-20260715-131139}, all remote, \SI{60}{\minute} of compute, none
failed. The bit-identity gate passed at exactly \num{170} hits with the
reconstructed concentration matching the on-record arm to \num{15} decimal places,
and a penalty-active gate was pre-registered as a silent-failure detector: it
\textbf{passed on three arms and FAILED --- reporting the penalty inert --- on
four}, which is the campaign's result.

\emph{Results.} Table~\ref{tab:nexttrack-eagermargin} is the batch and
Table~\ref{tab:nexttrack-eagermargin-paired} the paired statistics.
Figure~\ref{fig:nexttrack-eager-margin} shows the mechanism.

\begin{table}[ht]
  \centering
  \caption{The \texttt{eager\_beta} $\times$ margin batch, eight
  \texttt{seq-dualgru} runs on \dsid{seq-20260715-131139} ($n=\num{1431}$). Every
  run id is \protect\runid{run-20260730-} plus the value in the \emph{run} column.
  $m$ is the hinge margin, $\beta$ the penalty weight. H@10 is generation~3
  (Section~\ref{sec:nexttrack-power-generations}). ``best ep/ran'' is the best epoch over
  the epochs actually run; Recall@10 is \texttt{hits}$/\num{1431}$ throughout
  (C0 \num{0.118798}, A1 \num{0.053809}, A2 \num{0.013976}, B1 \num{0.114605}, and
  \num{0.118798} on every inert arm). \texttt{ild} is omitted for width and moves
  only on the arms that fire (\num{0.4506} on C0, \num{0.5572} on A1, \num{0.6762}
  on A2, \num{0.4530} on B1, \numrange{0.4504}{0.4506} on the inert arms).
  \texttt{mood}, mus@10, H@10 higher-better; \texttt{conc}, \texttt{adj}
  lower-better. The lowest-eagerness
  cells are bolded and belong to \emph{degenerate} arms, which is what the
  relevance floor exists to catch.
  \emph{Source: \texttt{2026-07-30-eager-margin-calibrated-refutation.md}.}}
  \label{tab:nexttrack-eagermargin}
  {\scriptsize\setlength{\tabcolsep}{2.8pt}\begin{tabular}{llllrlllllll}
    \toprule
    Arm & run & $m$ & $\beta$ & hits & conc & adj & mood & mus@10 & H@10 & ep & verdict \\
    \midrule
    C0 & \texttt{175402-cd006} & 0 & 0 & \best{170} & 0.370588 & 0.362963 & \best{0.479846} & \best{0.455347} & \best{0.023785} & 8/14 & CONTROL \\
    A1 & \texttt{181405-53088} & 0 & 3 & 77 & 0.240717 & 0.206010 & 0.346867 & 0.399997 & 0.013339 & 18/24 & DEGENERATE \\
    A2 & \texttt{181405-0bc1a} & 0 & 12 & 20 & \best{0.133535} & \best{0.107687} & 0.221756 & 0.358132 & 0.003569 & 37/40 & DEGENERATE \\
    B1 & \texttt{181405-84da4} & 0.40 & 9 & 164 & 0.367606 & 0.364151 & 0.463440 & 0.447386 & 0.021601 & 15/21 & no effect \\
    B2 & \texttt{175402-7c304} & 0.5843 & 19 & 170 & 0.370634 & 0.363103 & 0.479955 & 0.455210 & 0.023751 & 8/14 & INERT \\
    B3 & \texttt{184556-b08c1} & 0.70 & 35 & 170 & 0.370588 & 0.362963 & 0.479846 & 0.455347 & 0.023785 & 8/14 & INERT (bit-id.) \\
    B4 & \texttt{175402-0d310} & 0.8304 & 92 & 170 & 0.370588 & 0.362963 & 0.479846 & 0.455347 & 0.023785 & 8/14 & INERT (bit-id.) \\
    B5 & \texttt{184556-47d7b} & 0.5843 & 57 & 170 & 0.370743 & 0.363103 & 0.479955 & 0.455379 & 0.023771 & 8/14 & INERT \\
    \bottomrule
  \end{tabular}}
\end{table}

\begin{table}[ht]
  \centering
  \caption{Paired statistics against the in-batch control C0 (\num{2000}
  resamples, rng \num{1337}, $n=\num{1431}$). $n_{disc}$ counts sessions whose
  hit@10 changed. Measured half-widths on the \emph{active} arms
  (\numrange{0.0080}{0.0173}) match the design's predicted \num{0.0102}, so the
  instrument was correctly powered --- the axis failed anyway.
  \emph{Source: \texttt{2026-07-30-eager-margin-calibrated-refutation.md}.}}
  \label{tab:nexttrack-eagermargin-paired}
  {\footnotesize\setlength{\tabcolsep}{4.5pt}\begin{tabular}{lllrrr}
    \toprule
    Arm & $\Delta$\texttt{artist\_conc} [95\% CI] & $\Delta$\texttt{music@10} [95\% CI] & $\Delta H$ & $n_{disc}$ & mean $|\Delta$conc$|$ \\
    \midrule
    A1 & $-0.12987$ [$-0.14354$, $-0.11605$] & $-0.05535$ [$-0.06606$, $-0.04465$] & $-0.010445$ & 125 & 0.191304 \\
    A2 & $-0.23705$ [$-0.25493$, $-0.22026$] & $-0.09721$ [$-0.11107$, $-0.08333$] & $-0.020215$ & 166 & 0.282895 \\
    B1 & $-0.00298$ [$-0.01085$, $+0.00519$] & $-0.00796$ [$-0.01411$, $-0.00194$] & $-0.002183$ & 62 & 0.093734 \\
    B2 & $+0.00005$ [$-0.00028$, $+0.00039$] & $-0.00014$ [$-0.00037$, $+0.00000$] & $-0.000034$ & \best{0} & 0.000419 \\
    B3 & $+0.00000$ [$+0.00000$, $+0.00000$] & $+0.00000$ [$+0.00000$, $+0.00000$] & $+0.000000$ & \best{0} & \best{0.000000} \\
    B4 & $+0.00000$ [$+0.00000$, $+0.00000$] & $+0.00000$ [$+0.00000$, $+0.00000$] & $+0.000000$ & \best{0} & \best{0.000000} \\
    B5 & $+0.00016$ [$-0.00031$, $+0.00068$] & $+0.00003$ [$-0.00036$, $+0.00045$] & $-0.000014$ & \best{0} & 0.000901 \\
    \bottomrule
  \end{tabular}}
\end{table}

\begin{figure}[ht]
  \centering
  \includegraphics[width=0.95\textwidth]{nexttrack_eager_margin}
  \caption{\textbf{Left:} the campaign's central result. The hinge's realized
  penalty at initialization, reconstructed from epoch-1 loss inflation at shared
  initialization (identical seed, so identical init and batch order). The margins
  were calibrated from the \emph{item}$\to$\emph{item} transition cosine (median
  \num{0.262}, p75 \num{0.584}, p90 \num{0.830}), but the penalty acts on
  \emph{prediction}$\to$\emph{current}, whose ceiling over training sits between
  \num{0.5843} and \num{0.70}: at $m\ge\num{0.5843}$ the penalized quantity never
  enters the hinge's support and the loss trace is bit-identical to the control at
  all \num{14} epochs, at $\beta$ up to the schema cap. \textbf{Right:} where the
  knob \emph{does} fire, eagerness and relevance are the same mechanism. Pushing
  the prediction off the current track's latent also pushes it off the region where
  the true next track lives, so there is no window between real de-eagering and
  breaching the relevance floor. The dash-dotted line marks the free on-record
  comparator that no arm here matched inside the floor.}
  \label{fig:nexttrack-eager-margin}
\end{figure}

\emph{Interpretation.} \textbf{The refutation is expressive, not dosimetric, and
that is the whole finding: the calibration measured the wrong distribution.} The
margins were drawn from $\cos(x_t, x_{t+1})$ --- how far apart two consecutively
played tracks are --- but the hinge acts on $\cos(\mathrm{pred}_t, x_t)$, how close
the model's \emph{prediction} sits to the current track. Those distributions are
not comparable. Reconstructing the hinge from epoch-1 loss inflation puts the
ceiling of the penalized quantity between \num{0.5843} and \num{0.70}, so p75 and
p90 of the real-transition distribution lie \emph{outside the support of the
quantity being penalized}, and arms B3 and B4 reproduced the control's loss to
\num{16} significant digits at every one of \num{14} epochs. The design's hedge ---
that the realized penalty would \emph{exceed} its nominal value --- is empirically
inverted by orders of magnitude. The reason is structural: an InfoNCE-trained
prediction is a \emph{hedged near-average} over many plausible continuations, so its
absolute cosine to any single item is bounded low. \textbf{Eagerness is a ranking
property --- whether the current track's neighbours occupy the top slots --- and a
hinge on an absolute cosine cannot express it.}

\textbf{Where the knob does fire, eagerness and relevance turn out to be the same
mechanism.} At $m=0$ the lever works spectacularly on the facet --- A2 cuts
concentration \num{0.371}$\to$\num{0.134} and drives \texttt{ild} to \num{0.676}
--- while \texttt{music@10} collapses to \num{0.358} and recall to \num{20} of
\num{1431}. There is no separating window: \SI{27}{\percent} of real transitions
are negative-cosine, so pushing the prediction away from the current track's latent
also pushes it out of the region where the true next track lives, and
\texttt{mood\_coh} collapses too (\num{0.480}$\to$\num{0.222}). Already at
$\beta=3$ the absolute relevance floor is breached. The pre-registered failure mode
fired exactly as written.

\textbf{The prior is sharpened rather than retired, and the usable band is
identified.} The 2026-07-22 probe at $\beta \le \num{0.2}$, $m=0$ reported ``weak''
--- correct in band, wrong in interpretation. That setting sits just \emph{under}
the destruction threshold, and by $\beta=3$ the same setting is a relevance
destroyer. The entire usable dynamic range at $m=0$ therefore lies between
$\beta=\num{0.2}$ and $\beta=\num{3}$, a factor of \num{15}, unexplored, bounded
above by degeneration and below by no effect. The dose ladder run here was
\textbf{not underpowered; it was aimed past the target}.

\textbf{Two second-order results worth keeping.} First, demoting $H$ to secondary
was load-bearing: A2 has by far the best facets and the \emph{worst} $H$
(\num{0.003569}), because the grounding factor annihilates it exactly as the
$\num{5.46}\times$ derivative ratio predicted. Had $H$ been primary this campaign
would have read as a flat null instead of a mapped frontier. Second, the cheapest
real de-eagering on record remains an \emph{architecture} setting, not this knob:
the free on-record comparator at \texttt{fusion\_layers}${}=1$ gives
$\Delta$concentration $-\num{0.051}$ at a $-\num{0.016}$ recall cost, and no arm
here matched that inside the relevance floor --- the floor-respecting arms delivered
between $-\num{0.003}$ and $+\num{0.0002}$, and the arms that beat $-\num{0.051}$
paid $\numrange{4}{8}\times$ the recall.

\paragraph{Why the approved later phases did not run.} All three were gated on a
phase-A win, and there is none. Phase C's arms were additionally defined \emph{by
reference} to the winner's margin and $\beta$, so with no winner they are
undefined; applied literally to the arms that passed the floors, the selection rule
would have spent three runs transferring a \emph{provably inert} setting to another
family. The runner stopped rather than improvise, which is the correct outcome: the
next move is a re-aimed design, not the approved phase C. Concretely, re-specify
the penalty rather than re-tune it --- hinge on a \emph{ranking} quantity such as
$\cos(\mathrm{pred}, x_t) - \cos(\mathrm{pred}, x_{t+1})$, or an explicit
same-artist mass term over the top-$k$ --- and, by elimination, make the eval-time
re-rank reach the serving path, since with the training-time lever refuted the only
lever that works was at that moment one the serving path could not reproduce. That
last observation is what the next section acts on.
\emph{Source: \texttt{2026-07-30-eager-margin-calibrated-refutation.md}.}

\section{Next-track: the hard per-artist cap, and the second evaluation surface}
\label{sec:nexttrack-artistcap}
\emph{Rationale.} Two \emph{soft} anti-eager levers had been refuted in the
preceding day. The training-time regularizer is mis-specified
(Section~\ref{sec:nexttrack-eagermargin}): it hinges on an absolute cosine that a
hedged InfoNCE prediction never drives above ${\sim}\num{0.6}$, while eagerness is
a \emph{ranking} property. The eval-time MMR re-rank is real
(Section~\ref{sec:nexttrack-mmr}) but measures \textbf{inert on the pathological
case}: on a real user playlist the model ranks \num{59} of its top \num{60}
candidates as one artist and the next artist first appears at \textbf{rank
\num{75}}, while MMR's entire diversity budget there is worth
$\num{0.3}\times\num{0.26}\approx\num{0.078}$ of normalized score --- nowhere near
enough to pay for a 70-rank relevance drop. Both failures share one mechanism:
\textbf{a soft penalty in score space loses to a steep relevance gradient.} The
lever tried here does not compete on score at all. A \emph{hard} per-artist cap on
the top-$k$ \textbf{cannot be outscored} --- the same property the lab's journey
cooldown already relied on --- and it is the first de-eagering lever in this project
that works.

\paragraph{The second evaluation surface, in brief.}
\label{sec:nexttrack-walksurface}
Every next-track number in the sections above is a \emph{one-shot} top-10: score
the catalog once from a held-out session prefix and look for the next distinct
track. The shipped showcase, however, serves an \emph{autoregressive walk} --- it
appends its own prediction to the prefix and steps again, twenty times --- and a run
of individually legal steps can end in a completely different region of the latent
space while every single step looks fine. Drift is in fact invisible \emph{by
construction} on the one-shot surface, because the \texttt{mood\_coh} facet of
holisticness@10 measures each candidate's cosine to the \emph{prefix} centroid,
which on a walk travels with the model. A separate measurement line --- consolidated
in full in the companion next-track document and grounded in the user's own
\num{73632} consecutive real transitions, whose median step cosine is \num{0.261}
with \SI{27.0}{\percent} of steps \emph{negative} --- therefore added four
walk-level metrics: \texttt{vibe} (mean cosine of each generated track to the
\emph{fixed seed} centroid: does the journey hold its mood), \texttt{stride\_err}
($|$median step cosine $-\,\num{0.261}|$, lower better: does it move like the
listener it imitates), \texttt{drift} (end-of-journey displacement, nearer zero
better), \texttt{genres} (distinct genres visited, counted as a \emph{credit}) and
\texttt{ground} (a degeneracy floor: graded agreement with the real continuation).
Two deterministic retrieval-time controls tune the walk without touching the model
--- a seed \emph{anchor} $a$ and a \emph{stride} term $s$, one pulling toward a fixed
point and one pushing off the last point --- so a walk configuration is written as a
cell $a$/$s$. The deployed engine is a dual-tower GRU reading \texttt{latent} and
\texttt{cummean} views with a purely linear fusion, at cell
$a\num{0.8}$/$s\num{0.5}$; it is this corpus's only measured architecture
\emph{win}, and it won on the walk surface after three campaigns had refuted its
family on recall.

\paragraph{Implementation, and why it ports where MMR cannot.} The cap is applied
to the \emph{ordering}, after MMR, in both the training-time evaluation and the
serving ranker, so evaluation and serving agree by construction --- which avoids the
divergence class fixed earlier the same day. It is threaded through all five
declaring predictor families with a default of ``off'', so every pre-existing run
and model is byte-identical. It is \textbf{best-effort rather than absolute}
(unit-verified): the scan walks the whole order looking for admissible artists, but
where a prefix's candidate space is exhausted the remaining slots are backfilled
with deferred items instead of returning fewer than $k$; identity holds when the cap
is off or $\ge k$, and unknown-artist items are never capped. Crucially it needs
\emph{no sonic vectors}, so unlike MMR it ports to a thin client unchanged --- MMR
would require shipping a \SI{48}{\mega\byte} similarity dial against the
${\sim}\SI{6}{\mega\byte}$ one already baked.

\emph{Results.} Table~\ref{tab:nexttrack-artistcap} is the cap on the one-shot
surface, measured off-run: the scores are computed once per model and memoized, so
every cap arm sees \emph{bit-identical} score vectors and the pairing is exact
rather than merely matched. Figure~\ref{fig:nexttrack-artist-cap} plots the
trade-off.

\begin{table}[ht]
  \centering
  \caption{The hard per-artist cap on the one-shot surface, canonical split
  \dsid{seq-20260715-131139} ($n=\num{1431}$), on the deployed champion and on the
  champion plus the evidence gate of Section~\ref{sec:nexttrack-markovgate}. Paired
  on bit-identical cached scores (\num{2000} resamples, rng \num{1337}). H@10 is
  generation~3 (Section~\ref{sec:nexttrack-power-generations}). R@10, MRR, mus@10,
  \texttt{ild}, H@10 higher-better; \texttt{conc}, \texttt{adj} lower-better. Best
  cell per metric bolded within each block; the recommended deployment candidate
  (gate $+$ cap 3) is bolded in full.
  \emph{Source: \texttt{2026-07-30c-artist-cap-measurement.md}.}}
  \label{tab:nexttrack-artistcap}
  {\footnotesize\setlength{\tabcolsep}{4pt}\begin{tabular}{lccccccccc}
    \toprule
    cap & R@10 & hits & MRR & conc & adj & ild & mood & mus@10 & H@10 \\
    \midrule
    \multicolumn{10}{l}{\emph{Champion \Rp+M+\Cp\ (deployed, ungated)}} \\
    off & \best{0.211740} & \best{303} & \best{0.121985} & 0.691653 & 0.609713 & 0.406374 & \best{0.398222} & 0.453790 & 0.008603 \\
    5   & 0.178896 & 256 & 0.119693 & 0.252613 & 0.337247 & 0.563417 & 0.358768 & 0.474291 & 0.027258 \\
    3   & 0.160727 & 230 & 0.116915 & 0.122975 & 0.215304 & 0.613641 & 0.339311 & \best{0.476787} & 0.033862 \\
    2   & 0.138365 & 198 & 0.112313 & \best{0.061946} & \best{0.150384} & \best{0.639717} & 0.328459 & 0.474317 & \best{0.036331} \\
    \midrule
    \multicolumn{10}{l}{\emph{Champion $+$ \texttt{markov\_gate} (the deployable candidate)}} \\
    off & \best{0.232006} & \best{332} & \best{0.133521} & 0.598214 & 0.574773 & 0.433765 & \best{0.462264} & 0.493189 & 0.018699 \\
    5   & 0.198463 & 284 & 0.131137 & 0.231338 & 0.335849 & 0.580367 & 0.413910 & 0.505703 & 0.040056 \\
    \best{3} & \best{0.175402} & \best{251} & \best{0.127788} & \best{0.113440} & 0.217191 & 0.631939 & 0.386846 & \best{0.505762} & \best{0.048063} \\
    2   & 0.150943 & 216 & 0.122627 & \best{0.057365} & \best{0.152690} & \best{0.660122} & 0.368732 & 0.503767 & 0.051175 \\
    \bottomrule
  \end{tabular}}
\end{table}

\begin{figure}[ht]
  \centering
  \includegraphics[width=0.95\textwidth]{nexttrack_artist_cap}
  \caption{\textbf{Left:} artist concentration collapses,
  \num{0.692}$\to$\num{0.113} at cap 3 on the gated champion. \textbf{Middle:} the
  result that was \emph{not} predicted --- graded relevance \emph{rises}, with the
  paired CI above zero at every cap on both models, and it replicates on the
  disjoint \SI{64}{\percent}-cold split at $+\num{0.026}$. The pre-registered worry
  was that a cap would de-eager by degrading into generic filler; the opposite
  happens, because the items promoted from rank \num{75} and beyond are \emph{more}
  musically apt to the true continuation than the same-artist tracks they displace.
  \textbf{Right:} the price, on the recall $\times$ crown plane. The cap is
  \emph{priced}, not free: gate $+$ cap 3 costs \num{52} of \num{1431} exact hits
  against the deployed status quo, while landing a holisticness@10 about
  \SI{40}{\percent} above the incumbent crown at $\num{1.5}\times$ its recall. Every
  arm clears the \num{155}-hit recall floor, so none is an instance of the
  ``maximize holisticness by predicting worse'' pathology.}
  \label{fig:nexttrack-artist-cap}
\end{figure}

\emph{Interpretation.} \textbf{\texttt{music@10} goes up, with the CI above zero, at
every cap on both models} ($+\num{0.0105}$ to $+\num{0.0230}$). This is the
campaign's most important result and it was not predicted: \textbf{the cap trades
\emph{literal} relevance for \emph{graded} relevance}. Only exact Recall@10
regresses. At gate $+$ cap 5 against the deployed status quo, recall falls
\num{0.2117}$\to$\num{0.1985} ($-\num{19}$ hits) while MRR \emph{rises}
\num{0.1220}$\to$\num{0.1311}, concentration falls \num{0.692}$\to$\num{0.231},
\texttt{ild} rises \num{0.406}$\to$\num{0.580}, \texttt{mood\_coh} rises
\num{0.398}$\to$\num{0.414}, \texttt{music@10} rises \num{0.454}$\to$\num{0.506},
and holisticness@10 rises \num{0.0086}$\to$\num{0.0401}.

\textbf{The cap does not pay for itself out of the gate's gain, and that must be
stated plainly.} The gate is $+\num{29}$ hits; the mildest cap costs $-\num{48}$ on
the gated model. So gate $+$ cap is net $\mathbf{-52}$ hits against the deployed
status quo at cap 3. Anyone who cares about exact recall should read this as a
trade-off, not a free win; the project's stated objective is session holisticness,
where it is a $\num{5.6}\times$ gain. Every cap arm clears the \num{0.108} recall
floor (\numrange{0.151}{0.198}), so these are not instances of the pathology the
floor exists to reject.

\textbf{The cap binds only where there is a problem.} On an already-varied top-10
it is a literal no-op at every cap value. The spot-check on three real user
playlists makes this concrete, counting distinct artists in the top-10: on a
playlist the champion already serves with 8 distinct artists, the cap changes
\emph{nothing} at any value. On the pathological playlist --- the case that defeated
both soft levers outright --- it goes from \num{1} distinct artist to \num{5} at cap
3, i.e.\ from ten consecutive tracks by one artist to a genuine spread across five.
And it \textbf{exactly repairs the gate's only regression}: the gate alone takes one
playlist from 8 distinct artists to 4 (concentration \num{0.044}$\to$\num{0.467}),
and at cap 3 it is back to 7 (\num{0.089}), at cap 2 to 8 (\num{0.044}) ---
identical to the deployed champion. The gate's accuracy gain therefore becomes
available \emph{without} its diversity cost.

\paragraph{Confirmed by registered runs, then hardened on a second split.} Three
registered runs on the canonical split, launched from the gated champion's
registry-declared hyperparameters with only the cap varying, reproduced
\textbf{every reported figure to every digit}
(\runid{run-20260730-235524-0ec1d-seq-blend} at cap off,
\runid{run-20260730-235524-93d38-seq-blend} at cap 3,
\runid{run-20260730-234635-a6252-seq-blend} at cap 5), so the off-run measurement
transfers exactly and the plumbing is verified end to end. Four further runs on the
disjoint \SI{64}{\percent}-cold split \dsid{seq-20260718-211238} then tested two
pre-registered legs, \emph{both} required: $\Delta$concentration
$\le-\num{0.05}$ with the CI below zero, \emph{and} a $\Delta$\texttt{music@10}
lower bound \emph{above} zero --- graded relevance had to keep \emph{improving},
not merely avoid collapse. \textbf{Both hold} at cap 3 (concentration
$-\num{0.52699}$ [$-\num{0.54546}$, $-\num{0.50825}$]; \texttt{music@10}
$+\num{0.02644}$ [$+\num{0.01919}$, $+\num{0.03368}$]), with the recall cost
\emph{smaller} and the relevance gain \emph{larger} on the colder window, and with
absolute concentration at cap 3 within \num{0.008} of the canonical figure. The
same batch delivered a third independent reproduction of the evidence gate's paired
effect, matching the recorded $+\num{0.01957}$ exactly.

\paragraph{Three cautions, one of them a lesson about floors.} First, the canonical
split's \emph{absolute} \texttt{music@10} $\ge\num{0.43}$ floor does \textbf{not}
transfer: that split's own control sits at \num{0.4295} on a legitimately colder
window, so pre-registering the rule on the paired $\Delta$ rather than an absolute
threshold is the only reason this did not read as a false failure. \textbf{Derive
floors per baseline.} Second, the measurement is at $k=\num{10}$ only, and $k$
interacts with the cap: a \emph{fixed} cap 3 at $k=\num{50}$ drives concentration
to \num{0.021} --- five times past the hardened operating point --- and reaches base
rank p95 \num{278}, whereas scaling the cap to preserve the \emph{share} reaches
only p95 \num{156} and holds concentration at \num{0.118}. Scaled is both gentler
on relevance and faithful to what was validated, and the two are identical at
$k=\num{10}$; it is implemented as \SI{30}{\percent} of slots with a floor of 3, in
both the client and the server ranker, so the two paths agree by construction.
Third, nothing here was promoted, registered as a champion or deployed, and the
recommendation on record is \texttt{markov\_gate} with cap 3 as the next deployment
candidate (cap 5 as the conservative alternative, whose MRR is still above today's).
An open eval-versus-serving row-level divergence, documented separately, means
production output will be close to rather than identical with these numbers.

\paragraph{Two engineering findings that generalize.} First, a bit-identity gate
caught an \emph{invalid launch}: the initial control ran against a stale remote
checkout containing zero occurrences of either new flag, so both hyperparameters
were silently ignored and the run reported the \emph{ungated, uncapped} champion's
metrics under gated hyperparameters --- a wrong number, not a crash. The lesson is
sharper than ``resync the worker'': \textbf{checksum the predictor files for the
family being launched immediately before posting}, not at design time, because the
preceding campaign had verified two \emph{other} families' files. Second, the same
incident left an orphaned run row stuck at \texttt{running}, which neither the stop
nor the delete endpoint can clear (the former tracks only local runs, the latter
refuses runs claimed by a worker) --- a gap worth closing before it happens again.
Both levers were subsequently ported to the two showcase clients as pure
post-processing, with no model rebuild, no new vectors and no retraining, and with
parity tests pinning the ported ordering against the Python reference; in one client
the champion engine turns out to be \textbf{degenerate uncapped} --- thirty tracks by
one artist on an Argentine-rock seed --- so there the cap is load-bearing rather
than cosmetic.
\emph{Source: \texttt{2026-07-30c-artist-cap-measurement.md}.}

\section{Next-track: stage-stack topology: depth refuted on both surfaces}
\label{sec:nexttrack-stagestack}
\emph{Rationale.} The user's question was direct: does an explicit \emph{stack} of
per-step MLP and recurrent stages beat a single GRU, on a corpus where five
consecutive architecture campaigns have returned nulls? Four arrangements were
specified --- \texttt{ann+gru+ann+gru}, \texttt{ann+gru+gru+ann},
\texttt{ann+gru+gru+gru}, \texttt{gru+gru+gru+gru} --- plus a fifth,
\texttt{gru+ann+gru+gru}, chosen by the user to test \emph{why} a leading MLP hurt
at $N=1$: Section~\ref{sec:nexttrack-preencoder} attributed that loss to half-wave
rectification destroying signed prefix information that a GRU's gated linear input
transform preserves, so putting the MLP on a \emph{gated hidden state} instead of
raw signed PCA should make the harm vanish. This is also the first campaign in this
arc to be judged on \emph{two} surfaces by design rather than in retrospect, and it
sets the standard the campaign after it inherits.

\paragraph{Design, and the equivalence proof that made a control possible.} Seven
runs of a new predictor, \texttt{seq-stack}, on the canonical split
\dsid{seq-20260715-131139} (\num{1431} test sessions), \SI{90}{\minute} of compute,
zero failures, zero contingency runs spent. Capacity was equalised at
${\sim}\num{1282000}$ parameters across the five 4-stage arms (mutual spread
\SI{0.445}{\percent}), and all seven counts were re-verified by instantiating the
model rather than trusted from the design. Everything but the stage string and the
widths is held: dropout \num{0.1}, InfoNCE, $\tau$~\num{0.07}, \num{40} epochs,
patience \num{6}, lr \num{0.001}, batch \num{128}, seed \num{1337}, MMR off, artist
cap off.

Bit-identity with the family's historical \num{176} hits is \emph{not available}
here and its absence is not a failure: the trainer re-seeds on entry, so handing it
a pre-built model offsets the dropout stream even when the initial weights match.
The campaign therefore gated on a \textbf{functional-equivalence proof} plus an
\textbf{in-batch control}, and every verdict is stated against that control. The
proof was run \emph{before} the registry edit and the server restart, so a failure
would have cost nothing: at \texttt{stages=["gru"]} the new model computes the same
function as the incumbent single GRU, with $\max|\Delta|$ \textbf{\texttt{0.00e+00}}
on both \texttt{forward} and \texttt{predict\_next} at prefix lengths
$T=1,2,5,17,40$ and \num{394944} parameters on both sides. The in-batch control
then returned \textbf{\num{178} of \num{1431}}, inside the pre-registered band
$[\num{155},\num{198}]$ and only \num{2} hits off the historical \num{176} ---
exactly the small offset the re-seeding argument predicts. Two further gates
mattered: metric completeness was green on all seven runs (no worker-image key
stripping of the kind Section~\ref{sec:nexttrack-mmr} documents), and the walk instrument
was verified against the shipped engine, whose $\Delta$\texttt{vibe}
($-\num{0.0033}$) and $\Delta$\texttt{stride\_err} ($+\num{0.0137}$) both straddle
zero against the freshly-trained control. The walk harness additionally reproduced
\emph{all ten} previously recorded means exactly and stayed bit-stable as arms were
added, which is what licenses pairing across separate invocations.

\emph{Results, the one-shot surface.} Table~\ref{tab:nexttrack-stagestack-s1}.

\begin{table}[ht]
  \centering
  \caption{The stage-stack batch on the one-shot surface, seven runs on
  \dsid{seq-20260715-131139} ($n_{\mathrm{test}}=\num{1431}$), sorted by
  Recall@10. Run ids are \protect\runid{run-20260801-} plus the chunk time and hex
  suffix, in full: \texttt{003536-eaf81} (C0), \texttt{003536-bcedf} (C1),
  \texttt{003646-07994} (B), \texttt{003646-8cc2d} (E), \texttt{003646-dbe68} (C),
  \texttt{003536-d612e} (D), \texttt{003646-6e3bf} (A). Stage codes: \texttt{a}
  $=$ per-step MLP, \texttt{g} $=$ GRU. MRR is truncated at \num{10} and
  recomputed from the stored predictions (the stored MRR is over the full ranking
  and is therefore not pairable). Every H@10 was recomputed from
  \texttt{predictions.json} under the live definition and matches the stored value
  exactly. The paired $\Delta$ is against the in-batch control C0 (\num{2000}
  resamples, rng \num{1337}). Truncated MRR@10 and \texttt{mood\_coh} are omitted
  for width and track recall closely: MRR@10 \num{0.04555} (C0), \num{0.04579}
  (C1), then \numrange{0.02080}{0.02894} across the five 4-stage arms;
  \texttt{mood\_coh} \num{0.4788} (C0), \num{0.4856} (C1), then
  \numrange{0.3821}{0.3964}. R@10, H@10, mus@10
  higher-better; \texttt{conc}, \texttt{adj} lower-better --- and the
  lowest-eagerness cells belong to \emph{refuted} arms, which is what the recall
  floor is for. \texttt{ild} rises monotonically as the arms get worse
  (\num{0.4549} on C0, \num{0.4424} on C1, then \numrange{0.5112}{0.5373} across
  the five 4-stage arms) and is omitted here for width.
  \emph{Source: \texttt{2026-08-01-stage-stack-topology.md}.}}
  \label{tab:nexttrack-stagestack-s1}
  {\scriptsize\setlength{\tabcolsep}{2.0pt}\begin{tabular}{lllrllllll}
    \toprule
    Arm & id & stages / width & params & R@10 (hits) & $\Delta$R@10 [95\% CI] & H@10 & conc & adj & mus@10 \\
    \midrule
    C0 \emph{ctrl} & \texttt{eaf81} & \texttt{gru} h256 & \num{394944} & \best{0.12439 (178)} & --- & \best{0.024378} & 0.3589 & 0.3482 & \best{0.4552} \\
    C1 \emph{cap.} & \texttt{bcedf} & \texttt{gru} h537 & \num{1280937} & 0.12159 (174) & $-0.00280$ [$-.0119$, $+.0056$] & 0.023418 & 0.3773 & 0.3653 & 0.4556 \\
    B & \texttt{07994} & \texttt{agga} h299 & \num{1281407} & 0.08595 (123) & $-0.03843$ [$-.0538$, $-.0231$] & 0.017558 & 0.3176 & 0.2639 & 0.4264 \\
    E & \texttt{8cc2d} & \texttt{gagg} h259 & \num{1277321} & 0.08246 (118) & $-0.04193$ [$-.0566$, $-.0280$] & 0.021053 & 0.2704 & 0.2376 & 0.4310 \\
    C & \texttt{dbe68} & \texttt{aggg} h256 & \num{1283008} & 0.07966 (114) & $-0.04472$ [$-.0601$, $-.0301$] & 0.021442 & \best{0.2626} & 0.2375 & 0.4324 \\
    D & \texttt{d612e} & \texttt{gggg} h229 & \num{1282821} & 0.07687 (110) & $-0.04752$ [$-.0622$, $-.0328$] & 0.018529 & 0.2644 & 0.2326 & 0.4222 \\
    A & \texttt{6e3bf} & \texttt{agag} h299 & \num{1281407} & 0.07058 (101) & $-0.05381$ [$-.0685$, $-.0391$] & 0.017374 & 0.2730 & \best{0.2335} & 0.4168 \\
    \bottomrule
  \end{tabular}}
\end{table}

No arm meets any win clause. Every 4-stage arm loses recall \emph{and}
holisticness by paired CI, and all five sit \textbf{below the mandatory
\num{0.108} absolute recall floor} (\numrange{0.0706}{0.0860}). Derived
per-baseline from the control's own \texttt{music@10} of \num{0.45519}, every
4-stage arm is \SIrange{5.0}{8.4}{\percent} below the relevance floor too, so none
is even relevance-neutral. The losses are $\numrange{2.5}{3.7}\times$ the measured
half-width: this is not a noise-band question.

\emph{Results, the walk surface.} Table~\ref{tab:nexttrack-stagestack-s2}, at the matched
primary cell $a\num{0.4}$/$s\num{0.3}$, paired against the in-batch control on the
walk metrics of Section~\ref{sec:nexttrack-walksurface}. \texttt{vibe} is always reported with
\texttt{stride\_err} beside it, at matched stride, per the standing caveat that a
vibe gain bought by taking shorter steps is not a gain.

\begin{table}[ht]
  \centering
  \caption{The stage-stack batch on the generated walk: \num{40} held-out sessions
  $\times$ \num{20} steps at artist cap 1, paired against the \emph{in-batch}
  control C0 (\num{2000} resamples). The refuted clause is
  ``$\Delta$\texttt{vibe} CI below zero \emph{and} $\Delta$\texttt{stride\_err}
  not CI below zero''. \texttt{vibe}, \texttt{genres}, \texttt{ground}
  higher-better; \texttt{stride\_err} lower-better; \texttt{drift} nearer zero
  better. CI bounds are quoted where the interval clears zero.
  \emph{Source: \texttt{2026-08-01-stage-stack-topology.md}.}}
  \label{tab:nexttrack-stagestack-s2}
  {\footnotesize\setlength{\tabcolsep}{2.6pt}\begin{tabular}{lccccclc}
    \toprule
    Arm & vibe & stride\_err & drift & genres & ground & $\Delta$vibe [95\% CI] & V2 \\
    \midrule
    GRU shipped $a$0.4 $s$0.0 & \best{0.609} & 0.399 & $-0.089$ & 9.925 & 0.345 & $+0.0376$ CI$>$0 & reference \\
    shipped engine $a$0.4 $s$0.3 & 0.568 & 0.129 & $-0.080$ & \best{10.450} & \best{0.357} & $-0.0033$ straddles & instrument gate \\
    \textbf{C0} $a$0.4 $s$0.3 & 0.571 & 0.116 & $-0.080$ & 10.325 & 0.344 & base & base \\
    C1 \texttt{gru} h537 & 0.561 & 0.122 & $-0.082$ & 10.625 & 0.359 & $-0.0100$ straddles & indistinguishable \\
    D \texttt{gggg} & 0.557 & 0.122 & $-0.071$ & 10.225 & 0.319 & $-0.0140$ [$-.0255$, $-.0022$] & REFUTED $\times2$ \\
    C \texttt{aggg} & 0.548 & \best{0.113} & $-0.064$ & 10.250 & 0.340 & $-0.0231$ [$-.0381$, $-.0084$] & REFUTED \\
    B \texttt{agga} & 0.545 & 0.126 & \best{$-0.061$} & 9.675 & 0.335 & $-0.0258$ [$-.0421$, $-.0108$] & REFUTED \\
    E \texttt{gagg} & 0.540 & 0.109 & $-0.084$ & 9.225 & 0.335 & $-0.0316$ [$-.0494$, $-.0144$] & REFUTED \\
    A \texttt{agag} & 0.534 & 0.117 & $-0.085$ & 9.125 & 0.336 & $-0.0375$ [$-.0559$, $-.0214$] & REFUTED \\
    \bottomrule
  \end{tabular}}
\end{table}

Every 4-stage arm fires the refuted clause. D fires it \emph{twice}, also failing
the grounding floor ($\Delta$\texttt{ground} $-\num{0.0250}$ [$-\num{0.0474}$,
$-\num{0.0048}$]), and A additionally loses genre variety ($-\num{1.200}$
[$-\num{2.125}$, $-\num{0.300}$]). No arm approaches a win clause. Two robustness
checks close the escape routes. A pre-registered \textbf{stride probe} over
$s\in\{\num{0.2},\num{0.3},\num{0.4}\}$ was run for the only two arms with any
directional movement toward a walk win, and \emph{both peak at $s\num{0.3}$ --- the
cell they were already measured in} (C: vibe/stride\_err \num{0.586}/\num{0.183},
\num{0.548}/\num{0.113}, \num{0.517}/\num{0.153}; E: \num{0.576}/\num{0.155},
\num{0.540}/\num{0.109}, \num{0.516}/\num{0.174}); at $s\num{0.4}$ both lose vibe
with the CI below zero while \texttt{stride\_err} gets \emph{worse}. And the single
best arm was re-measured at the \textbf{deployed} cell $a\num{0.8}$/$s\num{0.5}$,
where it is refuted again ($\Delta$\texttt{vibe} $-\num{0.0158}$
[$-\num{0.0267}$, $-\num{0.0059}$], $\Delta$\texttt{stride\_err} straddling).

\begin{figure}[ht]
  \centering
  \includegraphics[width=0.95\textwidth]{nexttrack_stage_stack}
  \caption{\textbf{Left:} all five 4-stage arms fall \emph{through} the mandatory
  absolute recall floor, while the $\num{3.24}\times$ capacity control ties the
  anchor --- so the damage is depth, not capacity. Error bars are the paired
  bootstrap CI against the in-batch control, mapped onto the hit scale.
  \textbf{Middle:} the campaign's most interesting number. Arm B reaches the
  \emph{best} validation InfoNCE in the batch and retrieves \num{55} fewer hits
  than the control: the objective the whole family trains on and top-10 retrieval
  quality come apart at depth, so any architecture search selecting on validation
  loss would pick the wrong arm. \textbf{Right:} the two-surface closure. The walk
  read at the matched primary cell refutes the same five arms that the one-shot
  surface refutes; the capacity control ties on both. The axis does not merely fail
  to close on one surface --- it closes on two.}
  \label{fig:nexttrack-stage-stack}
\end{figure}

\emph{Interpretation.} \textbf{The axis is closed, and it is the first null in this
arc to close on both surfaces.} All five 4-stage arms are refuted by paired CI on
one-shot recall \emph{and} on the walk with the instrument gate green, and the
capacity control is indistinguishable on both. This is the sixth consecutive
architecture null on this corpus and, unlike the five before it, it satisfies the
standing rule that a one-shot verdict says nothing about the surface the product
serves.

\textbf{The user's mechanism hypothesis is refuted, and that is the most reusable
fact here.} Arm E put the MLP on a \emph{gated hidden state} rather than on raw
signed PCA, which the rectifier story predicts should remove the harm. It does not:
E reaches \num{118} of \num{1431}, ties two other arrangements, beats only the
worst, and comes nowhere near the control's \num{178}. So half-wave rectification
of the \emph{raw input} is not the mechanism --- whatever is wrong survives moving
the rectifier downstream of a gate. One run bought a clean refutation of a
plausible and cheap-to-believe story.

\textbf{The damage is depth itself, and three controls close the alternatives
simultaneously.} It is not capacity: the $\num{3.24}\times$ control ties on both
surfaces. It is not regularization or early stopping: \emph{every} 4-stage arm
trains \emph{longer} than the control (\numrange{21}{40} epochs against \num{14}),
and two of five reach a \emph{lower} best validation loss, so the pre-authorized
dropout-count contingency failed both of its trigger clauses and was not spent ---
the four-fold dropout application is exonerated. And it is not the MLP: a leading
per-step MLP at depth 4 is recall-neutral and \emph{holisticness-positive} against
the pure-recurrent arrangement ($\Delta H$ $+\num{0.00291}$, CI above zero). What
remains is the depth of the recurrent chain, and its cost is large ---
\numrange{55}{77} hits of \num{1431}.

\textbf{Bracketing beats interleaving at byte-identical capacity.}
\texttt{ann+gru+gru+ann} beats \texttt{ann+gru+ann+gru} by $+\num{0.01537}$ recall
with the CI above zero at \emph{the same parameter count by formula} --- the
tightest-controlled topology comparison the campaign could make. Interposing a
per-step MLP \emph{between} two recurrences is worse than putting MLPs on the
outside, which is consistent with the depth reading: breaking the recurrent chain
in the middle costs more than adding depth at the ends. Two further contrasts are
flat: the number of recurrent stages (2 versus 3) straddles on recall, and the
rectifier's position buys nothing beyond beating the worst arrangement.

\textbf{A leading per-step MLP is \emph{not} harmful at depth 4, which narrows
rather than confirms the $N=1$ result.} Section~\ref{sec:nexttrack-preencoder}'s ``leading
MLP is harmful'' was measured at $N=1$ against a single GRU; at depth 4 the same
addition is free-to-slightly-good. The pre-encoder pitfall should be scoped to the
shallow case.

\textbf{Objective and retrieval dissociate at depth, and this is the sharpest
instance yet.} Arm B achieves the \emph{lowest} validation InfoNCE in the entire
batch (\num{6.5378} against the control's \num{6.5832}) while retrieving
\textbf{\num{55} fewer hits} (\num{123} against \num{178}). Every prior
architecture null on this corpus was consistent with ``the extra machinery does
nothing''; this one shows the extra machinery does something measurable to the loss
and that it is the \emph{wrong} something. The operational consequence is blunt:
select on Recall@10 and the walk, never on the training objective. One caveat rides
with it --- B ran to the \num{40}-epoch cap with patience never triggering, so its
recall is a \emph{lower bound}, and a longer re-run is the clean test of the
dissociation (the prediction being a better validation loss and no better recall).

\textbf{The eagerness facets improve monotonically as the arms get worse, and the
floors caught it.} Every 4-stage arm de-eagers substantially (concentration
\num{0.359}$\to$\numrange{0.263}{0.318}) and raises \texttt{ild}; under a naive
holisticness-primary read some would look interesting. All are killed by the
absolute recall floor and by the per-baseline relevance floor --- the same mechanism
that has now saved four campaigns. Note that the canonical \emph{absolute}
\num{0.43} relevance figure would have passed several of these arms, so the floor
must be derived per baseline.

\paragraph{Two by-products, and the governance line.} A previously closed $N=2$
relative --- a dual tower with a pre-MLP on one side, closed on a one-shot recall
tie --- was walk-evaluated for the first time here and is \textbf{refuted on the
walk} ($\Delta$\texttt{vibe} $-\num{0.0255}$ [$-\num{0.0432}$, $-\num{0.0096}$]),
its one real gain being less drift. That is the same signature a recurrent-cell
variant showed earlier --- better stride and drift, worse vibe --- which now looks
like a general property of adding pre-recurrence machinery rather than anything
specific to either arm. Second, the grounding metric again discriminated almost
nothing, separating one arm of nine, consistent with its recorded arithmetic
(between-arm range \num{0.0060} against a per-session standard deviation of
\num{0.1417}, needing ${\sim}\num{2140}$ sessions): it is a floor that catches
degeneracy, not a ranking metric, and it was used only that way. Nothing was
promoted as champion, nothing was registered as a crown claim and nothing was
deployed; the seven arms were promoted \emph{for measurement only}, each carrying
the verbatim note ``promoted to enable walk evaluation; not a crown claim'', which
flooded the best-models group and was handed to the curation agent.

\paragraph{Two process findings worth keeping.} First, the predictor-checksum gate
fired \emph{on a file the design did not name for copying} --- it listed one module
to copy and a shared module only to \emph{verify}, and the verification failed. The
divergence turned out to be inert for this batch, but that is a \emph{post-hoc}
judgement that required reading the diff; the rule that caught it is the one worth
keeping. \textbf{Verify the checksum of every file in the family's import closure,
not just the ones you copied.} Second, the approved launch body was rejected by the
approved registry block: the design's JSON sent six hyperparameters the block does
not declare, so the first posts failed outright. It was resolved by dropping the
six keys \emph{after proving the drop is a no-op} --- the predictor's defaults supply
exactly those values and the resolved hyperparameter dict is identical either way,
verified by direct comparison --- and the registry block was left exactly as
approved.
\emph{Source: \texttt{2026-08-01-stage-stack-topology.md}.}

\section{Next-track: parallel tower banks closed on both surfaces}
\label{sec:nexttrack-towerbank}
This is the seventh architecture campaign on the next-track corpus and the one
that closes the \emph{parallel} branch of the axis. It is the second campaign
consolidated here judged on \textbf{both} evaluation surfaces, to the standard
Section~\ref{sec:nexttrack-stagestack} set, so it needs one paragraph of
orientation before the results.

\paragraph{The second surface, in brief.} Every next-track number in the
sections above is a \emph{one-shot} top-10: score the catalog once from a session
prefix and look for the held-out next track. The shipped showcase, however, serves
an \emph{autoregressive walk} --- it appends its own prediction to the prefix and
steps again, twenty times --- and a run of individually legal steps can end in a
completely different region of the latent space while every single step looks
fine. A separate measurement line (introduced in
Section~\ref{sec:nexttrack-walksurface} and consolidated in full in the companion
next-track document, and grounded in the user's own \num{73632} consecutive real
transitions, whose median step cosine is \num{0.261}) therefore added four
walk-level metrics: \texttt{vibe} (mean cosine of each generated track to the
\emph{seed} centroid --- does the journey hold its mood), \texttt{stride\_err}
($|$median step cosine $-\,\num{0.261}|$, lower better --- does it move like the
listener it imitates), \texttt{drift} (end-of-journey displacement, nearer zero
better), \texttt{genres} (distinct genres visited) and \texttt{ground} (a
degeneracy floor: agreement with the real continuation). Two deterministic
retrieval-time controls tune the walk without touching the model --- a seed
\emph{anchor} $a$ and a \emph{stride} term $s$ --- so a walk configuration is
written as a cell $a$/$s$. The deployed engine is a dual-tower GRU reading
\texttt{latent} and \texttt{cummean} views with a purely linear fusion
(\texttt{fusion\_layers}${}=0$) at cell $a\num{0.8}$/$s\num{0.5}$; it is this
corpus's only measured architecture \emph{win}, and it won on the walk surface
after three campaigns had refuted its family on recall.

\emph{Rationale.} That win came with a stated mechanism: the tower carries a
\textbf{fixed reference never overwritten by recurrence}, which is what holds a
mood. A mechanism of that shape makes a testable prediction --- it should be
\emph{dose-responsive}. If one fixed-reference tower buys a held vibe plus
$+\num{1.5}$ genres, then $N$ towers, enough of them on the raw \texttt{latent}
view to hold the recall floor and the rest on \texttt{cummean}, should buy
\emph{more} of the same. This campaign runs that dose curve.

The vehicle is a new predictor, \texttt{seq-bank}: $N$ \emph{parallel}
single-layer bare GRU towers over configurable causal views, their per-step hidden
states concatenated and read by \textbf{one linear head}, with tower count and
view multiset collapsed into a \emph{single} \texttt{views} hyperparameter so that
$N$ cannot drift out of sync with the views. What it cannot express matters as
much as what it can: it declares exactly \num{15} hyperparameters, and none of
\texttt{fusion\_layers}, \texttt{fusion\_hidden}, \texttt{pre\_hidden},
\texttt{pre\_layers}, \texttt{layers}, \texttt{num\_layers}, \texttt{arch},
\texttt{bidirectional}, \texttt{residual} or \texttt{eager\_beta} is among them.
Every mechanism the sections above already refuted --- the MLP fusion head
(Section~\ref{sec:nexttrack-dualtower}), the per-step pre-encoder
(Section~\ref{sec:nexttrack-preencoder}), the soft anti-eager regularizer
(Section~\ref{sec:nexttrack-holisticness}) --- is therefore \emph{structurally
unreachable} rather than merely unused, and a \texttt{parse\_views} guard rejects
any spec without a \texttt{latent} tower, putting the registered
\numrange{0.058}{0.087} \texttt{delta}/\texttt{cummean} pathologies out of reach
as well. Two hypotheses were pre-registered: the \emph{dose} hypothesis above,
and the independent sub-hypothesis that if the raw \texttt{latent} view
\emph{subsumes} its siblings (the dual-tower campaign's mechanism), then $N$
identical \texttt{latent} towers buy nothing at any $N$.

\paragraph{Design.} Ten runs on the canonical split
\dsid{seq-20260715-131139} (\num{1431} test sessions, cold-item rate
\num{0.5507}), \SI{84.2}{\minute} of compute, zero failures. Everything but
\texttt{views} / \texttt{hidden} / \texttt{dropout} is held: \texttt{loss}
InfoNCE, $\tau=\num{0.07}$, \num{40} epochs, patience \num{6}, lr \num{0.001},
batch \num{128}, seed \num{1337}, MMR off, artist cap off. \textbf{Capacity is
matched at \num{3.00}$\times$ the h256 baseline} for every bank arm \emph{and}
for the single-GRU control C1 (h512), with counts re-verified from the artifacts
rather than trusted from the design --- \texttt{seq\_bank} emits
\texttt{n\_params} into \texttt{metrics.json}, and the mutual spread across the
seven \num{3.00}$\times$ arms including C1 is \SI{0.405}{\percent}. The decision
rule was written on both surfaces before launch. On \textbf{S1} (one-shot
recall@10) a win requires a paired $\Delta$ above zero against \emph{both} the
h256 anchor C0 \emph{and} the parameter-matched C1, subject to an absolute recall
floor ($\ge\num{0.108}$, i.e.\ \num{155} of \num{1431}) and a per-baseline
relevance floor (paired $\Delta$\,music@10 vs C0 must not fall below zero). On
\textbf{S2} (the walk: \num{40} held-out sessions $\times$ \num{20} steps at
artist cap 1, paired against the \emph{in-batch} C0) the primary cell is
$a\num{0.4}$/$s\num{0.3}$, and a saved definition additionally requires
\textbf{the same sign at the deployed cell} $a\num{0.8}$/$s\num{0.5}$.

\paragraph{The instrument was pinned before anything was measured.} All ten
pre-registered gates came back green, and three are worth recording because they
make the comparison unusually tight. A \emph{functional-equivalence} gate showed
\texttt{seq-bank} to be a strict generalization of both incumbents rather than a
new family: at \texttt{views="latent"} it reproduces the single-GRU
\texttt{SeqNextLatent} and at \texttt{views="latent+latent"} the
\texttt{fusion\_layers}${}=0$ \texttt{DualTowerNextLatent}, with $\max|\Delta|$
\textbf{\texttt{0.00e+00}} on both \texttt{forward} and \texttt{predict\_next} at
prefix lengths $T=1,2,5,17,40$ and identical parameter counts on both sides; an
extra training-stream check matched per-epoch train \emph{and} validation loss to
the last digit over three epochs, with all ten state-dict tensors at
$\max|\Delta|$ \texttt{0.00e+00}. Both bit-identity gates then hit exactly: C0
reproduced its on-record recall@10 \num{0.12299091544374564} $=$ \num{176} of
\num{1431} --- a \textbf{10th} reproduction, with \texttt{music@10}, H@10 and MRR
matching to the last digit --- and R2 reproduced the dual-tower
\texttt{latent/latent} f0 arm at \num{0.1187980433263452} $=$ \num{170} of
\num{1431}, so the design's band fallback was never needed. The walk instrument
was checked against the shipped \texttt{gru.onnx}: $\Delta$\texttt{vibe}
$+\num{0.0028}$ and $\Delta$\texttt{stride\_err} $+\num{0.0005}$ both straddle
zero, so S2 verdicts against the in-batch C0 are not an artefact of the base.

\emph{Results, S1.} Table~\ref{tab:nexttrack-towerbank-s1} is the batch.
\textbf{No arm is bolded as a winner because there is none}: no arm achieves even
the first leg of the win rule, all nine non-anchor arms are ties against C0, and
\emph{five of the eight bank arms lose to the parameter-matched single GRU C1 with
the CI below zero}. Both floors pass for every arm --- the worst is V1 at
\num{162} hits, clear of the \num{155} floor, and the most negative relevance
$\Delta$ is V2's $-\num{0.006109}$, straddling.

\begin{table}[ht]
  \centering
  \caption{The tower-bank batch on S1, ten runs on \dsid{seq-20260715-131139}
  (\num{1431} test sessions), sorted by hits. The batch was launched in four
  chunks; every run id is \protect\runid{run-20260802-} plus the chunk time and
  the hex suffix in the id column, in full: \texttt{182050-91b54},
  \texttt{182050-2e6f2}, \texttt{183404-91e31}, \texttt{183404-62604},
  \texttt{183404-67185}, \texttt{190415-1d662}, \texttt{190415-90866},
  \texttt{190415-b5928}, \texttt{192627-17502}, \texttt{192627-836b1}. Paired
  $\Delta$ is a per-session paired bootstrap (\num{2000} resamples, rng
  \num{1337}) computed from the stored per-row \texttt{predictions.json}; a CI
  straddling zero is a \emph{tie}, and \textsuperscript{$\ast$} marks a $\Delta$
  vs C1 whose CI lies entirely \emph{below} zero. d\num{0.1}\,/\,d\num{0.0} is
  the dropout of the byte-identical D2\,/\,R2 pair. Capacity: the seven arms C1,
  B2, B3, B4, V1, V2 and V3 span \numrange{2.996}{3.007}$\times$ the h256
  baseline (mutual spread \SI{0.405}{\percent}), D2 and R2 sit at exactly
  \num{2.000}$\times$ and C0 at \num{1.000}$\times$. View codes:
  lat\,$=$\,\texttt{latent}, cum\,$=$\,\texttt{cummean},
  del\,$=$\,\texttt{delta}. H@10 is the \emph{generation-3} holisticness
  (Section~\ref{sec:nexttrack-power-generations}), in raw next-track units. R@10,
  H@10 and mus@10 higher-better; a\_conc is lower-better but display-only and is
  deliberately \emph{not} bolded --- no arm's eagerness is a claim in this
  campaign. Best cell per metric bolded; C1 and B4 tie exactly on R@10 and both
  are bolded.
  \emph{Source: \texttt{2026-08-02-tower-bank-parallel-towers.md}.}}
  \label{tab:nexttrack-towerbank-s1}
  {\scriptsize\setlength{\tabcolsep}{1.9pt}\begin{tabular}{lllrlllccc}
    \toprule
    Arm & id & views / width & params & R@10 (raw) & $\Delta$ vs C0 [95\% CI] & $\Delta$ vs C1 & H@10 & mus@10 & a\_conc \\
    \midrule
    C1 \emph{ctrl} & \texttt{91e31} & 1$\times$GRU h512 & \num{1182912} & \best{0.12998 (186)} & $+0.0070$ [$-.002$, $+.016$] & --- & 0.02380 & 0.4581 & 0.3825 \\
    B4 & \texttt{1d662} & lat$\times$4 h211 & \num{1187700} & \best{0.12998 (186)} & $+0.0070$ [$-.002$, $+.017$] & $+0.0000$ & 0.02446 & \best{0.4594} & 0.3716 \\
    D2 & \texttt{62604} & lat$\times$2 h256 d\num{0.1} & \num{789696} & 0.12369 (177) & $+0.0007$ [$-.008$, $+.008$] & $-0.0063$ & 0.02426 & 0.4567 & 0.3698 \\
    C0 \emph{base} & \texttt{91b54} & 1$\times$GRU h256 & \num{394944} & 0.12299 (176) & --- & $-0.0070$ & \best{0.02454} & 0.4558 & 0.3559 \\
    B3 & \texttt{90866} & lat$\times$3 h256 & \num{1184448} & 0.12020 (172) & $-0.0028$ [$-.010$, $+.006$] & $\mathbf{-0.0098}$\textsuperscript{$\ast$} & 0.02358 & 0.4554 & 0.3764 \\
    R2 & \texttt{2e6f2} & lat$\times$2 h256 d\num{0.0} & \num{789696} & 0.11880 (170) & $-0.0042$ [$-.013$, $+.004$] & $\mathbf{-0.0112}$\textsuperscript{$\ast$} & 0.02378 & 0.4554 & 0.3706 \\
    B2 & \texttt{67185} & lat$\times$2 h334 & \num{1186560} & 0.11880 (170) & $-0.0042$ [$-.013$, $+.003$] & $\mathbf{-0.0112}$\textsuperscript{$\ast$} & 0.02336 & 0.4554 & 0.3799 \\
    V3 & \texttt{17502} & lat$+$cum$+$cum h256 & \num{1184448} & 0.11880 (170) & $-0.0042$ [$-.016$, $+.006$] & $-0.0112$ & 0.02320 & 0.4497 & 0.3581 \\
    V2 & \texttt{836b1} & lat$+$cum$+$del h256 & \num{1184448} & 0.11600 (166) & $-0.0070$ [$-.020$, $+.005$] & $\mathbf{-0.0140}$\textsuperscript{$\ast$} & 0.02245 & 0.4497 & 0.3547 \\
    V1 & \texttt{b5928} & lat$\times$2$+$cum h256 & \num{1184448} & 0.11321 (162) & $-0.0098$ [$-.022$, $+.001$] & $\mathbf{-0.0168}$\textsuperscript{$\ast$} & 0.02299 & 0.4500 & 0.3572 \\
    \bottomrule
  \end{tabular}}
\end{table}

\paragraph{The sharpest number in the batch is a zero.} B4 --- four independent
h211 \texttt{latent} towers, \num{1187700} parameters --- and C1 --- one h512
GRU, \num{1182912} parameters --- return the \textbf{identical \num{186} hits of
\num{1431}}, a paired $\Delta$ of precisely $\num{0.000000}$
$[-\num{0.00839},+\num{0.00839}]$. At a matched ${\sim}\num{3.00}\times$ budget,
\emph{topology and width are interchangeable on this corpus}: whatever four
towers with different random initializations discover, a single wider recurrence
already had. Figure~\ref{fig:nexttrack-tower-bank} (left) plots the $N$-axis at
matched capacity, and the line through $\num{186}/\num{170}/\num{172}/\num{186}$
is flat inside the split's own measured resolution
(Section~\ref{sec:nexttrack-power-stats}).

This is also the cleanest demonstration on record of \emph{why the
parameter-matched control leg is mandatory}. Read against the h256 anchor alone,
B4's $+\num{10}$ hits would have looked like a topology gain and the five arms
that lose to C1 would all have read as ties --- the anchor cannot distinguish
``the bank helped'' from ``the extra parameters helped, and the bank spent them
worse than width would have''. The C1 leg does exactly that, and it is what turns
a batch of nine ties into a closure. Note what the flatness costs the \emph{other}
direction: C1's $+\num{10}$ hits over C0 is itself not significant
($+\num{0.006988}$ $[-\num{0.00210},+\num{0.01607}]$), the \textbf{seventh}
parameter-matched width tie on record, so the campaign cannot claim h512 as an
improvement either despite it holding the highest single-GRU hit count measured
(\num{186} against h454's \num{179}). The single-GRU width points on PCA-192 now
read h256\,$\to$\,\num{176}, h297\,$\to$\,\num{162}, h425\,$\to$\,\num{173},
h454\,$\to$\,\num{179}, h512\,$\to$\,\num{186}, h537\,$\to$\,\num{174} ---
non-monotone, spanning \num{0.0168}, with no trend.

\emph{Results, S2.} Table~\ref{tab:nexttrack-towerbank-s2} is the walk. Two arms
--- B2 (\texttt{latent}$+$\texttt{latent} h334) and V3
(\texttt{latent}$+$\texttt{cummean}$+$\texttt{cummean}) --- reached a genuine
\textbf{TRADE} at the primary cell, each buying a drift improvement with the CI
above zero while holding vibe. Both then \textbf{flip} at the deployed cell: the
drift gain evaporates into a straddle and \texttt{vibe} --- the non-negotiable
line --- goes below zero for both. Under the pre-agreed rule, \textbf{no
definition was saved}.

\begin{table}[ht]
  \centering
  \caption{The tower-bank batch on S2, the generated walk: \num{40} held-out
  sessions $\times$ \num{20} steps at artist cap 1, paired against the
  \emph{in-batch} C0 (\num{2000} resamples). The upper block is the primary cell
  $a\num{0.4}$ / $s\num{0.3}$; the lower block re-measures the two TRADE arms at
  the \emph{deployed} cell $a\num{0.8}$ / $s\num{0.5}$, where the save clause
  requires the same sign. \texttt{vibe}, $\Delta$drift, $\Delta$genres and
  $\Delta$ground higher-better; \texttt{stride\_err} lower-better. CI bounds are
  quoted only where the interval clears zero, as in the source report, and
  \textsuperscript{$\ast$} marks every such interval --- a \emph{gain} on
  $\Delta$drift, a \emph{refutation} on $\Delta$vibe and $\Delta$genres.
  \emph{Source: \texttt{2026-08-02-tower-bank-parallel-towers.md}.}}
  \label{tab:nexttrack-towerbank-s2}
  {\scriptsize\setlength{\tabcolsep}{2.6pt}\begin{tabular}{lccllllc}
    \toprule
    Arm & vibe & stride\_err & $\Delta$vibe & $\Delta$stride\_err & $\Delta$drift & $\Delta$genres & $\Delta$ground \\
    \midrule
    \multicolumn{8}{l}{\emph{Primary cell $a\num{0.4}$ / $s\num{0.3}$}} \\
    C0 (base) & 0.566 & 0.130 & --- & --- & --- & --- & --- \\
    \textbf{B2} \emph{TRADE} & 0.567 & 0.139 & $+0.0017$ & $+0.0099$ & $\mathbf{+0.0339}$ [$+.007$, $+.063$]\textsuperscript{$\ast$} & $-0.225$ & $+0.0044$ \\
    \textbf{V3} \emph{TRADE} & 0.556 & 0.124 & $-0.0094$ & $-0.0045$ & $\mathbf{+0.0286}$ [$+.001$, $+.059$]\textsuperscript{$\ast$} & $-0.450$ & $-0.0024$ \\
    C1 & 0.563 & 0.130 & $-0.0020$ & $+0.0015$ & $+0.0194$ & $-0.375$ & $+0.0008$ \\
    D2 & 0.559 & 0.129 & $-0.0065$ & $+0.0005$ & $+0.0104$ & $-0.750$ & $+0.0003$ \\
    B4 & 0.566 & 0.139 & $+0.0006$ & $+0.0103$ & $+0.0164$ & $-0.775$ & $+0.0004$ \\
    B3 & 0.557 & 0.135 & $-0.0078$ & $+0.0060$ & $+0.0146$ & $\mathbf{-1.025}$\textsuperscript{$\ast$} & $+0.0154$ \\
    V2 \emph{refuted} & 0.551 & 0.127 & $\mathbf{-0.0144}$ [$-.029$, $-.001$]\textsuperscript{$\ast$} & $-0.0019$ & $+0.0259$ & $-0.025$ & $-0.0148$ \\
    V1 \emph{refuted} & 0.548 & 0.119 & $\mathbf{-0.0174}$ [$-.035$, $-.001$]\textsuperscript{$\ast$} & $-0.0099$ & $+0.0227$ & $\mathbf{-1.725}$\textsuperscript{$\ast$} & $+0.0017$ \\
    \midrule
    \multicolumn{8}{l}{\emph{Deployed cell $a\num{0.8}$ / $s\num{0.5}$ --- both TRADEs collapse}} \\
    B2 \emph{refuted} & --- & --- & $\mathbf{-0.0064}$ [$-.012$, $-.001$]\textsuperscript{$\ast$} & $-0.0107$ & $-0.0052$ & $-0.175$ & $+0.0006$ \\
    V3 \emph{refuted} & --- & --- & $\mathbf{-0.0112}$ [$-.017$, $-.006$]\textsuperscript{$\ast$} & $-0.0074$ & $+0.0043$ & $-0.400$ & $-0.0039$ \\
    \bottomrule
  \end{tabular}}
\end{table}

\begin{figure}[ht]
  \centering
  \includegraphics[width=0.95\textwidth]{nexttrack_tower_bank}
  \caption{The parallel-tower-bank axis, closed on both surfaces.
  \textbf{Left:} the $N$-axis at matched capacity (${\sim}\num{3.00}\times$ the
  h256 baseline). Error bars are the paired bootstrap CI against C0 mapped onto
  the hit scale; the grey band is the split's measured resolution
  ($\pm\num{0.0124}$ recall@10 $=$ $\pm\num{17.7}$ hits,
  Figure~\ref{fig:nexttrack-power}), and every arm sits inside it. Four h211
  towers (B4) and one h512 GRU (C1) return the same \num{186} hits, paired
  $\Delta$ $\num{0.000000}$. \textbf{Middle:} the two arms that earned a walk
  TRADE at the primary cell both flip at the deployed cell --- the drift gain that
  bought the TRADE decays to a straddle while \texttt{vibe} goes below zero. Grey
  markers are point estimates whose CI straddles (the source report quotes bounds
  only for intervals clearing zero). \textbf{Right:} the campaign's one positive
  result, and it is not about topology --- at byte-identical parameters
  (\num{789696} both, \texttt{model.pt} \num{3162277} bytes both) dropout
  \num{0.1} beats dropout \num{0.0}, CI above zero.}
  \label{fig:nexttrack-tower-bank}
\end{figure}

\emph{Interpretation.} \textbf{The dose hypothesis is refuted, and it fails in
the direction opposite to the prediction.} Adding fixed-reference
\texttt{cummean} towers did not buy more of what $N=2$ bought; it monotonically
destroyed the walk. Ranked by $\Delta$\texttt{vibe} at the primary cell the
ordering is B4/B2 (all-\texttt{latent}, ${\approx}0$) $>$ B3 $>$ V3 (one
\texttt{cummean}, $-\num{0.0094}$) $>$ V2 ($-\num{0.0144}$, CI below zero) $>$
V1 ($-\num{0.0174}$, CI below zero), and the same ordering holds on S1, where the
three \texttt{cummean}-bearing arms occupy the \emph{bottom three} places
(\num{162}, \num{166}, \num{170} hits) and are the only arms whose
$\Delta$\,music@10 point estimates fall past $-\num{0.005}$. The irony is worth
stating plainly: \textbf{the tower added in order to hold a mood is the tower
that made the walk's mood measure worse} --- the only two arms whose
\texttt{vibe} is refuted outright are the two carrying the most
\texttt{cummean}. Whatever the deployed $N=2$ \texttt{latent}/\texttt{cummean}
engine is doing, ``more \texttt{cummean} is more mood-holding'' is not it, and the
structural reading that motivated the design does not survive being dosed.

\textbf{The sub-hypothesis is confirmed and extends cleanly from $N=2$ to
$N=4$.} B2 (\num{170}), B3 (\num{172}) and B4 (\num{186}) all straddle C0 and
show no trend distinguishable from width noise. This is the same mechanism the
dual-tower scan found (Section~\ref{sec:nexttrack-dualtower}) --- the raw
\texttt{latent} view subsumes its siblings --- now measured at three tower counts
instead of two, which is what closing an axis by \emph{dose} rather than by a
single point requires.

\textbf{The one positive finding is a by-product of the null, and it lands on the
deployed engine.} D2 and R2 differ in \emph{one} hyperparameter: parameters are
byte-identical (\num{789696} both, \texttt{model.pt} \num{3162277} bytes both),
only \texttt{dropout} moves, and $\Delta$recall@10 is $+\num{0.004892}$
$[+\num{0.000699},+\num{0.009783}]$, \textbf{CI above zero} --- dropout
\num{0.1} \emph{beats} dropout \num{0.0} (\num{177} against \num{170} hits).
Read carefully this does not overturn the record's three-times-confirmed
``\texttt{fusion\_layers}${}=0$ beats
\texttt{fusion\_layers}${}=1$'' result; it does something more
awkward. Dropout enters \texttt{seq\_dualgru} only through a pre-MLP, an
inter-layer gap or the fusion MLP, and at \texttt{fusion\_layers}${}=0$ with one
layer and no pre-encoder \emph{all three are absent} --- the f0 path constructs
\textbf{zero} dropout modules. So every f0 arm on record, \textbf{including the
deployed showcase engine}, trained under-regularized and was leaving
${\approx}+\num{0.005}$ recall on the table; and the three recorded f0-vs-f1 gaps
($+\num{0.016}$ / $+\num{0.020}$ / $+\num{0.0266}$,
Table~\ref{tab:nexttrack-fusion-depth}) were measured against an f1 that
\emph{had} regularization its f0 comparator lacked, so they are not like-for-like
and the true topology advantage of f0 is if anything \emph{larger} than recorded.
Wherever those three figures are quoted as clean topology measurements they
should be re-stated as not like-regularized. The cheapest actionable follow-up on
the whole board falls out of this: retrain the deployed engine at
\texttt{dropout}${}=\num{0.1}$ and re-read both surfaces --- a free improvement,
not an architecture bet.

\textbf{The validation-loss dissociation gets a fourth confirmation, now as a
rank correlation} (Figure~\ref{fig:nexttrack-bank-valloss}). Across the ten arms,
Spearman(best validation InfoNCE, hits) $=-\num{0.2954}$, $p=\num{0.407}$ --- no
usable predictive power, and if anything the wrong sign. The sharp instances are
brutal: B4 and B2 differ by \num{0.0003} in validation loss and by \num{16}
hits; V1 holds the fourth-best validation loss in the batch and the \emph{worst}
hit count; C0 holds the \emph{worst} validation loss and beats six of the eight
bank arms; and R2 reaches a better validation loss than D2 while retrieving seven
fewer hits, which makes the dissociation and the dropout finding two faces of one
coin. The decision rule never read validation loss, and this is why.

\begin{figure}[ht]
  \centering
  \includegraphics[width=0.8\textwidth]{nexttrack_bank_valloss}
  \caption{The training objective does not rank retrieval --- fourth
  confirmation, and its sharpest instance yet. Ten arms of the tower-bank batch:
  best validation InfoNCE against hits of \num{1431}. Spearman
  $=-\num{0.2954}$, $p=\num{0.407}$. B4 and B2 are \num{0.0003} apart in
  validation loss and \num{16} hits apart; C0, the worst fit in the batch, beats
  six of the eight banks. Every arm early-stopped (epoch \num{14} seven times,
  then \num{21} and \num{24} twice), so no arm's recall is a lower bound imposed
  by the epoch cap.}
  \label{fig:nexttrack-bank-valloss}
\end{figure}

\textbf{Both walk TRADEs were cell-local, which is a methodological result in its
own right.} B2 and V3 each earned a real $\Delta$drift above zero at
$a\num{0.4}$/$s\num{0.3}$ with vibe held, and both flipped to a vibe below zero
at $a\num{0.8}$/$s\num{0.5}$ with the drift gain gone. Had the design not
pre-registered the secondary-cell requirement, this campaign would have reported
two positive walk results and saved two definitions. \textbf{The secondary-cell
clause is the clause that did the work here}, and it should stay mandatory in
every future walk design. The pre-registered stride probe points the same way: at
$a\num{0.4}$ over $s\in\{\num{0.2},\num{0.3},\num{0.4}\}$ all three
\texttt{cummean} arms peak at $s\num{0.3}$, the cell they were already measured
in --- the third campaign in a row in which re-tuning the stride rescues nothing
--- and at $s\num{0.2}$ their $\Delta$\texttt{vibe} turns positive-straddling
only by taking visibly shorter steps ($\Delta$\texttt{stride\_err} above zero for
all three) and shedding genre variety ($-\num{1.90}$ / $-\num{1.78}$ /
$-\num{1.00}$, all below zero). That is precisely the trade the walk win clause
exists to reject.

\textbf{The failure signature is the opposite of the preceding architecture
campaign's, and that is diagnostic.} No arm here was refuted by a floor: the
recall floor and the relevance floor both pass for all ten runs, and eight of the
nine non-anchor arms sit at \emph{higher} artist concentration than the anchor
(\numrange{0.3547}{0.3799} against \num{0.3559}), i.e.\ the arms did not
de-eager. The floors were therefore not the operative mechanism --- these arms
are not degenerate, they are simply \textbf{flat}. That is a different failure
mode from the immediately preceding depth-in-series campaign
(Section~\ref{sec:nexttrack-stagestack}), in which all five arms fell
\emph{through} the recall floor, and it is why this axis closes as a null rather
than as a harm.

\paragraph{The scope of the closure, precisely.} What closes is
``$N\ge3$ parallel bare-GRU tower banks at \texttt{fusion\_layers}${}=0$'', and
only that, to the standard Section~\ref{sec:nexttrack-stagestack} set for a
closure on this corpus: both surfaces, walk instrument verified. It does
\textbf{not} re-close the dual-tower family, which stays reopened on the walk
surface (Section~\ref{sec:nexttrack-walksurface}) and remains the deployed engine. Nothing
was promoted as champion, no definition was saved, and no dataset was built; the
eight measurement arms were promoted for measurement only, each carrying the
verbatim note ``promoted to enable walk evaluation; not a crown claim'', and the
resulting flooding of the best-models group (to \num{15} entries) was disclosed up
front and handed to the curation agent. Zero of the three pre-authorized
contingency runs fired, and the epoch trace says why: the dropout-confound
extension needed R2 to early-stop later than D2 and both stopped at epoch
\num{14}; the epoch-cap contingency needed an arm to reach the \num{40}-epoch cap
and every arm early-stopped; the width-bracket contingency needed C1 to lose, and
it did not. One limitation is inherited rather than introduced: the walk-surface
reopening this campaign was built to follow up has \emph{no campaign report of its
own} and survives only as a facts roll-up, a git commit and the transcription in
the companion next-track document --- the harness-reproduction gate re-derived all
ten of its recorded means and its deltas exactly, so the \emph{measurement} is
reproducible and the instrument has not moved, but the primary write-up is
missing and should be backfilled.
\emph{Source: \texttt{2026-08-02-tower-bank-parallel-towers.md}.}

\section{Conclusions}
\label{sec:conclusions}
Thirteen campaigns (Jun--Jul 2026) on the pivoted rotation target converge
on a single story, with one late-breaking caveat (point 12). (1) Across two feature campaigns and three model-family scans (two
corpora), every dense modality and every dense distance/margin/dense-net family
loses to categorical identity one-hots read by gradient-boosted trees --- the
song-AE latent dilutes (Section~\ref{sec:ae-blend}), kernel SVM is
worse-than-mean (Section~\ref{sec:svm}), and on the new full-64-d corpus the two
tree families lead while the dense pyramid-MLP is last
(Section~\ref{sec:bakeoff-p64}). (2) The choice of metric is load-bearing on a
binary target: MAE is degenerate (a two-level $\epsilon$-tube predictor posts a
deceptively low MAE while being worse than the mean), and a grid-by-MAE blend
even mis-selects such a predictor; the honest classification metrics (AUC,
log-loss, Brier) restore a fair comparison (Section~\ref{sec:auc}). (3) The
campaigns delivered two measured noise bands --- $2\sigma=0.0086$ MAE and
$2\sigma=0.0126$ AUC --- and the project's first registered, promotable rotation
champion, \texttt{xgboost-classifier-rotation} at AUC 0.7429 on DS-META, plus a
saved new-corpus xgboost definition at AUC 0.7204. (4) The dense neural net was
diagnosed in full (Section~\ref{sec:anns}): an architecture scan refuted the
budget hypothesis (every geometry self-stops at 73--100 iterations and the
converged net reproduces the under-trained AUC to all digits), so the binding
constraint is the raw one-hot \emph{representation}, not optimization --- and the
net's worse-than-random log-loss beside a genuine AUC marks it as confidently
miscalibrated, not merely weak. (5) That diagnosis was then tested directly: the
Rust/\texttt{burn} \texttt{burn-deep-classifier}, which feeds the identity
one-hots \emph{raw}, repaired the calibration completely (log-loss
$2.0\to\sim 0.62$, into the tree regime) and lifted AUC $\sim$2 bands to
\num{0.7038} (Section~\ref{sec:anns-burn}) --- confirming the pathology was a
representation artifact, not an intrinsic property of nets. But the gap to the
trees only \emph{narrowed} (from $\sim$3.3 bands to $\sim$1), it did not close:
the net still lands below \texttt{catboost} and \texttt{xgboost}, the plain
\texttt{mlp} beat both the embedding and wide-deep topologies (entity embeddings
are starved at the corpus's low identity cardinality), and nothing cleared the
$\sim$0.71 save bar. (6) That standing lever was then pulled. A paired re-run on a
\emph{high-vocab} matrix --- identical to the 60/40 dataset but with the artist
cap raised $60\to300$ and genre $40\to120$ --- \emph{flipped} the topology
ranking (Section~\ref{sec:anns-highvocab}): \texttt{embeddings}, starved and
last on the 60/40 matrix, now \emph{beats} \texttt{mlp} (a $+0.0130$ swing that
reverses the prior result and confirms the starvation hypothesis directionally),
and the saved net reaches the in-dataset tree tier at AUC \num{0.7192} (within
$\sim$1 band of \texttt{catboost}) --- the first ANN to clear the rotation save
bar, saved as \texttt{burn-deep-embeddings-highvocab}. The convergent finding
still stands at the top (\texttt{xgboost} leads every matrix, and higher
resolution lifts \emph{every} family, not just the embedding), but the
ANN-pitfalls thread is now closed: a representation-aware net that consumes raw
one-hots is calibrated in the tree regime, and given enough identity resolution
its embedding topology reaches the tree tier --- identity \emph{resolution} was
the lever all along.
(7) A final probe then tested the inverse of the lever --- whether a learned representation \emph{derived from} the tree, rather than fed alongside it, could beat the bare tree --- and refuted it: the Facebook GBDT$\to$LR trick generalized to GBDT$\to$MLP, with a fresh head trained on the frozen champion's one-hot leaf indices, caps \num{1.8}--\num{2.4} bands below the bare xgboost champion (best head AUC \num{0.6978}) and is badly mis-calibrated on every head, because the leaf one-hots are a strictly lossier view of a tree whose additive leaf-value scoring already weights the leaves optimally (Section~\ref{sec:leaf-head}). The dense head cannot even recover the tree that generated its features --- the seventh convergent confirmation that identity resolution, not new heads or representations layered on the trees, is the durable lever.
(8) Finally, the standing identity-resolution lever was mapped to its edge: a vocab-cap scan on the bare GBDT --- frozen champion config, only \texttt{vocab\_top\_n} moved --- shows the lever is \emph{largely spent} (Section~\ref{sec:vocab-cap}). The first cardinality jump ($60/40\to300/120$) carries the gain; past 300/120 every step is sub-band, the curve plateaus inside the noise envelope, xgboost reaches only \num{0.7349} at 1000/300 ($+0.005$, $\sim$0.4 band over the anchor, so no save fired), and \texttt{catboost} tracks it --- two tree families agreeing the marginal high-cardinality columns carry little independent signal. Not saturated and not reversed, but diminishing: raw cap height is no longer the durable lever, and future identity work should target resolution \emph{quality} (follower/popularity interactions, genre hierarchy) rather than bigger caps --- with the full-vocabulary point (6633 cols) still unmeasured against a host memory ceiling.

(9) A tenth campaign characterized a \emph{second} representation-aware net and reached the same ceiling: the two-phase \texttt{burn-ae-classifier}, which pretrains an autoencoder over the assembled feature vector and transplants the encoder into a classifier, is flat across 13 capacity/head configs (a \num{0.0054}-wide strip inside the band), ties the plain \texttt{burn-deep} \texttt{mlp}, and stays $\sim$1.8 bands below the tree champion; freezing the MSE-pretrained encoder under a deep head is a documented collapse mode (Section~\ref{sec:anns-ae}). The unsupervised reconstruction latent is task-agnostic --- a third confirmation that a learned representation downstream of the identity one-hots buys no ranking power the trees lack.
(10) Two interpretability campaigns opened the nets with a per-model sparse autoencoder (Sections~\ref{sec:anns-sae}, \ref{sec:anns-capacity}): at equal AUC, a tapering pyramid \emph{consolidates} its interpretable concepts into a compressed bottleneck while a uniform net keeps a broader set, and at high vocab an AUC-tied embeddings/mlp pair captures the \emph{same} $\sim$0.14-decodable signal with different concept allocation --- ``same signal, different layout,'' invisible to the leaderboard.
(11) A capacity scan on the high-vocab matrix confirmed the dense net is at its ceiling there: neither widening nor deepening beats the $[128,64]$ control by a band, a third embedding layer \emph{hurts} ($-\num{0.0152}$, dilution not starvation), and the best \texttt{mlp} ties the best \texttt{embeddings} --- capacity, not the embeddings mechanism, was the lever, and it is spent (Section~\ref{sec:anns-capacity}).
(12) Finally, the honest forward question was asked --- and it reframes the whole leaderboard. A first-class chronological (train-earlier / test-later) split drops the champion from AUC \num{0.735} to \num{0.619}, a $\sim$0.11 collapse that is \emph{universal} across xgboost, catboost and logistic and is genuine taste non-stationarity, not right-censoring; a walk-forward read lands at \num{0.59}--\num{0.67} every year, never near the random baseline (Section~\ref{sec:drift}). Every on-record AUC is therefore an upper bound on same-era ranking; real forward ``will I replay a song I discover next'' skill is $\sim$0.62--0.63. The mechanism is identity churn --- unseen future artists fall into \texttt{\_\_other\_\_} and the model must lean on the weaker $\sim$0.60 intrinsic floor --- which is precisely why higher vocab caps and bigger nets, both same-era memorization levers, do not transfer forward. This does not overturn any random-split number; it clarifies what they measure and makes richer generalizable intrinsic signal the only remaining lever for forward performance.

(13) A parallel line of work --- \emph{next-track} sequence recommendation, whose task and metrics are unrelated to the rotation boards above (Sections~\ref{sec:nexttrack}--\ref{sec:nexttrack-phase2}) --- was consolidated across twenty-one campaigns and one research memo. A session GRU alone only \emph{ties} the first-order Markov co-listening bar and loses on MRR, and neither loss-tuning nor any topology (depth/width/residual/direction) beats it. But a \emph{rank-fusion} R+M z-blend (GRU $\times$ Markov) beats the bar decisively ($+61\%$ R@10, three-seed confirmed) because the two legs rank \emph{different} correct tracks; a learned XGB stacker underperforms the fixed z-blend and adding redundant neural legs does not help --- the win is signal diversity, not model count. The one genuinely additive new signal is \emph{content}: an item-representation sweep found latent \emph{dimension} (not the compressor) is the lever and peaks at PCA-192, and on that rich space a three-leg R+M+C z-blend is the current best-on-record row at R@10 \num{0.1859}, MRR \num{0.1025} (paired $\Delta$ over the prior champion $+\num{0.0133}$, CI$>$0). Two caveats bound the crown: the reproduction confirm had zero spread because \texttt{seq-nexttrack} is deterministic on this CPU (so it verified determinism, not seed-robustness), and the paired-$\Delta$ lower bound clears zero by only $+\num{0.0007}$ on a single split. Every next-track ranking predictor \emph{was} train-only and off-server until the projection arc: a supervised \emph{learned content projection} (\Rp+M+\Cp) --- an InfoNCE metric map of the PCA-192 item space in which the GRU and content legs retrieve --- then superseded the R+M+C crown at R@10 \num{0.2117} (paired $\Delta$ $+\num{0.0259}$ CI$>$0), resolving both R+M+C caveats (it is seed-robust $\times$3 and cold-reaching) and hardening on a second disjoint 0.64-cold split (paired $\Delta$ $+\num{0.0426}$ CI$>$0, wider on the colder split); it is now \emph{registered, promoted and servable} (\texttt{blend-gru-markov-content-proj}, best-models rank 1, Sections~\ref{sec:nexttrack-projection}--\ref{sec:nexttrack-projection-scan}), closing the productization gap, while every borrowed literature-fit method (session-kNN, rank fusion, learned combiners, Markov smoothing, GRU loss tricks) closed. An item-vector representation exploration then confirmed the equal-weighted PCA-192 geometry handicaps a \emph{raw} model but the fix is \emph{absorbed} by the projection (Section~\ref{sec:nexttrack-representation}), and a continuous ``sounds-alike'' \emph{music@10} metric --- correlated with Recall@10 ($r=\num{0.609}$) but crowning the content-regressing \texttt{seq-nexttrack} as musically-closest --- was added as a display-only column (Section~\ref{sec:nexttrack-music}). Finally, a
measurement campaign added a display-only \emph{session-holisticness} suite
(artist\_adj@10 / mood\_coh@10 / ild@10): the eagerness symptom \emph{tracks}
Recall@10 inside the blend family (the champion is the most artist-eager model at
artist\_adj@10 \num{0.610}) but is sourced from the \emph{Markov} bigram leg, not
the RNN (single GRU / LSTM artist\_adj@10 $\approx$ \num{0.345}), and the GRU and
LSTM are statistically tied (the ``GRU wins because more eager'' hypothesis is
refuted); of two anti-eager levers the eval-time MMR re-rank is the strong knob
($\lambda{=}\num{0.9}$ near-free) and the training-time $\beta$ regularizer is weak,
nothing was promoted, and $\lambda\approx\num{0.9}$ is flagged for a future user
holisticness decision (Section~\ref{sec:nexttrack-holisticness}).

(14) Two closing next-track architecture campaigns then shut the last open
architecture axes, and both are nulls with mechanism. A jointly-trained
\emph{dual-tower} GRU, crossing the complete
$\{$\texttt{latent},\texttt{delta},\texttt{cummean}$\}^2$ view matrix against a
\SI{0.098}{\percent} parameter-matched single tower, fails on all eight arms
(Section~\ref{sec:nexttrack-dualtower}) --- the \emph{fourth} independent
learned-combiner null and the decisive one, because it answers the standing
rebuttal to the earlier three (``of course a bolted-on combiner fails --- train
the fusion end-to-end''): trained end-to-end at the representation level, it
still fails. The mechanism is that the raw \texttt{latent} view already
\emph{subsumes} the others (a GRU over raw latents integrates movement and
running mean internally; the one genuinely complementary pair,
\texttt{delta}$+$\texttt{cummean}, beats both its parents yet is hopeless in
absolute terms). Two by-products are durable: the MLP fusion head is actively
\emph{harmful} (removing it wins in every regime measured, three times over), and
single-GRU width is now closed on PCA-192 as it already was on AE-64. A learned
per-step \emph{pre-encoder} then fails harder
(Section~\ref{sec:nexttrack-preencoder}): the decisive single-tower ablation
\emph{loses} outright ($-\num{0.0140}$, CI below zero, \num{156} hits against
\num{176}), the dose curve is flat to twice the input dimension, and placement is
irrelevant. Because a \emph{linear} pre-encoder is provably expressivity-neutral
here, the rectifier is the only candidate cause --- and a counterfactual
interpretability read confirmed it, measuring the affine map neutral at
${\sim}0.8\sigma$ and the ReLU costing ${\sim}5.9\sigma$, with
\texttt{frac\_exact\_zero} \num{0.5327} matching
\texttt{frac\_negative\_coords} \num{0.5327} to four digits: half-wave
rectification of a zero-centred latent deletes exactly its negative half. Across
five architecture nulls the pattern is hard to dismiss --- adding representation
machinery \emph{around} a single GRU over a good latent space does not help on
this corpus, while every real gain came from improving the latent space itself or
from the fixed z-blend. It is the same shape of finding as the rotation line's
own arc (points 1, 7, 9): learned machinery layered on an already-good
representation buys nothing, and \emph{resolution of the representation} is the
durable lever.
(15) Finally, two facts about the \emph{instrument} rather than any model reframe
the architecture arc (Section~\ref{sec:nexttrack-power}). The nulls are
\textbf{power-limited}: at $n=\num{1431}$ the paired-bootstrap CI half-width is
\num{0.0124} on Recall@10 and ${\sim}\num{0.0018}$ on holisticness@10, while a
whole architecture family spans only \num{0.028} and \num{0.0031} --- so of eight
adjacent arms roughly two tiers are separable on recall and one on holisticness,
and resolving a \num{0.005} difference would need ${\sim}\num{8800}$ sessions.
The honest reading is that \emph{large harms were detected reliably while
moderate wins were never detectable at all}, which is a different claim from
``these architectures do not work'' and yields a prioritization rule: further
architecture arms are a poor use of runs on this split whatever they return,
because the axis's whole dynamic range is barely twice the measurement error.
And \emph{holisticness@10} exists in three mutually non-comparable code
generations (3-factor; 4-factor grounded; grounded plus artist-concentration),
whose differences dwarf every model difference and under all of which the
servable next-track champion ranks \emph{last} --- so the next-track board stays
denominated in Recall@10 and every holisticness column is tagged with its
generation (Table~\ref{tab:nexttrack-hgen}). Note the parallel with point (12):
both lines of work end at a measurement caveat rather than a modelling one ---
rotation's random-split AUC is an upper bound on forward skill, and next-track's
architecture nulls are bounded by test-set size.
(16) A seventh next-track architecture campaign then closed the \emph{parallel}
branch of the axis, and it is the second verdict here to rest on \textbf{two}
evaluation surfaces (point (19) is the first) --- the familiar one-shot top-10,
and the autoregressive
\emph{walk} the showcase actually serves (Section~\ref{sec:nexttrack-towerbank}).
A bank of $N$ parallel bare-GRU towers under one linear head, matched at
\num{3.00}$\times$ the h256 baseline and dosed from $N=2$ to $N=4$, wins on
neither surface: every arm \emph{ties} the baseline on one-shot recall while
\textbf{five of eight lose to a parameter-matched single GRU} with the CI below
zero, and the two arms that earn a walk trade at the primary cell both flip to
\texttt{vibe} below zero at the deployed cell, so no definition was saved. The
decisive datapoint is a zero: four h211 towers (\num{1187700} params) and one
h512 GRU (\num{1182912} params) return the \emph{identical} \num{186} hits of
\num{1431}, paired $\Delta$ \num{0.000000} --- \emph{topology and width are
interchangeable at this budget}, which is also the cleanest demonstration on
record of why the parameter-matched control leg is mandatory: against the anchor
alone, five losing arms would have read as ties. The campaign was built to dose
the walk win's own stated mechanism (a fixed reference never overwritten by
recurrence) and the dose is refuted \emph{ironically} --- the \texttt{cummean}
towers added in order to hold a mood are the worst arms in the batch on both
surfaces and the only two whose walk \texttt{vibe} is refuted outright. The one
CI-clear result in the whole batch is a by-product about \emph{regularization},
not topology: at byte-identical parameters dropout \num{0.1} beats dropout
\num{0.0} by $+\num{0.0049}$ recall@10. That is consequential because the
\texttt{fusion\_layers}${}=0$ path constructs \emph{zero} dropout modules, so
every f0 arm on record --- including the deployed showcase engine --- trained
under-regularized, the three recorded f0-vs-f1 gaps are not like-regularized, and
the cheapest improvement on the board is a single retrain. It is the same shape as
points (7), (9) and (14): machinery layered on an already-good representation
buys nothing, and the durable levers are the representation and the training
regime.

(17) Six further next-track campaigns then complete that line, and the two
\emph{positive} results in them are both score-time changes rather than models.
The eval-time MMR $\lambda$ frontier (Section~\ref{sec:nexttrack-mmr}) is the first
confirmed holisticness lever on record --- eight of eight arms lift it with the CI
above zero, three single-GRU arms clear the prior all-time bar at a recall tie, and
three of the four crown factors improve together --- and it resolves its
pre-registered mechanism question unanimously: sonic diversification
de-concentrates \emph{artists} $\numrange{1.59}{3.06}\times$ faster than it reduces
seed-echo in 8 of 8 arms, so the explicit artist-penalty sequel is cancelled. Its
by-product is a measurement rule: shared-weights pairing is
$\numrange{3}{33}\times$ sharper than the architecture-family noise floor of point
(15), so eval-time levers must budget their minimum detectable effect from the
pairing. Then \texttt{markov\_gate}
(Section~\ref{sec:nexttrack-markovgate}) delivers the largest single recall gain in
the whole arc from nine lines and zero new parameters: making the Markov leg
\emph{abstain} where the last prefix item has no train-bigram support
(\SI{55.6}{\percent} of queries) lifts Recall@10 to \num{0.23201} (\num{332} of
\num{1431}), paired $+\num{0.02027}$ canonically and $+\num{0.01957}$ on a disjoint
\SI{64}{\percent}-cold split. The defect is the blend's \emph{scale-invariant}
normalizer re-inflating an evidence-free back-off to a full one-third vote, and the
bucket decomposition --- whole effect in the zero-support bucket, both warm buckets
\textbf{byte-for-byte identical} --- mechanically proves that no global static leg
weight could have found it, which \emph{narrows} the four-times-confirmed
learned-combiner pitfall to \emph{learned score-level fusion} rather than per-query
evidence gating. It is a recommendation, not a promotion.
(18) The negative results in the same window are equally load-bearing. The
training-time anti-eager regularizer is refuted \emph{expressively}, not
dosimetrically (Section~\ref{sec:nexttrack-eagermargin}): its margins were
calibrated from the item$\to$item transition tail while the hinge acts on
prediction$\to$current, whose ceiling lies below the quantiles used, so four arms
reproduce the control's loss to \num{16} significant digits and the arms that do
fire destroy relevance in the same motion that de-eagers. What \emph{does} work is
the first lever in this project that cannot be outscored: a \textbf{hard per-artist
cap} on the top-$k$ (Section~\ref{sec:nexttrack-artistcap}), which takes
\texttt{artist\_conc@10} \num{0.692}$\to$\num{0.113} and generation-3
holisticness@10 \num{0.0086}$\to$\num{0.0481} while \texttt{music@10} \emph{rises}
with the CI above zero at every setting on both models --- it trades \emph{literal}
relevance for \emph{graded} relevance --- hardened on a disjoint
\SI{64}{\percent}-cold split on two pre-registered legs, at a stated price of
\num{52} of \num{1431} exact hits. A purpose-built rollout metric
(Section~\ref{sec:nexttrack-rollout}) returns a sixth architecture null in the
deployment regime and finds per-seed variance dominating model identity; and a
research memo (Section~\ref{sec:nexttrack-memo}), reported at a lower register
because nothing in it was run on the server, converges from five independent
directions on the thesis that every surviving proposal is an
\emph{evidence-conditioning} change rather than a capacity change --- the same shape
as points (1), (7), (9) and (14) on the rotation side.
(19) A stage-stack depth campaign (Section~\ref{sec:nexttrack-stagestack}) finally
closes the \emph{serial} branch of the architecture axis on \emph{both} evaluation
surfaces, the standard the tower-bank campaign of point (16) then inherits. Five
capacity-equalised arrangements of per-step MLP and recurrent stages all fall
\emph{through} the mandatory recall floor and are refuted again on the walk, while a
$\num{3.24}\times$ capacity control ties on both --- so the damage is depth itself,
and not capacity, regularization or the MLP (a leading per-step MLP is
recall-neutral and holisticness-positive at depth 4, which narrows the $N=1$
pre-encoder result of point (14) to the shallow case). Two facts outlive the null:
bracketing beats interleaving at byte-identical capacity ($+\num{0.01537}$, CI
above zero), and the training objective \textbf{dissociates} from retrieval --- the
arm with the best validation InfoNCE in the batch retrieves \num{55} fewer hits
than the control, so any architecture search on this family that selects on
validation loss selects the wrong arm.

\section{Source reports}
This synthesis consolidates thirty-five source documents from \texttt{experiments/}
--- thirty-four campaign reports and one research memo, the latter flagged as such
because nothing in it was run on the server:
\begin{itemize}
  \item \texttt{2026-06-09-ae-xgboost-blend.md} --- autoencoder-latent +
    XGBoost blend on rotation (Section~\ref{sec:ae-blend}).
  \item \texttt{2026-06-09-svm-kernel-rotation.md} --- kernel SVM scan and
    conditional blend; first measured MAE noise band (Section~\ref{sec:svm}).
  \item \texttt{2026-06-09-rotation-classification.md} --- AUC bake-off, first
    measured AUC noise band, first registered rotation champion
    (Section~\ref{sec:auc}).
  \item \texttt{2026-06-12-classifier-bakeoff-p64.md} --- four-family
    classifier bake-off and xgboost depth/lr sweep on the new full-64-d corpus
    (Sections~\ref{sec:bakeoff-p64} and~\ref{sec:anns}).
  \item \texttt{2026-06-12-pyramid-mlp-arch-scan.md} --- six-geometry
    pyramid-MLP architecture scan, refuting the budget hypothesis
    (Sections~\ref{sec:mlp-scan} and~\ref{sec:anns}).
  \item \texttt{2026-06-12-burn-deep-topology.md} --- the Rust/\texttt{burn}
    \texttt{burn-deep-classifier} topology bake-off; resolves the
    representation-aware-net question, fixing calibration but not the AUC gap
    (Section~\ref{sec:anns-burn}).
  \item \texttt{2026-06-12-high-vocab-topology.md} --- the paired topology
    re-run on a higher-identity-cardinality matrix; flips
    \texttt{embeddings}$>$\texttt{mlp} and lifts the saved net into the
    in-dataset tree tier, the first ANN to clear the rotation save bar
    (Section~\ref{sec:anns-highvocab}).
  \item \texttt{2026-06-15-gbdt-leaf-head-ceiling-probe.md} --- the Facebook
    GBDT$\to$LR trick generalized to GBDT$\to$MLP; every leaf-embedding head
    caps below the bare xgboost champion with broken calibration, the seventh
    convergent refutation (Section~\ref{sec:leaf-head}).
  \item \texttt{2026-06-15-vocab-cap-scan-gbdt.md} --- the identity-resolution vocab-cap scan on the bare GBDT; maps the cap$\to$AUC curve and finds the lever largely spent past 300/120 (sub-band gains, no save fired), with the all-vocab point unmeasured at the host memory ceiling (Section~\ref{sec:vocab-cap}).
  \item \texttt{2026-06-15-burn-ae-classifier-capacity-scan.md} --- first
    characterization of the two-phase \texttt{burn-ae-classifier}; 13 configs in
    one sub-band strip, ties the plain \texttt{burn-deep} \texttt{mlp}, and a
    frozen-deep-head collapse mode (Section~\ref{sec:anns-ae}).
  \item \texttt{2026-07-12-pyramid-vs-rect-sae.md} --- per-model SAE
    representation read; at equal AUC the pyramid consolidates concepts into its
    bottleneck while the uniform net keeps a broader set
    (Section~\ref{sec:anns-sae}).
  \item \texttt{2026-07-12-high-vocab-capacity-scan.md} --- the
    \texttt{burn-deep-classifier} capacity scan on the high-vocab matrix
    (CEILING) plus a cross-family SAE allocation read
    (Section~\ref{sec:anns-capacity}).
  \item \texttt{2026-07-12-taste-drift-temporal-split.md} --- the chronological
    versus random split; a universal $\sim$0.11 AUC drop reframing the
    leaderboard as an upper bound on same-era ranking (Section~\ref{sec:drift}).
  \item \texttt{2026-07-15-nexttrack-compression-representation-bakeoff.md} ---
    the first \emph{next-track} sequence-recommendation campaign folded in: an
    item-representation compression bake-off (PCA vs autoencoder $\times$ latent
    dim 32/64/128) under a frozen GRU readout. PCA-128 is the single-GRU winner
    (paired-$\Delta$ $+0.0189$, CI$>$0) but the win does not carry into the R+M
    blend, so nothing was promoted (Section~\ref{sec:nexttrack-compression}).
  \item \texttt{2026-07-13-nexttrack-sequence-sweep.md} --- the founding
    next-track campaign: a GRU/LSTM sequence sweep against the first-order
    Markov bar; the best GRU-infonce ties on R@10 but loses on MRR
    (Section~\ref{sec:nexttrack-sweep}).
  \item \texttt{2026-07-13-nexttrack-topology-scan.md} --- depth / width /
    residual / bidirectional topology scan; architecture is a spent lever, all
    arms cap below the Markov bar (Section~\ref{sec:nexttrack-topology}).
  \item \texttt{2026-07-14-nexttrack-blend-and-ensemble.md} --- the
    GRU$\times$Markov R+M z-blend beats the bar $+61\%$ R@10 (the first
    next-track champion); no richer ensemble improves on it and the learned XGB
    stacker underperforms (Section~\ref{sec:nexttrack-blend}).
  \item \texttt{2026-07-14-nexttrack-graded-relevance-and-new-legs.md} ---
    first-class graded-relevance metrics (artist@k / genre@k / artist\_mrr) and
    new leak-free base learners; content-kNN is the only additive signal, tying
    the champion on AE-64 (Section~\ref{sec:nexttrack-graded}).
  \item \texttt{2026-07-15b-nexttrack-compression-dim-extension.md} --- the
    PCA/AE dimension sweep extended to 192/256; the curve peaks at PCA-192 and
    turns over, so PCA-128 stays the Phase-2 base (Section~\ref{sec:nexttrack-dim}).
  \item \texttt{2026-07-15c-nexttrack-phase2-model-sweep.md} --- the
    model-family sweep on the rich PCA spaces; the fixed z-blend beats the
    learned stacker, LSTM is the best single model, and R+M+C on PCA-192 is
    crowned the new best-on-record row (Section~\ref{sec:nexttrack-modelsweep}).
  \item \texttt{2026-07-18-nexttrack-literature-fit-campaign.md} --- a
    bibliography-driven exploration; a supervised learned content projection
    (\Rp+M+\Cp) is a new best-on-record row (R@10 \num{0.2117}, paired $\Delta$ over
    R+M+C $+\num{0.0259}$ CI$>$0, seed-robust, cold-reaching), and every borrowed
    method (session-kNN, rank fusion, learned combiner, Markov smoothing, GRU
    importance negatives) closes (Section~\ref{sec:nexttrack-projection-win}).
  \item \texttt{2026-07-18b-nexttrack-projection-registration-and-scan.md} --- the
    projection reproduced on-server, registered $+$ promoted $+$ servable as
    \texttt{blend-gru-markov-content-proj} (best-models rank 1); an objective / rank
    / $\tau$ scan (InfoNCE $\gg$ BPR; default held) and second-split hardening on a
    disjoint 0.64-cold split (paired $\Delta$ $+\num{0.0426}$ CI$>$0)
    (Section~\ref{sec:nexttrack-projection-scan}).
  \item \texttt{2026-07-18c-nexttrack-representation-exploration.md} --- an
    item-vector representation audit; the equal-weighted PCA-192 geometry handicaps
    the raw GRU (six of eight alt spaces lift it CI$>$0) but the gain is absorbed by
    the projection (Section~\ref{sec:nexttrack-representation}).
  \item \texttt{2026-07-20-musical-distance-metric.md} --- the continuous
    ``sounds-alike'' music@10 metric and its 53-run offline backfill; it correlates
    with Recall@10 ($r=\num{0.609}$) but crowns a different winner, the
    content-regressing \texttt{seq-nexttrack} (Section~\ref{sec:nexttrack-music}).
  \item \texttt{2026-07-22-session-holisticness-and-antieager-levers.md} ---
    a display-only session-holisticness metric suite (artist\_adj@10 /
    mood\_coh@10 / ild@10) and two anti-eager levers (eval-time MMR
    \texttt{mmr\_lambda}, training-time \texttt{eager\_beta}); the champion is
    the most artist-eager model, the eagerness is sourced from the Markov leg,
    GRU\,$\approx$\,LSTM, and MMR $\lambda$ is the strong knob
    (Section~\ref{sec:nexttrack-holisticness}).
  \item \texttt{2026-07-25-nexttrack-dual-tower-fusion-scan.md} --- the
    dual-tower view-pairing $\times$ fusion-depth scan against a
    \SI{0.098}{\percent} parameter-matched single tower; \textbf{REFUTED} (the
    fourth learned-combiner null, closing the ``train the fusion end-to-end''
    rebuttal), with the raw-latent view shown to subsume
    \texttt{delta}/\texttt{cummean}, the MLP fusion head shown actively harmful,
    single-GRU width closed on PCA-192, and holisticness/recall shown
    anti-correlated across the batch (Section~\ref{sec:nexttrack-dualtower}).
  \item \texttt{2026-07-26-nexttrack-pre-encoder-scan.md} --- the learned
    per-step pre-encoder scan (presence $\times$ placement $\times$ dose) with the
    decisive single-tower ablation; \textbf{REFUTED and actively harmful}
    ($-\num{0.0140}$ Recall@10, CI below zero), the rectifier identified
    analytically and then measured (affine map neutral at ${\sim}0.8\sigma$, ReLU
    ${\sim}5.9\sigma$, \texttt{frac\_exact\_zero} $=$
    \texttt{frac\_negative\_coords} $=$ \num{0.5327}), plus the measured power
    analysis and the three holisticness generations
    (Sections~\ref{sec:nexttrack-preencoder} and~\ref{sec:nexttrack-power}).
  \item \texttt{2026-07-26-nexttrack-mmr-lambda-frontier.md} --- the native
    eval-time MMR $\lambda$ frontier on both anchors; \textbf{CONFIRMED} as the
    first genuine holisticness lever on record (three TIER-1 ``free'' single-GRU
    arms above the prior all-time bar at a recall tie, best $H$ \num{0.030355}),
    with the pre-registered mechanism question resolved unanimously to the
    responsive hypothesis (8 of 8 arms, ratio \numrange{1.59}{3.06}), the binding
    facet shown to switch as $\lambda$ falls, the candidate-pool axis closed as a
    fine-tuning knob, and the shared-weights pairing measured
    $\numrange{3}{33}\times$ sharper than the architecture-family noise floor
    (Section~\ref{sec:nexttrack-mmr}).
  \item \texttt{2026-07-27-rollout-hold-metric.md} --- \texttt{hold@}$N$, a
    generation metric for the autoregressive rollout regime (per-axis fade
    calibrated on \num{2135} real train sessions, with the anti-gaming term inside
    the metric), plus five sweeps and \num{240} journeys; \textbf{REFUTED} as a
    model discriminator --- a sixth architecture null, obtained in the deployment
    regime --- with per-seed variance shown to dominate model identity and three
    retracted positives documented as a catalogue of generation-metric pitfalls
    (Section~\ref{sec:nexttrack-rollout}).
  \item \texttt{2026-07-27b-nexttrack-model-topology-research-memo.md} ---
    \textbf{a research memo, not a campaign report}: nothing was run on the
    server, nothing promoted, nothing registered, and its numbers come from
    offline screens over frozen checkpoints. Consolidated at a lower register,
    with its own verified/provisional partition preserved and its test-set
    screening caveat stated; its reusable content is the convergent thesis that
    every surviving proposal is an \emph{evidence-conditioning} change rather than
    a capacity change, the mechanical explanation of the learned-combiner anomaly,
    and the graveyard (Section~\ref{sec:nexttrack-memo}).
  \item \texttt{2026-07-27c-nexttrack-markov-gate-confirm.md} --- the native
    two-split confirm of \texttt{markov\_gate}; \textbf{CONFIRMED} and the largest
    single recall gain in the arc (R@10 \num{0.23201} $=$ \num{332} of \num{1431},
    paired $+\num{0.02027}$ canonically and $+\num{0.01957}$ on a disjoint
    \SI{64}{\percent}-cold split, MRR up on both), with the warm evidence buckets
    \emph{byte-for-byte identical}, an unpredicted crown gain in which all four
    factors move the right way at once, a first recall-neutral crown arm when
    composed with MMR, and the learned-combiner pitfall narrowed accordingly
    (Section~\ref{sec:nexttrack-markovgate}).
  \item \texttt{2026-07-30-eager-margin-calibrated-refutation.md} --- the
    training-time eagerness regularizer at margins calibrated from the measured
    transition tail; \textbf{REFUTED across the whole axis, with a mechanism}: the
    margins were drawn from the item$\to$item cosine while the hinge acts on
    prediction$\to$current, whose ceiling sits between \num{0.5843} and \num{0.70},
    so four arms are provably inert (loss traces bit-identical at $\beta$ up to the
    schema cap) and the arms that fire destroy relevance before they help
    (Section~\ref{sec:nexttrack-eagermargin}).
  \item \texttt{2026-07-30c-artist-cap-measurement.md} --- the \emph{hard
    per-artist cap}, the first de-eagering lever that works, and it works because
    it cannot be outscored: \texttt{artist\_conc@10}
    \num{0.692}$\to$\num{0.113} with \texttt{music@10} \emph{rising} (CI above
    zero) at every setting on both models, hardened on a disjoint
    \SI{64}{\percent}-cold split on two pre-registered legs, at a stated price of
    \num{52} of \num{1431} exact hits; it is also where this document introduces
    the second evaluation surface, the autoregressive walk
    (Section~\ref{sec:nexttrack-artistcap}).
  \item \texttt{2026-08-01-stage-stack-topology.md} --- the stage-stack
    topology campaign (arrangements of per-step MLP $\times$ recurrence at
    equalised capacity, via the new \texttt{seq-stack} predictor);
    \textbf{REFUTED on both surfaces} and the first closure on this corpus to that
    standard: all five 4-stage arms fall \emph{through} the recall floor and are
    refuted again on the walk, the $\num{3.24}\times$ capacity control ties on
    both, the user's rectifier-position hypothesis is refuted, bracketing beats
    interleaving at byte-identical capacity, and the best-fitting arm in the batch
    retrieves \num{55} fewer hits than the control
    (Section~\ref{sec:nexttrack-stagestack}).
  \item \texttt{2026-08-02-tower-bank-parallel-towers.md} --- the
    parallel-tower-bank dose campaign ($N$ bare-GRU towers under one linear head,
    via the new \texttt{seq-bank} predictor) and the first campaign here judged on
    \emph{two} surfaces; \textbf{CLOSED on both}: no arm wins, five of eight lose
    to a parameter-matched single GRU, four h211 towers and one h512 GRU return the
    identical \num{186} hits (paired $\Delta$ \num{0.000000}), the
    \texttt{cummean} dose makes the walk strictly worse, both walk trades are
    cell-local, and the batch's one CI-clear result is that dropout \num{0.1}
    beats dropout \num{0.0} at byte-identical parameters --- which implies every
    \texttt{fusion\_layers}${}=0$ arm on record, the deployed engine included,
    trained under-regularized (Section~\ref{sec:nexttrack-towerbank}).
\end{itemize}

\end{document}
